\section{Irreducible supports: topology, conductor and equality}
\label{sec:II:night-sign}
The following proofs retain every Part~4 result of the night run. The final classification is Theorem~\ref{night:P4.58}; its exact normalization maps and its local exclusions are proved below.
\paragraph{Standing torus-fibration setting.} The ambient space is a smooth compact complex threefold $Y$ with $b_1(Y)=b_2(Y)=0$, and $f:Y\to B$ is proper and surjective with connected fibres and smooth general fibre a complex two-torus. Every statement below concerning an irreducible support uses its actual full fibre $mS$, normalization and resolution. The global assertions expressly retaining the hypothesis that every support is irreducible or has normal crossings are used with that hypothesis. Constructions and comparison examples with different Betti numbers state their own assumptions.
\subsection{The explicit threefold with identity monodromy}

Throughout this part we orient complex spaces by their complex structures.
All homology groups in the example have coefficients in $\Z$. A meridian is
counterclockwise in its indicated holomorphic coordinate; at infinity this
coordinate is $1/t$. We write
\[
 E=\C/(\Z+i\Z)=\C^*/q^\Z,\qquad q=e^{-2\pi},\qquad
 A=E^2,\qquad \Lambda=\Z e_1\oplus\Z e_2\oplus\Z e_3\oplus\Z e_4,
\]
where $(e_1,e_2,e_3,e_4)=((1,0),(i,0),(0,1),(0,i))$. The comparison
$E\longrightarrow\C^*/q^\Z$ is $z\mapsto e^{2\pi iz}$.

\begin{proposition}[P4.1]\label{night:P4.1}\label{stmt:P4.1}
\StatementRecord{night:P4.1}{proved}{Historical events: GPT-6 Astra (OpenAI), ChatGPT Pro mode, effort Pro, S6\_06\_BLIND, 16--18 November 2026: targeted proof check. Also checked in this writing. Pre-continuation writing events: GPT-6 Astra (OpenAI), ChatGPT Pro mode, effort Pro, S6\_10\_PAPERS, 18 November 2026: targeted proof check; used and rechecked in this run. Continuation metadata events: GPT-6 Astra (OpenAI), ChatGPT Pro mode, effort Pro, S6\_10\_PAPERS continuation, 19 November 2026: used and rechecked in this run; status-vocabulary reconciliation; metadata reconciliation only; no new mathematical check. \emph{Pre-continuation recorded proof scope.} The night-run re-derivation was by GPT-6 Astra (OpenAI, ChatGPT Pro mode, effort Pro), 16--18 November 2026. Also checked in this writing by GPT-6 Astra (OpenAI), ChatGPT Pro mode, effort Pro, 18 November 2026. The two actual blowup charts and line frames give lambda\_B=v lambda\_A, and the canonical section gives the punctured product map lambda\_A/u=lambda\_B/s. The scalar quotient is free, proper and Hausdorff. Tautological evaluation identifies the exceptional nonzero line with C2 minus zero equivariantly. The other component has primitive circle Chern class, b1=3 and trivial canonical line, matching the currently opened BHPV classification table. The map uv gives the exact reduced SNC central divisor}{\ref{night:P4.3}}

Let $V=\operatorname{Bl}_{(0,0)}(B\times\Delta)$, where $B=E$, let $R$ be
the exceptional curve, and put $L=\OO_V(R)$. The holomorphic quotient
$U=L^\times/q^\Z$ is a smooth Hausdorff threefold, proper over $\Delta$.
Its noncentral fibres are $E\times B$. Its central fibre is the reduced
union of a primary Kodaira surface and the scalar Hopf surface
$(\C^2\setminus\{0\})/q^\Z$, meeting transversally in one elliptic curve.
\end{proposition}
\begin{proof}
\noindent\textbf{Step 1.} We use coordinates $(z,t)$ at the blown-up point. The blowup has charts
\[
 (z,t)=(u,uv),\qquad (z,t)=(rs,s),\qquad r=v^{-1},\quad s=uv.
\]
The frames $u^{-1}$ and $s^{-1}$ for $L$ give fibre coordinates related by
$\lambda_B=v\lambda_A$. Each nonzero fibre coordinate, modulo multiplication
by $q$, gives a product chart with fibre $E$. The transitions are holomorphic
translations on $E$. The scalar action is free and properly discontinuous:
the logarithm of the norm of a nonzero fibre coordinate changes by the
nonzero constant $\log q$, and compact subsets meet only finitely many
translates. Thus the quotient is Hausdorff and smooth. These product charts
have compact fibre $E$, so $U\to V$ is proper. The blowup and $B\times\Delta
\to\Delta$ are proper, proving properness of their composite.

\noindent\textbf{Step 2.} The canonical section of $\OO(R)$ has coefficients $u$ and $s$. Over
$t\ne0$ it never vanishes. The ratios $\lambda_A/u=\lambda_B/s$ therefore
give the specified holomorphic product trivialization over $\Delta^*$.

\noindent\textbf{Step 3.} The central divisor in $V$ is $B'\cup R$, with $R^2=-1$ and $R\cdot B'=1$.
Thus $L|_R=\OO_{\Pp^1}(-1)$ and $L|_{B'}$ has degree one. Over $B'$ the
quotient is a principal elliptic bundle with primitive nonzero circle Chern
class. Its canonical bundle is trivial: the base differential wedged with
$d\lambda/\lambda$ is invariant under the line-bundle transition, because
the additional term is pulled back from the one-dimensional base. The
circle-bundle calculation in P4.3 gives $b_1=3$. The surface classification
The classification \cite[Chapter~VI, Theorem~1.1 and Table~10, pp.~243--244]{BHPV2004E} consequently identifies this surface as
a primary Kodaira surface.

\noindent\textbf{Step 4.} The tautological evaluation morphism
$\OO_{\Pp^1}(-1)^\times\to\C^2\setminus\{0\}$ is a biholomorphism: a
nonzero vector uniquely determines its line. It intertwines both scalar
actions, and hence identifies the other component with the stated Hopf
surface. Over $B'\cap R$ the components meet in $E$. In local product
coordinates the map is $(u,v,c)\mapsto uv$; this proves transversality,
multiplicity one, and the absence of other critical points.
\end{proof}

\begin{proposition}[P4.2]\label{night:P4.2}\label{stmt:P4.2}
\StatementRecord{night:P4.2}{proved}{Historical events: GPT-6 Astra (OpenAI), ChatGPT Pro mode, effort Pro, S6\_06\_BLIND, 16--18 November 2026: targeted proof check. Also checked in this writing. Pre-continuation writing events: GPT-6 Astra (OpenAI), ChatGPT Pro mode, effort Pro, S6\_10\_PAPERS, 18 November 2026: targeted proof check. Continuation metadata events: GPT-6 Astra (OpenAI), ChatGPT Pro mode, effort Pro, S6\_10\_PAPERS continuation, 19 November 2026: used and rechecked in this run; status-vocabulary reconciliation; metadata reconciliation only; no new mathematical check. \emph{Pre-continuation recorded proof scope.} The night-run re-derivation was by GPT-6 Astra (OpenAI, ChatGPT Pro mode, effort Pro), 16--18 November 2026. Also checked in this writing by GPT-6 Astra (OpenAI), ChatGPT Pro mode, effort Pro, 18 November 2026. The logarithm divided by 2pi i defines a single-valued elliptic-valued transition because its increment is the integral lattice generator. The three displayed holomorphic group isomorphisms, their determinant-one real lattice matrices, local properness and point separation give the smooth compact Hausdorff total space. Their punctured trivializations make the remaining return an actual translation, which acts identically on integral H1}{\ref{night:P4.1}}

The three markings and the principal-bundle clutch in the prompt define a
smooth compact complex threefold $f:Y\to\Pp^1$ with connected fibres,
exactly three singular fibres of the form in P4.1, and identity monodromy
on $H_1$ of the regular fibre.
\end{proposition}
\begin{proof}
\noindent\textbf{Step 1.} The map $\ell(t)=(\log t)/(2\pi i)$ with values in $E$ is single-valued and
holomorphic on an annulus: changing the logarithm adds an integer.
Consequently $a_{\rm out}=a_{\rm in}+(\ell(t),0)$ is a holomorphic transition
for a principal $A$-bundle on $\Pp^1$. Product charts with compact fibre
prove that this bundle is Hausdorff and compact.

\noindent\textbf{Step 2.} We choose disjoint discs about $0,1,\infty$ and replace their product models
by P4.1 using the holomorphic group isomorphisms
\[
 C_i\times B_i\longrightarrow A:\qquad
 (z,w)\mapsto(w,z),\quad(iz,iz+w),\quad(z,iz+w),
\]
respectively. Each lattice matrix $S_i$ displayed below has determinant
one. These maps, $\lambda_B=v\lambda_A$, the canonical punctured-disc
trivialization in P4.1, and the clutch $\ell$ specify all transitions.
Different base points separate in the base. Points over the same base
point lie in one Hausdorff local model, which separates them. The gluing is
therefore Hausdorff. Properness holds over each base chart and hence over
the base; compactness follows.

\noindent\textbf{Step 3.} The two central components have connected elliptic fibres over connected
curves and meet each other. All remaining fibres equal $A$. The equations
$uv=t$ give exactly the indicated critical values. The punctured-disc
trivializations identify all fibre lattices with $\Lambda$, and the remaining
transition is a translation on $A$. Translation induces the identity on
$H_1(A;\Z)$, giving the asserted monodromy representation.
\end{proof}

\begin{proposition}[P4.3]\label{night:P4.3}\label{stmt:P4.3}
\StatementRecord{night:P4.3}{proved}{Historical events: GPT-6 Astra (OpenAI), ChatGPT Pro mode, effort Pro, S6\_06\_BLIND, 16--18 November 2026: targeted proof check. Also checked in this writing. Pre-continuation writing events: GPT-6 Astra (OpenAI), ChatGPT Pro mode, effort Pro, S6\_10\_PAPERS, 18 November 2026: targeted proof check. Continuation metadata events: GPT-6 Astra (OpenAI), ChatGPT Pro mode, effort Pro, S6\_10\_PAPERS continuation, 19 November 2026: used and rechecked in this run; status-vocabulary reconciliation; metadata reconciliation only; no new mathematical check. \emph{Pre-continuation recorded proof scope.} The night-run re-derivation was by GPT-6 Astra (OpenAI, ChatGPT Pro mode, effort Pro), 16--18 November 2026. Also checked in this writing by GPT-6 Astra (OpenAI), ChatGPT Pro mode, effort Pro, 18 November 2026. The blowup is the explicit Hopf-boundary homotopy pushout B wedge S2. The circle Chern class is zero on a regular B and minus one on the exceptional sphere, yielding the pushout T3 from S1 to S3. After its product with the norm circle, the only cellular boundaries are r to a and rc to a wedge c. This proves the two precise inclusion quotients and free ranks. The canonical section trivializes the stated meridian bounding chains, including every mixed product}{the hypotheses and explicit local calculation stated here}

Order the local fibre lattice as $(a,c,b,d)$, where $(a,c)$ are the cycles
$(1,i)$ in $C_i$. The regular-fibre inclusions induce the quotient maps
\[
 H_1(A)=\Lambda\longrightarrow\Lambda/\Z a=H_1(U),\qquad
 H_2(A)=\bigwedge^2\Lambda\longrightarrow
 \bigwedge^2\Lambda/\Z(a\wedge c)=H_2(U).
\]
Both target groups are free. The meridian and every product of the
meridian with a fibre one-cycle map to zero in these degrees.
\end{proposition}
\begin{proof}
\noindent\textbf{Step 1.} A Hermitian metric gives a diffeomorphism
$L^\times/q^\Z\cong S(L)\times S^1_c$, with factors the circle bundle and
the logarithm of the norm modulo $\log q$. We first specify a homotopy
model for $V$. Blowing up a four-ball replaces it by the disc bundle of
$\OO(-1)$ along its Hopf boundary. The complement of the ball in
$B\times\Delta$ has homotopy type $B\vee S^3$; the boundary sphere is
the $S^3$ summand, since its map to $B$ is nullhomotopic. The resulting
homotopy pushout is $B\vee S^2$. Equivalently, on simply connected
covers the replacement adds one sphere at each lift of the centre and
no other homology; the comparison is an isomorphism on $\pi_1$ and on
the homology of those covers, and hence a homotopy equivalence. We may
choose the $B$ summand to be a regular $B_t$ and the sphere to be $R$,
joining their base points by a path.

\noindent\textbf{Step 2.} The circle Chern class is zero on $B_t$ and $-1$ on $R$. Therefore
$S(L)$ is represented by the homotopy pushout
\[
 T^3_{(a,b,d)}\longleftarrow S^1_a\longrightarrow S^3.
\]
The second map is a Hopf fibre and is nullhomotopic. Thus this pushout
has the homotopy type $(T^3\cup_aD^2)\vee S^3$. We multiply by $S^1_c$.
The resulting integral cellular complex starts with the exterior-cell
complex of $T^4$, ordered by $(a,c,b,d)$, and adds cells $r$ in degree
two, $rc,h$ in degree three, and $hc$ in degree four. Its only nonzero
differentials are
\[
 \partial r=a,\qquad \partial(rc)=a\wedge c.
\]
We choose the orientation of $r$ to give the first equality and use the
product orientation for $rc$. The complete differential matrices are
printed below. The complex has free homology of ranks $(1,3,5,5,2)$;
the inclusion of the torus subcomplex gives exactly the two quotient maps.

\noindent\textbf{Step 3.} We represent a base point, and both base generators $b,d$, away from
$0\in B$. Their products with the disc miss $R$, so the canonical section
trivializes $L$ over the resulting chains. The meridian and the products
with $b,d$ consequently bound. The products with $a,c$ bound over one
such base point. This proves all the remaining assertions on the maps.
\end{proof}

The lattice matrices, with columns $(a_i,c_i,b_i,d_i)$ in the common basis,
are
\[
S_1=\begin{pmatrix}0&0&1&0\\0&0&0&1\\1&0&0&0\\0&1&0&0\end{pmatrix},\quad
S_2=\begin{pmatrix}0&-1&0&0\\1&0&0&0\\0&-1&1&0\\1&0&0&1\end{pmatrix},\quad
S_3=\begin{pmatrix}1&0&0&0\\0&1&0&0\\0&-1&1&0\\1&0&0&1\end{pmatrix}.
\]
Put $Q_i=S_i^{-1}$. Deleting the first row of $Q_i$ gives $R_{1,i}$.
The row indices of $\bigwedge^2Q_i$ are $(ac,ab,ad,cb,cd,bd)$ and its
column indices are $(12,13,14,23,24,34)$; deleting its first row gives
$R_{2,i}$. Thus these are the integer matrices of the local quotient maps,
with target bases $(c_i,b_i,d_i)$ and $(a_ib_i,a_id_i,c_ib_i,c_id_i,b_id_i)$.

\begin{proposition}[P4.4]\label{night:P4.4}\label{stmt:P4.4}
\StatementRecord{night:P4.4}{proved}{Historical events: GPT-6 Astra (OpenAI), ChatGPT Pro mode, effort Pro, S6\_06\_BLIND, 16--18 November 2026: targeted proof check. Also checked in this writing. Pre-continuation writing events: GPT-6 Astra (OpenAI), ChatGPT Pro mode, effort Pro, S6\_10\_PAPERS, 18 November 2026: targeted proof check. GPT-6 Astra (OpenAI), ChatGPT Pro mode, effort Pro, S6\_10\_PAPERS, 18 November 2026 05:30:45.837959+00:00: computation or formal certificate. Continuation metadata events: GPT-6 Astra (OpenAI), ChatGPT Pro mode, effort Pro, S6\_10\_PAPERS continuation, 19 November 2026: used and rechecked in this run; status-vocabulary reconciliation; metadata reconciliation only; no new mathematical check. \emph{Pre-continuation recorded proof scope.} The night-run re-derivation was by GPT-6 Astra (OpenAI, ChatGPT Pro mode, effort Pro), 16--18 November 2026. Also checked in this writing by GPT-6 Astra (OpenAI), ChatGPT Pro mode, effort Pro, 18 November 2026. Reconstructed F1 and F2 directly from the three holomorphic marking matrices and logarithmic clutch, including every e\_j wedge g\_i mixed term and the two base meridian rows. The current independent replay gives determinant F1=-1, fifteen unit Smith factors for F1, twenty-nine for F2, and reduced relation determinants -1 and 1. The Mayer-Vietoris exact sequence and explicit van Kampen relations then give integral H1=H2=0 and trivial pi1. The computation scope is recorded separately in analysis/sign\_P44\_integral\_replay.json}{the hypotheses and explicit local calculation stated here}

The threefold $Y$ has $H_1(Y;\Z)=H_2(Y;\Z)=0$ and $\pi_1(Y)=1$.
\end{proposition}
\begin{proof}
\noindent\textbf{Step 1.} Let $P$ be the complement of three small open base discs. The remaining
bundle $W=f^{-1}(P)$ is trivial. We use $W=A\times P$, let $s_1,s_2$
generate $H_1(P)$, and set $s_3=-s_1-s_2$. The overlaps are
$N_i=A\times S^1$. Positive continuation of $\ell$ around zero adds $e_1$.
With
\[
 g_1=e_1,\quad g_2=g_3=0,\qquad
 v_1=\binom10,\quad v_2=\binom01,\quad v_3=\binom{-1}{-1},
\]
the overlap maps into $W$ are
\[
 x\mapsto x,\quad s_i\mapsto g_i+v_i,\qquad
 x\wedge y\mapsto x\wedge y,\quad
 x\wedge s_i\mapsto x\wedge g_i+x\otimes v_i.
\]
These are the maps represented by the products of the cycles, or by the
induced maps into the aspherical space $A\times P$; in particular they
retain the degree-two contribution of the clutch.

\noindent\textbf{Step 2.} We use the ordered basis $(e_1,e_2,e_3,e_4,s_1,s_2)$ for $H_1(W)$ and
append the three bases $(c_i,b_i,d_i)$ for the local groups. On the source
we concatenate the three bases $(e_1,e_2,e_3,e_4,s_i)$. The resulting
$15\times15$ Mayer--Vietoris matrix $F_1$ has, in its $i$th five-column
block, upper block
\[
 \begin{pmatrix}I_4&g_i\\0_{2\times4}&v_i\end{pmatrix}
\]
and local block $(-R_{1,i},0)$ in the $i$th summand; its other local
blocks vanish. The sign is the usual difference of the two inclusions.

\noindent\textbf{Step 3.} For $H_2(W)$ we use
\[
 (12,13,14,23,24,34,e_1s_1,\ldots,e_4s_1,e_1s_2,\ldots,e_4s_2)
\]
and append the three local five-element bases. The source is the
concatenation of $(12,13,14,23,24,34,e_1s_i,\ldots,e_4s_i)$.
The $29\times30$ matrix $F_2$ has top six rows
$[I_6,(e_j\wedge g_i)_{j=1}^4]$ in its $i$th block; its next eight rows
have the coefficients of $e_j\otimes v_i$ in the displayed order, and
its $i$th local block is $[-R_{2,i},0]$. All other blocks vanish.
Finally
\[
 F_0=\begin{pmatrix}1&1&1\\-1&0&0\\0&-1&0\\0&0&-1\end{pmatrix}
\]
has zero kernel. Thus the required exact sequence is
\[
 \Z^{30}\xrightarrow{F_2}\Z^{29}\longrightarrow H_2(Y)
 \longrightarrow\Z^{15}\xrightarrow{F_1}\Z^{15}\longrightarrow H_1(Y)
 \longrightarrow\Z^3\xrightarrow{F_0}\Z^4.
\]

\noindent\textbf{Step 4.} We verify integral surjectivity by hand. Eliminating the local $H_1$
groups and two independent base meridians leaves $\Lambda$ modulo
$e_1,a_1,a_2,a_3$. The relation matrix is
\[
 R_{H_1}=\begin{pmatrix}1&0&0&1\\0&0&1&0\\0&1&0&0\\0&0&1&1\end{pmatrix},
 \qquad\det R_{H_1}=-1.
\]
Hence $F_1$ is an integral isomorphism. In degree two the fibre
$\bigwedge^2\Lambda$ columns eliminate the fifteen local generators,
leaving the three relations $\eta_i=a_i\wedge c_i$. The second meridian
columns kill $\Lambda\otimes s_2$, the third then kill
$\Lambda\otimes s_1$, and the first leave $e_1\wedge\Lambda$. The six
relations, in the order $e_{12},e_{13},e_{14},\eta_1,\eta_2,\eta_3$, have
matrix
\[
R_{H_2}=\begin{pmatrix}
1&0&0&0&1&1\\0&1&0&0&0&-1\\0&0&1&0&1&0\\
0&0&0&0&-1&0\\0&0&0&0&0&-1\\0&0&0&1&1&1
\end{pmatrix},\qquad \det R_{H_2}=1.
\]
This proves surjectivity of $F_2$ over $\Z$ and excludes torsion. The
executed independent Smith calculation gives fifteen unit factors for
$F_1$ and twenty-nine unit factors for $F_2$, and computes $\det F_1=-1$.

\noindent\textbf{Step 5.} For the fundamental group the local pushout gives
$\pi_1(U_i)=\Lambda/\Z a_i$, with trivial meridian. Van Kampen starts
from $\pi_1(W)=\Lambda\times F(s_1,s_2)$ and imposes
\[
 a_1=a_2=a_3=0,\qquad s_1+e_1=0,\quad s_2=0,\quad s_1+s_2=0.
\]
Both base generators and $e_1$ vanish. The same unimodular matrix
$R_{H_1}$ then kills the fibre lattice, proving $\pi_1(Y)=1$.
\end{proof}

The program \texttt{code/part4/example\_mayer\_vietoris.py} recomputes all
entries, differentials, determinants and Smith forms exactly. Its output
and machine-readable matrices are supplied alongside it. For completeness
we print every integer matrix, including the zero differentials:
% Generated by code/part4/matrices_to_tex.py from example_matrices.json.
\begingroup
\setlength{\arraycolsep}{1.7pt}
\renewcommand{\arraystretch}{0.88}
\par\smallskip\noindent
$S_1$ has size $4\times 4$; its full integer matrix is
\[
\resizebox{!}{11.2mm}{$\displaystyle S_1=\begin{pmatrix}
0 & 0 & 1 & 0\\
0 & 0 & 0 & 1\\
1 & 0 & 0 & 0\\
0 & 1 & 0 & 0\\
\end{pmatrix}$}
\]
\par\smallskip\noindent
$S_2$ has size $4\times 4$; its full integer matrix is
\[
\resizebox{!}{11.2mm}{$\displaystyle S_2=\begin{pmatrix}
0 & -1 & 0 & 0\\
1 & 0 & 0 & 0\\
0 & -1 & 1 & 0\\
1 & 0 & 0 & 1\\
\end{pmatrix}$}
\]
\par\smallskip\noindent
$S_3$ has size $4\times 4$; its full integer matrix is
\[
\resizebox{!}{11.2mm}{$\displaystyle S_3=\begin{pmatrix}
1 & 0 & 0 & 0\\
0 & 1 & 0 & 0\\
0 & -1 & 1 & 0\\
1 & 0 & 0 & 1\\
\end{pmatrix}$}
\]
\par\smallskip\noindent
$R_{1,1}$ has size $3\times 4$; its full integer matrix is
\[
\resizebox{!}{8.399999999999999mm}{$\displaystyle R_{1,1}=\begin{pmatrix}
0 & 0 & 0 & 1\\
1 & 0 & 0 & 0\\
0 & 1 & 0 & 0\\
\end{pmatrix}$}
\]
\par\smallskip\noindent
$R_{1,2}$ has size $3\times 4$; its full integer matrix is
\[
\resizebox{!}{8.399999999999999mm}{$\displaystyle R_{1,2}=\begin{pmatrix}
-1 & 0 & 0 & 0\\
-1 & 0 & 1 & 0\\
0 & -1 & 0 & 1\\
\end{pmatrix}$}
\]
\par\smallskip\noindent
$R_{1,3}$ has size $3\times 4$; its full integer matrix is
\[
\resizebox{!}{8.399999999999999mm}{$\displaystyle R_{1,3}=\begin{pmatrix}
0 & 1 & 0 & 0\\
0 & 1 & 1 & 0\\
-1 & 0 & 0 & 1\\
\end{pmatrix}$}
\]
\par\smallskip\noindent
$R_{2,1}$ has size $5\times 6$; its full integer matrix is
\[
\resizebox{!}{14.0mm}{$\displaystyle R_{2,1}=\begin{pmatrix}
0 & -1 & 0 & 0 & 0 & 0\\
0 & 0 & 0 & -1 & 0 & 0\\
0 & 0 & -1 & 0 & 0 & 0\\
0 & 0 & 0 & 0 & -1 & 0\\
1 & 0 & 0 & 0 & 0 & 0\\
\end{pmatrix}$}
\]
\par\smallskip\noindent
$R_{2,2}$ has size $5\times 6$; its full integer matrix is
\[
\resizebox{!}{14.0mm}{$\displaystyle R_{2,2}=\begin{pmatrix}
1 & 0 & 0 & 1 & 0 & 0\\
0 & 0 & 0 & 0 & 1 & 0\\
0 & -1 & 0 & 0 & 0 & 0\\
1 & 0 & -1 & 0 & 0 & 0\\
1 & 0 & -1 & 1 & 0 & 1\\
\end{pmatrix}$}
\]
\par\smallskip\noindent
$R_{2,3}$ has size $5\times 6$; its full integer matrix is
\[
\resizebox{!}{14.0mm}{$\displaystyle R_{2,3}=\begin{pmatrix}
1 & 1 & 0 & 0 & 0 & 0\\
0 & 0 & 1 & 0 & 0 & 0\\
0 & 0 & 0 & 1 & 0 & 0\\
1 & 0 & 0 & 0 & 1 & 0\\
1 & 1 & 0 & 0 & 1 & 1\\
\end{pmatrix}$}
\]
\par\smallskip\noindent
$\partial_1$ has size $1\times 4$; its full integer matrix is
\[
\resizebox{!}{2.8mm}{$\displaystyle \partial_1=\begin{pmatrix}
0 & 0 & 0 & 0\\
\end{pmatrix}$}
\]
\par\smallskip\noindent
$\partial_2$ has size $4\times 7$; its full integer matrix is
\[
\resizebox{!}{11.2mm}{$\displaystyle \partial_2=\begin{pmatrix}
0 & 0 & 0 & 0 & 0 & 0 & 1\\
0 & 0 & 0 & 0 & 0 & 0 & 0\\
0 & 0 & 0 & 0 & 0 & 0 & 0\\
0 & 0 & 0 & 0 & 0 & 0 & 0\\
\end{pmatrix}$}
\]
\par\smallskip\noindent
$\partial_3$ has size $7\times 6$; its full integer matrix is
\[
\resizebox{!}{19.599999999999998mm}{$\displaystyle \partial_3=\begin{pmatrix}
0 & 0 & 0 & 0 & 1 & 0\\
0 & 0 & 0 & 0 & 0 & 0\\
0 & 0 & 0 & 0 & 0 & 0\\
0 & 0 & 0 & 0 & 0 & 0\\
0 & 0 & 0 & 0 & 0 & 0\\
0 & 0 & 0 & 0 & 0 & 0\\
0 & 0 & 0 & 0 & 0 & 0\\
\end{pmatrix}$}
\]
\par\smallskip\noindent
$\partial_4$ has size $6\times 2$; its full integer matrix is
\[
\resizebox{!}{16.799999999999997mm}{$\displaystyle \partial_4=\begin{pmatrix}
0 & 0\\
0 & 0\\
0 & 0\\
0 & 0\\
0 & 0\\
0 & 0\\
\end{pmatrix}$}
\]
\par\smallskip\noindent
$F_0$ has size $4\times 3$; its full integer matrix is
\[
\resizebox{!}{11.2mm}{$\displaystyle F_0=\begin{pmatrix}
1 & 1 & 1\\
-1 & 0 & 0\\
0 & -1 & 0\\
0 & 0 & -1\\
\end{pmatrix}$}
\]
\par\smallskip\noindent
$F_1$ has size $15\times 15$; its full integer matrix is
\[
\resizebox{0.94\linewidth}{!}{$\displaystyle F_1=\begin{pmatrix}
1 & 0 & 0 & 0 & 1 & 1 & 0 & 0 & 0 & 0 & 1 & 0 & 0 & 0 & 0\\
0 & 1 & 0 & 0 & 0 & 0 & 1 & 0 & 0 & 0 & 0 & 1 & 0 & 0 & 0\\
0 & 0 & 1 & 0 & 0 & 0 & 0 & 1 & 0 & 0 & 0 & 0 & 1 & 0 & 0\\
0 & 0 & 0 & 1 & 0 & 0 & 0 & 0 & 1 & 0 & 0 & 0 & 0 & 1 & 0\\
0 & 0 & 0 & 0 & 1 & 0 & 0 & 0 & 0 & 0 & 0 & 0 & 0 & 0 & -1\\
0 & 0 & 0 & 0 & 0 & 0 & 0 & 0 & 0 & 1 & 0 & 0 & 0 & 0 & -1\\
0 & 0 & 0 & -1 & 0 & 0 & 0 & 0 & 0 & 0 & 0 & 0 & 0 & 0 & 0\\
-1 & 0 & 0 & 0 & 0 & 0 & 0 & 0 & 0 & 0 & 0 & 0 & 0 & 0 & 0\\
0 & -1 & 0 & 0 & 0 & 0 & 0 & 0 & 0 & 0 & 0 & 0 & 0 & 0 & 0\\
0 & 0 & 0 & 0 & 0 & 1 & 0 & 0 & 0 & 0 & 0 & 0 & 0 & 0 & 0\\
0 & 0 & 0 & 0 & 0 & 1 & 0 & -1 & 0 & 0 & 0 & 0 & 0 & 0 & 0\\
0 & 0 & 0 & 0 & 0 & 0 & 1 & 0 & -1 & 0 & 0 & 0 & 0 & 0 & 0\\
0 & 0 & 0 & 0 & 0 & 0 & 0 & 0 & 0 & 0 & 0 & -1 & 0 & 0 & 0\\
0 & 0 & 0 & 0 & 0 & 0 & 0 & 0 & 0 & 0 & 0 & -1 & -1 & 0 & 0\\
0 & 0 & 0 & 0 & 0 & 0 & 0 & 0 & 0 & 0 & 1 & 0 & 0 & -1 & 0\\
\end{pmatrix}$}
\]
\par\smallskip\noindent
$F_2$ has size $29\times 30$; its full integer matrix is
\[
\resizebox{0.94\linewidth}{!}{$\displaystyle F_2=\begin{pmatrix}
1 & 0 & 0 & 0 & 0 & 0 & 0 & -1 & 0 & 0 & 1 & 0 & 0 & 0 & 0 & 0 & 0 & 0 & 0 & 0 & 1 & 0 & 0 & 0 & 0 & 0 & 0 & 0 & 0 & 0\\
0 & 1 & 0 & 0 & 0 & 0 & 0 & 0 & -1 & 0 & 0 & 1 & 0 & 0 & 0 & 0 & 0 & 0 & 0 & 0 & 0 & 1 & 0 & 0 & 0 & 0 & 0 & 0 & 0 & 0\\
0 & 0 & 1 & 0 & 0 & 0 & 0 & 0 & 0 & -1 & 0 & 0 & 1 & 0 & 0 & 0 & 0 & 0 & 0 & 0 & 0 & 0 & 1 & 0 & 0 & 0 & 0 & 0 & 0 & 0\\
0 & 0 & 0 & 1 & 0 & 0 & 0 & 0 & 0 & 0 & 0 & 0 & 0 & 1 & 0 & 0 & 0 & 0 & 0 & 0 & 0 & 0 & 0 & 1 & 0 & 0 & 0 & 0 & 0 & 0\\
0 & 0 & 0 & 0 & 1 & 0 & 0 & 0 & 0 & 0 & 0 & 0 & 0 & 0 & 1 & 0 & 0 & 0 & 0 & 0 & 0 & 0 & 0 & 0 & 1 & 0 & 0 & 0 & 0 & 0\\
0 & 0 & 0 & 0 & 0 & 1 & 0 & 0 & 0 & 0 & 0 & 0 & 0 & 0 & 0 & 1 & 0 & 0 & 0 & 0 & 0 & 0 & 0 & 0 & 0 & 1 & 0 & 0 & 0 & 0\\
0 & 0 & 0 & 0 & 0 & 0 & 1 & 0 & 0 & 0 & 0 & 0 & 0 & 0 & 0 & 0 & 0 & 0 & 0 & 0 & 0 & 0 & 0 & 0 & 0 & 0 & -1 & 0 & 0 & 0\\
0 & 0 & 0 & 0 & 0 & 0 & 0 & 1 & 0 & 0 & 0 & 0 & 0 & 0 & 0 & 0 & 0 & 0 & 0 & 0 & 0 & 0 & 0 & 0 & 0 & 0 & 0 & -1 & 0 & 0\\
0 & 0 & 0 & 0 & 0 & 0 & 0 & 0 & 1 & 0 & 0 & 0 & 0 & 0 & 0 & 0 & 0 & 0 & 0 & 0 & 0 & 0 & 0 & 0 & 0 & 0 & 0 & 0 & -1 & 0\\
0 & 0 & 0 & 0 & 0 & 0 & 0 & 0 & 0 & 1 & 0 & 0 & 0 & 0 & 0 & 0 & 0 & 0 & 0 & 0 & 0 & 0 & 0 & 0 & 0 & 0 & 0 & 0 & 0 & -1\\
0 & 0 & 0 & 0 & 0 & 0 & 0 & 0 & 0 & 0 & 0 & 0 & 0 & 0 & 0 & 0 & 1 & 0 & 0 & 0 & 0 & 0 & 0 & 0 & 0 & 0 & -1 & 0 & 0 & 0\\
0 & 0 & 0 & 0 & 0 & 0 & 0 & 0 & 0 & 0 & 0 & 0 & 0 & 0 & 0 & 0 & 0 & 1 & 0 & 0 & 0 & 0 & 0 & 0 & 0 & 0 & 0 & -1 & 0 & 0\\
0 & 0 & 0 & 0 & 0 & 0 & 0 & 0 & 0 & 0 & 0 & 0 & 0 & 0 & 0 & 0 & 0 & 0 & 1 & 0 & 0 & 0 & 0 & 0 & 0 & 0 & 0 & 0 & -1 & 0\\
0 & 0 & 0 & 0 & 0 & 0 & 0 & 0 & 0 & 0 & 0 & 0 & 0 & 0 & 0 & 0 & 0 & 0 & 0 & 1 & 0 & 0 & 0 & 0 & 0 & 0 & 0 & 0 & 0 & -1\\
0 & 1 & 0 & 0 & 0 & 0 & 0 & 0 & 0 & 0 & 0 & 0 & 0 & 0 & 0 & 0 & 0 & 0 & 0 & 0 & 0 & 0 & 0 & 0 & 0 & 0 & 0 & 0 & 0 & 0\\
0 & 0 & 0 & 1 & 0 & 0 & 0 & 0 & 0 & 0 & 0 & 0 & 0 & 0 & 0 & 0 & 0 & 0 & 0 & 0 & 0 & 0 & 0 & 0 & 0 & 0 & 0 & 0 & 0 & 0\\
0 & 0 & 1 & 0 & 0 & 0 & 0 & 0 & 0 & 0 & 0 & 0 & 0 & 0 & 0 & 0 & 0 & 0 & 0 & 0 & 0 & 0 & 0 & 0 & 0 & 0 & 0 & 0 & 0 & 0\\
0 & 0 & 0 & 0 & 1 & 0 & 0 & 0 & 0 & 0 & 0 & 0 & 0 & 0 & 0 & 0 & 0 & 0 & 0 & 0 & 0 & 0 & 0 & 0 & 0 & 0 & 0 & 0 & 0 & 0\\
-1 & 0 & 0 & 0 & 0 & 0 & 0 & 0 & 0 & 0 & 0 & 0 & 0 & 0 & 0 & 0 & 0 & 0 & 0 & 0 & 0 & 0 & 0 & 0 & 0 & 0 & 0 & 0 & 0 & 0\\
0 & 0 & 0 & 0 & 0 & 0 & 0 & 0 & 0 & 0 & -1 & 0 & 0 & -1 & 0 & 0 & 0 & 0 & 0 & 0 & 0 & 0 & 0 & 0 & 0 & 0 & 0 & 0 & 0 & 0\\
0 & 0 & 0 & 0 & 0 & 0 & 0 & 0 & 0 & 0 & 0 & 0 & 0 & 0 & -1 & 0 & 0 & 0 & 0 & 0 & 0 & 0 & 0 & 0 & 0 & 0 & 0 & 0 & 0 & 0\\
0 & 0 & 0 & 0 & 0 & 0 & 0 & 0 & 0 & 0 & 0 & 1 & 0 & 0 & 0 & 0 & 0 & 0 & 0 & 0 & 0 & 0 & 0 & 0 & 0 & 0 & 0 & 0 & 0 & 0\\
0 & 0 & 0 & 0 & 0 & 0 & 0 & 0 & 0 & 0 & -1 & 0 & 1 & 0 & 0 & 0 & 0 & 0 & 0 & 0 & 0 & 0 & 0 & 0 & 0 & 0 & 0 & 0 & 0 & 0\\
0 & 0 & 0 & 0 & 0 & 0 & 0 & 0 & 0 & 0 & -1 & 0 & 1 & -1 & 0 & -1 & 0 & 0 & 0 & 0 & 0 & 0 & 0 & 0 & 0 & 0 & 0 & 0 & 0 & 0\\
0 & 0 & 0 & 0 & 0 & 0 & 0 & 0 & 0 & 0 & 0 & 0 & 0 & 0 & 0 & 0 & 0 & 0 & 0 & 0 & -1 & -1 & 0 & 0 & 0 & 0 & 0 & 0 & 0 & 0\\
0 & 0 & 0 & 0 & 0 & 0 & 0 & 0 & 0 & 0 & 0 & 0 & 0 & 0 & 0 & 0 & 0 & 0 & 0 & 0 & 0 & 0 & -1 & 0 & 0 & 0 & 0 & 0 & 0 & 0\\
0 & 0 & 0 & 0 & 0 & 0 & 0 & 0 & 0 & 0 & 0 & 0 & 0 & 0 & 0 & 0 & 0 & 0 & 0 & 0 & 0 & 0 & 0 & -1 & 0 & 0 & 0 & 0 & 0 & 0\\
0 & 0 & 0 & 0 & 0 & 0 & 0 & 0 & 0 & 0 & 0 & 0 & 0 & 0 & 0 & 0 & 0 & 0 & 0 & 0 & -1 & 0 & 0 & 0 & -1 & 0 & 0 & 0 & 0 & 0\\
0 & 0 & 0 & 0 & 0 & 0 & 0 & 0 & 0 & 0 & 0 & 0 & 0 & 0 & 0 & 0 & 0 & 0 & 0 & 0 & -1 & -1 & 0 & 0 & -1 & -1 & 0 & 0 & 0 & 0\\
\end{pmatrix}$}
\]
\par\smallskip\noindent
$R_{H_1}$ has size $4\times 4$; its full integer matrix is
\[
\resizebox{!}{11.2mm}{$\displaystyle R_{H_1}=\begin{pmatrix}
1 & 0 & 0 & 1\\
0 & 0 & 1 & 0\\
0 & 1 & 0 & 0\\
0 & 0 & 1 & 1\\
\end{pmatrix}$}
\]
\par\smallskip\noindent
$R_{H_2}$ has size $6\times 6$; its full integer matrix is
\[
\resizebox{!}{16.799999999999997mm}{$\displaystyle R_{H_2}=\begin{pmatrix}
1 & 0 & 0 & 0 & 1 & 1\\
0 & 1 & 0 & 0 & 0 & -1\\
0 & 0 & 1 & 0 & 1 & 0\\
0 & 0 & 0 & 0 & -1 & 0\\
0 & 0 & 0 & 0 & 0 & -1\\
0 & 0 & 0 & 1 & 1 & 1\\
\end{pmatrix}$}
\]
\endgroup


\begin{proposition}[P4.5]\label{night:P4.5}\label{stmt:P4.5}
\StatementRecord{night:P4.5}{proved}{Historical events: GPT-6 Astra (OpenAI), ChatGPT Pro mode, effort Pro, S6\_06\_BLIND, 16--18 November 2026: targeted proof check. Also checked in this writing. Pre-continuation writing events: GPT-6 Astra (OpenAI), ChatGPT Pro mode, effort Pro, S6\_10\_PAPERS, 18 November 2026: targeted proof check. Continuation metadata events: GPT-6 Astra (OpenAI), ChatGPT Pro mode, effort Pro, S6\_10\_PAPERS continuation, 19 November 2026: used and rechecked in this run; status-vocabulary reconciliation; metadata reconciliation only; no new mathematical check. \emph{Pre-continuation recorded proof scope.} The night-run re-derivation was by GPT-6 Astra (OpenAI, ChatGPT Pro mode, effort Pro), 16--18 November 2026. Also checked in this writing by GPT-6 Astra (OpenAI), ChatGPT Pro mode, effort Pro, 18 November 2026. Stratified Euler additivity gives zero for the actual smooth torus fibres and Kodaira-Hopf central supports, and Poincare duality with the computed integral H1 and H2 gives the asserted torsion-free cohomology and b3=2. An effective horizontal polar divisor would restrict to a nonzero effective divisor of zero real Chern class on a Kahler torus. Excluding it, removing vertical poles by local powers of the base coordinate, and connected proper descent identify the meromorphic function field with that of P1}{\ref{night:P4.4}}

For this example,
\[
 a(Y)=1,\qquad c_3[Y]=0,\qquad b_3(Y)=2,
\]
and the integral homology is free of ranks $(1,0,0,2,0,0,1)$.
\end{proposition}
\begin{proof}
\noindent\textbf{Step 1.} The regular fibre has Euler number zero. Both components of each special
fibre are elliptic bundles over compact curves, and their intersection is
elliptic, so additivity gives zero for each special support. Stratifying
the base and applying multiplicativity to the smooth bundle gives
$e(Y)=0$. The Euler class of the real bundle underlying a complex
rank-three tangent bundle is its top Chern class, with the complex
orientation. Hence $c_3[Y]=0$. Duality and P4.4 give $b_4=b_5=0$ and
$e(Y)=2-b_3$, so $b_3=2$. Integral duality and the universal-coefficient
theorem identify the torsion of $H_3$ with that of $H_2$, which is zero.
All other groups have already been determined.

\noindent\textbf{Step 2.} We also have $H^2(Y;\Z)=0$. A horizontal effective divisor would restrict
to a nonzero effective divisor on the Kähler torus $A$ with zero first
Chern class. Its pairing with a Kähler class would be positive by
integration over the divisor and zero by the Chern-class computation.
Thus every divisor on $Y$ is vertical. The polar divisor of a meromorphic
function is consequently vertical; on a general fibre the function is
holomorphic and therefore constant. Local holomorphic sections give a
descended meromorphic function on the regular part of the base. Near
each exceptional base value a sufficiently high power of the base
coordinate removes the vertical poles. Connected fibres then give
holomorphic descent there, and dividing by that power gives meromorphic
descent. Therefore $\C(Y)=f^*\C(\Pp^1)$ and $a(Y)=1$.
\end{proof}

\subsection{The normal-crossings sign theorem}

\begin{proposition}[P4.6]\label{night:P4.6}\label{stmt:P4.6}
\StatementRecord{night:P4.6}{proved}{Historical events: GPT-6 Astra (OpenAI), ChatGPT Pro mode, effort Pro, S6\_06\_BLIND, 16--18 November 2026: targeted proof check. Also checked in this writing. Pre-continuation writing events: GPT-6 Astra (OpenAI), ChatGPT Pro mode, effort Pro, S6\_10\_PAPERS, 18 November 2026: targeted proof check. Continuation metadata events: GPT-6 Astra (OpenAI), ChatGPT Pro mode, effort Pro, S6\_10\_PAPERS continuation, 19 November 2026: used and rechecked in this run; status-vocabulary reconciliation; metadata reconciliation only; no new mathematical check. \emph{Pre-continuation recorded proof scope.} The night-run re-derivation was by GPT-6 Astra (OpenAI, ChatGPT Pro mode, effort Pro), 16--18 November 2026. Also checked in this writing by GPT-6 Astra (OpenAI), ChatGPT Pro mode, effort Pro, 18 November 2026. The branchwise normal derivative gives the actual line isomorphism N\_i=nu\_i*O(F) tensor O(-D\_i), also for self-intersecting global components. Its tangent Chern sequence gives c2=K2 and K numerically minus D. The unramified branch-choice double cover and three conductor nodes at each triple point give e(Gamma)=3T. The alternating normalization complex carries the degree-zero sign line on its double term; coherent HRR for the Cartier sequence gives T=2 sum chi(O\_Si). Topological inclusion-exclusion then gives e(F)=T and parity, without a multiplicity assumption}{the hypotheses and explicit local calculation stated here}

Let $X$ be a compact complex threefold with $b_2(X)=0$, and let $F$ be a
reduced compact normal-crossings divisor. We allow different local
branches of one global irreducible component to meet. If $T(F)$ is the
number of ordinary triple points, then
\[
 e(F)=T(F),\qquad T(F)\text{ is even}.
\]
This assertion concerns the reduced support, regardless of the
multiplicities with which it occurs in a fibre.
\end{proposition}
\begin{proof}
\noindent\textbf{Step 1.} Write $\nu:S=\coprod S_i\to X$ for the normalization immersion, and $D_i$
for the conductor divisor on $S_i$. The surfaces $S_i$ are smooth and
the curves $D_i$ are reduced nodal divisors. The quotient
$N_i=\nu_i^*T_X/T_{S_i}$ is a line bundle. A local equation of $F$ gives
\[
 N_i\simeq \nu_i^*\OO_X(F)\otimes\OO_{S_i}(-D_i).
\]
Indeed, for the coordinate equation $x_1\cdots x_k=0$ the normal
derivative on $x_1=0$ is $x_2\cdots x_k$, whose zero divisor is $D_i$.
These are morphisms of holomorphic vector bundles, including when
$S_i$ is not embedded globally.

\noindent\textbf{Step 2.} Since $b_2(X)=0$, both $c_1(T_X)$ and $c_1(\OO_X(F))$ vanish in real
cohomology. Duality gives $b_4(X)=0$, so $c_2(T_X)$ also vanishes. The
tangent exact sequence yields
\[
 K_{S_i}\equiv-D_i,\qquad c_2(S_i)=K_{S_i}^2.
\]
If $a=\sum_i\chi(\OO_{S_i})$, the surface Riemann--Roch formula
$12\chi(\OO_{S_i})=K_{S_i}^2+e(S_i)$ consequently gives
$\sum_i e(S_i)=6a$.

\noindent\textbf{Step 3.} Let $\Gamma$ be the normalization of the double-curve locus, with one
sheet for each unordered pair of local branches. The normalization of
$\coprod D_i$ is an unramified degree-two cover of $\Gamma$, recording
the choice of branch. Each triple point contributes three nodes to
$\coprod D_i$. Adjunction gives
$\chi(\OO_{D_i})=-D_i\cdot(D_i+K_{S_i})/2=0$. The normalization
sequence for these nodal curves therefore gives
\[
 0=2\chi(\OO_\Gamma)-3T(F)=e(\Gamma)-3T(F).
\]
The normalization complex of $F$ is the exact alternating restriction
complex of the local coordinate hyperplanes. Permuting the branches
globally replaces its double-curve term by $\OO_\Gamma$ tensored with
the sign line of the two-sheeted branch cover. That line has order at
most two, and hence degree zero on every component; it has the same
coherent Euler characteristic as $\OO_\Gamma$. The triple term has
length one at each triple point. It follows that
\[
 \chi(\OO_F)=a-\chi(\OO_\Gamma)+T(F)=a-T(F)/2.
\]
The exact sequence for the Cartier divisor $F$ and analytic
Hirzebruch--Riemann--Roch give $\chi(\OO_F)=0$, since
$\operatorname{ch}(\OO_X(-F))=\operatorname{ch}(\OO_X)$ in rational
cohomology. Thus $T(F)=2a$. Finally ordinary Euler additivity gives
\[
 e(F)=\sum_i e(S_i)-e(\Gamma)+T(F)=6a-2T(F)=T(F).
\]
The alternating count is also verified directly at a point with one,
two or three branches. This proves the identity and parity, without
using fibre multiplicities or a base change.
\end{proof}

\begin{proposition}[P4.7]\label{night:P4.7}\label{stmt:P4.7}
\StatementRecord{night:P4.7}{proved}{Historical events: GPT-6 Astra (OpenAI), ChatGPT Pro mode, effort Pro, S6\_06\_BLIND, 16--18 November 2026: targeted proof check. Also checked in this writing. Pre-continuation writing events: GPT-6 Astra (OpenAI), ChatGPT Pro mode, effort Pro, S6\_10\_PAPERS, 18 November 2026: targeted proof check. Continuation metadata events: GPT-6 Astra (OpenAI), ChatGPT Pro mode, effort Pro, S6\_10\_PAPERS continuation, 19 November 2026: used and rechecked in this run; status-vocabulary reconciliation; metadata reconciliation only; no new mathematical check. \emph{Pre-continuation recorded proof scope.} The night-run re-derivation was by GPT-6 Astra (OpenAI, ChatGPT Pro mode, effort Pro), 16--18 November 2026. Also checked in this writing by GPT-6 Astra (OpenAI), ChatGPT Pro mode, effort Pro, 18 November 2026. The connected-fibre constant-sheaf edge injects H1 of the base. Smooth-fibre adjunction gives c2=0 from the ambient H4 vanishing. Properness makes the critical-value set finite, and Ehresmann plus Euler additivity gives the exact sum over supports. Poincare duality and the alternating middle pairing force the nonnegative sum to lie in \{0,2\}. No torus hypothesis is used in this numerical argument}{\ref{night:P4.6}}

Let $f:X\to B$ be a proper surjective holomorphic map with connected
fibres from a compact complex threefold satisfying $b_1(X)=b_2(X)=0$
to a smooth compact curve. Suppose all reduced fibre supports have
normal crossings. Then
\[
 B=\Pp^1,\qquad
 c_3[X]=\sum_{b\in B}T((f^{-1}(b))_{\mathrm{red}})\in\{0,2\},\qquad
 b_3(X)=2-c_3[X].
\]
The numerical conclusion does not require the general fibre to be a
torus.
\end{proposition}
\begin{proof}
\noindent\textbf{Step 1.} Connected fibres give $f_*\R=\R$. The beginning of the constant-sheaf
Leray sequence injects $H^1(B;\R)$ into $H^1(X;\R)$, so $B$ has genus
zero. A smooth fibre $G$ has trivial normal line bundle. Its tangent
sequence gives $c_2(G)=c_2(X)|_G=0$, hence $e(G)=0$ for any such fibre.
The critical values are finite. Ehresmann's theorem gives a smooth
bundle over their complement. Euler additivity and P4.6 give the
displayed sum. Duality gives $e(X)=2-b_3(X)$, and the nondegenerate
alternating cup product on $H^3(X;\R)$ makes $b_3$ even. The
nonnegative sum is therefore zero or two.
\end{proof}

\begin{proposition}[P4.8]\label{night:P4.8}\label{stmt:P4.8}
\StatementRecord{night:P4.8}{proved}{Historical events: GPT-6 Astra (OpenAI), ChatGPT Pro mode, effort Pro, S6\_06\_BLIND, 16--18 November 2026: targeted proof check. Also checked in this writing. Pre-continuation writing events: GPT-6 Astra (OpenAI), ChatGPT Pro mode, effort Pro, S6\_10\_PAPERS, 18 November 2026: targeted proof check; used and rechecked in this run. Continuation metadata events: GPT-6 Astra (OpenAI), ChatGPT Pro mode, effort Pro, S6\_10\_PAPERS continuation, 19 November 2026: used and rechecked in this run; status-vocabulary reconciliation; metadata reconciliation only; no new mathematical check. \emph{Pre-continuation recorded proof scope.} The night-run re-derivation was by GPT-6 Astra (OpenAI, ChatGPT Pro mode, effort Pro), 16--18 November 2026. Also checked in this writing by GPT-6 Astra (OpenAI), ChatGPT Pro mode, effort Pro, 18 November 2026. The equality e=K2 gives k=2(1-g) in the ruled case and minimality with b2=0 in the non-Kahler classVII case. The actual effective conductor excludes Inoue covers and, for irreducible supports, its torsion adjunction kills every positive plurigenus. Pullback of an adjacent component excludes Kahler normalization in a reducible support. The zero triple count makes the conductor curves disjoint elliptic curves with the unramified branch maps. The normalized root cover of a smooth multiple support is a proper submersion and supplies the finite unramified torus cover. BHPV, BBK, OVV, Khaled-Teleman and the earlier Catanese exact source loci are now matched; the printed Remark3.4 pointer is corrected to Remark3.5}{\ref{night:P4.12}; \ref{night:P4.6}; \ref{night:P4.7}}

For a torus fibration in P4.7 whose reduced supports have normal
crossings, an irreducible singular support with
rational normalization has Euler number two; every other irreducible
support has Euler number zero. In the zero case the supports have the
forms stated in the prompt: a smooth support has a finite unramified
torus cover; a singular irreducible support has smooth minimal
elliptic-ruled or Hopf normalization; a reducible support has smooth
minimal Hopf or non-Kähler elliptic normalization components, with
smooth elliptic double curves and finite unramified branch maps.
Thus $c_3[X]=0$ precisely when every fibre belongs to these zero cases.
\end{proposition}
\begin{proof}
\noindent\textbf{Step 1.} We apply \cite[Chapter~VI, Theorem~1.1 and Table~10, pp.~243--244]{BHPV2004E} to the smooth compact
normalization components, using the identity $e(S_i)=K_{S_i}^2$ proved
in P4.6. The relevant non-Kähler classification is stated in
\cite[Theorem~4.2 and proof, pp.~17--21; Theorem~5.1 and proof, pp.~27--29]{OVV2020}: a minimal non-Kähler
surface is elliptic or of class VII, and the reduced fibres of a
non-Kähler elliptic surface are smooth elliptic curves.

\noindent\textbf{Step 2.} For a component with nonempty conductor, a Kähler form $\omega$ would
give $K_{S_i}\cdot\omega=-D_i\cdot\omega<0$. No positive power of
the canonical bundle can then have a nonzero section. The classification
makes a Kähler component rational or ruled. If it is obtained by $k$
blowups of a ruled surface over a genus-$g$ curve, then
\[
 e=4(1-g)+k,\qquad K^2=8(1-g)-k.
\]
Equality gives $k=2(1-g)$, so $g\le1$ and the genus-one case is minimal.
A rational component has $\chi(\OO)=1$, hence $e=6$ here, including
the case with minimal model $\Pp^2$.

\noindent\textbf{Step 3.} A minimal non-Kähler elliptic surface has $e=K^2=0$, by its smooth
elliptic reduced fibres. A minimal class-VII surface has $b_1=1$ and
$p_g=0$, so $\chi(\OO)=0$ and $K^2=-e$, with $e=b_2\ge0$.
Blowups preserve $\chi(\OO)$ and change $(e,K^2)$ by $(1,-1)$.
Consequently $e=K^2$ forces minimality and $e=0$ in all these
non-Kähler cases. The class-VII case has $b_2=0$. Theorem 1 of
Buonerba--Bogomolov--Kurnosov \cite[Theorem~1, p.~1]{BBK2019} supplies a finite unramified Hopf or
Inoue cover. The conductor excludes an Inoue cover: Inoue surfaces
have no compact curves, by Inoue's theorem as stated and credited in
\cite[Theorem~3.1, p.~303]{KhaledTeleman2026Journal}.
The type-III family has the explicit type-II unramified double cover
of \cite[proof of Remark~3.5, p.~305]{KhaledTeleman2026Journal}. A finite unramified quotient of a Hopf surface has universal
cover $\C^2\setminus\{0\}$ and is therefore a Hopf surface.

\noindent\textbf{Step 4.} For an irreducible singular support there is one normalization
component. P4.6 gives $T=2\chi(\OO_{S_1})$, which is two in the
rational case and zero in the other cases. P4.12 proves that the
dualizing sheaf of an irreducible support is torsion. Adjunction on
the normalization therefore makes $K_{S_1}+D_1$ torsion. Since
$D_1\ne0$, every positive plurigenus of $S_1$ vanishes: multiplying
a pluricanonical section by an appropriate positive power of the
conductor section would give a section of a torsion line bundle with
a zero, whose sufficiently high power is a nonzero constant. Thus a
non-Kähler normalization in this irreducible case is of class VII,
and hence is the Hopf case already obtained.

\noindent\textbf{Step 5.} For a reducible connected fibre, each normalization component meets
another global irreducible component. Pulling back the divisor of
that component gives a nonzero effective divisor on $S_i$ with zero
real first Chern class. A Kähler form would pair positively with it,
contradicting this zero class. Hence every component is non-Kähler,
has $\chi(\OO)=0$, and $T=0$. The conductor components are disjoint
and smooth, because all their nodes would come from triple points.
For any conductor component $C$, adjunction and $K+D\equiv0$ give
$2g(C)-2=(K+C)\cdot C=0$, since the other conductor components
are disjoint. Thus the double curves are smooth elliptic, and their
branch maps are the unramified maps already described in P4.6.

\noindent\textbf{Step 6.} Finally let the support $G$ of a multiple fibre $mG$ be smooth. The
normalized cyclic base change $t=u^m$ is locally a disjoint union of
smooth charts: the original equation is $t=z^m$ times a holomorphic
unit, whose root can be extracted locally. These charts show that
the new central fibre maps finitely and unramifiedly to $G$, and
that the new map to the $u$-disc is a proper submersion. Its nearby
fibres are the original tori. Catanese's theorem
\cite[Theorem~2.1]{Catanese2004} states that every
deformation of a complex torus is a complex torus; its definition
uses precisely a proper holomorphic submersion over a connected
disc, with no Kähler hypothesis on the total space. We have verified
all these hypotheses. The new central fibre is therefore a torus;
Ehresmann's theorem also shows that it is connected. This gives
the required finite unramified cover of $G$.
\end{proof}

\begin{proposition}[P4.9]\label{night:P4.9}\label{stmt:P4.9}
\StatementRecord{night:P4.9}{proved}{Historical events: GPT-6 Astra (OpenAI), ChatGPT Pro mode, effort Pro, S6\_06\_BLIND, 16--18 November 2026: targeted proof check. Also checked in this writing. Pre-continuation writing events: GPT-6 Astra (OpenAI), ChatGPT Pro mode, effort Pro, S6\_10\_PAPERS, 18 November 2026: targeted proof check. Continuation metadata events: GPT-6 Astra (OpenAI), ChatGPT Pro mode, effort Pro, S6\_10\_PAPERS continuation, 19 November 2026: used and rechecked in this run; status-vocabulary reconciliation; metadata reconciliation only; no new mathematical check. \emph{Pre-continuation recorded proof scope.} The night-run re-derivation was by GPT-6 Astra (OpenAI, ChatGPT Pro mode, effort Pro), 16--18 November 2026. Also checked in this writing by GPT-6 Astra (OpenAI), ChatGPT Pro mode, effort Pro, 18 November 2026. Followed the actual local irreducible sign, rational and classVII equality statements, zero-Euler normalization maps, and positive-ruled exclusions through P4.26, P4.32, P4.35-37, P4.54-58 and P4R.18-20. The global conclusion retains the hypothesis that all supports are irreducible or NC. This assembly review adds no new individual proof credit to those separately reviewed inputs}{\ref{night:P4.26}; \ref{night:P4.32}; \ref{night:P4.35}; \ref{night:P4.37}; \ref{night:P4.54}; \ref{night:P4.57}; \ref{night:P4.58}; \ref{night:P4R.18}; \ref{night:P4R.20}}

In the requested irreducible and normal-crossing fibre classes,
the complete sign, exact Euler values and normalization equality
classification hold as stated in P4.58. In particular every
irreducible singular support with rationally resolved normalization
has Euler number two, and every other irreducible support has
Euler number zero.
\end{proposition}
\begin{proof}
\noindent\textbf{Step 1.} P4.26 proves the unrestricted irreducible numerical sign;
P4.35--P4.37 give the rational and class-VII exact values.
P4.32 proves the zero-Euler classification. The positive ruled
alternative is excluded by the complete sequence P4.54--P4.57
and P4R.18--P4R.20, synthesized in P4.58.
\end{proof}
We checked CDP 2020, Lemma~4.6, Propositions~4.7,
4.11 and~7.3, deriving the sign and classification independently
from its local conductor identities. P4.14 records the precise
incorrect curve-classification inference in Proposition~7.3.
Earlier proof obligations and unsuccessful reductions are retained
in the working notes; the final classification is P4.58.
The original stated-value row also claimed a sign for arbitrary
reducible supports without normal crossings. Such supports are
outside the two classes explicitly requested in Part~4.2; the
proved global assertion retains the stated irreducible-or-NC
support hypothesis.

\subsection{Explicit checks on the hypotheses}

\begin{proposition}[P4.10]\label{night:P4.10}\label{stmt:P4.10}
\StatementRecord{night:P4.10}{proved}{Historical events: GPT-6 Astra (OpenAI), ChatGPT Pro mode, effort Pro, S6\_06\_BLIND, 16--18 November 2026: targeted proof check. Also checked in this writing. Pre-continuation writing events: GPT-6 Astra (OpenAI), ChatGPT Pro mode, effort Pro, S6\_10\_PAPERS, 18 November 2026: targeted proof check; used and rechecked in this run. Continuation metadata events: GPT-6 Astra (OpenAI), ChatGPT Pro mode, effort Pro, S6\_10\_PAPERS continuation, 19 November 2026: used and rechecked in this run; status-vocabulary reconciliation; metadata reconciliation only; no new mathematical check. \emph{Pre-continuation recorded proof scope.} The night-run re-derivation was by GPT-6 Astra (OpenAI, ChatGPT Pro mode, effort Pro), 16--18 November 2026. Also checked in this writing by GPT-6 Astra (OpenAI), ChatGPT Pro mode, effort Pro, 18 November 2026. The two Weierstrass charts have twelve simple nodal discriminant values, smooth total surface and relative dualizing line O(1), so the displayed direct-image computation makes the surface rational. Blowing up the actual diagonal ideal (x-u,y-v) in xy=uv gives both displayed smooth charts and reduced normal-crossing central equations. No new surface component appears. Each exceptional P1 changes the fibre Euler number from one to two. The extended section and two factor surface maps kill the actual complement generators of pi1. The result has c3=24 and a nonzero Kahler class in H2}{the hypotheses and explicit local calculation stated here}

There is a smooth projective threefold $Z$ with $b_1(Z)=0$ and a
connected torus fibration $Z\to\Pp^1$ whose reduced special fibres
are irreducible normal-crossings surfaces, but $c_3[Z]=24$.
Consequently the hypothesis $b_2=0$ cannot be omitted, even in the
intersection of the two classes of fibres in the question.
\end{proposition}
\begin{proof}
\noindent\textbf{Step 1.} On the two affine charts of $\Pp^1$, with $s=t^{-1}$, we take
\[
 y^2z=x^3-3xz^2+t^6z^3,\qquad
 y_\infty^2z=x_\infty^3-3s^4x_\infty z^2+z^3,
\]
glued by $x_\infty=s^2x$ and $y_\infty=s^3y$. These are the
Weierstrass cubics in $\Pp(\OO\oplus\OO(2)\oplus\OO(3))$ for
the line bundle $\OO_{\Pp^1}(1)$. They define an elliptic surface
$\pi:S\to\Pp^1$ with zero section $[0:1:0]$. The discriminant
is a nonzero constant times $t^{12}-4$, or $1-4s^{12}$ on the
other chart. There are twelve simple zeros and a smooth fibre at
infinity. At a discriminant zero the cubic has a double root
$x=1$ or $x=-1$ and a distinct third root, so the fibre has an
ordinary node. Since $6t^5\ne0$ there, the total surface is smooth.
At the zero section the derivative in $z$ is nonzero. The map near
each node is therefore $xy=t$ after holomorphic coordinate changes
and translation of the base.

\noindent\textbf{Step 2.} The relative dualizing differential $dx/y$ has transitions
$\pi^*\OO(1)$ and generates the dualizing sheaf also at the nodes.
Thus $\omega_S=\pi^*\OO(-1)$. Relative duality gives
$R^1\pi_*\OO_S=\OO(-1)$, and connected fibres give
$\pi_*\OO_S=\OO$. Leray yields $h^1(\OO_S)=0$, while all
positive plurigenera vanish. The classification \cite[Chapter~VI, Theorem~1.1, p.~243]{BHPV2004E} makes $S$
rational; in the ruled case the base genus equals $h^1(\OO_S)=0$.
In particular $\pi_1(S)=1$, since the minimal rational models
$\Pp^2$ and the Hirzebruch surfaces are simply connected and a
point blowup preserves the fundamental group.

\noindent\textbf{Step 3.} The projective fibre product $W=S\times_{\Pp^1}S$ has exactly
twelve singularities, at the pair of nodes in each nodal fibre.
Their equations are $xy=uv$, with base function $xy=uv$. We blow
up the coherent ideal of the diagonal $\Delta_S\subset W$. This
divisor is Cartier away from the twelve points, so the blowup is
the identity there. Put $r=x-u$, $h=y-v$ at a singular point.
The equation becomes $ry+uh=0$ and the diagonal ideal is $(r,h)$.
The blowup has charts
\[
 h=ar,\quad y=-au,\quad x=u+r,\quad f=-aux;
 \qquad
 r=bh,\quad u=-by,\quad x=-bv,\quad f=-byv.
\]
The independent coordinates are $(a,u,x)$ and $(b,y,v)$.
Both charts are smooth, with $b=a^{-1}$ on the overlap, and the
exceptional fibre is $\Pp^1$. Thus the resulting $Z$ is smooth
and projective. The displayed equations make every special fibre
reduced and normal crossings. Its support remains irreducible:
the original product of two irreducible nodal cubics is irreducible,
and the exceptional locus has dimension one, so introduces no
surface component. Its normalization is rational, by the
birational map from $\Pp^1\times\Pp^1$.

\noindent\textbf{Step 4.} A nodal cubic has Euler number $2-1=1$, so its self-product has
Euler number one. The resolution replaces one point by $\Pp^1$;
each resolved special support therefore has Euler number two.
All other fibres are $E_t\times E_t$ with Euler number zero.
Hence $e(Z)=12\cdot2=24=c_3[Z]$.

\noindent\textbf{Step 5.} For completeness we compute $\pi_1(Z)$. Every loop can be moved
off the fibre divisors, so the fundamental group of their complement
surjects onto $\pi_1(Z)$. The complement is a bundle with fibre
$E\times E$ and a section over a twelve-punctured sphere. The
section extends over $\Pp^1$, killing the base generators. The
two maps
\[
 S\longrightarrow Z:\quad p\mapsto(p,O(\pi(p))),\qquad
 p\mapsto(O(\pi(p)),p)
\]
are defined because their images miss every conifold point, and
extend the two factor inclusions of $E$. Since $S$ is simply
connected they kill both fibre fundamental groups. Thus
$\pi_1(Z)=1$ and $b_1(Z)=0$. A projective Kähler class is nonzero
in $H^2(Z;\R)$, because its cube has positive integral. This
identifies the omitted hypothesis: $b_2(Z)>0$.
\end{proof}

\begin{proposition}[P4.11]\label{night:P4.11}\label{stmt:P4.11}
\StatementRecord{night:P4.11}{proved}{Historical events: GPT-6 Astra (OpenAI), ChatGPT Pro mode, effort Pro, S6\_06\_BLIND, 16--18 November 2026: targeted proof check. Also checked in this writing. Pre-continuation writing events: GPT-6 Astra (OpenAI), ChatGPT Pro mode, effort Pro, S6\_10\_PAPERS, 18 November 2026: targeted proof check. Continuation metadata events: GPT-6 Astra (OpenAI), ChatGPT Pro mode, effort Pro, S6\_10\_PAPERS continuation, 19 November 2026: used and rechecked in this run; status-vocabulary reconciliation; metadata reconciliation only; no new mathematical check. \emph{Pre-continuation recorded proof scope.} The night-run re-derivation was by GPT-6 Astra (OpenAI, ChatGPT Pro mode, effort Pro), 16--18 November 2026. Also checked in this writing by GPT-6 Astra (OpenAI), ChatGPT Pro mode, effort Pro, 18 November 2026. Freeness and properness of the exact complex action, its norm-level S3 times S3 model, and the induced elliptic pencil with Hopf fibres were re-derived. The fibre fundamental group is Z, whereas an etale quotient of a two-torus has a finite-index Z4 subgroup. The P4.7 dependency is used only for the final numerical comparison}{\ref{night:P4.7}}

There is a smooth compact complex threefold $C$ with
$b_1(C)=b_2(C)=0$, $e(C)=0$, and a holomorphic submersion
$C\to\Pp^1$ whose fibres are primary Hopf surfaces. No fibre has
a finite unramified torus cover. Thus the torus hypothesis is
necessary for the equality classification as stated in the prompt.
P4.7 proves it unnecessary for the numerical normal-crossings
sign theorem.
\end{proposition}
\begin{proof}
\noindent\textbf{Step 1.} The additive group $\C$ acts on
$V=(\C^2\setminus\{0\})\times(\C^2\setminus\{0\})$ by
\[
 \tau\cdot(z,w)=(e^\tau z,e^{i\tau}w).
\]
For $\tau=a+ib$, the logarithms of the two Euclidean norms change
by $a$ and $-b$. The action is free. It is proper: when both a
point and its translate lie in prescribed compact subsets of $V$,
the two norm ratios bound $a$ and $b$. Each orbit has a unique
representative with both norms one, obtained with
$a=-\log\|z\|$ and $b=\log\|w\|$. This gives an explicit
diffeomorphism $C=V/\C\cong S^3\times S^3$ and proves compactness.
Holomorphic transverse slices, for example $z_j=1$ after choosing
a local logarithm, give quotient charts with holomorphic
transitions. Hence $C$ is a smooth compact complex threefold.
The product cell complex gives free integral homology of ranks
$(1,0,0,2,0,0,1)$ and therefore the stated Betti and Euler numbers.

\noindent\textbf{Step 2.} The invariant map $[z,w]\mapsto[z_1:z_2]$ is a holomorphic
submersion to $\Pp^1$. In an affine base chart write
$z=\lambda(1,\xi)$ and set $\lambda=1$ using the action.
The choices of logarithm differ by $2\pi ik$, changing $w$ by
$e^{-2\pi k}$. Thus we obtain the explicit local fibre bundle
with fibre
$H=(\C^2\setminus\{0\})/\langle w\mapsto e^{-2\pi}w\rangle$.
Its universal cover retracts onto the simply connected $S^3$, so
$\pi_1(H)=\Z$. A connected finite unramified torus cover would
inject $\Z^4$ as a finite-index subgroup of $\Z$. Every subgroup
of $\Z$ is cyclic, whereas $\Z^4$ has rank four, which is
impossible. This proves the failure of the stated smooth-support
classification using the exact covering morphism.
\end{proof}

The executed program \texttt{code/part4/necessity\_examples.py}
checks the degree and squarefreeness of the discriminant, both
resolution substitutions, the Euler count $24$, and the product
cell ranks used above. Its complete output is supplied.

\subsection{Irreducible supports: canonical torsion and the normal case}

\begin{proposition}[P4.12]\label{night:P4.12}\label{stmt:P4.12}
\StatementRecord{night:P4.12}{proved}{Historical events: GPT-6 Astra (OpenAI), ChatGPT Pro mode, effort Pro, S6\_06\_BLIND, 16--18 November 2026: targeted proof check. Also checked in this writing. Pre-continuation writing events: GPT-6 Astra (OpenAI), ChatGPT Pro mode, effort Pro, S6\_10\_PAPERS, 18 November 2026: targeted proof check; used and rechecked in this run. Continuation metadata events: GPT-6 Astra (OpenAI), ChatGPT Pro mode, effort Pro, S6\_10\_PAPERS continuation, 19 November 2026: used and rechecked in this run; status-vocabulary reconciliation; metadata reconciliation only; no new mathematical check. \emph{Pre-continuation recorded proof scope.} The night-run re-derivation was by GPT-6 Astra (OpenAI, ChatGPT Pro mode, effort Pro), 16--18 November 2026. Also checked in this writing by GPT-6 Astra (OpenAI), ChatGPT Pro mode, effort Pro, 18 November 2026. The rank-one direct-image generator has divisor aS with 0<=a<m, so adjunction and O(mS)=O give torsion omega\_S. On the minimal resolution the principal-function intersection equation and the scaled quadratic identity prove negative definiteness; minimal-resolution adjunction gives discrepancies at most zero. Formal-functions comparison gives p=h0(R1 pi\_*O)>=h1(O\_Z) for every exceptional divisor Z. The canonical-trivial, rational, class-VII and positive-genus ruled alternatives each give the stated contradiction unless the normal support is smooth. For a nonnormal support the finite-duality conductor and the same matrix calculation give an effective nonzero pole divisor, hence Kodaira dimension minus infinity. Mumford's original section I p.6 and the CDP20 conductor locus are opened; no unrestricted class-VII classification is invoked}{\ref{night:P4.6}}

In the standing torus-fibration setting of this section, let $f^{-1}(b)=mS$
have irreducible reduced support. Then
\[
 \omega_S^{\otimes m}\simeq\OO_S,\qquad
 \chi(\OO_S)=0,\qquad h^1(\OO_S)\in\{1,2\}.
\]
Every normal such support is smooth. For a singular, necessarily
nonnormal support, a minimal resolved model of its normalization
has Kodaira dimension $-\infty$.
\end{proposition}
\begin{proof}
\noindent\textbf{Step 1. Canonical torsion.}
Choose a disc $\Delta$ with no other critical value, with coordinate
$t$ vanishing at $b$, and write $U=f^{-1}(\Delta)$. The coherent
direct image $f_*\omega_U$ is torsion-free of rank one: nonzero
base functions act injectively on sections, and a nearby torus has
one-dimensional canonical space. It is therefore a line bundle
on the disc. A generator gives a section $v$ of $\omega_U$,
nonzero and nowhere vanishing on every noncentral fibre. Its
divisor is $aS$ for an integer $a\ge0$. If $a\ge m$, then $v/t$
would be holomorphic, contradicting primitivity of the generator
at zero. We have $0\le a<m$ and
\[
 \omega_U\simeq\OO_U(aS),\qquad\OO_U(mS)\simeq\OO_U.
\]
Adjunction gives $\omega_S\simeq\OO_U((a+1)S)|_S$, proving
$\omega_S^{\otimes m}\simeq\OO_S$. The ambient
Riemann--Roch computation of P4.6, applied to $S$, gives
$\chi(\OO_S)=0$. Reduced irreducibility gives $h^0(\OO_S)=1$.
A nonzero section of a torsion line bundle has a power equal to
a nonzero constant, and hence is nowhere zero; its space of
sections consequently has dimension at most one. Serre duality
on the Gorenstein surface $S$ gives
$h^2(\OO_S)=h^0(\omega_S)\in\{0,1\}$, proving the assertion
about $h^1$.

\noindent\textbf{Step 2. Exceptional curves and discrepancies.}
For a minimal resolution $\pi:T\to S$ of a normal Gorenstein
surface write $K_T=\pi^*K_S+\sum a_jE_j$. Its exceptional
intersection matrix $M$ has nonnegative off-diagonal entries
and is negative definite. We recall the local analytic proof.
Pulling back a generic holomorphic function vanishing at a
contracted point gives $H+\sum m_jE_j$, with $m_j>0$ and $H$
the strict transform. Intersecting gives $Mm\le0$, strictly in
at least one entry of each connected block. After scaling by
$\operatorname{diag}(m_j)$, the negative quadratic form is a
sum of nonnegative terms $c_{ij}(z_i-z_j)^2$ and $d_i z_i^2$,
with connected support and some $d_i>0$ in every block. It is
positive definite. This is the quadratic identity of
\cite[\S I, unnumbered assertion and proof, p.~6]{Mumford1961}.
The displayed local holomorphic function and exceptional blocks give
all its inputs for these analytic germs.

Adjunction gives
$K_T\cdot E_j=2p_a(E_j)-2-E_j^2\ge0$: the sole negative
possibility is a smooth rational $(-1)$-curve, excluded by
minimality of the resolution. The inverse of a negative definite
matrix with nonnegative off-diagonal entries has nonpositive
entries. Indeed, put $Q=-M$ and choose $c$ larger than all its
eigenvalues and diagonal entries; then
$Q^{-1}=c^{-1}\sum_{k\ge0}(I-Q/c)^k$ is entrywise nonnegative.
The discrepancy equations $Ma=(K_T\cdot E_j)_j$ imply
$a_j\le0$. Hence $K_T=\pi^*K_S-E$ with $E$ an effective
exceptional integral divisor.

Put $p=\dim H^0(S,R^1\pi_*\OO_T)$. Leray, with this sheaf
supported at finitely many points, gives
\begin{equation}\label{eq:P4chi-resolution}
 \chi(\OO_T)=\chi(\OO_S)-p=-p.
\end{equation}
For any effective exceptional cycle $Z$, the higher direct image
$R^2\pi_*I_Z$ vanishes because all fibres of $\pi$ have dimension
at most one. The exact sequence of $Z$ gives a surjection
$R^1\pi_*\OO_T\to R^1\pi_*\OO_Z$, and therefore
\begin{equation}\label{eq:P4exceptional-genus}
 p\ge h^1(Z,\OO_Z).
\end{equation}

\noindent\textbf{Step 3. Exclusion of normal singular supports.}
Suppose first that $E=0$. Then $K_T$ is torsion and the
classification gives Kodaira dimension zero and nonnegative
holomorphic Euler characteristic. Equation
\eqref{eq:P4chi-resolution} forces $p=0$ and $\chi(\OO_T)=0$.
The remaining surfaces are tori, bielliptic surfaces and Kodaira
surfaces. They have no rational curves: a map from $\Pp^1$ to
a torus is constant, bielliptic surfaces have finite unramified
torus covers, and a primary Kodaira surface has an elliptic
fibration with smooth elliptic reduced fibres over an elliptic
curve. A map from $\Pp^1$ to a secondary Kodaira surface lifts
to its finite unramified primary Kodaira cover. On a torsion
canonical surface a negative curve is a smooth rational
$(-2)$-curve by adjunction. Thus there are no exceptional
curves, and $S$ is smooth.

Suppose next that $E\ne0$. Since $K_T+E$ is torsion, all
positive plurigenera vanish, so $T$ is rational, ruled, or a
blowup of a class-VII surface. Rationality contradicts
\eqref{eq:P4chi-resolution}, since $\chi(\OO_T)=1$. In class
VII that equation gives $p=0$, whereas
\[
 \chi(\OO_E)=-E\cdot(E+K_T)/2=0.
\]
The nonempty effective compact divisor $E$ has $h^0(\OO_E)\ge1$,
so $h^1(\OO_E)\ge1$, contradicting
\eqref{eq:P4exceptional-genus}.

Finally suppose $T$ is ruled, after point blowups, over a curve
$C$ of genus $g\ge1$. Equation \eqref{eq:P4chi-resolution}
gives $p=g-1$. A general ruling fibre $F$ has $K_T\cdot F=-2$,
so $E\cdot F=2$. An exceptional component $E_j$ therefore
maps onto $C$. Its normalization has genus at least $g$ by
Riemann--Hurwitz, so $h^1(\OO_{E_j})\ge g$. Taking $Z=E_j$
in \eqref{eq:P4exceptional-genus} gives $p\ge g$, a contradiction.
This exhausts the normal case.

\noindent\textbf{Step 4. The resolved normalization of a nonnormal support.}
Write $\nu:N\to S$ for the finite normalization and $\pi:T\to N$
for its minimal resolution. The residue comparison of dualizing
sheaves gives a pole divisor along the conductor; its strict
transform $D$ on $T$ is nonzero and effective. For completeness,
pull back a local generator of $\omega_S$ as a meromorphic
two-form. Its canonical divisor has the form $-D+\sum a_jE_j$.
The equations for the exceptional coefficients have right side
$(K_T+D)\cdot E_j\ge0$, because both $K_T\cdot E_j$ and
$D\cdot E_j$ are nonnegative. The same inverse-matrix argument
gives $a_j\le0$. Taking the $m$th power to trivialize the pulled
back $\omega_S$ consequently gives
\[
 mK_T+D'\sim0
\]
for a nonzero effective integral divisor $D'$. This is also the
conductor formula of \cite[Lemma~4.6, p.~7 of the attached-source edition]{CDP20}; its local proof does
not use the disputed global restriction conclusion. Any
positive pluricanonical section on $T$, multiplied by the
appropriate effective section and raised to a power, would give
a nonzero holomorphic function on compact $T$ with a zero.
It follows that $\kappa(T)=-\infty$. The classification \cite[Chapter~VI, Theorem~1.1, p.~243]{BHPV2004E}
gives precisely the stated alternatives for the minimal model.
\end{proof}

The ruled-genus and nonnormal-conductor arguments needed after
this reduction are proved below and completed in P4.58.

\begin{proposition}[P4.13]\label{night:P4.13}\label{stmt:P4.13}
\StatementRecord{night:P4.13}{proved}{Historical events: GPT-6 Astra (OpenAI), ChatGPT Pro mode, effort Pro, S6\_06\_BLIND, 16--18 November 2026: targeted proof check. Also checked in this writing. Pre-continuation writing events: GPT-6 Astra (OpenAI), ChatGPT Pro mode, effort Pro, S6\_10\_PAPERS, 18 November 2026: targeted proof check; used and rechecked in this run. Continuation metadata events: GPT-6 Astra (OpenAI), ChatGPT Pro mode, effort Pro, S6\_10\_PAPERS continuation, 19 November 2026: used and rechecked in this run; status-vocabulary reconciliation; metadata reconciliation only; no new mathematical check. \emph{Pre-continuation recorded proof scope.} The night-run re-derivation was by GPT-6 Astra (OpenAI, ChatGPT Pro mode, effort Pro), 16--18 November 2026. Also checked in this writing by GPT-6 Astra (OpenAI), ChatGPT Pro mode, effort Pro, 18 November 2026. The exceptional multiplicity bound, explicit local monodromy traces on every normal-crossing stratum, spectral-sequence alternating trace, quasi-unipotence and determinant sign and bound were re-derived. A'Campo1975 section2 pp237-243 and section3 pp243-245 were subsequently opened in primary-host indexed text, including the exact analytic hypotheses and stalk-trace proof. Direct visual-scan retrieval remains separately unavailable}{the hypotheses and explicit local calculation stated here}

Let a proper holomorphic degeneration of two-dimensional complex
tori have smooth total threefold and central fibre $mS$ with
irreducible reduced support. If $A$ is the monodromy on $H_1$ of
a nearby torus, then
\[
 m e(S_{\mathrm{reg}})=\det(I-A^m),\qquad
 0\le e(S_{\mathrm{reg}})\le16/m.
\]
This assertion requires no Betti-number hypothesis on a compactification.
\end{proposition}
\begin{proof}
\noindent\textbf{Step 1.} We resolve the pair by a proper holomorphic modification
$\pi:U'\to U$ which is an isomorphism away from $S$ and over
$S_{\mathrm{reg}}$. The total central divisor has normal crossings.
The strict transform has multiplicity $m$, and each exceptional
component has multiplicity at least $2m$. Indeed, its first centre
lies over a singular point of the reduced divisor, whose order is
at least two; subsequent centres either retain such an exceptional
multiplicity or add further positive multiplicities.

\noindent\textbf{Step 2.} We spell out the trace comparison. At a point on $k$ branches,
with local equation $z_1^{n_1}\cdots z_k^{n_k}=t$, the nearby
fibre retracts onto a disjoint union of real $(k-1)$-tori.
Monodromy permutes these components and translates any component
that it preserves. For $k\ge2$ the alternating cohomological
trace is zero. On a single branch of multiplicity $n$ there are
$n$ contractible nearby components in a cyclic permutation;
the $d$th iterate has trace $n$ when $n\mid d$, and zero otherwise.
Alternating traces in the nearby-cycle spectral sequence, with a
finite stratification of the compact central divisor, therefore give
\[
 L(A^d;H^*(T^4;\C))=\sum_{n_j\mid d}n_j e(E_j^\circ).
\]
This is the global form of A'Campo 1975, Theorem 1; his proper
analytic construction in Section 2 and spectral-sequence proof in
Section 3 give precisely this comparison. The explicit local
calculation verifies the hypotheses without a Kähler assumption.
For $d=m$ the strict transform is the sole contribution, and its
open stratum is $S_{\mathrm{reg}}$.

\noindent\textbf{Step 3.} These same charts prove quasi-unipotence. A common multiple $N$
of the $n_j$ makes monodromy act trivially on all local
nearby-cohomology sheaves, hence on the $E_2$ page. Its action
on the filtered abutment is consequently unipotent. Every
eigenvalue of $A$ is a root of unity. Since torus cohomology is
the exterior algebra on $H_1$, its Lefschetz number is
$\det(I-A^m)$. Complex conjugate eigenvalues contribute
$|1-\lambda^m|^2$; real eigenvalues are $1$ or $-1$ and
contribute zero or two. Thus this determinant is nonnegative
and at most $2^4$, as required.
\end{proof}

\begin{proposition}[P4.14]\label{night:P4.14}\label{stmt:P4.14}
\StatementRecord{night:P4.14}{proved}{Historical events: GPT-6 Astra (OpenAI), ChatGPT Pro mode, effort Pro, S6\_06\_BLIND, 16--18 November 2026: targeted proof check. Also checked in this writing. Pre-continuation writing events: GPT-6 Astra (OpenAI), ChatGPT Pro mode, effort Pro, S6\_10\_PAPERS, 18 November 2026: targeted proof check; used and rechecked in this run. Continuation metadata events: GPT-6 Astra (OpenAI), ChatGPT Pro mode, effort Pro, S6\_10\_PAPERS continuation, 19 November 2026: used and rechecked in this run; status-vocabulary reconciliation; metadata reconciliation only; no new mathematical check. \emph{Pre-continuation recorded proof scope.} The night-run re-derivation was by GPT-6 Astra (OpenAI, ChatGPT Pro mode, effort Pro), 16--18 November 2026. Also checked in this writing by GPT-6 Astra (OpenAI), ChatGPT Pro mode, effort Pro, 18 November 2026. For y(y-x\textasciicircum{}2)=0, the normalization map is (a,b) with b-a divisible by t\textasciicircum{}2; its inverse expression a+y(b-a)/x\textasciicircum{}2 gives delta two and conductor order two on both branches. For xy(x-y)=0, equality of constants gives two conditions and the linear relation gives one more; the displayed degree-k monomials separate every coefficient for k>=2, giving delta three and order-two conductor on all three branches. Plane-cubic adjunction proves trivial dualizing sheaf and arithmetic genus one. These are actual reducible curves contradicting the sentence in CDP20 Proposition7.3, attached-source p.16, without asserting a counterexample to the global threefold sign theorem}{the hypotheses and explicit local calculation stated here}

The reducible-curve inference in \cite[proof of Proposition~7.3,
p.~16 of the attached-source edition]{CDP20} omits the two plane cubics
\[
 Z(Y^2-XZ)=0,\qquad XY(X-Y)=0.
\]
They are reduced Gorenstein curves of arithmetic genus one with
trivial dualizing sheaf and rational normalization, but have
respectively two tangent components and three concurrent
components. Their conductor order is two on every normalization
branch. These curves refute the asserted nodal-cycle
classification at that step; they are not counterexamples to
the threefold sign assertion.
\end{proposition}
\begin{proof}
\noindent\textbf{Step 1.} Plane adjunction gives $\omega_C=\OO_C(3-3)=\OO_C$ and
$p_a(C)=1$. The first cubic has a line tangent to a smooth
conic at $[1:0:0]$. Its local equation is $y(y-x^2)=0$.
The normalization restrictions are $(p(t,0),p(t,t^2))$.
Their image is exactly the pairs $(a,b)$ with $b-a$ divisible
by $t^2$: the inverse expression is
$p(x,y)=a(x)+y(b(x)-a(x))/x^2$. Thus the normalization
quotient has length two and the conductor exponent is two on
each branch.

\noindent\textbf{Step 2.} For the second cubic the branches are $(t,0),(0,t),(t,t)$.
All constant coefficients must agree, and the third linear
coefficient is the sum of the first two. These three conditions
are sufficient. In every degree $k\ge2$, the monomials
$x^k-x^{k-1}y$, $y^k-xy^{k-1}$ and $x^{k-1}y$ independently
realize the coefficients on the three branches. Hence the
normalization quotient has length three and each conductor
exponent is again two. The normalization Euler characteristics
are two and three; subtracting these quotient lengths gives
$\chi(\OO_C)=0$ in both cases.

\noindent\textbf{Step 3.} The program \texttt{code/part4/genus\_one\_conductor.py}
executes the exact jet restriction matrices through orders
$2,3,4,6,8$, computing the ranks and conductor exponents.
When these branches vary in a surface, their permutation may
assemble them into a single normalization curve covering the
image conductor with degree two or three. Consequently a
unique horizontal conductor component with multiplicity two
does not alone force a degree-one map or generically reduced
image conductor. An additional argument is needed at that
precise point of the general irreducible case.
\end{proof}

The branch permutations in P4.14 have explicit local realizations.
Take an elliptic curve $C$ with a free translation of order two or
three. The map $(x,y)\mapsto(x,x^2-y)$ exchanges the two tacnode
branches and preserves $y(y-x^2)$. The order-three map
$(x,y)\mapsto(-y,x-y)$ cycles the three branches and preserves
$xy(x-y)$. The diagonal action on $C$ times a small invariant
neighbourhood of the origin is free. Its quotient is a smooth
threefold; the divisor is irreducible because the action permutes
its branches transitively. Its normalization conductor is one
elliptic curve mapping to $C/\langle\tau\rangle$ with degree two
or three. These specify the finite morphisms in P4.14 explicitly;
they make no assertion of a compact torus smoothing.

\begin{proposition}[P4.15]\label{night:P4.15}\label{stmt:P4.15}
\StatementRecord{night:P4.15}{proved}{Historical events: GPT-6 Astra (OpenAI), ChatGPT Pro mode, effort Pro, S6\_06\_BLIND, 16--18 November 2026: targeted proof check. Also checked in this writing. Pre-continuation writing events: GPT-6 Astra (OpenAI), ChatGPT Pro mode, effort Pro, S6\_10\_PAPERS, 18 November 2026: targeted proof check. Continuation metadata events: GPT-6 Astra (OpenAI), ChatGPT Pro mode, effort Pro, S6\_10\_PAPERS continuation, 19 November 2026: used and rechecked in this run; status-vocabulary reconciliation; metadata reconciliation only; no new mathematical check. \emph{Pre-continuation recorded proof scope.} The night-run re-derivation was by GPT-6 Astra (OpenAI, ChatGPT Pro mode, effort Pro), 16--18 November 2026. Also checked in this writing by GPT-6 Astra (OpenAI), ChatGPT Pro mode, effort Pro, 18 November 2026. The actual derivative is a sheaf injection of rank two into rank three. The depth lemma excludes a zero-dimensional torsion subsheaf in its cokernel, yielding I\_Z tensor L and e(T)=K\_T\textasciicircum{}2-length(Z). Substitution in each ruled and class-VII Chern formula gives the stated bounds and minimal zero cases. Adjunction on a minimal elliptic ruled surface makes every coefficient-one horizontal conductor smooth elliptic and excludes vertical poles; the doubled-section case uses the independently proved nonnegative Euler number of the open support. On a Hopf surface rational curves are excluded by their lift to C\textasciicircum{}2 minus zero. At a finite immersion every local pair of surface branches intersects along the same smooth conductor curve, so the actual branch maps are finite unramified elliptic maps and the support Euler number is zero}{\ref{night:P4.12}; \ref{night:P4.13}; \ref{night:P4.8}}

Let $mS$ be an irreducible singular fibre in the setting of 4.2.
Suppose its normalization $T$ is smooth and $d\nu:T_T\to\nu^*T_X$
has rank two away from finitely many points. Let $\ell$ be the
length of its maximal-minor scheme. Then
\[
 e(T)=K_T^2-\ell.
\]
The alternatives for $T$ are a rational surface with
$e(T)=6-\ell/2$ and $0\leq\ell\leq6$, a minimal ruled surface
over an elliptic curve with $\ell=0$, or a Hopf surface with
$\ell=0$. In either of the last two alternatives $e(S)=0$.
The reduced conductor and singular locus are disjoint unions of
smooth elliptic curves, and their normalization maps are finite
and unramified.
\end{proposition}
\begin{proof}
\noindent\textbf{Step 1.} The differential is injective as a sheaf map: its kernel vanishes
on a dense open set, and a locally free sheaf is torsion-free.
Its rank-one cokernel $Q$ is locally free away from finitely many
points. The depth lemma applied to the length-one locally free
resolution gives depth at least one at those points, excluding
a finite-length subsheaf. All torsion would be supported there,
so $Q$ is torsion-free. On the smooth surface we obtain
\[
 Q=\mathcal I_Z\otimes L,\qquad
 L=K_T\otimes\nu^*(\det T_X),\qquad \operatorname{length}(Z)=\ell.
\]
The two-column presentation by $d\nu$ identifies the local ideal
of $Z$ with its three maximal minors. The classes $c_1(T_X)$ and
$c_2(T_X)$ vanish in real cohomology, since $b_2(X)=b_4(X)=0$.
The second Chern class of the exact sequence therefore gives
\[
 0=e(T)-K_T^2+\ell.
\]
P4.12 gives $\kappa(T)=-\infty$. For a ruled surface over a
genus-$g$ curve, blown up $k$ times, we have
\[
 e(T)=4(1-g)+k,\qquad K_T^2=8(1-g)-k.
\]
Consequently $4(g-1)+2k+\ell=0$. Thus $g\leq1$, and $g=1$
forces $k=\ell=0$. In the class-VII case, including blowups,
$K_T^2=-e(T)$ and $e(T)=b_2(T)\geq0$. Hence
$2e(T)+\ell=0$, and the surface is minimal with $b_2=0$.
Its nonempty conductor, the finite-cover classification, and the
absence of curves on Inoue surfaces proved in P4.8 force it to
be Hopf. For rational $T$, Noether's formula gives
$2e(T)+\ell=12$. The inequality $e(T)\geq3$ gives the asserted
range. When $\ell=0$, the cokernel is a line bundle and $d\nu$
has rank two everywhere.

\noindent\textbf{Step 2.} We now compute the image Euler number in the two latter cases.
The conductor divisor $D$ satisfies $K_T+D\equiv0$. First let
$T$ be a minimal ruled surface over an elliptic curve $E$;
it is a $\mathbb P^1$-bundle. For a ruling fibre $F$, $D\cdot F=2$.
If a horizontal component $H$ occurs with multiplicity one,
adjunction gives
\[
 2p_a(H)-2=(K_T+H)H=-(D-H)H\leq0.
\]
The map to $E$ gives geometric genus at least one. Equality
therefore holds, $H$ is smooth elliptic, and $H$ is disjoint
from every other component of $D$. There is no vertical
component because every ruling fibre meets $H$. The possibilities
are one elliptic double-section or two disjoint elliptic sections.

\noindent\textbf{Step 3.} The other possibility is $D=2H+\sum_j a_jF_j$, where $H$ is a
section, $F_j$ are distinct ruling fibres and $a_j>0$.
The reduced conductor has Euler number equal to the number $k$
of these fibres: $e(H)=0$, $e(F_j)=2$, and each meets $H$ once.
The immersive normalization identifies $T\setminus D_{\mathrm{red}}$
with $S_{\mathrm{reg}}$. Hence
\[
 e(S_{\mathrm{reg}})=e(T)-e(D_{\mathrm{red}})=-k.
\]
P4.13 forces $k=0$. Thus in every case the reduced conductor is
a disjoint union of smooth elliptic curves.

\noindent\textbf{Step 4.} On a Hopf surface every irreducible curve $H$ satisfies
$H^2=K_TH=0$, so $p_a(H)=1$. A rational curve would lift from
$\mathbb P^1$ to the universal cover $\mathbb C^2\setminus\{0\}$,
where both coordinate functions must be constant. Thus $H$
is smooth elliptic. Distinct curves have intersection zero and
are disjoint, proving the same conductor description.

\noindent\textbf{Step 5.} To check the image, at a singular point the finite immersive
normalization presents the germ as distinct smooth surface
branches $S_i$. The reduced conductor on each branch is a single
smooth curve germ $C_i$. Any two branches intersect in a curve
in the smooth ambient threefold, and its support lies in both
$C_i$ and $C_j$. Their irreducibility forces those curve germs
to agree. Thus the image singular locus is one smooth curve
germ, and the conductor maps locally isomorphically to it.
Globally this map is finite and unramified. Riemann--Hurwitz
then makes every image component elliptic. Finally
\[
 e(S)=e(T)-e(D_{\mathrm{red}})+e(\operatorname{Sing}S)=0.
\]
\end{proof}

\begin{proposition}[P4.16]\label{night:P4.16}\label{stmt:P4.16}
\StatementRecord{night:P4.16}{proved}{Historical events: GPT-6 Astra (OpenAI), ChatGPT Pro mode, effort Pro, S6\_06\_BLIND, 16--18 November 2026: targeted proof check. Also checked in this writing. Pre-continuation writing events: GPT-6 Astra (OpenAI), ChatGPT Pro mode, effort Pro, S6\_10\_PAPERS, 18 November 2026: targeted proof check. Continuation metadata events: GPT-6 Astra (OpenAI), ChatGPT Pro mode, effort Pro, S6\_10\_PAPERS continuation, 19 November 2026: used and rechecked in this run; status-vocabulary reconciliation; metadata reconciliation only; no new mathematical check. \emph{Pre-continuation recorded proof scope.} The night-run re-derivation was by GPT-6 Astra (OpenAI, ChatGPT Pro mode, effort Pro), 16--18 November 2026. Also checked in this writing by GPT-6 Astra (OpenAI), ChatGPT Pro mode, effort Pro, 18 November 2026. The explicit cover of x\textasciicircum{}2-zy\textasciicircum{}n=1 has U=C* times C and two contractible graph neighborhoods, with each overlap C* times a disc. Van Kampen and Mayer-Vietoris give pi1=0 and H2=Z. The compactification performs n blowups over each of (1,0) and (-1,0); its rational-tree boundary makes the compact class primitive. Orthogonality to that boundary gives the class sum E\_plus minus sum E\_minus, of square -2n. Positive weights (n+2,2,4) carry compact representatives and transverse companions into an actual local fibre. The weighted holomorphic unit-root map removes a unit exactly, and the full-fibre root comparison embeds the resulting spherical class into a nearby aspherical torus, where its nonzero square is impossible. Archived numerical runs are not counted as current computations}{the hypotheses and explicit local calculation stated here}

For every integer $n\geq1$, the reduced hypersurface germ
$x^2-zy^n=0$ cannot occur in a fibre of a proper holomorphic
degeneration of complex two-tori with smooth total threefold,
even with any positive common multiplicity on that germ.
In particular the Whitney umbrella ($n=2$) is excluded.
No hypothesis on $b_2$ is needed.
\end{proposition}
\begin{proof}
\noindent\textbf{Step 1.} Put $F_n=\{x^2-zy^n=1\}\subset\mathbb C^3$. We exhibit its
compact spherical homology class and compute its square.
The open subset $U=\{y\ne0\}$ is
$\mathbb C^*\times\mathbb C$ in coordinates $(y,x)$.
For $0<\delta<1$, take the two disjoint graph neighbourhoods
\[
 N_\pm:\quad x=\pm\sqrt{1+zy^n},\quad |zy^n|<\delta.
\]
Their $(y,z)$ domains contract by simultaneous real scaling.
Each $N_\pm\cap U$ has coordinates $(y,zy^n)$ in
$\mathbb C^*\times\{|w|<\delta\}$. Its map to $U$ takes
the generating $y$-loop to the generating $y$-loop. These three
sets cover $F_n$. Van Kampen gives $\pi_1(F_n)=0$.
Mayer--Vietoris identifies $H_2(F_n;\mathbb Z)$ with the
kernel of $(a,b)\mapsto a+b$ on the two intersection
$H_1$ groups, hence with $\mathbb Z$. Hurewicz represents its
generator by a map $S^2\to F_n$.

\noindent\textbf{Step 2.} We compute the intersection form by compactification. Blow up
$\mathbb P^1_x\times\mathbb P^1_y$ successively $n$ times above
each point $(1,0)$ and $(-1,0)$. Near either point let
$a=x^2-1$, $b=y$. The successive charts at the centres are
\[
 a=b^j a_j,\qquad z=\frac{a_j}{b^{n-j}}.
\]
For $j<n$, blow up the next centre $a_j=b=0$.
At $j=n$, $z=a_n$ is regular. The final exceptional curve
minus its point on the preceding boundary is precisely the
line $y=0,x=\pm1$ in $F_n$. Every earlier exceptional curve
is a pole divisor of $z$. Consequently $F_n=T_n\setminus B$,
where $T_n$ is this smooth projective blown-up surface and
$B$ consists of the two infinity divisors, the strict transform
of $y=0$, and the first $n-1$ exceptional curves on each chain.
Both infinity divisors are removed, so no extension of the
rational function at their intersection is asserted or needed.

\noindent\textbf{Step 3.} In the total-transform basis $H_x,H_y,E_{\pm,j}$ the nonzero
products are $H_xH_y=1$ and $E_{\pm,j}^2=-1$. The boundary
classes are
\[
 H_x,\ H_y,\ H_y-E_{+,1}-E_{-,1},\quad
 E_{+,j}-E_{+,j+1},\ E_{-,j}-E_{-,j+1}\quad(1\leq j<n).
\]
They form a connected tree of rational curves; thus $H_1(B)=0$.
Compact-support duality and the sequence of $(T_n,B)$ identify
the image of $H_2(F_n)$ in $H_2(T_n)$ with the kernel of
intersection against these boundary classes. Orthogonality to
$H_x,H_y$ removes both ruling coefficients. Orthogonality to
the exceptional differences makes the coefficients constant
along each chain. The strict transform of $y=0$ makes the two
constants opposite. Thus the kernel is generated by the
primitive integral vector
\[
 v=\sum_{j=1}^n E_{+,j}-\sum_{j=1}^n E_{-,j},\qquad v^2=-2n.
\]
The sequence introduces no multiple: $H_1(B)=0$ injects
$H^2(T_n,B)$ into $H^2(T_n)$ with precisely the stated kernel.
Therefore the compact spherical generator has square $-2n$.
The exact matrices, primitive kernels, boundary graphs and
blowup identities for $n=1,\ldots,12$ are independently checked
by \texttt{code/part4/milnor\_intersection.py}; its output and
JSON data are included. The proof is uniform in $n$.

\noindent\textbf{Step 4.} The defining polynomial is weighted homogeneous of degree
$d=2n+4$ with positive weights $(n+2,2,4)$. Scaling by
\[
 (x,y,z)\longmapsto
 (\varepsilon^{n+2}x,\varepsilon^2y,\varepsilon^4z)
\]
identifies $F_n$ with the fibre of value $\varepsilon^{2n+4}$.
It carries any compact representative of $v$, together with a
compact transverse representative for its square, into an
arbitrarily small neighbourhood of the origin.

\noindent\textbf{Step 5.} We also check the unit and multiplicity in an actual defining
map. Take an analytic $m$-th root of its unit. For a unit $h$,
the local analytic isomorphism
\[
 \Phi(x,y,z)=
 \bigl(h^{(n+2)/d}x,h^{2/d}y,h^{4/d}z\bigr)
\]
satisfies $g\circ\Phi=hg$, where $g=x^2-zy^n$; its derivative
at the origin is invertible. Thus the actual local map to the
base is $(g\circ\Phi)^m$. The small local fibre of $g$ embeds
in the corresponding smooth compact torus fibre.

\noindent\textbf{Step 6.} The spherical representative is null-homotopic in that torus,
whose universal cover is the contractible space $\mathbb C^2$.
It is therefore zero in $H_2$, and its self-intersection there
is zero. The two compact representatives already chosen compute
the same intersection in the open local fibre and in the torus.
Their intersection is $-2n$, which is a contradiction.
\end{proof}

P4.16 excludes the displayed pinching and ramification germs.
The subsequent arguments retain singular normalizations and
nonimmersive normalization maps and complete their classification
in P4.58.

\begin{proposition}[P4.17]\label{night:P4.17}\label{stmt:P4.17}
\StatementRecord{night:P4.17}{proved}{Historical events: GPT-6 Astra (OpenAI), ChatGPT Pro mode, effort Pro, S6\_06\_BLIND, 16--18 November 2026: targeted proof check. Also checked in this writing. Pre-continuation writing events: GPT-6 Astra (OpenAI), ChatGPT Pro mode, effort Pro, S6\_10\_PAPERS, 18 November 2026: targeted proof check. Continuation metadata events: GPT-6 Astra (OpenAI), ChatGPT Pro mode, effort Pro, S6\_10\_PAPERS continuation, 19 November 2026: used and rechecked in this run; status-vocabulary reconciliation; metadata reconciliation only; no new mathematical check. \emph{Pre-continuation recorded proof scope.} The night-run re-derivation was by GPT-6 Astra (OpenAI, ChatGPT Pro mode, effort Pro), 16--18 November 2026. Also checked in this writing by GPT-6 Astra (OpenAI), ChatGPT Pro mode, effort Pro, 18 November 2026. With a graph neighborhoods the same exact cover gives H2=ker(Z\textasciicircum{}a to Z) and a simply connected fibre. The a exceptional chains in the compactification yield the primitive differences v\_i=sum E\_i,j-sum E\_i+1,j; their intersection form is minus n times the Cartan matrix A\_(a-1). The positive-weight and unit-root maps give the local torus contradiction. The finite normalization map (u,v) to (u\textasciicircum{}n v,u\textasciicircum{}a,v\textasciicircum{}a) has its exact mu\_a action and degree-a function-field extension, so the normal invariant ring is the normalization. For (a,n)=(2,3) it is UW=V\textasciicircum{}2 with x=UV,y=U,z=W; for n=a the conductor map is v to v\textasciicircum{}a. This proves the exclusions through the displayed morphisms}{\ref{night:P4.16}}

For all $a\geq2$ and $n\geq1$, the germ $x^a-zy^n=0$ is
excluded in the setting of P4.16. Its affine nearby fibre has
$\pi_1=0$ and compact second-homology intersection lattice
$A_{a-1}(-n)$. This includes the order-three cyclic-ramification
germ $x^3-zy^3$ and examples with singular normalization.
\end{proposition}
\begin{proof}
\noindent\textbf{Step 1.} Replace the two roots of $x^2=1$ in P4.16 by the $a$ distinct
roots of $x^a=1$. The open set $y\ne0$ is still
$\mathbb C^*\times\mathbb C$. There are $a$ contractible graph
neighbourhoods, each meeting it in $\mathbb C^*\times\Delta$
with the same generating $y$-loop. Thus van Kampen gives
$\pi_1=0$, and Mayer--Vietoris gives
\[
 H_2(F_{a,n};\mathbb Z)=\ker\bigl(\mathbb Z^a
   \xrightarrow{(c_i)\mapsto\sum_i c_i}\mathbb Z\bigr).
\]
At each root, $x^a-1$ is a local coordinate. The same $n$
successive blowup charts give a compactification with $a$
exceptional chains and boundary classes
\[
 H_x,\ H_y,\ H_y-\sum_i E_{i,1},\qquad
 E_{i,j}-E_{i,j+1}\quad(1\leq i\leq a,\ 1\leq j<n).
\]
This boundary is a tree of rational curves. The exact sequence
of the pair identifies compact $H_2$ integrally with the kernel
of intersection against these classes. Its primitive basis is
\[
 v_i=\sum_{j=1}^n E_{i,j}-\sum_{j=1}^n E_{i+1,j}
       \quad(1\leq i<a).
\]
The products are $v_i^2=-2n$, $v_iv_{i+1}=n$, with all other
off-diagonal products zero. This is $A_{a-1}(-n)$.
Hurewicz makes every class spherical. Positive weights
$(n+1,a,a)$, of degree $a(n+1)$, scale compact representatives
into arbitrarily small nearby fibres. The explicit unit removal
and torus argument of P4.16 now exclude $v_1$, whose square
is $-2n$. The program independently checks the complete
matrices and primitive bases for $2\leq a\leq6$, $1\leq n\leq6$.

\noindent\textbf{Step 2.} We specify the normalization morphism. Let $\mu_a$ act on
$\mathbb C^2_{u,v}$ by $(u,v)\mapsto(\zeta u,\zeta^{-n}v)$.
Then
\[
 x=u^nv,\qquad y=u^a,\qquad z=v^a
\]
gives a finite map to $x^a=zy^n$: both $u$ and $v$ satisfy
monic equations over its coordinate ring. The function-field
degree is $a$, with the displayed faithful $\mu_a$ action.
Consequently the induced quotient map is finite and birational.
The invariant ring is integrally closed, so this is the
normalization. For $(a,n)=(2,3)$ it has coordinates
$U=u^2,V=uv,W=v^2$ with $UW=V^2$, and its map is
$(x,y,z)=(UV,U,W)$. Thus it has an $A_1$ singularity.
For $n=a$ the quotient is smooth with coordinates $(y,v)$,
and its conductor map is $v\mapsto v^a$. This verifies
explicitly the cyclic branch-permutation description.
\end{proof}

\begin{proposition}[P4.18]\label{night:P4.18}\label{stmt:P4.18}
\StatementRecord{night:P4.18}{proved}{Historical events: GPT-6 Astra (OpenAI), ChatGPT Pro mode, effort Pro, S6\_06\_BLIND, 16--18 November 2026: targeted proof check. Also checked in this writing. Pre-continuation writing events: GPT-6 Astra (OpenAI), ChatGPT Pro mode, effort Pro, S6\_10\_PAPERS, 18 November 2026: targeted proof check. Continuation metadata events: GPT-6 Astra (OpenAI), ChatGPT Pro mode, effort Pro, S6\_10\_PAPERS continuation, 19 November 2026: used and rechecked in this run; status-vocabulary reconciliation; metadata reconciliation only; no new mathematical check. \emph{Pre-continuation recorded proof scope.} The night-run re-derivation was by GPT-6 Astra (OpenAI, ChatGPT Pro mode, effort Pro), 16--18 November 2026. Also checked in this writing by GPT-6 Astra (OpenAI), ChatGPT Pro mode, effort Pro, 18 November 2026. Writing b=dB,c=dC and Bp+Cq=1 gives the literal punctured-fibre coordinates y=h\textasciicircum{}p t\textasciicircum{}C,z=h\textasciicircum{}q t\textasciicircum{}-B, with t=y\textasciicircum{}q z\textasciicircum{}-p. Over a disc joining two simple zeros of x\textasciicircum{}a-1, the two graph-neighborhood loop maps have matrix [[B,C],[q,-p]], determinant -1, and kill both generators. The finite map (X,y,z) to (X,y,z\textasciicircum{}C) has degree C; pullback of the compact Poincare-dual class of square -2B has square -2BC. Hurewicz, positive weights and the exact unit-root comparison place a nonzero spherical square in an actual nearby torus, giving the contradiction. The separate mixed example normalizes by (u,v) to (u\textasciicircum{}3 v\textasciicircum{}2,v,u\textasciicircum{}2); its monomial ring gives conductor u\textasciicircum{}2 v\textasciicircum{}2 and its derivative minors have common factor u, verifying the asserted positive-dimensional rank drop}{\ref{night:P4.15}; \ref{night:P4.16}}

For $a\geq2$ and $b,c\geq1$, the reduced germ
$x^a-y^bz^c=0$ cannot occur with a common positive multiplicity
in a torus degeneration with smooth total threefold. Its affine
nearby fibre contains a spherical class of nonzero square.
In particular this excludes $x^2-z^3y^4$, whose smooth
normalization has a divisorial differential rank drop.
\end{proposition}
\begin{proof}
\noindent\textbf{Step 1.} Write $d=\gcd(b,c)$, $b=dB$, $c=dC$, so $\gcd(B,C)=1$.
The nearby fibre $F=\{x^a-y^bz^c=1\}$ maps to the smooth
affine curve $A=\{x^a-w^d=1\}$ by $w=y^Bz^C$.
The function $w$ on $A$ has $a$ simple zeros. Choose an
embedded arc between two zeros avoiding the others, and a
small open disc $D$ around it. In a holomorphic coordinate
$\xi$ on $D$, write $w=u(\xi)(\xi-\alpha)(\xi-\beta)$
with $u$ a unit. Taking a $B$-th root of $u$, rescaling $y$,
and making an affine coordinate change identifies the preimage
of $D$ with
\[
 W=\{y^Bz^C=X^2-1,\ X\in D'\},
\]
where $D'$ is a disc containing both $-1$ and $1$.

\noindent\textbf{Step 2.} Choose $Bp+Cq=1$. With $h=X^2-1$, the explicit trivialization
away from the two special fibres is
\[
 y=h^pt^C,\quad z=h^qt^{-B},\quad t=y^qz^{-p}.
\]
Its fundamental group is $F(\beta_-,\beta_+)\times\mathbb Z\tau$.
The neighbourhood of each special fibre is the graph
$X=\pm\sqrt{1+y^Bz^C}$ on $|y^Bz^C|<\delta$, a contractible
domain by real scaling. Its punctured part has two generating
loops in $y,z$, whose images in $\langle\beta_\pm,\tau\rangle$
are the columns of
\[
 \begin{pmatrix}B&C\\q&-p\end{pmatrix},\qquad \det=-1.
\]
Thus adjoining each neighbourhood kills $\beta_\pm$ and $\tau$.
Van Kampen gives $\pi_1(W)=0$.

\noindent\textbf{Step 3.} The map
\[
 \pi:W\longrightarrow W'=\{y^BZ=X^2-1,\ X\in D'\},
 \qquad (X,y,z)\longmapsto(X,y,z^C)
\]
is proper finite of degree $C$. The compact class $v$ in
P4.16, of square $-2B$, is already supported over $D'$.
Indeed, the cover just given with $C=1$ has the identical
Mayer--Vietoris maps before and after enlarging $D'$ to the
whole plane: both punctured bases have the same two meridians,
and the same two graph neighbourhoods are attached. The
inclusion is therefore an isomorphism on $H_2$ and preserves
the primitive class and its square.
For its compact-support dual $\alpha\in H_c^2(W';\mathbb Z)$,
properness and degree give
\[
 \int_W(\pi^*\alpha)^2=C\int_{W'}\alpha^2=-2BC.
\]
One may verify this by smooth compactly supported forms:
outside the measure-zero branch divisor the map is an oriented
$C$-sheeted covering. The dual of $\pi^*\alpha$ is spherical
by $\pi_1(W)=0$ and Hurewicz. Its compact representatives
embed in $F$ with the same nonzero square.

\noindent\textbf{Step 4.} The polynomial has positive weights $(b+c,a,a)$. Weighted
scaling places these representatives in arbitrarily small
nearby fibres. P4.16's explicit unit removal handles the unit
and common multiplicity. The spherical class becomes zero
in a torus, contradicting its square $-2BC$.

\noindent\textbf{Step 5.} For the mixed example the normalization is the finite
birational map
\[
 (u,v)\longmapsto(u^3v^2,v,u^2).
\]
Finiteness follows from $u^2=z$, and $u=x/(zy^2)$ gives
birationality. The image ring contains all even powers of $u$
and the odd monomials $u^iv^j$ with $i\geq3,j\geq2$.
Its conductor is $u^2v^2\mathbb C\{u,v\}$, giving divisor
$2\{u=0\}+2\{v=0\}$. Its differential minors are
$3u^2v^2,-4u^4v,-2u$, with common divisorial factor $u$.
Thus this example explicitly lies outside P4.15's finite
rank-drop assumption. The 64 exponent matrices for
$1\leq b,c\leq8$, transfer squares and normalization
identities are executed in
\texttt{code/part4/binomial\_milnor.py}; its output and
JSON data are saved alongside it.
\end{proof}

\begin{proposition}[P4.19]\label{night:P4.19}\label{stmt:P4.19}
\StatementRecord{night:P4.19}{proved}{Historical events: GPT-6 Astra (OpenAI), ChatGPT Pro mode, effort Pro, S6\_06\_BLIND, 16--18 November 2026: targeted proof check. Also checked in this writing. Pre-continuation writing events: GPT-6 Astra (OpenAI), ChatGPT Pro mode, effort Pro, S6\_10\_PAPERS, 18 November 2026: targeted proof check. Continuation metadata events: GPT-6 Astra (OpenAI), ChatGPT Pro mode, effort Pro, S6\_10\_PAPERS continuation, 19 November 2026: used and rechecked in this run; status-vocabulary reconciliation; metadata reconciliation only; no new mathematical check. \emph{Pre-continuation recorded proof scope.} The night-run re-derivation was by GPT-6 Astra (OpenAI, ChatGPT Pro mode, effort Pro), 16--18 November 2026. Also checked in this writing by GPT-6 Astra (OpenAI), ChatGPT Pro mode, effort Pro, 18 November 2026. The plane and Whitney umbrella normalization maps, conductor ideals x squared and v squared, ramification at the origin, and the actual nonzero-fibre isomorphism with C-star times C were re-derived. The comparison cites the separately stated sphere obstructions and does not assert a compact counterexample}{\ref{night:P4.16}; \ref{night:P4.18}}

The germ $y(x^2-zy^2)=0$ has ramified conductor, of multiplicity
two on each normalization component, but its affine nearby
fibre is $\mathbb C^*\times\mathbb C$. Ramification alone
therefore does not give the spherical intersection obstruction
of P4.16--P4.18. This is a local statement, not a compact
counterexample to 4.2.
\end{proposition}
\begin{proof}
\noindent\textbf{Step 1.} One normalization component is the plane $y=0$; the other is
the Whitney normalization $(u,v)\mapsto(uv,v,u^2)$.
The singular locus is the smooth $z$-axis. The plane conductor
is $x^2=0$, obtained by restricting the other component's
equation. On the Whitney component the intrinsic conductor
is $v\mathbb C\{u,v\}$; multiplication by the added plane
equation $y=v$ gives $v^2\mathbb C\{u,v\}$. Directly, the
union ring is the fibre product of the component rings over
$\mathbb C\{x,z\}/(x^2)$, giving exactly these conductor ideals.
The reduced conductor cover is the identity sheet from the
plane together with $u\mapsto u^2$ from the Whitney component.
The latter is ramified.

\noindent\textbf{Step 2.} On $y(x^2-zy^2)=1$ we have $y\ne0$, and the inverse maps
\[
 (x,y,z)\longmapsto(y,x),\qquad
 (y,x)\longmapsto\left(x,y,\frac{yx^2-1}{y^3}\right)
\]
identify this fibre with $\mathbb C^*\times\mathbb C$.
Its $\pi_2$ and $H_2$ vanish. Hence every compact spherical
class in an open subset has zero square after inclusion,
and therefore has zero square already in that open subset.
\end{proof}

\begin{proposition}[P4.20]\label{night:P4.20}\label{stmt:P4.20}
\StatementRecord{night:P4.20}{proved}{Historical events: GPT-6 Astra (OpenAI), ChatGPT Pro mode, effort Pro, S6\_06\_BLIND, 16--18 November 2026: targeted proof check. Also checked in this writing. Pre-continuation writing events: GPT-6 Astra (OpenAI), ChatGPT Pro mode, effort Pro, S6\_10\_PAPERS, 18 November 2026: targeted proof check. Continuation metadata events: GPT-6 Astra (OpenAI), ChatGPT Pro mode, effort Pro, S6\_10\_PAPERS continuation, 19 November 2026: used and rechecked in this run; status-vocabulary reconciliation; metadata reconciliation only; no new mathematical check. \emph{Pre-continuation recorded proof scope.} The night-run re-derivation was by GPT-6 Astra (OpenAI, ChatGPT Pro mode, effort Pro), 16--18 November 2026. Also checked in this writing by GPT-6 Astra (OpenAI), ChatGPT Pro mode, effort Pro, 18 November 2026. Boundary-preserving deformations, the finite flat critical-rank scheme, the exact crosscap normal form, and transport into the actual torus were re-derived. For immersive sheets, dependent gradients give an explicit nonzero critical value and an arbitrarily small quadratic Morse perturbation, excluding non-NC intersections. The final global applications retain P4.6/P4.8 hypotheses}{\ref{night:P4.16}; \ref{night:P4.6}; \ref{night:P4.8}}

Let $mS$ be an irreducible fibre of a proper holomorphic
degeneration of complex two-tori with smooth total threefold.
Suppose the normalization of $S$ is smooth and, outside
finitely many points, $S$ is smooth or has the ordinary-double
germ $xy=0$. Then $S$ has normal crossings everywhere.
In the global setting of 4.2, P4.6--P4.8 apply to this subclass.
\end{proposition}
\begin{proof}
\noindent\textbf{Step 1.} We give the boundary and transport arguments explicitly.
Take a small ball $B$ around an exceptional point, with
boundary transverse to the smooth and double strata away
from that point. Let $g$ be a reduced equation. An analytic
$m$-th root of the unit in the actual defining function lets
us write that function as $g^m$.

\noindent\textbf{Step 2.} Consider a small holomorphic deformation $G_s$ of $g$ whose
zero set near $\partial B$ retains its smoothly moving smooth
and ordinary-double strata, transverse to $\partial B$.
There is a uniform $\delta>0$ such that
$G_s|_{\partial B}$ is submersive over $0<|t|<\delta$.
At a boundary point on a double curve, choose coordinates
$a,b,c$ with $G_s=ab$, absorbing a unit in $a$.
Transversality of the double curve to $\partial B$ allows one
real $c$-coordinate to be solved from the boundary equation.
The complex coordinates $a,b$ remain free there, so $ab$ is
submersive when nonzero. At a smooth point use $G_s=a$.
Compactness gives uniform bounds; points outside these
neighbourhoods stay away from zero.

\noindent\textbf{Step 3.} For each fixed small $s$ there are finitely many nonzero
critical values in this disc. After shrinking, the analytic
critical set in a neighbourhood of the compact closed ball
has finitely many relevant local components. The function is constant on
each analytic critical component. A component at a nonzero
small value cannot meet $\partial B$, so it is compact;
the maximum principle excludes a positive-dimensional
compact analytic subset of a ball. Thus the relevant critical
points are isolated and finite. Paths avoiding these values
join a small regular value to one fixed regular value $t_0$.
Proper Ehresmann transport on the closed ball, with the
verified boundary submersion, identifies the fibres along
such paths. For sufficiently small deformation the fibres
$G_s=t_0$ also form a smooth proper family to $s=0$, by
compactness. Hence every compact spherical class of nonzero
square in a local nearby fibre of $G_s$ transports into
$g=t_0$. Its inclusion in the actual torus fibre is impossible.

\noindent\textbf{Step 4.} At a normalization preimage write the map germ as
$\phi:(\mathbb C^2,0)\to(\mathbb C^3,0)$. Away from the
exceptional point it is immersive with transverse multiple
images. If $d\phi$ drops rank, its maximal minors define a
nonempty finite scheme. Small generic polynomial
perturbations have transverse rank-one points near it.
We check both genericity and conservation of their positive
number. Varying linear coefficients makes the universal
derivative map submersive onto the $3$-by-$2$ matrix space.
The rank-one locus is smooth of codimension two; the rank-zero
locus has codimension six. Generic perturbations are transverse
to the first and avoid the second. The height-two maximal-minor
ideal has its Hilbert--Burch resolution. Its one-parameter
zero scheme is Cohen--Macaulay of dimension one; the finite
special fibre excludes a vertical component. Thus it is flat
and preserves the positive length. Further generic polynomial
coefficients keep the finitely many rank-one images off other
sheets: incidence of a zero-dimensional critical set with a
two-dimensional image in a three-dimensional target has
negative expected dimension.

\noindent\textbf{Step 5.} Each transverse rank-one point is a Whitney crosscap, as we
now verify. Choose its first target coordinate as a source
coordinate $u$, and write $\phi=(u,F(u,v),H(u,v))$.
Transversality of the equations $F_v=H_v=0$ gives a linear
combination with nonzero second $v$-derivative. The
holomorphic Morse lemma with parameter $u$, followed by
subtracting the critical value, makes the second coordinate
$v^2$. Subtract the even part of the third coordinate to
write it as $vA(u,v^2)$. Transversality gives $A_u(0,0)\ne0$.
The target change $u\mapsto A(u,Y)$ and the matching source
change give exactly $(u,v)\mapsto(u,v^2,uv)$.
Its image is the crosscap of P4.16, with a compact spherical
class of square $-4$ in its nearby fibre.

\noindent\textbf{Step 6.} The deformed images have holomorphic equations $G_s$ with
$G_0=g$. To see this, a generic projection of the original
finite normalization to two target coordinates is finite;
Weierstrass preparation preserves this after shrinking the
deformation. The total image is an analytic hypersurface
in $\mathbb C^3\times\Delta$. Its central equation has
multiplicity one, as seen at a smooth point of $S$; a unit
adjustment gives $G_0=g$. Outside a smaller ball the original
map is a transverse immersion, and this compact boundary
data persist under small deformation. We therefore have
the boundary property already proved. Transporting the
crosscap class contradicts its torus inclusion. The
normalization is consequently immersive everywhere.

\noindent\textbf{Step 7.} Write its image germ as a unit times $g_1\cdots g_r$, where
each branch is smooth and $v_i=dg_i(0)\ne0$. Absorb the
unit in one factor. Independent gradients give $r\leq3$
and normal-crossing coordinates by the inverse function
theorem. Suppose instead $\sum_i c_iv_i=0$ is a nonzero
relation, and let $I$ be its nonzero support. Set
$a_i=\varepsilon/c_i$ for $i\in I$ and $a_j=\eta$ for
$j\notin I$, where $0<|\varepsilon|\ll|\eta|\ll1$.
Perturb constants so that $h_i(0)=a_i$. Keep all gradients
except one $i_0\in I$, and change that gradient by
\[
 \Delta v_{i_0}=-a_{i_0}\sum_{j\notin I}\frac{v_j}{a_j}.
\]
This is arbitrarily small and makes $\sum_i dh_i(0)/a_i=0$.
Thus $H=\prod_i h_i$ is critical at the origin and has
nonzero critical value $\prod_i a_i$. Its Hessian divided
by that value is
\[
 \sum_i\frac{\operatorname{Hess}(h_i)}{a_i}
 -\sum_i\frac{dh_i\otimes dh_i}{a_i^2}.
\]
An arbitrarily small generic quadratic change to $h_{i_0}$
makes this matrix invertible: its determinant is a nonzero
polynomial in the freely variable symmetric Hessian.
The branch functions remain smooth and retain their
ordinary-double/smooth boundary strata. We have constructed
a holomorphic Morse point off the divisor.

\noindent\textbf{Step 8.} Its nearby fibre contains the real vanishing sphere in
$w_1^2+w_2^2+w_3^2=\epsilon$. Its normal bundle is
$iTS^2$, with the opposite orientation to $TS^2$ under the
complex surface orientation. Its self-intersection is
therefore $-e(S^2)=-2$. Boundary transport again puts this
class in the original torus fibre, a contradiction. Hence
the original branch gradients are independent everywhere,
which proves the normal-crossing assertion.
\end{proof}

The rank-one count is also stated in Gim\'enez Conejero--Pint\'er,
arXiv:2401.14951v2, Lemma 2.3; Section 5.3, p.~33,
discusses the stable-image/nearby-fibre comparison.
The proof above verifies the boundary transport required here.
It does not extend that argument to cuspidal, tacnodal or
triple-curve transverse types.

\paragraph{Exact source locations used here.}
The classification inputs are \cite[Chapter~VI, Theorem~1.1 and
Table~10, pp.~243--244]{BHPV2004E},
\cite[Theorems~4.2 and~5.1, with proofs, pp.~17--21 and~27--29]{OVV2020},
\cite[Theorem~1, p.~1]{BBK2019}, and
\cite[Theorem~2.1]{Catanese2004}.
The no-curve theorem is credited to Inoue by
\cite[Theorem~3.1, p.~303]{KhaledTeleman2026Journal}; the actual
unramified double cover is in its proof of Remark~3.5, p.~305.
The exceptional intersection identity is
\cite[\S I, p.~6]{Mumford1961}; the proof above supplies its
local analytic function and its exact exceptional-divisor map.
The current source records distinguish theorem-statement openings
from complete source-proof readings. The sign and equality assertions
are proved through the displayed normalization, conductor, filling
and intersection maps, with the stated irreducible-or-NC scope for
the global conclusion.

\subsection{The remaining smooth ruled normalizations}

\begin{proposition}[P4.21]\label{night:P4.21}\label{stmt:P4.21}
\StatementRecord{night:P4.21}{proved}{Historical events: GPT-6 Astra (OpenAI), ChatGPT Pro mode, effort Pro, S6\_06\_BLIND, 16--18 November 2026: targeted proof check. Also checked in this writing. Pre-continuation writing events: GPT-6 Astra (OpenAI), ChatGPT Pro mode, effort Pro, S6\_10\_PAPERS, 18 November 2026: targeted proof check. Continuation metadata events: GPT-6 Astra (OpenAI), ChatGPT Pro mode, effort Pro, S6\_10\_PAPERS continuation, 19 November 2026: used and rechecked in this run; status-vocabulary reconciliation; metadata reconciliation only; no new mathematical check. \emph{Pre-continuation recorded proof scope.} The night-run re-derivation was by GPT-6 Astra (OpenAI, ChatGPT Pro mode, effort Pro), 16--18 November 2026. Also checked in this writing by GPT-6 Astra (OpenAI), ChatGPT Pro mode, effort Pro, 18 November 2026. On a smooth rational normalization, torsion adjunction gives an actual anticanonical divisor. Its complete cubic and ruled minimal-model cases have the stated elliptic, cycle or simply connected rational-support topology. Under each allowed blowup the pole coefficient is multiplicity minus one; collapsing the inserted rational tree gives the explicit fundamental-group isomorphism. Every initially nontrivial reduced nodal case remains reduced nodal and gives the NC conclusion by P4.20. In the other cases both normalization and conductor are simply connected, so the conductor pushout identifies pi1(S) with the curve graph free group. Virtual abelianness bounds its rank by one and gives positive curve Euler number; the open-support bound completes the positive sign}{\ref{night:P4.12}; \ref{night:P4.13}; \ref{night:P4.20}; \ref{night:P4.6}; \ref{night:P4.8}; \ref{night:P4R.1}; \ref{night:P4R.3}}

In the standing torus-fibration setting of this section, an irreducible singular support $S$
with smooth rational normalization $T$ has $e(S)>0$.
\end{proposition}
\begin{proof}
\noindent\textbf{Step 1.} Let $D$ be the conductor divisor and
$C=\operatorname{Sing}(S)_{\mathrm{red}}$. By P4.12,
$K_T+D\equiv0$. Contract $T$ to a minimal rational surface
$T_0$ and push $D$ forward to $D_0$. At every reversed
blowup the exceptional coefficient is $\mu-1$, where
$\mu$ is the multiplicity of the current conductor at the
centre. In particular every centre lies on the conductor.

\noindent\textbf{Step 2.} We give the complete minimal-surface case split. On
$\mathbb P^2$, the divisor $D_0$ is a plane cubic. An
irreducible cubic is smooth elliptic, rational nodal, or
rational cuspidal. A reducible cubic is a smooth conic
and a line, three lines, a double line and another line,
or a triple line. The conic and line meet at two
transverse points or at one tangency; three distinct
lines form a triangle or are concurrent. Thus either
$D_0$ is reduced nodal and has fundamental group
$\mathbb Z$ or $\mathbb Z^2$, or its reduced support is
simply connected and all normalized components are
rational.

\noindent\textbf{Step 3.} On a Hirzebruch surface, write $F$ for a ruling fibre.
Since $D_0\cdot F=2$, its horizontal part is a double
section, two distinct reduced sections, or a reduced
irreducible bisection. A double section together with
vertical fibres is a rational tree. For the other cases
write
\[
 D_0=H_1+H_2+\sum_j a_jF_j
 \quad\hbox{or}\quad
 D_0=H+\sum_j a_jF_j.
\]
In the first case adjunction on either section gives
\[
 H_1\cdot H_2+\sum_j a_j=2.
\]
The support is therefore a rational tree or a single
rational cycle. A cycle occurs only with all coefficients
one and ordinary double intersection points. This
includes the cases of two section-intersection points,
one section-intersection point and one vertical fibre,
or two vertical fibres. Coincident intersections or a
vertical fibre through the section-intersection point
instead give trees.

\noindent\textbf{Step 4.} For a bisection, adjunction gives
\[
 p_a(H)=1-\sum_j a_j.
\]
There is either no vertical component and $H$ has
arithmetic genus one, or one coefficient-one vertical
fibre and a smooth rational bisection. In the former
case $H$ is smooth elliptic, rational nodal, or rational
cuspidal. In the latter it meets the fibre at two distinct
points or at one tangency. Again the only nontrivial
fundamental groups are $\mathbb Z$ and $\mathbb Z^2$,
and those cases have reduced nodal conductor.

\noindent\textbf{Step 5.} We follow these alternatives through the blowups. At a
point of conductor multiplicity at least two, the new
conductor contains the exceptional $\mathbb P^1$.
Its contraction is the quotient which collapses that
connected simply connected subcomplex to the old
point. Van Kampen therefore makes the induced map on
fundamental groups an isomorphism. At a point of
multiplicity one the conductor is smooth and reduced,
and its strict transform is isomorphic to it. Thus
$\pi_1(D_{\mathrm{red}})$ is unchanged throughout.
If the initial group was nontrivial, the conductor was
reduced nodal. A smooth-point blowup adds no conductor
component, and a node blowup inserts one component of
coefficient one. It remains reduced nodal.

\noindent\textbf{Step 6.} In this latter case the generic transverse conductor
exponent of every normalization branch is one. The
exponent is its intrinsic conductor exponent plus the
intersection multiplicities with the other branches.
A singular branch has intrinsic exponent at least two.
Exponent one therefore means two smooth branches
meeting transversely. Consequently $S$ has ordinary
double curves outside finitely many points. P4.20 makes
$S$ normal crossing everywhere, and P4.6--P4.8 give
$e(S)=2$ for this rational normalization.

\noindent\textbf{Step 7.} In the remaining case $D_{\mathrm{red}}$ is connected
and simply connected, and all its normalized components
are rational. The conductor pushout verified in P4R.3
gives
\[
 \pi_1(S)=\pi_1(T)*_{\pi_1(D_{\mathrm{red}})}\pi_1(C)
          =\pi_1(C).
\]
Every component of $C$ has rational normalization,
because it is the finite nonconstant image of a
rational component of $D_{\mathrm{red}}$. P4R.1 makes
$\pi_1(C)$ virtually abelian. For a connected curve
with rational normalized components, this is the free
group $F_b$ of its point-identification graph, so
$b\le1$. If $c\ge1$ is its number of components, then
\[
 e(C)=c+1-b\ge c\ge1.
\]
P4.13 gives $e(S_{\mathrm{reg}})\ge0$. Additivity proves
$e(S)=e(S_{\mathrm{reg}})+e(C)>0$.
\end{proof}

\begin{proposition}[P4.22]\label{night:P4.22}\label{stmt:P4.22}
\StatementRecord{night:P4.22}{proved}{Historical events: GPT-6 Astra (OpenAI), ChatGPT Pro mode, effort Pro, S6\_06\_BLIND, 16--18 November 2026: targeted proof check. Also checked in this writing. Pre-continuation writing events: GPT-6 Astra (OpenAI), ChatGPT Pro mode, effort Pro, S6\_10\_PAPERS, 18 November 2026: targeted proof check. Continuation metadata events: GPT-6 Astra (OpenAI), ChatGPT Pro mode, effort Pro, S6\_10\_PAPERS continuation, 19 November 2026: used and rechecked in this run; status-vocabulary reconciliation; metadata reconciliation only; no new mathematical check. \emph{Pre-continuation recorded proof scope.} The night-run re-derivation was by GPT-6 Astra (OpenAI, ChatGPT Pro mode, effort Pro), 16--18 November 2026. Also checked in this writing by GPT-6 Astra (OpenAI), ChatGPT Pro mode, effort Pro, 18 November 2026. Adjunction and Riemann-Hurwitz force coefficient-one horizontal curves on a minimal elliptic ruled surface to be smooth elliptic, unramified and disjoint. Their subsequent blowups have zero exceptional pole coefficient, so P4.15 forces the minimal zero case. In the doubled-section alternative the actual conductor inclusion induces the normalization fundamental-group isomorphism, and the pushout identifies the support group with that of the conductor image. The graph Euler formula gives nonnegativity. Equality leaves a single elliptic image; finiteness and Riemann-Hurwitz exclude rational poles and therefore all blowups. The transverse conductor order two has exactly the cusp, tangent-double, and ordinary-triple possibilities, each giving the claimed zero case through its actual normalization map}{\ref{night:P4.13}; \ref{night:P4.15}; \ref{night:P4R.1}; \ref{night:P4R.3}}

In the standing torus-fibration setting of this section, an irreducible singular support
$S$ whose smooth normalization $T$ is birationally ruled
over an elliptic curve has $e(S)\ge0$. Equality holds
exactly when $T$ is minimal.
\end{proposition}
\begin{proof}
\noindent\textbf{Step 1.} Push the effective conductor to the minimal ruled
surface $T_0\to B$, where $B$ is elliptic. We have
$K_{T_0}+D_0\equiv0$. If a horizontal component $H$
has coefficient one, adjunction gives $p_a(H)\le1$,
whereas its nonconstant map to $B$ gives geometric
genus at least one. Thus $H$ is smooth elliptic and
$(D_0-H)\cdot H=0$. There is no vertical component.
The horizontal part is one elliptic bisection or two
disjoint elliptic sections, all with coefficient one.
All subsequent blowups have smooth coefficient-one
conductor centres; their exceptional curves are absent
from the conductor. The normalization differential
has no divisorial rank drop. P4.15 makes $T$ minimal,
makes the differential immersive everywhere, and
gives $e(S)=0$.

\noindent\textbf{Step 2.} The other possibility is
\[
 D_0=2H+\sum_j a_jF_j,
\]
where $H$ is a smooth elliptic section. As in P4R.3,
$D_{\mathrm{red}}$ is $H$ with rational trees attached,
and its inclusion in $T$ induces an isomorphism on
fundamental groups. The conductor pushout and P4R.1
make $\pi_1(C)$ virtually abelian. The reduced-curve
calculation in P4R.3 gives $e(C)\ge0$, and P4.13 then
gives $e(S)\ge0$.

\noindent\textbf{Step 3.} If $e(S)=0$, the curve calculation forces $C$ to have
one component, genus-one smooth normalization, and
no point-identification graph cycles. Each conductor
component maps finitely and nonconstantly onto $C$.
Riemann--Hurwitz excludes rational conductor
components. Thus $D_{\mathrm{red}}=H$ and no blowup
occurred: the first blowup on $2H$ would create a
vertical conductor component of coefficient one.
Therefore $T$ is minimal.

\noindent\textbf{Step 4.} Conversely suppose $T$ is minimal. In the double-section
case let $v$ be the number of distinct vertical fibres
in the conductor. Since $e(T)=e(H)=0$ and each vertical
$\mathbb P^1$ meets $H$ once,
\[
 e(S_{\mathrm{reg}})=e(T\setminus D_{\mathrm{red}})=-v.
\]
P4.13 forces $v=0$, hence $D=2H$. At a general point
of its image curve, every normalization branch has
conductor exponent two. The branch exponent equals
its intrinsic conductor exponent plus the sum of its
intersections with the other branches. A singular
branch together with another branch would have
exponent greater than two. Thus either there is one
ordinary cusp branch, or all branches are smooth:
two with intersection multiplicity two, or three
with pairwise intersection multiplicity one. Four
smooth branches would have exponent at least three,
so the list is exhaustive.

\noindent\textbf{Step 5.} In the cusp case $H\to C$ has generic degree one.
The normalization of $C$ is therefore $H$, an elliptic
curve. Virtual abelianness of $\pi_1(C)$ excludes any
additional point-identification graph cycle, so
$e(C)=0$. Additivity gives $e(S)=0$. In the other two
cases the normalization differential is immersive at
the generic point of $H$. There is no other conductor
component, so its rank-drop set is finite. P4.15
applies and again gives $e(S)=0$.
\end{proof}


\subsection{Class-VII minimal models with positive second Betti number}

\begin{proposition}[P4.23]\label{night:P4.23}\label{stmt:P4.23}
\StatementRecord{night:P4.23}{proved}{Historical events: GPT-6 Astra (OpenAI), ChatGPT Pro mode, effort Pro, S6\_06\_BLIND, 16--18 November 2026: targeted proof check. Also checked in this writing. Pre-continuation writing events: GPT-6 Astra (OpenAI), ChatGPT Pro mode, effort Pro, S6\_10\_PAPERS, 18 November 2026: targeted proof check. Continuation metadata events: GPT-6 Astra (OpenAI), ChatGPT Pro mode, effort Pro, S6\_10\_PAPERS continuation, 19 November 2026: used and rechecked in this run; status-vocabulary reconciliation; metadata reconciliation only; no new mathematical check. \emph{Pre-continuation recorded proof scope.} The night-run re-derivation was by GPT-6 Astra (OpenAI, ChatGPT Pro mode, effort Pro), 16--18 November 2026. Also checked in this writing by GPT-6 Astra (OpenAI), ChatGPT Pro mode, effort Pro, 18 November 2026. The effective numerically anticanonical divisor supplies the exact hypothesis of the opened Dloussky theorem. On an Inoue-Hirzebruch surface the negative-definite cycle matrix forces the reduced divisor; on a parabolic Inoue surface the numerically trivial cycle coefficient is one because an effective nonzero divisor cannot be torsion. Their reduced nodal transforms give the NC Chern contradiction at positive b2. In the remaining cycle-with-trees case the full conductor pushout is <t,c\_1,...,c\_b | t\textasciicircum{}k=w>. The injective cyclic-amalgam case embeds F\_b; the zero-word cases are checked as the corresponding actual free products. When k=0 and w is nontrivial, any nontrivial quotient factor produces a free subgroup on t and g t g\textasciicircum{}-1, so virtual abelianness forces b<=1. The exact curve Euler count and open-support bound then give strict positivity}{\ref{night:P4.12}; \ref{night:P4.13}; \ref{night:P4.20}; \ref{night:P4.21}; \ref{night:P4.6}; \ref{night:P4R.1}}

In the standing torus-fibration setting of this section, let an irreducible singular support
$S$ have smooth normalization $T$. If its minimal model $T_0$
belongs to class VII and $b_2(T_0)>0$, then $e(S)>0$.
\end{proposition}
\begin{proof}
\noindent\textbf{Step 1.} Let $D$ be the conductor. P4.12 makes $K_T+D$ torsion;
pushing down the point blowups gives $K_{T_0}+D_0$ torsion,
with $D_0$ effective and nonzero. This is a numerically
anticanonical divisor in the sense of Dloussky,
Definition~2.20 of \cite{Dloussky2006}. His Theorem~0.2
therefore gives a global spherical shell. Proposition~2.27
and the square-zero-divisor paragraph after Lemma~2.22
leave the parabolic Inoue or Inoue--Hirzebruch cases,
or a connected rational cycle with attached rational
trees. The latter configuration and $\pi_1(T_0)=\mathbb Z$
are stated explicitly in \cite[\S1, pp.~284--285]{DOT2003}.
We verify the consequences needed
for our conductor.

\noindent\textbf{Step 2.} For an Inoue--Hirzebruch surface, the reduced cycle or
union of two cycles is numerically anticanonical. The
curve intersection matrix is negative definite. Subtracting
that divisor from $D_0$ and intersecting with every curve
therefore makes every coefficient difference zero.
Thus $D_0$ is reduced nodal. In the parabolic Inoue case
there is an elliptic curve $E$ and a disjoint rational
cycle $\Gamma$, with
\[
 K_{T_0}+E+\Gamma\sim0,\qquad \Gamma\equiv0.
\]
These relations are stated after Dloussky's Lemma~2.22.
Intersections with $E$ and the rational curves give
$D_0=E+k\Gamma$: the cycle matrix has kernel generated
by the sum of its components. Torsion of $K_{T_0}+D_0$
then makes $\mathcal O((k-1)\Gamma)$ torsion. A nonzero
effective divisor cannot have torsion associated line
bundle on a compact connected reduced complex space:
a power of its defining section would be a nonzero
constant and could have no zero. Hence $k=1$ and
$D_0$ is again reduced nodal.

\noindent\textbf{Step 3.} The blowup coefficient $\mu-1$ preserves the reduced
nodal property, as in P4.21. Thus $S$ has ordinary
double curves outside finitely many points. P4.20
makes it normal crossing everywhere. Its normalization
is then immersive and P4.6 gives $e(T)=K_T^2$.
For a class-VII surface, however,
\[
 e(T)=b_2(T)>0,\qquad K_T^2=-b_2(T),
\]
a contradiction. These two exceptional minimal-model
cases are excluded.

\noindent\textbf{Step 4.} In the remaining case the support of the conductor is
a connected rational cycle with attached rational trees.
It is the whole connected curve configuration: Dloussky's
Lemma~2.22(2) puts every curve meeting the anticanonical
support in that support, and Proposition~2.27 gives its
connectedness. Its fundamental group is $\mathbb Z$.
The blowups preserve this group by the quotient argument
of P4.21, and preserve the surface fundamental group.
Consequently
\[
 \pi_1(D_{\mathrm{red}})=\mathbb Z,\qquad
 \pi_1(T)=\mathbb Z.
\]
Every normalized component of
$C=\operatorname{Sing}(S)_{\mathrm{red}}$ is rational,
being the finite nonconstant image of a rational conductor
component. The curve is connected. Write $\pi_1(C)=F_b$
for the free group of its point-identification graph.
The conductor pushout gives
\[
 \pi_1(S)=
 \langle t,c_1,\ldots,c_b\mid t^k=w(c_1,\ldots,c_b)\rangle.
\]
Here $k$ records the conductor-cycle map to $\pi_1(T)$,
and $w$ records its map to $\pi_1(C)$. Neither map is
assumed to be an isomorphism.

\noindent\textbf{Step 5.} Virtual abelianness from P4R.1 forces $b\le1$.
If $k\ne0$ and $w\ne1$, both cyclic amalgam maps are
injective; alternating nonidentity coset representatives
give the reduced normal form and show that $F_b$ injects
in the amalgam. If $w=1$, the group is
$(\mathbb Z/k\mathbb Z)*F_b$, with
$\mathbb Z/0\mathbb Z=\mathbb Z$, and $F_b$ again
injects. A nonabelian free subgroup is impossible in a
virtually abelian group. Finally, if $k=0$ and $w\ne1$,
the group is
\[
 \mathbb Z*\bigl(F_b/\langle\!\langle w\rangle\!\rangle\bigr).
\]
A free product of an infinite cyclic group and a nontrivial
group contains a nonabelian free subgroup: for a nontrivial
$g$ in the second factor, $t$ and $gtg^{-1}$ have no
nonempty reduced relation. Thus the second factor must
be trivial. Its abelianization has rank at least $b-1$,
which again gives $b\le1$.

\noindent\textbf{Step 6.} For the number $c\ge1$ of components of $C$, we obtain
\[
 e(C)=c+1-b\ge c\ge1.
\]
By P4.13, $e(S_{\mathrm{reg}})\ge0$, so
$e(S)=e(S_{\mathrm{reg}})+e(C)>0$.
\end{proof}

\subsection{Minimality of a smooth Hopf normalization}

\begin{proposition}[P4.24]\label{night:P4.24}\label{stmt:P4.24}
\StatementRecord{night:P4.24}{proved}{Historical events: GPT-6 Astra (OpenAI), ChatGPT Pro mode, effort Pro, S6\_06\_BLIND, 16--18 November 2026: targeted proof check. Also checked in this writing. Pre-continuation writing events: GPT-6 Astra (OpenAI), ChatGPT Pro mode, effort Pro, S6\_10\_PAPERS, 18 November 2026: targeted proof check; used and rechecked in this run. Continuation metadata events: GPT-6 Astra (OpenAI), ChatGPT Pro mode, effort Pro, S6\_10\_PAPERS continuation, 19 November 2026: used and rechecked in this run; status-vocabulary reconciliation; metadata reconciliation only; no new mathematical check. \emph{Pre-continuation recorded proof scope.} The night-run re-derivation was by GPT-6 Astra (OpenAI, ChatGPT Pro mode, effort Pro), 16--18 November 2026. Also checked in this writing by GPT-6 Astra (OpenAI), ChatGPT Pro mode, effort Pro, 18 November 2026. The primary Hopf normal form, compact-curve monomial calculation, nowhere-zero foliation, blowup saddle charts, saturated derivative Chern identity, integral quadratic bound, and actual smooth elliptic conductor maps were re-derived. Khaled Theorem 2.1(1) and formula (1) were opened in the primary preprint}{the hypotheses and explicit local calculation stated here}

In the standing torus-fibration setting of this section, let $S$ be an irreducible
singular fibre support with smooth normalization $T$.
If the minimal model of $T$ is a Hopf surface, then
$T$ is already minimal. The reduced conductor on $T$
and its reduced image in $S$ are disjoint unions of
smooth elliptic curves. Each map between their
components is a finite unramified covering, and
$e(S)=0$.
\end{proposition}
\begin{proof}
\noindent\textbf{Step 1.} Write $\nu:T\to S\hookrightarrow Y$ for the finite
holomorphic normalization map followed by the closed
embedding. We use the effective integral conductor
divisor $D$ with $K_T+D$ torsion. Let $T\to T_0$
be the sequence of $k$ point blowups giving the
minimal Hopf model. At each step the conductor
coefficient on the exceptional curve is $\mu-1$,
where $\mu$ is the multiplicity of the conductor
at the centre. Effectiveness therefore puts every
centre on the current conductor.

\noindent\textbf{Step 2.} There is a finite unramified primary Hopf cover
$T'_0\to T_0$. Pulling back every blowup gives a
finite unramified cover $T'\to T$ and a blowup
presentation of $T'$ over $T'_0$. A centre has
exactly the degree of the covering many lifts,
so the new number of blowups is that degree times
$k$. The composite $T'\to Y$ is finite and is
immersive away from the pulled-back conductor.
The rational Chern classes of its pulled-back
ambient tangent bundle vanish. Indeed $b_2(Y)=0$
gives $c_1(T_Y)=c_2(T_Y)=0$ in rational cohomology,
using Poincar\'e duality in degree four. The
argument below applies to this finite map, without
requiring it to be a normalization. Proving
minimality of $T'$ proves $k=0$. We therefore use
$T,T_0,\nu,k$ for the covering objects during the
Chern calculation.

\noindent\textbf{Step 3.} We first construct on $T_0$ a nonsingular
holomorphic foliation tangent to every compact
curve. Khaled, \emph{Real structures on primary
Hopf surfaces}, Theorem~2.1(1) and formula~(1),
gives
\[
 T_0=(\mathbb C^2\setminus\{0\})/\langle F\rangle,
 \qquad F(z,w)=(\alpha z+\lambda w^n,\beta w),
\]
where $0<|\alpha|\le|\beta|<1$ and
$\lambda(\alpha-\beta^n)=0$. We verify the needed
curve statement from these coordinates. Every
compact irreducible curve is smooth elliptic:
adjunction and $b_2(T_0)=0$ give arithmetic genus
one, and a rational normalization would lift to
$\mathbb C^2\setminus\{0\}$ and have constant
coordinate functions.

\noindent\textbf{Step 4.} The map from the fundamental group of such an
elliptic curve to $\pi_1(T_0)=\mathbb Z$ has
nonzero image, since zero image would give the
same impossible compact lift. Its full inverse
image has finitely many connected components.
The Remmert--Stein extension across the origin
is an analytic curve germ; a power $F^h$ fixes
one of its branches. On the normalization of
that branch the induced contraction can be
written $t\mapsto\eta t$, $0<|\eta|<1$.
For completeness, if the initial coordinate
has return map $g(t)=\eta t+O(t^2)$, choose
$|\eta|<\rho<\sqrt{|\eta|}$. On a sufficiently
small disc, $|g(t)|\le\rho|t|$, and the
successive differences in $g^j(t)/\eta^j$
are bounded by a constant times
$(\rho^2/|\eta|)^j|t|^2$. Their convergent sum
gives the asserted linearizing coordinate.

\noindent\textbf{Step 5.} When $\lambda=0$, comparison of coefficients in
$z(\eta t)=\alpha^h z(t)$ and
$w(\eta t)=\beta^h w(t)$ shows that each
nonzero coordinate is a single monomial.
Thus every branch away from the coordinate axes
has parametrization $(a t^p,b t^q)$. Its ratio
$p/q=\log|\alpha|/\log|\beta|$ is independent
of the branch. If such a branch exists, take
coprime positive integers $p,q$ in this ratio
and the vector field
$p z\partial_z+q w\partial_w$.
If no such branch exists, use
$z\partial_z+w\partial_w$.
Both are invariant under $F$, have no zero on
$\mathbb C^2\setminus\{0\}$, and are tangent
to every compact curve.

\noindent\textbf{Step 6.} When $\lambda\ne0$, we have $\alpha=\beta^n$ and
\[
 F^h(z,w)=
 (\alpha^h z+h\lambda\alpha^{h-1}w^n,\beta^h w).
\]
If a normalized invariant branch had
$w(t)=b t^q\ne0$, the coefficient of $t^{nq}$
in the first coordinate identity would give
$0=h\lambda\alpha^{h-1}b^n$, a contradiction.
Thus the only curve is the image of $w=0$.
The vector field $n z\partial_z+w\partial_w$
is invariant, nowhere zero, and tangent to it.
This proves the foliation assertion.

\noindent\textbf{Step 7.} Pull this foliation through the point blowups
and saturate its tangent sheaf. It is a line
bundle $L\subset T_T$. Every strict transform
of a conductor component and every exceptional
curve is invariant. The only singularities of
the foliation are simple saddles. To verify
this assertion and its multiplicities, at a
regular point use coordinates $(x,y)$ with
leaves $y=\text{constant}$. In the chart
$x=u,y=uv$ the saturated generator is
$u\partial_u-v\partial_v$; there is one
saddle, and the other blowup chart is regular
away from it. At an existing saddle use
$b x\partial_x-a y\partial_y$, with positive
integers $a,b$. The two charts have generators
\[
 b u\partial_u-(a+b)v\partial_v,\qquad
 (a+b)u\partial_u-a v\partial_v.
\]
One saddle is replaced by two. Each has
zero ideal $(u,v)$ and length one.
Consequently there are exactly $k$ saddles.

\noindent\textbf{Step 8.} Let $V$ be the divisor measuring the
divisorial zero of the pulled-back transverse
one-form before saturation. Numerically
\[
 c_1(L)=-K_T+V.
\]
The omitted pullback line bundle is topologically
trivial because it comes from $T_0$.
The preceding charts give, at every blowup $\mu$,
\[
 K'=\mu^*K+E,\qquad
 V'=\mu^*V+\varepsilon E,\qquad
 \varepsilon=
 \begin{cases}0&\text{at a regular point},\\
 1&\text{at a saddle}.
 \end{cases}
\]
Starting with $V=0$ and $K_{T_0}^2=0$, we obtain
\[
 K_T^2=-k,\qquad K_TV-V^2=0.
\]
The second identity is preserved because its
change is $\varepsilon^2-\varepsilon=0$.

\noindent\textbf{Step 9.} Put $B=\nu^*T_Y$ and compose the holomorphic
maps $L\to T_T\to B$. This map has no
divisorial zero. Away from the conductor this
follows from immersivity. At the generic point
of a conductor component the foliation is
tangent to that curve, and finiteness of
$\nu$ makes its differential along the curve
nonzero. A nonconstant holomorphic map of
curves has nonzero derivative at the generic
point. Thus the cokernel $Q=B/L$ is a
torsionfree sheaf of rank two. Its double
dual $Q^{**}$ is locally free, and
\[
 z_0=\operatorname{length}(Q^{**}/Q)
\]
is finite. At a point, write the map $L\to B$
as a column $(f_1,f_2,f_3)$ in trivializations.
If this column vanishes, the ideal
$(f_1,f_2,f_3)$ has finite colength. Dualizing
$0\to\mathcal O\to\mathcal O^3\to Q\to0$
gives $\mathcal E xt^1(Q,\mathcal O)
=\mathcal O/(f_1,f_2,f_3)$.
Dualizing $0\to Q\to Q^{**}\to Q^{**}/Q\to0$
and using the two-dimensional regular local
ring identifies this with
$\mathcal E xt^2(Q^{**}/Q,\mathcal O)$,
which has the same length. At each foliation
saddle all three coefficients vanish.
It follows that
\[
 z_0\ge k.
\]

\noindent\textbf{Step 10.} Let $E_B\subset B$ be the saturation of the
generic rank-two image of $d\nu:T_T\to B$.
Its determinant differs from that of $T_T$
by an effective integral divisor $R$, the
divisorial rank-drop divisor:
\[
 c_1(E_B)=-K_T+R.
\]
The support of $R$ is contained in the
conductor. Saturate $E_B/L$ in $Q^{**}$.
The resulting rank-one reflexive subsheaf
is a line bundle $J$, with
$c_1(J)=R-V$. Its torsionfree rank-one
quotient has the form
$\mathcal I_{Z_1}\otimes J_1$ with
$c_1(J_1)=K_T-R$ and
$z_1=\operatorname{length}(Z_1)\ge0$.
All these assertions concern sheaves on the
smooth surface $T$; determinants are computed
off finite sets and therefore determine their
line-bundle classes.

\noindent\textbf{Step 11.} The identities $c_1(B)=c_2(B)=0$ give
\[
 c_2(Q)=(-K_T+V)^2,\qquad
 c_2(Q^{**})=(-K_T+V)^2-z_0.
\]
The exact sequence with $J$ gives another
expression for the same number:
\[
 (-K_T+V)^2-z_0=(R-V)(K_T-R)+z_1.
\]
Using $K_T^2=-k$ and $K_TV-V^2=0$, this becomes
\[
 G(R):=R^2-K_TR-VR=k+z_0+z_1\ge2k.
\]

\noindent\textbf{Step 12.} We compare this expression with the actual
blowup sequence. Push $R$ and $V$ down at
each step; $R$ has integral coefficients.
Writing $R'=\mu^*R+cE$ gives $c\in\mathbb Z$,
without any positivity requirement on $c$.
The intersection calculation is
\[
 G(R')-G(R)=-c^2+(1+\varepsilon)c.
\]
For $\varepsilon=0$ this is at most zero
for every integer $c$; for $\varepsilon=1$
it is $1-(c-1)^2\le1$. On $T_0$ every
divisor is numerically zero, so $G(R_0)=0$.
Consequently
\[
 G(R)\le\#\{\text{blowups at saddles}\}\le k.
\]
The lower bound $G(R)\ge2k$ forces
$k=z_0=z_1=0$. Returning through the finite
unramified cover proves that the original
normalization $T$ is minimal.

\noindent\textbf{Step 13.} It remains to identify the conductor maps,
rather than only their Euler number.
The reduced conductor on a minimal Hopf
surface consists of disjoint smooth elliptic
curves. Since $z_0=0$ on the primary cover,
the differential of $\nu$ along each such
curve is nowhere zero, also before that
cover. Thus every branch of its image is
smooth. At an image point, two different
curve germs would have to come from two
different normalization points. Their
surface branches in the smooth threefold
intersect in a curve germ. This germ is
contained in the conductor image of both
branches; the unique conductor germ at
each normalization point forces the two
image germs to be equal. Therefore the
image conductor is locally a single smooth
curve. Its components are disjoint, and
each conductor map is finite unramified.
Riemann--Hurwitz makes every image component
elliptic. Finally, normalization is an
isomorphism off these curves, so
\[
 e(S)=e(T)-e(D_{\mathrm{red}})
                 +e(\operatorname{Sing}S)=0.
\]
\end{proof}


\subsection{The reduced resolved-conductor case}

\begin{proposition}[P4.25]\label{night:P4.25}\label{stmt:P4.25}
\StatementRecord{night:P4.25}{proved}{Historical events: GPT-6 Astra (OpenAI), ChatGPT Pro mode, effort Pro, S6\_06\_BLIND, 16--18 November 2026: targeted proof check. Also checked in this writing. Pre-continuation writing events: GPT-6 Astra (OpenAI), ChatGPT Pro mode, effort Pro, S6\_10\_PAPERS, 18 November 2026: targeted proof check; used and rechecked in this run. Continuation metadata events: GPT-6 Astra (OpenAI), ChatGPT Pro mode, effort Pro, S6\_10\_PAPERS continuation, 19 November 2026: used and rechecked in this run; status-vocabulary reconciliation; metadata reconciliation only; no new mathematical check. \emph{Pre-continuation recorded proof scope.} The night-run re-derivation was by GPT-6 Astra (OpenAI, ChatGPT Pro mode, effort Pro), 16--18 November 2026. Also checked in this writing by GPT-6 Astra (OpenAI), ChatGPT Pro mode, effort Pro, 18 November 2026. The actual normalized boundary gives the log-canonical resolution inequalities. Liu-Rollenske2012's main theorem, Table1, Definition2.3, and Propositions3.1/4.1 were opened at their exact author-preprint loci. Every listed non-NC germ was deformed with its full multiplicity and boundary Hessian retained, reducing to the proved crosscap or tangent-sheet obstruction. The final global applications retain P4.6/P4.8 hypotheses}{\ref{night:P4.16}; \ref{night:P4.20}; \ref{night:P4.6}; \ref{night:P4.8}}

Every semi-log-canonical reduced fibre support in
Part~4.2 has normal crossings. In particular, let
$\rho:T\to N\to S$ be a resolution of the normalization
of an irreducible support, and let $D$ be the integral
divisor defined by
$\omega_T(D)\simeq\rho^*\omega_S$. When this actual
divisor is reduced and has normal crossings, $S$
has normal crossings, its normalization is smooth,
and P4.6--P4.8 apply.
\end{proposition}
\begin{proof}
\noindent\textbf{Step 1.} We first explain why the stated resolved-conductor
case is semi-log-canonical. At a generic conductor
point, an exponent one on every normalization
branch forces the transverse plane curve to be an
ordinary node. Indeed the conductor exponent on
branch $i$ is
$2\delta_i+\sum_{j\ne i}I_{ij}$; positivity of
intersection multiplicities leaves exactly two
smooth transverse branches. The support $S$ is a
Cartier hypersurface, hence is Cohen--Macaulay and
has invertible dualizing sheaf. On its normal
normalization $N$, the pair consisting of $N$
and its reduced conductor is log canonical.
To check this last assertion, pull its canonical
divisor to $T$. The resulting boundary is exactly
$D$, by the displayed isomorphism and the
divisorial orders of a local generator. A smooth
surface with reduced normal-crossing boundary is
log canonical: blowing up a point gives new
boundary coefficient $\mu-1\le1$, and this
bound persists under successive point blowups.
Discrepancy comparison therefore gives the
assertion on $N$. This is precisely the
semi-log-canonical definition of
Liu--Rollenske, Definition~2.3 and Lemma~2.2,
which explicitly applies to complex analytic
spaces.

\noindent\textbf{Step 2.} We now treat all local semi-log-canonical
hypersurface germs. We use the complete analytic
classification in Liu--Rollenske,
\emph{Two-dimensional semi-log-canonical
hypersurfaces}, main Theorem, Table~1,
Propositions~3.1 and~4.1. The normal singular
germs are isolated. Any such germ is excluded
in our family: a small generic linear
perturbation of its reduced equation has
Morse critical points, the number counted
with multiplicity being the positive
colength of its gradient ideal. The equality
of this number before and after perturbation
follows by conservation of the local
intersection multiplicity of the three
gradient equations. A constant perturbation
puts their critical values away from zero.
The original zero divisor is smooth near
a small transverse boundary sphere, so it
retains that boundary property under these
perturbations. The proper boundary
transport proved in P4.20 carries a
vanishing sphere of square $-2$ into a
nearby torus fibre, which is impossible.

\noindent\textbf{Step 3.} For the nonnormal germs we give the exact
deformations and check their boundary strata.
The ordinary double point and $xyz=0$
already have normal crossings. A pinch point
$x^2+y^2z=0$ is excluded by P4.16.
The remaining double-point equations are
\[
 x^2+y^2z^2,\qquad
 x^2+y^2(z^2+y^{q-2}),\quad q\ge3.
\]
For the first equation the two normalization
branches are the smooth graphs $x=\pm i yz$.
Their tangent planes coincide at the origin,
whereas they meet transversely off that point
along the two axes. The local proof of P4.20
therefore supplies a vanishing sphere after
a deformation retaining its boundary
ordinary-double strata.

\noindent\textbf{Step 4.} For finite $q$ put
\[
 G_s=x^2+y^2(z^2+y^{q-2}+s).
\]
For $s\ne0$ small, at each point
$(0,0,z_0)$ with $z_0^2=-s$ the function
$Z=z^2+y^{q-2}+s$ is a local coordinate:
its $z$ derivative is $2z_0\ne0$.
The actual equation there is exactly
$x^2+y^2Z$, a pinch point. At the original
boundary double axis $x=y=0$, $z\ne0$,
the transverse Hessian determinant is
$4(z^2+s)$, which stays nonzero on the
compact boundary when $s$ is small.
All boundary smooth points remain smooth.
Thus the zero divisor retains the same
boundary smooth and double strata.
The spherical class of square $-4$
from the pinch-point nearby fibre transports,
by the detailed P4.20 construction, to
the original nearby torus fibre.
This excludes every finite $q$.

\noindent\textbf{Step 5.} The remaining triple-point equations are
\[
 xyz+x^p,\quad p\ge3,\qquad
 xyz+x^p+y^q,\quad q\ge p\ge3.
\]
For the first put
$G_s=xyz+x^p+s x^2
=x(yz+x^{p-1}+s x)$.
When $s\ne0$, its two local normalization
branches are smooth: the derivative in $x$
of the second factor at the origin is $s$.
Their tangent planes coincide, so the
origin is not a normal crossing. Along
the boundary $y$ axis the transverse
Hessian determinant is $-y^2$, and along
the boundary $z$ axis it is $-z^2$.
Both axes persist as ordinary-double
strata. The local argument of P4.20 on
the deformed germ supplies a compact
spherical class of nonzero square, and
the same boundary transport gives the
contradiction.

\noindent\textbf{Step 6.} For the second equation put
\[
 G_s=xyz+x^p+y^q+s x^2+s y^3,
 \qquad 0<|s|\ll1.
\]
Its transverse Hessian along the boundary
$z$ axis has determinant $-z^2$,
so its sole boundary double stratum
persists. We check the smooth normalization
of its germ at the origin explicitly.
Since $\partial^2G_s/\partial x^2(0)=2s$,
the critical equation in $x$ has a unique
analytic solution $x=\xi(y,z)$ with
$\xi$ divisible by $yz$. The parametric
holomorphic Morse change in $x$ writes
\[
 G_s=X^2+y^2H(y,z).
\]
Divisibility follows directly on substituting
$\xi$: each of $\xi yz,\xi^p,s\xi^2$
is divisible by $y^2$, as are the two
pure $y$ terms. Moreover
\[
 H_y(0,0)=
 \begin{cases}1+s&q=3,\\s&q>3,\end{cases}
 \qquad H(0,z)=-z^2/(4s).
\]
The finite
birational map
\[
 \{v^2+H(y,z)=0\}\longrightarrow\{G_s=0\},
 \qquad (v,y,z)\longmapsto(X=yv,y,z),
\]
followed by the inverse Morse coordinate
change, is the normalization. Its source
is smooth because $H_y(0,0)\ne0$,
and the map is finite because $v$
satisfies the displayed monic equation.
It is birational on $y\ne0$ and is
therefore the normalization by normality
of its source. At the origin its
differential has rank one, so this germ
is not a normal crossing. P4.20 again
produces a spherical class, with the
already verified boundary transport.

\noindent\textbf{Step 7.} These calculations exhaust the analytic
classification. Thus only the smooth,
ordinary-double, and triple-normal-crossing
germs remain. Their local equations give
the claimed normal-crossing support and
smooth normalization.
\end{proof}


\subsection{Numerical sign for arbitrary irreducible supports}

\begin{proposition}[P4.26]\label{night:P4.26}\label{stmt:P4.26}
\StatementRecord{night:P4.26}{proved}{Historical events: GPT-6 Astra (OpenAI), ChatGPT Pro mode, effort Pro, S6\_06\_BLIND, 16--18 November 2026: targeted proof check. Also checked in this writing. Pre-continuation writing events: GPT-6 Astra (OpenAI), ChatGPT Pro mode, effort Pro, S6\_10\_PAPERS, 18 November 2026: targeted proof check. Continuation metadata events: GPT-6 Astra (OpenAI), ChatGPT Pro mode, effort Pro, S6\_10\_PAPERS continuation, 19 November 2026: used and rechecked in this run; status-vocabulary reconciliation; metadata reconciliation only; no new mathematical check. \emph{Pre-continuation recorded proof scope.} The night-run re-derivation was by GPT-6 Astra (OpenAI, ChatGPT Pro mode, effort Pro), 16--18 November 2026. Also checked in this writing by GPT-6 Astra (OpenAI), ChatGPT Pro mode, effort Pro, 18 November 2026. The rational pole-support transfers, conductor pushout and group presentations, class-VII alternatives, local ruled and Hopf comparisons, normal cover case, Euler sum and parity were re-derived with the source's exact separation between the single-support conclusion and the global all-support hypothesis}{\ref{night:P4.12}; \ref{night:P4.13}; \ref{night:P4.21}; \ref{night:P4.23}; \ref{night:P4.25}; \ref{night:P4.6}; \ref{night:P4.7}; \ref{night:P4.8}; \ref{night:P4R.1}; \ref{night:P4R.6}; \ref{night:P4R.7}; \ref{night:P4R.8}; \ref{night:P4R.9}}

In Part~4.2 every irreducible reduced fibre support
has nonnegative topological Euler number. For a
singular irreducible support, rational resolution
of the normalization gives strictly positive
Euler number; so does a resolution whose minimal
model is of class VII with positive second
Betti number. These assertions permit singular
normalization.

Every normal-crossing support has the
nonnegative even Euler number of P4.6.
Consequently, for a fibration all of whose
reduced supports are irreducible or have
normal crossings,
\[
 B=\mathbb P^1,\qquad c_3(Y)\in\{0,2\},
 \qquad b_3(Y)=2-c_3(Y).
\]
This last conclusion does not assume a
classification of the equality cases for
irreducible nonnormal supports.
\end{proposition}
\begin{proof}
\noindent\textbf{Step 1.} For an irreducible singular support use the
minimal resolution of its normalization
$\rho:T\to N\to S$, and put
$A=\rho^{-1}(\operatorname{Sing}S)_{\rm red}$,
$C=\operatorname{Sing}(S)_{\rm red}$.
P4R.6 gives an effective integral divisor
$D$ with support $A$ and $K_T+D$ torsion.
The resolution $T$ has Kodaira dimension
$-\infty$, by P4.12. We explain the two
remaining transfers from smooth
normalization to this resolution.

\noindent\textbf{Step 2.} Suppose first that $T$ is rational.
Push $D$ to a minimal rational surface.
The complete anticanonical-support calculation
in P4.21 has the following exact alternatives.
The reduced support is connected; either
its fundamental group is trivial and all
its normalized components are rational,
or the divisor itself is reduced with
normal crossings. On $\mathbb P^2$ these
are the plane cubic configurations
enumerated there. On a ruled minimal
surface they are the degree-two horizontal
configurations there. Those arguments use
only the effective integral divisor and
$K+D\equiv0$, so apply to the present
resolved divisor.

\noindent\textbf{Step 3.} We check their preservation through the
actual blowup morphisms. A blowup at
conductor multiplicity $\mu$ gives
exceptional coefficient $\mu-1$.
For $\mu=1$, the centre is a smooth
point of the reduced support and its
strict transform is isomorphic to that
support. For $\mu\ge2$, the whole
exceptional $\mathbb P^1$ belongs to
the new support. Its contraction is
the blowup map on supports. Collapsing
this connected simply connected
subcomplex does not change the
fundamental group, by van Kampen.
The reduced normal-crossing case stays
reduced normal crossing: a smooth
point gives coefficient zero and a
double point gives coefficient one.
Thus the same alternatives hold for
the actual divisor $D$ on $T$.

\noindent\textbf{Step 4.} In the reduced normal-crossing alternative,
P4.25 makes $S$ normal crossing. Its
normalization is smooth, so this minimal
normalization resolution is the
normalization itself. The rational case
of P4.6--P4.8 gives $e(S)=2$.
In the other alternative $\pi_1(A)=1$
and $\pi_1(T)=1$. The proper-map pushout
$T\cup_A C\to S$ is a homeomorphism:
it is a continuous compact-to-Hausdorff
bijection, since $\rho$ is an isomorphism
off $A$ and $A$ is the full inverse
image of $C$. Compatible triangulation
makes $A\to T$ a cofibration.
Van Kampen therefore gives
$\pi_1(S)\simeq\pi_1(C)$.
Every normalized component of $C$ is
rational: a nonconstant map from a
normalized component of $A$ onto it
is a finite map from $\mathbb P^1$.
Contracted components of $A$ contribute
no exception to this statement.
The curve $C$ is connected, and its
fundamental group is a free group
$F_b$ on the cycles of its
point-identification graph. P4R.1
implies $b\le1$. If $c\ge1$ is the
number of its components, then
\[
 e(C)=c+1-b\ge1.
\]
P4.13 gives $e(S_{\rm reg})\ge0$,
so $e(S)>0$.

\noindent\textbf{Step 5.} Now suppose that the minimal model
$T_0$ is of class VII with
$b_2(T_0)>0$. The divisor pushed to
$T_0$ satisfies the hypotheses of
Dloussky's numerically anticanonical
theorem exactly as in P4.23.
The verified sources and their
normalization-independent hypotheses
are specified there. The
Inoue--Hirzebruch and parabolic Inoue
alternatives give a reduced
normal-crossing divisor, preserved
under the blowups just checked.
P4.25 then makes $S$ normal crossing
and $T=N$ smooth. The immersive
normalization Chern formula would
give $e(T)=K_T^2$, whereas
$e(T)=b_2(T)>0$ and
$K_T^2=-b_2(T)$. These alternatives
cannot occur.

\noindent\textbf{Step 6.} In the intermediate alternative
$A$ is a connected rational cycle
with attached rational trees, and
\[
 \pi_1(A)=\mathbb Z,\qquad
 \pi_1(T)=\mathbb Z.
\]
Both statements survive the actual
blowups; they concern the resolution
and do not require $N$ to be smooth.
Every normalized component of $C$
is again rational and $C$ is connected.
With $\pi_1(C)=F_b$, the pushout
presentation is exactly
\[
 \pi_1(S)=
 \langle t,c_1,\ldots,c_b
 \mid t^k=w(c_1,\ldots,c_b)\rangle
\]
for the actual two conductor-group
homomorphisms. A component contracted
to a point merely induces the zero
homomorphism on its group. The
amalgam and free-product normal-form
argument in P4.23, which allowed
arbitrary $k$ and $w$, gives $b\le1$
from virtual abelianness.
Thus $e(C)=c+1-b\ge1$, and
$e(S)>0$ follows again from P4.13.

\noindent\textbf{Step 7.} This exhausts these two resolution
classes. Ruled resolutions of positive
base genus are covered by P4R.7
and P4R.9, and Hopf minimal models
by P4R.8. The class-VII minimal
model with $b_2=0$ is Hopf or Inoue,
by the verified classification in
P4.8. An Inoue model has no curves,
so it cannot carry the nonempty
effective divisor pushed down from
$D$. Therefore it does not occur
here. All irreducible singular
supports have now been covered.

\noindent\textbf{Step 8.} A normal irreducible support is
smooth by P4.12. The cyclic
base-change comparison and torus
deformation theorem in P4.8 show
that it has a finite unramified
torus cover, hence Euler number
zero. This proves the full
irreducible assertion.

\noindent\textbf{Step 9.} For the final global statement,
$B=\mathbb P^1$ follows from
$b_1(Y)=0$ and connected fibres,
as proved in P4.7. Euler
additivity over the finitely many
exceptional base values gives
\[
 e(Y)=\sum_{b\in B}
 e\bigl((f^{-1}(b))_{\rm red}\bigr),
\]
because a smooth torus has Euler
number zero and the complement
is a proper smooth fibre bundle.
Every summand is nonnegative
under the stated support
hypothesis. Poincar\'e duality
and $b_1=b_2=0$ give
$e(Y)=2-b_3(Y)$; the nondegenerate
skew-symmetric cup-product
pairing on $H^3(Y;\mathbb Q)$
makes $b_3(Y)$ even. Thus
$0\le e(Y)\le2$ and $e(Y)$ is
even, which proves the asserted
values and their Chern-number
interpretation.
\end{proof}


\subsection{Equality in the Hopf-resolution case}

\begin{proposition}[P4.27]\label{night:P4.27}\label{stmt:P4.27}
\StatementRecord{night:P4.27}{proved}{Historical events: GPT-6 Astra (OpenAI), ChatGPT Pro mode, effort Pro, S6\_06\_BLIND, 16--18 November 2026: targeted proof check. Also checked in this writing. Pre-continuation writing events: GPT-6 Astra (OpenAI), ChatGPT Pro mode, effort Pro, S6\_10\_PAPERS, 18 November 2026: targeted proof check. Continuation metadata events: GPT-6 Astra (OpenAI), ChatGPT Pro mode, effort Pro, S6\_10\_PAPERS continuation, 19 November 2026: used and rechecked in this run; status-vocabulary reconciliation; metadata reconciliation only; no new mathematical check. \emph{Pre-continuation recorded proof scope.} The night-run re-derivation was by GPT-6 Astra (OpenAI, ChatGPT Pro mode, effort Pro), 16--18 November 2026. Also checked in this writing by GPT-6 Astra (OpenAI), ChatGPT Pro mode, effort Pro, 18 November 2026. The exact Hopf pole graph has r elliptic roots and v rational vertices after k blowups, so equality forces k=v and zero conductor-image Euler number. In the all-elliptic image case every rational pole contracts and no elliptic root does. The resolution-to-minimal-model and resolution-to-normalization maps then have exactly the same connected fibres; normalizing their graph image gives proper finite birational projections to normal surfaces, hence a holomorphic isomorphism of the normalization with the minimal Hopf surface. In the sole remaining graph count, the homomorphism from the Hopf group to Z has finite kernel and each elliptic-root map has infinite image by the actual universal-cover lift argument. Quotienting each elliptic vertex by the saturated image of its rank-one kernel gives a surjection to a graph of cyclic groups. Removing its unique trivial-edge cycle leaves a tree amalgam containing Z; restoring it yields a free product with Z containing F2, contradicting virtual abelianness. Thus the equality normalization is smooth minimal Hopf}{\ref{night:P4.24}; \ref{night:P4R.1}; \ref{night:P4R.6}; \ref{night:P4R.8}}

Let \(S\) be an irreducible singular reduced fibre support in Part 4.2.
Suppose that the minimal resolution \(T\to N\) of its normalization has a
Hopf minimal model \(T_0\). Then
\[
 e(S)=0\quad\Longleftrightarrow\quad
 N\text{ is a smooth minimal Hopf surface}.
\]
In this case the reduced conductor and its image are disjoint unions of
smooth elliptic curves, and all component maps are finite unramified covers.
\end{proposition}
\begin{proof}
\noindent\textbf{Step 1.} Write \(\rho:T\to N\to S\), \(A=\rho^{-1}(\operatorname{Sing}S)_{\rm red}\),
and \(C=\operatorname{Sing}(S)_{\rm red}\). Let \(k\) be the number of
blowups in \(\pi:T\to T_0\). P4R.6 and P4R.8 give an effective integral pole
divisor whose reduced support \(A\) has \(r\) connected components. Each is
a smooth elliptic root with a rational tree attached. Let \(v\) be the
number of its rational components. Every such component is one of the
\(k\) blowup exceptional curves. We have
\[
 e(S)=k-v+e(C),\qquad k-v\ge0,\qquad e(C)\ge0.
\]
Thus equality forces \(k=v\) and \(e(C)=0\).

\noindent\textbf{Step 2.} Let \(c,s,g,b\) denote the numbers of irreducible components of \(C\),
connected components of \(C\), elliptic normalized components of \(C\), and
the sum of its point-identification graph cycle ranks. P4R.8 gives
\[
 e(C)=c+s-2g-b,\qquad
 b+r-s\le1,\qquad c\ge g,\quad r\ge g.
\]
Consequently \(e(C)=0\) implies \(c+r-2g\le1\).

\noindent\textbf{Step 3.} First suppose \(c=g\). Every normalized component of \(C\) is elliptic.
Every rational component of \(A\) is contracted by \(T\to N\), since a
finite nonconstant map from \(\mathbb P^1\) to an elliptic normalization
does not exist. A root cannot be contracted. Indeed all rational curves
attached to it are already contracted, so its whole rooted connected
component would be a connected exceptional block disjoint from the strict
conductor. P4R.6 excludes precisely such a block. Hence all roots are
noncontracted.

\noindent\textbf{Step 4.} The equality \(k=v\) now says that every blowup exceptional curve of
\(\pi:T\to T_0\) belongs to \(A\) and is contracted by \(T\to N\).
Conversely every exceptional curve of \(T\to N\) belongs to \(A\), by
P4R.6, and cannot be a root. The two proper birational maps therefore have
exactly the same exceptional curves. Their fibres are identical: each
nontrivial fibre is a connected component of the exceptional locus, and
their fibres are connected by Stein factorization and normality of the
targets.

\noindent\textbf{Step 5.} Here is the holomorphic comparison, rather than only a topological
identification. Take the image of
\[
 (\pi,\rho_N):T\longrightarrow T_0\times N
\]
and then its normalization \(G\). The two projections \(G\to T_0,N\)
are proper birational maps with zero-dimensional fibres, because the
original two maps have identical fibres. They are finite. A finite
birational morphism to a normal complex space is an isomorphism, as follows
on local rings from integral closedness. Thus \(N\simeq T_0\)
biholomorphically. It is a smooth minimal Hopf surface.

\noindent\textbf{Step 6.} The only other possibility is \(c=g+1,r=g\). There is one rational
normalized component of \(C\), and all \(g=r\) roots are noncontracted and
map to distinct elliptic normalized components of \(C\). Each such map is
an isogeny. Also
\[
 b+r-s=1.
\]
We exclude this possibility by a quotient of the exact fundamental-group
pushout; the quotient is constructed explicitly below.

\noindent\textbf{Step 7.} Represent \(C\) as its disjoint normalized smooth curves with a finite
graph of point identifications. This is a graph of spaces with elliptic
vertex groups \(\mathbb Z^2\), trivial rational vertex groups, and trivial
edge groups. Its graph has \(s\) connected components and total cycle rank
\(b\). Add a vertex for \(T\), and add one root edge for each connected
component of \(A\). The root edge group is \(\mathbb Z^2\), since its
rational tree is simply connected. Its map to the \(C\) side factors
through the elliptic normalization receiving that root; its map to the
\(T\) side is the actual inclusion. This graph-of-spaces presentation is
the ordinary van Kampen presentation of the homotopy pushout
\(T\cup_A C=S\). It has total graph cycle rank \(b+r-s=1\).

\noindent\textbf{Step 8.} Let \(\Gamma=\pi_1(T)=\pi_1(T_0)\). It is virtually infinite cyclic, and
\(b_1(T)=1\). The primitive map from its abelianization to its free part
therefore gives a surjection
\[
 h:\Gamma\longrightarrow\mathbb Z
\]
with finite kernel. Indeed a cyclic subgroup of finite index maps
nontrivially to \(\mathbb Z\), so the kernel meets that subgroup trivially
and is finite.

\noindent\textbf{Step 9.} For root \(H_i\), write
\[
 a_i:\pi_1(H_i)=\mathbb Z^2\longrightarrow\Gamma,\qquad
 q_i:\pi_1(H_i)=\mathbb Z^2\longrightarrow
       \pi_1(\widetilde C_i)=\mathbb Z^2.
\]
The second map is injective of finite index because it is induced by an
isogeny. The composite \(h a_i\) is nonzero. To see this, a finite image
of \(a_i\) would become trivial after a finite cover of \(H_i\); that
compact elliptic cover would lift holomorphically to the universal Hopf
cover \(\mathbb C^2\setminus\{0\}\), where its coordinate functions are
constant. Thus \(a_i\) has infinite image; since \(\ker h\) is finite,
\(h a_i\ne0\).

\noindent\textbf{Step 10.} Put \(K_i=\ker(h a_i)\), a primitive rank-one subgroup of
\(\mathbb Z^2\). Quotient the \(i\)-th elliptic vertex group by the
saturation of \(q_i(K_i)\). The resulting vertex group is \(\mathbb Z\).
Quotient the root edge group by \(K_i\); it too becomes \(\mathbb Z\).
Its two induced maps into the \(T\) vertex \(\mathbb Z\) and the elliptic
vertex \(\mathbb Z\) are nonzero multiplications and hence injective.
The latter assertion follows because an element outside \(K_i\) cannot
map under the injective rational map \(q_i\otimes\mathbb Q\) into the
rational span of \(q_i(K_i)\). This verifies compatibility of the chosen
quotients with both edge homomorphisms.

\noindent\textbf{Step 11.} We have constructed a surjective homomorphism from \(\pi_1(S)\) onto the
fundamental group \(G\) of a graph of groups. Its nontrivial vertex groups
are infinite cyclic, its root edges are infinite cyclic with injective
maps, and all point-identification edges are trivial. The root edges
form a star: the roots map to distinct elliptic normalizations. They
therefore contain no graph cycle. The unique graph cycle of the complete
graph contains a point-identification edge, whose edge group is trivial.

\noindent\textbf{Step 12.} Delete one such edge on the cycle. The remaining graph is a connected
tree, and the normal form for iterated amalgamated products along
injective edge maps embeds its \(T\) vertex group \(\mathbb Z\) into its
fundamental group \(G_0\). Restoring the removed edge is an HNN extension
over the trivial group, hence
\[
 G\simeq G_0*\mathbb Z.
\]
Choosing a nontrivial element of the embedded \(\mathbb Z\subset G_0\)
and the new stable letter gives an embedded free group \(F_2\), by reduced
free-product normal form. But every quotient of a virtually abelian group
is virtually abelian, and a virtually abelian group has no \(F_2\)
subgroup. This contradicts P4R.1. The case \(c=g+1,r=g\) is impossible.

\noindent\textbf{Step 13.} We have proved that equality makes \(N\) smooth minimal Hopf. Conversely
P4.24 applies to such an actual normalization and proves the unramified
elliptic-conductor description and \(e(S)=0\). This completes both
directions.
\end{proof}


\subsection{Cuspidal transverse type and the ruled equality case}

\begin{proposition}[P4.28]\label{night:P4.28}\label{stmt:P4.28}
\StatementRecord{night:P4.28}{proved}{Historical events: GPT-6 Astra (OpenAI), ChatGPT Pro mode, effort Pro, S6\_06\_BLIND, 16--18 November 2026: targeted proof check. Also checked in this writing. Pre-continuation writing events: GPT-6 Astra (OpenAI), ChatGPT Pro mode, effort Pro, S6\_10\_PAPERS, 18 November 2026: targeted proof check. Continuation metadata events: GPT-6 Astra (OpenAI), ChatGPT Pro mode, effort Pro, S6\_10\_PAPERS continuation, 19 November 2026: used and rechecked in this run; status-vocabulary reconciliation; metadata reconciliation only; no new mathematical check. \emph{Pre-continuation recorded proof scope.} The night-run re-derivation was by GPT-6 Astra (OpenAI, ChatGPT Pro mode, effort Pro), 16--18 November 2026. Also checked in this writing by GPT-6 Astra (OpenAI), ChatGPT Pro mode, effort Pro, 18 November 2026. The transverse Hessian along the smooth critical curve is treated over the one-variable DVR: a minimal diagonal entry divides the off-diagonal entry, so the square completion is holomorphic. The rank-one case reduces to X\textasciicircum{}2+y\textasciicircum{}3 h; a simple-zero perturbation of h gives the excluded X\textasciicircum{}2+y\textasciicircum{}3 Z through an actual coordinate change. In the rank-zero case a small cubic perturbation, Weierstrass division and a Tschirnhaus change give Y\textasciicircum{}3+x\textasciicircum{}2 H; a simple-zero perturbation of H gives the excluded Y\textasciicircum{}3+x\textasciicircum{}2 Z. On the two nested boundaries the A2 normal form gives the required real submersion, and proper boundary transport carries the compact nonzero-square cycles to the original local fibre. The two coordinate perturbations are used on their specified nested balls. The remaining unit case is the actual product cusp, whose normalization is (u,z) to (u\textasciicircum{}3,-u\textasciicircum{}2,z), with conductor coordinate z unchanged}{\ref{night:P4.16}; \ref{night:P4.18}; \ref{night:P4.20}}

Let $p$ lie on a reduced fibre support of Part~4.2.
Suppose its reduced singular locus is a smooth curve
near $p$, and its generic transverse plane-curve germ
is the cusp $a^2+b^3=0$. Then the actual germ at $p$
is analytically isomorphic to
$\{a^2+b^3=0\}\times(\mathbb C,0)$.
In particular its normalization is smooth.
\end{proposition}
\begin{proof}
\noindent\textbf{Step 1.} We extend the boundary argument of P4.20 in the
precise form used here. Suppose a deformation retains
the smooth critical curve near a fixed transverse
boundary sphere and retains transverse type $A_2$.
Near a boundary point, the holomorphic Morse lemma
in the nondegenerate transverse variable followed
by extraction of a cube root of a unit gives
coordinates with defining function $a^2+b^3$.
The third coordinate parametrizes the critical
curve. Its transversality to the boundary lets
one real third-coordinate variable be solved
from the boundary equation; the complex
coordinates $a,b$ remain free. Thus the
restriction to the boundary is submersive
at every nonzero level. Compactness and
the smooth part of the boundary give the
same uniform nonzero-level submersion as
in P4.20. Its finite-critical-value and
proper Ehresmann transport proof applies
unchanged. A compact spherical class of
nonzero square in a deformed nearby
fibre would therefore lie in an actual
nearby torus fibre, which is impossible.
All deformations below retain this
boundary transverse type.

\noindent\textbf{Step 2.} Choose holomorphic coordinates $(x,y,z)$
so the smooth critical curve is $x=y=0$.
The transverse Hessian along it is a
symmetric rank-one matrix over the
discrete valuation ring $\mathbb C\{z\}$.
It is not the zero matrix generically.
It can be diagonalized holomorphically
to $h(z)x^2$. Here is the elementary
ring calculation. The determinant is
zero. Choose a diagonal entry of
minimal order; one exists since the
square of the off-diagonal entry is
the product of the two diagonal entries.
The ratio of the off-diagonal entry
to the chosen entry is holomorphic.
Completing its square gives the desired
invertible coordinate change. The
coefficient $h$ is not identically zero.

\noindent\textbf{Step 3.} First suppose that the Hessian has
rank one at $p$. The Morse lemma with
parameters $(y,z)$ writes the reduced
equation as
\[
 g=X^2+y^3h_1(y,z).
\]
Indeed after removing the $X$ direction,
the remaining function and its first
two $y$ derivatives vanish on the
critical curve; the third derivative
is not identically zero there. If
$h_1(0,0)\ne0$, extracting its cube
root in the $y$ coordinate proves the
asserted product form.

\noindent\textbf{Step 4.} Otherwise $h_1(0,z)$ has a zero of
finite positive order. Use
$g_\eta=X^2+y^3(h_1+\eta)$.
For arbitrarily small generic $\eta$,
the function $h_1(0,z)+\eta$ has a
simple zero $z_\eta$ near zero.
This follows directly from one-variable
Weierstrass preparation, avoiding the
finite critical values of $h_1(0,z)$.
At $(0,0,z_\eta)$ the function
$Z=h_1(y,z)+\eta$ is a coordinate,
since its $z$ derivative is nonzero.
The actual germ is
$X^2+y^3Z$, which is excluded by
P4.16: its nearby fibre has a
spherical class of square $-6$.
On a fixed sufficiently small boundary,
$h_1(0,z)$ is nonzero at the critical
curve, and this persists for small
$\eta$. The deformed transverse type
is therefore still $A_2$ there.
The preceding boundary transport gives
the contradiction.

\noindent\textbf{Step 5.} It remains to exclude Hessian rank
zero at $p$, so $h(0)=0$. In the
diagonalizing coordinates add
$\varepsilon y^3$ to $g$. The
critical curve and its transverse
Hessian do not change. The cubic
coefficient in the transverse kernel
remains nonzero near the original
compact boundary for small
$\varepsilon$, so this deformation
has the boundary property just proved.
Choose $\varepsilon$ so that the
coefficient of $y^3$ at the origin
is nonzero.

\noindent\textbf{Step 6.} On a smaller neighbourhood, Weierstrass
preparation in $y$ gives a unit times
\[
 P=y^3+a(x,z)y^2+b(x,z)y+c(x,z).
\]
Restriction to $x=0$ is a unit times
$y^3$, so all three coefficients
vanish at $x=0$. The critical-curve
equation $g_x(0,0,z)=0$ first gives
$c_x(0,z)=0$. In the identity
$g=U P$ the mixed derivative on
the critical curve is then
$g_{xy}=U b_x$, since the terms
containing $c,c_x,b$ vanish there.
The zero $xy$ Hessian entry therefore
gives $b_x(0,z)=0$. Thus
$x^2\mid b$, $x^2\mid c$, and
$x\mid a$.
The Tschirnhaus change $y=Y-a/3$
therefore gives exactly
\[
 P=Y^3+x^2H(x,Y,z).
\]
The coefficient $H(0,0,z)$ is a unit
multiple of the original Hessian
coefficient $h(z)$. It is not
identically zero but vanishes at
zero. Perturb it to $H+\eta$.
At a simple root of
$H(0,0,z)+\eta$, the coordinates
$(x,Y,Z=H+\eta)$ give
$Y^3+x^2Z$. P4.18 excludes this
binomial germ by its compact
spherical class of nonzero square.

\noindent\textbf{Step 7.} There are two nested transport
neighbourhoods in this argument.
First fix the original ball for
$g\to g+\varepsilon y^3$.
Then choose a smaller ball in the
Weierstrass coordinate neighbourhood,
where $g+\varepsilon y^3$ has only
the original critical curve and its
boundary has transverse type $A_2$.
The second perturbation transports
its spherical class into a nearby
fibre on this smaller ball. Inclusion
puts it in a fibre on the original
ball; the first boundary transport
puts it in a fibre of $g$.
This does not require extending the
second perturbation outside its
coordinate neighbourhood.
Units in the reduced equations are
retained during the perturbations.
At a binomial special point a unit
can be absorbed by the positive
weighted coordinate scaling of
P4.16--P4.18.

\noindent\textbf{Step 8.} We have ruled out both a zero
Hessian and a vanishing cubic
coefficient. The product form
follows. Its normalization is
the finite map
\[
 (u,z)\longmapsto (a=u^3,b=-u^2,z),
\]
whose source is smooth, whose
inverse rational parameter off
$b=0$ is $u=-a/b$, and whose
restriction to the conductor is
the identity on the $z$ coordinate.
\end{proof}

\begin{proposition}[P4.29]\label{night:P4.29}\label{stmt:P4.29}
\StatementRecord{night:P4.29}{proved}{Historical events: GPT-6 Astra (OpenAI), ChatGPT Pro mode, effort Pro, S6\_06\_BLIND, 16--18 November 2026: targeted proof check. Also checked in this writing. Pre-continuation writing events: GPT-6 Astra (OpenAI), ChatGPT Pro mode, effort Pro, S6\_10\_PAPERS, 18 November 2026: targeted proof check. Continuation metadata events: GPT-6 Astra (OpenAI), ChatGPT Pro mode, effort Pro, S6\_10\_PAPERS continuation, 19 November 2026: used and rechecked in this run; status-vocabulary reconciliation; metadata reconciliation only; no new mathematical check. \emph{Pre-continuation recorded proof scope.} The night-run re-derivation was by GPT-6 Astra (OpenAI, ChatGPT Pro mode, effort Pro), 16--18 November 2026. Also checked in this writing by GPT-6 Astra (OpenAI), ChatGPT Pro mode, effort Pro, 18 November 2026. The zero-Euler cusp graph, descent of the ruling, rationality of normalization singularities, Cohen-Macaulay reduced conductor, and vanishing conductor delta were re-derived. The product cusp normal form then gives the actual smooth minimal normalization}{\ref{night:P4.12}; \ref{night:P4.13}; \ref{night:P4.22}; \ref{night:P4.28}; \ref{night:P4R.7}; \ref{night:P4R.9}}

Let $S$ be an irreducible singular support
in Part~4.2 whose normalization has ruled
resolution, and suppose $e(S)=0$.
In the double-horizontal-divisor case,
when the generic transverse type is
cuspidal $A_2$, the normalization is
a smooth minimal elliptic ruled surface.
The image conductor is a smooth elliptic
curve and every point along it has the
product cusp germ of P4.28.
\end{proposition}
\begin{proof}
\noindent\textbf{Step 1.} Use the notation of P4R.7 and P4R.9.
The actual resolved pole support is
a section $H$ with rational trees,
and its inclusion in the ruled
resolution $T$ induces an isomorphism
on fundamental groups. Thus
$\pi_1(S)=\pi_1(C)$, where
$C=\operatorname{Sing}(S)_{\rm red}$.
The virtually abelian curve-group
calculation gives $e(C)\ge0$, and
P4.13 gives $e(S_{\rm reg})\ge0$.
Equality makes both numbers zero.
The same connected-curve calculation
shows that $C$ has exactly one
normalized component, of genus one,
and has no point-identification
cycles.

\noindent\textbf{Step 2.} All rational components of the
resolved pole support are contracted
to points: a nonconstant proper map
from a rational curve to the
elliptic normalization of $C$ is
impossible. Hence the strict
conductor consists only of $H$.
The generic cusp has one
normalization branch, so
$H\to C$ has generic degree one
and is the normalization of $C$.
In particular $H$ and the ruling
base both have genus one.

\noindent\textbf{Step 3.} Every exceptional curve of
$T\to N$ is now rational and
vertical for the ruling, since
$H$ is not contracted. The ruling
$r:T\to H$ descends to a holomorphic
map $\bar r:N\to H$. For an explicit
descent, the image of
$(\rho_N,r):T\to N\times H$
projects properly and bijectively
to $N$, because $r$ is constant
on each exceptional fibre. Its
normalized graph projects finitely
and birationally to the normal
space $N$, hence is isomorphic
to $N$.

\noindent\textbf{Step 4.} The normalization singularities
are rational. Indeed a sequence
of point blowups of a
$\mathbb P^1$-bundle has
$R^1r_*\mathcal O_T=0$.
For a blowup this follows from
the two affine charts, or from
the exact sequence of the
exceptional $\mathbb P^1$ with
$H^1(\mathcal O_{\mathbb P^1})=0$;
the assertion for the whole
sequence follows by Leray.
Since $\rho_{N*}\mathcal O_T
=\mathcal O_N$, the low-degree
Leray sequence for
$r=\bar r\circ\rho_N$ gives
\[
 \bar r_*R^1\rho_{N*}\mathcal O_T=0.
\]
Here $R^2\bar r_*\mathcal O_N=0$
because $\bar r$ has one-dimensional
fibres. The sheaf
$R^1\rho_{N*}\mathcal O_T$ is
supported at finitely many
normalization singular points,
and direct image of a nonzero
such sheaf to $H$ is nonzero.
Thus it vanishes. In particular
$\chi(\mathcal O_N)
=\chi(\mathcal O_T)=0$.

\noindent\textbf{Step 5.} We next check the scheme structure
of the conductor, to prove that
its image curve is smooth.
Write $I$ for the conductor ideal
on $S$ and on $N$, and let
$C_{\rm sch}\subset S$ and
$D_{\rm sch}\subset N$ be its
two zero schemes. Locally
$\mathcal O_S$ is Gorenstein
of dimension two and
$\nu_*\mathcal O_N$ is a
maximal Cohen--Macaulay module.
Finite duality identifies
its dual module with $I$,
which is again maximal
Cohen--Macaulay. The depth
lemma in
$0\to I\to\mathcal O_S
\to\mathcal O_{C_{\rm sch}}\to0$
therefore gives depth one
for the conductor curve.
At a generic cusp its local
length in the transverse
direction is one. A
one-dimensional Cohen--Macaulay
scheme which is generically
reduced is reduced: any
nilpotent ideal would have
zero-dimensional support,
contradicting the absence of
embedded associated points.
Consequently $C_{\rm sch}=C$.

\noindent\textbf{Step 6.} Put $L=\nu^*\omega_S$, a
torsion line bundle by P4.12.
Finite duality on $N$ gives
\[
 I=\omega_N\otimes L^{-1}.
\]
Serre duality on the compact
Cohen--Macaulay surface $N$
gives $\chi(I)=\chi(L)$.
The rational-singularity
comparison just proved and
the projection formula give
$\chi(L)=\chi(\rho_N^*L)$.
Riemann--Roch on the smooth
surface $T$, with the rational
Chern class of this torsion
line bundle equal to zero,
gives $\chi(\rho_N^*L)
=\chi(\mathcal O_T)=0$.
Therefore
$\chi(\mathcal O_{D_{\rm sch}})
=\chi(\mathcal O_N)-\chi(I)=0$.

\noindent\textbf{Step 7.} The exact conductor square is
\[
 0\longrightarrow\mathcal O_S
 \longrightarrow
 \nu_*\mathcal O_N\oplus\mathcal O_C
 \longrightarrow\nu_*\mathcal O_{D_{\rm sch}}
 \longrightarrow0.
\]
Taking Euler characteristics and
using $\chi(\mathcal O_S)=0$
from P4.12 yields
$\chi(\mathcal O_C)=0$.
Its normalization is the
smooth elliptic curve $H$,
so the normalization sequence
gives $\chi(\mathcal O_C)
 =-\sum_p\delta(C,p)$.
All delta invariants vanish;
thus $C$ is smooth.

\noindent\textbf{Step 8.} Every point of $C$ now satisfies
P4.28, so $S$ has the product
cusp germ there and $N$ is
smooth along the conductor.
It was already smooth away
from the conductor. P4.22
therefore applies and makes
this elliptic ruled
normalization minimal.
\end{proof}


\subsection{The conormal calculation and ordinary-triple equality}

\begin{proposition}[P4.30]\label{night:P4.30}\label{stmt:P4.30}
\StatementRecord{night:P4.30}{proved}{Historical events: GPT-6 Astra (OpenAI), ChatGPT Pro mode, effort Pro, S6\_06\_BLIND, 16--18 November 2026: targeted proof check. Also checked in this writing. Pre-continuation writing events: GPT-6 Astra (OpenAI), ChatGPT Pro mode, effort Pro, S6\_10\_PAPERS, 18 November 2026: targeted proof check. Continuation metadata events: GPT-6 Astra (OpenAI), ChatGPT Pro mode, effort Pro, S6\_10\_PAPERS continuation, 19 November 2026: used and rechecked in this run; status-vocabulary reconciliation; metadata reconciliation only; no new mathematical check. \emph{Pre-continuation recorded proof scope.} The night-run re-derivation was by GPT-6 Astra (OpenAI, ChatGPT Pro mode, effort Pro), 16--18 November 2026. Also checked in this writing by GPT-6 Astra (OpenAI), ChatGPT Pro mode, effort Pro, 18 November 2026. Differentiation on the elliptic normalization has image Omega\_E(-R) and a primitive rank-two kernel K of degree deg R. The symbolic conormal Q=P/P\textasciicircum{}(2) injects into j\_*K with finite quotient of length tau. At a unibranch germ of multiplicity m the derivative common order is m-1. The Pluecker minors of the gradient submodule bound the colength below by 2m in embedding dimension three and by m in the plane case; a smooth germ contributes zero. Hence tau>=deg R and chi(Q)<=0. The exact symbolic-square sequence gives chi(O/P\textasciicircum{}(2))<=-sum delta. Equality forces every delta to vanish, so the conductor curve is smooth. The bound is derived from the explicit gradient and normalization maps}{the hypotheses and explicit local calculation stated here}

Let $C$ be a compact irreducible reduced curve in the smooth
threefold $Y$ of Part~4.2. Suppose its normalization is an elliptic
curve $E$, and every point of $C$ is unibranch. Write $P$ for its
ideal in $Y$ and $P^{(2)}$ for the ideal of functions vanishing to
order at least two at the generic point of $C$. Then
\[
 \chi(\mathcal O_Y/P^{(2)})\le0.
\]
Equality forces $C$ to be smooth. The assertion uses
$c_1(TY)\cdot C=0$, which follows from $b_2(Y)=0$.
\end{proposition}
\begin{proof}
\noindent\textbf{Step 1.} Let $j:E\to Y$ be the normalization map followed by inclusion.
The derivative map
\[
 j^*\Omega_Y\longrightarrow\Omega_E
\]
has image $\Omega_E(-R)$ for an effective divisor $R$.
Its kernel $K$ is locally free of rank two: dividing the derivative
row by its common zero gives a surjection onto a line bundle.
Since $\deg j^*\Omega_Y=0$ and $\deg\Omega_E=0$, we have
$\deg K=\deg R$.

\noindent\textbf{Step 2.} Differentiation induces an injection
\[
 Q:=P/P^{(2)}\longrightarrow j_*K.
\]
Indeed its kernel consists precisely of functions with zero first
normal derivative at a generic smooth point of $C$, which is the
defining order-two condition for $P^{(2)}$. Both sides agree at the
generic point. Their quotient has finite length, denoted by $\tau$.
Riemann--Roch on the smooth elliptic curve gives
\[
 \chi(Q)=\deg R-\tau.
\]
We prove $\tau\ge\deg R$ by local equations, including at a
singular point of $C$.

\noindent\textbf{Step 3.} Parametrize its sole branch by $t$. Choose ambient coordinates
vanishing at the point, and let
$m=\min_i\operatorname{ord}_t x_i(t)$ be the multiplicity.
The common order of the three derivative entries is $m-1$,
so the coefficient of $R$ is $m-1$. The saturated kernel $K$
is a direct summand of the free rank-three module over
$\mathbb C\{t\}$. Consequently a generator of $\bigwedge^2K$
has at least one unit Pl\"ucker coordinate. If $L$ is the
$\mathbb C\{t\}$-span of the differentiated defining equations,
then the colength of $L$ in $K$ is the minimum order of their
two-by-two minors. This is the diagonal elementary-divisor
calculation for a rank-two submodule over a discrete valuation
ring. The local length of $j_*K/Q$ is at least this colength,
since $Q\subset L\subset K$.

\noindent\textbf{Step 4.} At a smooth point $m=1$ and there is nothing to prove.
At a singular point of embedding dimension three, every defining
equation belongs to the square of the ambient maximal ideal.
Every entry of its gradient therefore has order at least $m$
on the parametrization, and every two-by-two minor has order
at least $2m$. Hence the local quotient length is at least
$2m\ge m-1$.

\noindent\textbf{Step 5.} At a singular point of embedding dimension two, an invertible
ambient coordinate change writes the ideal as $(z,h(x,y))$,
where $h$ is a reduced irreducible plane equation without a
linear term. The conormal image has a free $dz$ summand.
Its other generator is $dh$; both its plane coefficients have
order at least $m$. The primitive rank-one kernel of the plane
derivative row has a unit coefficient, so $dh$ is its generator
multiplied by a function of order at least $m$. The quotient
length is therefore at least $m\ge m-1$. Embedding dimension
one is the smooth case. These cases prove the local bound and
therefore $\chi(Q)\le0$.

\noindent\textbf{Step 6.} The normalization exact sequence gives
$\chi(\mathcal O_C)=-\sum_p\delta(C,p)$.
Finally the exact sequence
\[
 0\longrightarrow P/P^{(2)}\longrightarrow
 \mathcal O_Y/P^{(2)}\longrightarrow\mathcal O_C
 \longrightarrow0
\]
gives
\[
 \chi(\mathcal O_Y/P^{(2)})
 =-\sum_p\delta(C,p)+\deg R-\tau\le0.
\]
Equality forces all delta invariants to vanish, so the
normalization is an isomorphism and $C$ is smooth.
\end{proof}

\begin{proposition}[P4.31]\label{night:P4.31}\label{stmt:P4.31}
\StatementRecord{night:P4.31}{proved}{Historical events: GPT-6 Astra (OpenAI), ChatGPT Pro mode, effort Pro, S6\_06\_BLIND, 16--18 November 2026: targeted proof check. Also checked in this writing. Pre-continuation writing events: GPT-6 Astra (OpenAI), ChatGPT Pro mode, effort Pro, S6\_10\_PAPERS, 18 November 2026: targeted proof check. Continuation metadata events: GPT-6 Astra (OpenAI), ChatGPT Pro mode, effort Pro, S6\_10\_PAPERS continuation, 19 November 2026: used and rechecked in this run; status-vocabulary reconciliation; metadata reconciliation only; no new mathematical check. \emph{Pre-continuation recorded proof scope.} The night-run re-derivation was by GPT-6 Astra (OpenAI, ChatGPT Pro mode, effort Pro), 16--18 November 2026. Also checked in this writing by GPT-6 Astra (OpenAI), ChatGPT Pro mode, effort Pro, 18 November 2026. The graph and ruled-pole bounds give a single unibranch conductor image with elliptic normalization and a degree-three horizontal section of genus g. The rational-exceptional contraction descends the ruling to the normalization and gives chi(O\_N)=1-g. The conductor image is pure Cohen-Macaulay with one associated prime; its generic ordinary-triple calculation is P\textasciicircum{}2, so primary localization and saturation identify its global ideal with P\textasciicircum{}(2). Finite duality, torsion adjunction and Serre duality give chi(D\_conductor)=0 and chi(C\_conductor)=g-1. P4.30 forces g=1 and a smooth image; the degree-three section map is unramified. Distinct section points cannot be joined by a vertical exceptional block. Each of the three hypersurface branches is normal, and its pairwise intersection with another branch is the reduced smooth Cartier curve; this forces each branch regular and pairwise transverse. The explicit coordinate changes take the third equation a x+b y to X-Y, yielding XY(X-Y), and the normalization is consequently smooth minimal elliptic ruled}{\ref{night:P4.22}; \ref{night:P4.29}; \ref{night:P4.30}; \ref{night:P4R.6}; \ref{night:P4R.7}; \ref{night:P4R.9}}

Let $S$ be an irreducible singular support in Part~4.2 whose
normalization has ruled resolution, and suppose $e(S)=0$.
In the double-horizontal-divisor case, when the generic
transverse type is the ordinary triple point $xy(x-y)=0$,
the normalization is a smooth minimal elliptic ruled surface.
Its conductor image is a smooth elliptic curve, and its
reduced conductor is an unramified cover of that curve of
degree three. Every germ along the image is analytically
the product of an ordinary plane triple point with a disc.
\end{proposition}
\begin{proof}
\noindent\textbf{Step 1.} Use the minimal resolution $T\to N\to S$ and the section $H$
of P4R.7 and P4R.9. Put $g=g(H)$ and
$C=\operatorname{Sing}(S)_{\rm red}$.
The same fundamental-group and Euler argument as in P4.29
shows that $C$ is connected, has just one normalized
component, that this normalization $E$ has genus one, and
that its point-identification graph has no cycles. Thus
every point of $C$ is unibranch. All rational components
of the resolved pole support are contracted, and the
only strict conductor component is $H$. Its finite map
to $C$ has degree three, since an ordinary transverse
triple point has three normalization branches.
In particular $g\ge1$ by Riemann--Hurwitz for $H\to E$.

\noindent\textbf{Step 2.} All exceptional curves of $T\to N$ are rational and
vertical for the ruling. Exactly the normalized-graph
descent and Leray calculation of P4.29 give a morphism
$N\to H$ and
\[
 R^1\rho_{N*}\mathcal O_T=0,\qquad
 \chi(\mathcal O_N)=\chi(\mathcal O_T)=1-g.
\]
Thus the normalization singularities are rational; no
smoothness assertion is used here.

\noindent\textbf{Step 3.} Let $I$ be the conductor ideal, and let $C_{\rm sch}$
and $D_{\rm sch}$ be its zero schemes in $S$ and $N$.
Finite duality and the depth argument of P4.29 show that
$\mathcal O_{C_{\rm sch}}$ is Cohen--Macaulay of pure
dimension one. If $P\subset\mathcal O_Y$ is the reduced
ideal of $C$, the inverse image $J$ of $I$ in
$\mathcal O_Y$ therefore has no embedded associated
prime. Its only associated prime locally is $P$.
The generic ordinary triple point has conductor equal
to the square of its maximal ideal in the transverse
plane. To check this directly, normalize its three
lines. Constant functions on them have a common value;
their linear terms are the evaluations of one linear
form on three distinct directions, a two-dimensional
space. Quadratic forms evaluate surjectively on those
three directions. All terms of degree at least two
can therefore be chosen independently on the three
branches, and a nonzero constant or linear term cannot
belong to the conductor. This proves the stated square
ideal, including its scheme structure.

\noindent\textbf{Step 4.} Consequently $J_P=P_P^2$, and in fact
\[
 J=P^{(2)}.
\]
For completeness, if $a\in P^{(2)}$, some $s\notin P$
satisfies $sa\in J$. Multiplication by $s$ is injective
on $\mathcal O_Y/J$, because its only associated prime
is $P$, so $a\in J$. The reverse inclusion follows by
localizing at $P$. Thus the conductor image is the
actual symbolic double curve to which P4.30 applies.

\noindent\textbf{Step 5.} Put $L=\nu^*\omega_S$, which is torsion. Finite duality
gives $I=\omega_N\otimes L^{-1}$. Serre duality, the
rational-singularity comparison and Riemann--Roch on
$T$ give
\[
 \chi(I)=\chi(L)=\chi(\rho_N^*L)
   =\chi(\mathcal O_T)=1-g.
\]
Hence $\chi(\mathcal O_{D_{\rm sch}})=0$.
Taking Euler characteristics in the exact conductor
square, and using $\chi(\mathcal O_S)=0$, yields
\[
 \chi(\mathcal O_{C_{\rm sch}})=g-1.
\]
P4.30 now gives $g-1\le0$. Since $g\ge1$, equality
holds, $g=1$, and $C$ is smooth. Riemann--Hurwitz then
makes the degree-three map $H\to C$ unramified.

\noindent\textbf{Step 6.} We finish by proving normalization smoothness from
the actual local morphisms. Over a point of $C$ the
three distinct points of $H$ remain distinct in $N$:
every exceptional block is vertical for the ruling
$T\to H$ and therefore cannot join different points
of $H$. The finite normalization consequently gives
three local irreducible branches of $S$. Each branch
is a hypersurface in $Y$ and is smooth at the generic
point of $C$. Its singularities are isolated, so the
hypersurface Cohen--Macaulay property and regularity
in codimension one make it normal. The corresponding
component of $N$ is that branch itself.

\noindent\textbf{Step 7.} Write its three local equations as $f_1,f_2,f_3$.
The intersection of any two branches is supported on
$C$ and has generic multiplicity one, because the
generic transverse triple point is ordinary. On a
normal branch the other equation is a nonzero divisor.
Its quotient is a one-dimensional Cohen--Macaulay
ring, generically reduced with support $C$, so it is
reduced and is exactly $\mathcal O_C$. Thus $C$ is a
smooth Cartier curve in each branch. A local ring of
dimension two with a nonzero divisor whose quotient
is a regular one-dimensional local ring has embedding
dimension at most two: lift its one maximal-ideal
generator and add the nonzero divisor. It is regular.
All three branches are therefore smooth.

\noindent\textbf{Step 8.} Any two branches meet transversely along $C$.
Choose their equations as coordinates $x,y$, with
$z$ a coordinate on $C$. The third equation is
$a(x,y,z)x+b(x,y,z)y$, where both $a$ and $b$ are
units by the pairwise transversality just proved.
The invertible coordinate change $X=ax$, $Y=-by$
writes the reduced union exactly as
$XY(X-Y)=0$. This is the asserted product triple
point, and its normalization is smooth. Away from
the conductor, normalization singularities were
already excluded by P4R.6. Finally P4.22 applies to
this smooth elliptic ruled normalization and proves
its minimality.
\end{proof}


\subsection{Complete zero-Euler classification in the requested classes}

\begin{proposition}[P4.32]\label{night:P4.32}\label{stmt:P4.32}
\StatementRecord{night:P4.32}{proved}{Historical events: GPT-6 Astra (OpenAI), ChatGPT Pro mode, effort Pro, S6\_06\_BLIND, 16--18 November 2026: targeted proof check. Also checked in this writing. Pre-continuation writing events: GPT-6 Astra (OpenAI), ChatGPT Pro mode, effort Pro, S6\_10\_PAPERS, 18 November 2026: targeted proof check. Continuation metadata events: GPT-6 Astra (OpenAI), ChatGPT Pro mode, effort Pro, S6\_10\_PAPERS continuation, 19 November 2026: used and rechecked in this run; status-vocabulary reconciliation; metadata reconciliation only; no new mathematical check. \emph{Pre-continuation recorded proof scope.} The night-run re-derivation was by GPT-6 Astra (OpenAI, ChatGPT Pro mode, effort Pro), 16--18 November 2026. Also checked in this writing by GPT-6 Astra (OpenAI), ChatGPT Pro mode, effort Pro, 18 November 2026. The zero-Euler normalization and conductor alternatives were checked against their stated local inputs. The local conclusion uses the single support, ambient b1=b2=0, and nearby tori; the final global assertion retains the source hypothesis that every support is irreducible or NC}{\ref{night:P4.12}; \ref{night:P4.22}; \ref{night:P4.24}; \ref{night:P4.26}; \ref{night:P4.27}; \ref{night:P4.29}; \ref{night:P4.31}; \ref{night:P4.8}; \ref{night:P4R.11}; \ref{night:P4R.7}; \ref{night:P4R.9}}

An irreducible support in Part~4.2 has Euler number zero exactly
in the following cases: it is smooth with a finite unramified
torus cover, or it is singular with smooth minimal elliptic ruled
normalization, or it is singular with smooth minimal Hopf
normalization. In the ruled double-horizontal case the conductor
image is a smooth elliptic curve; the local germs are product
cusps, product tacnodes, or product ordinary plane triple points,
and the reduced conductor maps have degrees one, two, or three
respectively and are unramified.

For a fibration whose supports are all irreducible or normal
crossing, $c_3(Y)=0$ exactly when every irreducible support has
one of these forms and every reducible support has the
zero-Euler normal-crossing form of P4.8.
\end{proposition}
\begin{proof}
\noindent\textbf{Step 1.} Normal irreducible supports are smooth with the stated finite
unramified cover by P4.8 and P4.12. For a singular irreducible
support, use its minimal normalization resolution $T$.
P4.26 excludes rational resolutions and class-VII minimal
models with positive second Betti number from equality.
The Inoue minimal model has no conductor curves and is
excluded there as well. A Hopf minimal model is covered
completely by P4.27: equality makes the actual normalization
a smooth minimal Hopf surface, and the converse is P4.24.

\noindent\textbf{Step 2.} It remains to use the positive-genus ruled resolution.
The reduced-horizontal alternatives of P4R.9 already give
smooth minimal elliptic ruled normalization. In the remaining
case the pole support is a section $H$ with rational trees,
and the coefficient of $H$ is two. Its inclusion in $T$
induces an isomorphism on fundamental groups. The conductor
pushout and the curve-group calculation of P4R.7--P4R.9 give
$\pi_1(S)=\pi_1(C)$ and $e(C)\ge0$ for the reduced image
conductor $C$. Equality in
$e(S)=e(S_{\rm reg})+e(C)$ makes $C$ a connected curve
with elliptic normalization and no point-identification
cycles. Every rational pole component is contracted, since
there is no finite nonconstant map from a rational curve
to an elliptic curve. Thus $H$ is the sole strict conductor
component, and every branch of a generic transverse
plane-curve singularity has conductor exponent two.

\noindent\textbf{Step 3.} We enumerate those transverse germs explicitly. For a branch
$B_i$ of a reduced plane curve, its conductor exponent is
\[
 c_i=2\delta(B_i)+\sum_{j\ne i}I(B_i,B_j).
\]
If a branch is singular, $c_i=2$ forces
$\delta(B_i)=1$ and no other branch. The one-gap
parametrization of a delta-one irreducible plane curve
has orders two and three and gives the ordinary cusp.
If all branches are smooth, two branches have intersection
multiplicity two and give a tacnode, or three branches
have all pairwise intersection multiplicities one and
give an ordinary triple point. Four branches would give
each conductor exponent at least three. One smooth branch
is nonsingular and has exponent zero. These exhaust the
possibilities.

\noindent\textbf{Step 4.} P4.29, P4R.11, and P4.31 respectively treat the cusp,
tacnode, and ordinary-triple cases, without assuming that
the initial normalization or image conductor was smooth.
They give the asserted smooth minimal elliptic ruled
normalization and the precise local conductor maps.
The converse for a smooth minimal elliptic ruled
normalization is P4.22. This proves the entire irreducible
zero-Euler classification.

\noindent\textbf{Step 5.} Finally P4.26 expresses $c_3(Y)$ as a finite sum of
nonnegative fibre Euler numbers. Its vanishing is therefore
equivalent to vanishing of every summand. The irreducible
summands have just been classified, and P4.8 gives the
complete reducible normal-crossing zero case. This proves
the final assertion.
\end{proof}


\subsection{Normalization cohomology for the remaining positive cases}

\begin{proposition}[P4.33]\label{night:P4.33}\label{stmt:P4.33}
\StatementRecord{night:P4.33}{proved}{Historical events: GPT-6 Astra (OpenAI), ChatGPT Pro mode, effort Pro, S6\_06\_BLIND, 16--18 November 2026: targeted proof check. Also checked in this writing. Pre-continuation writing events: GPT-6 Astra (OpenAI), ChatGPT Pro mode, effort Pro, S6\_10\_PAPERS, 18 November 2026: targeted proof check. Continuation metadata events: GPT-6 Astra (OpenAI), ChatGPT Pro mode, effort Pro, S6\_10\_PAPERS continuation, 19 November 2026: used and rechecked in this run; status-vocabulary reconciliation; metadata reconciliation only; no new mathematical check. \emph{Pre-continuation recorded proof scope.} The night-run re-derivation was by GPT-6 Astra (OpenAI, ChatGPT Pro mode, effort Pro), 16--18 November 2026. Also checked in this writing by GPT-6 Astra (OpenAI), ChatGPT Pro mode, effort Pro, 18 November 2026. Finite duality, the coherent conductor Euler relation, the two possibilities for the horizontal component, descent of the ruling, and the Leray length calculation were re-derived}{\ref{night:P4.12}; \ref{night:P4.26}; \ref{night:P4.29}}

Let $S$ be an irreducible singular support in Part~4.2, and let
$\rho:T\to N\to S$ be its minimal normalization resolution.
Then $H^2(N,\mathcal O_N)=0$. If
$p=\dim H^0(N,R^1\rho_{N*}\mathcal O_T)$ is the sum of the
geometric genera of the normalization singularities, then
\[
 p=h^1(T,\mathcal O_T)-h^1(N,\mathcal O_N).
\]
In particular a rational resolution forces every normalization
singularity to be rational and gives $\chi(\mathcal O_N)=1$.
For the scheme-theoretic image conductor $C_{\rm sch}$ we have
\[
 \chi(\mathcal O_{C_{\rm sch}})
   =-\chi(\mathcal O_N)=h^1(N,\mathcal O_N)-1.
\]

In the ruled double-horizontal case, with section $H$ of genus
$g\ge1$, a contracted $H$ gives $h^1(N,\mathcal O_N)=0$ and
$p=g$. A noncontracted $H$ gives $p=0$ and
$h^1(N,\mathcal O_N)=g$.
\end{proposition}
\begin{proof}
\noindent\textbf{Step 1.} Write $I_S\subset\mathcal O_S$ for the conductor ideal.
Finite duality gives
\[
 \nu_*\omega_N=I_S\otimes\omega_S.
\]
The dualizing line bundle $\omega_S$ is torsion by P4.12.
A nonzero section of a torsion line bundle on a compact
irreducible reduced complex space has no zero: raising it
to a trivializing positive power gives a nonzero holomorphic
function, hence a nonzero constant. Such a section cannot
belong to the proper conductor ideal along the nonempty
conductor. Therefore $H^0(N,\omega_N)=0$. Serre duality
on the compact Cohen--Macaulay normal surface gives
$H^2(N,\mathcal O_N)=0$.

\noindent\textbf{Step 2.} The resolution is birationally ruled or of class VII and
has $H^2(T,\mathcal O_T)=0$, by P4.12 and the classification
already used in P4.26. The low-degree Leray sequence,
using $\rho_{N*}\mathcal O_T=\mathcal O_N$ and the
finite support of $R^1\rho_{N*}\mathcal O_T$, is therefore
\[
 0\longrightarrow H^1(N,\mathcal O_N)
 \longrightarrow H^1(T,\mathcal O_T)
 \longrightarrow H^0(N,R^1\rho_{N*}\mathcal O_T)
 \longrightarrow0.
\]
This proves the formula for $p$. If $T$ is rational,
$h^1(\mathcal O_T)=0$, so $p=0$. Each nonnegative local
geometric genus is zero, giving the asserted rational
singularities and Euler characteristic.

\noindent\textbf{Step 3.} For completeness the conductor Euler identity does not
require rational singularities. Put $L=\nu^*\omega_S$,
and write $I_N=\omega_N\otimes L^{-1}$ for the conductor
on $N$. Serre duality gives $\chi(I_N)=\chi(L)$.
Riemann--Roch on $T$, applied to the torsion line bundle
$\rho_N^*L$, gives
$\chi(\rho_N^*L)=\chi(\mathcal O_T)$.
The projection formula makes
$R^1\rho_{N*}\rho_N^*L=L\otimes R^1\rho_{N*}\mathcal O_T$;
its finite length is still $p$. Leray therefore gives
$\chi(L)=\chi(\mathcal O_N)$.
Consequently the conductor scheme $D_{\rm sch}\subset N$
has $\chi(\mathcal O_{D_{\rm sch}})=0$.
The exact conductor square, together with
$\chi(\mathcal O_S)=0$, proves
$\chi(\mathcal O_{C_{\rm sch}})=-\chi(\mathcal O_N)$.
Since $N$ is compact connected and $H^2(\mathcal O_N)=0$,
the remaining equality follows.

\noindent\textbf{Step 4.} Finally suppose the ruled pole support is $H$ with
rational vertical trees. If $H$ is contracted, the
composition $H\to T\to N$ is constant. Pullback from
$H^1(N,\mathcal O_N)$ to $H^1(H,\mathcal O_H)$ is zero.
But restriction
$H^1(T,\mathcal O_T)\to H^1(H,\mathcal O_H)$ is an
isomorphism: the ruling has section $H$, a
$\mathbb P^1$-bundle has zero first higher direct image
of its structure sheaf, and the actual point blowups
preserve these groups. The Leray injection above
therefore forces $H^1(N,\mathcal O_N)=0$, so $p=g$.
If $H$ is not contracted, every exceptional curve is
vertical. The ruling descends by the normalized-graph
comparison of P4.29. Its Leray calculation gives $p=0$,
and the same exact sequence gives $h^1(\mathcal O_N)=g$.
\end{proof}


\subsection{The value one in the rational ordinary-triple sector}

\begin{proposition}[P4.34]\label{night:P4.34}\label{stmt:P4.34}
\StatementRecord{night:P4.34}{proved}{Historical events: GPT-6 Astra (OpenAI), ChatGPT Pro mode, effort Pro, S6\_06\_BLIND, 16--18 November 2026: targeted proof check. Also checked in this writing. Pre-continuation writing events: GPT-6 Astra (OpenAI), ChatGPT Pro mode, effort Pro, S6\_10\_PAPERS, 18 November 2026: targeted proof check. Continuation metadata events: GPT-6 Astra (OpenAI), ChatGPT Pro mode, effort Pro, S6\_10\_PAPERS continuation, 19 November 2026: used and rechecked in this run; status-vocabulary reconciliation; metadata reconciliation only; no new mathematical check. \emph{Pre-continuation recorded proof scope.} The night-run re-derivation was by GPT-6 Astra (OpenAI, ChatGPT Pro mode, effort Pro), 16--18 November 2026. Also checked in this writing by GPT-6 Astra (OpenAI), ChatGPT Pro mode, effort Pro, 18 November 2026. The symbolic double-curve Euler formula, the planar defect 5 delta, the nonplanar valuation bounds, and the exclusion of the apparent equality-six unibranch case were re-derived. The final exclusion of Euler one uses exactly the stated global upper bound}{\ref{night:P4.26}; \ref{night:P4.30}; \ref{night:P4.31}; \ref{night:P4.33}; \ref{night:P4.8}}

An irreducible singular support with rationally resolved
normalization and generic transverse ordinary triple point
cannot have Euler number one. For a fibration all of whose
supports are irreducible or normal crossing, its Euler number
is therefore two.
\end{proposition}
\begin{proof}
\noindent\textbf{Step 1.} Suppose its Euler number were one. The rational
normal-crossing case has Euler number two by P4.8, so we
are in the simply connected resolved-conductor alternative
of P4.26. That proof gives $\pi_1(S)=\pi_1(C)=F_b$,
$b\le1$, and
$e(C)=c+1-b$ for the number $c$ of normalized rational
components. Since $e(S_{\rm reg})\ge0$, the value one
forces $c=b=1$ and $e(S_{\rm reg})=0$.
Thus $C$ has normalization $E=\mathbb P^1$, exactly one
point with two analytic branches, and only unibranch
singular points elsewhere.

\noindent\textbf{Step 2.} P4.33 proves that $N$ has rational singularities and
that its image conductor scheme has
$\chi(\mathcal O_{C_{\rm sch}})=-1$.
The generic ordinary-triple conductor and pure
Cohen--Macaulay argument of P4.31 identify this scheme
with the symbolic double curve $\mathcal O_Y/P^{(2)}$.
At a point with two analytic branches, the symbolic ideal
means order at least two along each branch. The same
argument applies: a nonzero finite-support defect in the
conductor quotient would violate depth one.

\noindent\textbf{Step 3.} We compute its Euler characteristic and the possible
local defects. On $E$ use the bundles $K$ and the
ramification divisor $R$ of P4.30. This time
$\deg K=2+\deg R$, since $E=\mathbb P^1$.
Write $\tau_p$ for the local length of the conormal
quotient $j_*K/(P/P^{(2)})$ and $r_p$ for the sum of
the coefficients of $R$ over $p$. Normalization and
Riemann--Roch give
\[
 \chi(\mathcal O_Y/P^{(2)})
 =5-\sum_p d_p,
 \qquad d_p=\delta(C,p)+\tau_p-r_p.
\]
We prove the precise gap needed here: at an ordinary
node $d_p=5$; at every planar singularity
$d_p=5\delta(C,p)$; and at every nonplanar singularity
having at most two branches $d_p\ge7$.

\noindent\textbf{Step 4.} First consider a planar curve, with ideal $(z,h(x,y))$.
Its conormal is free over $\mathcal O_C$, generated by
$dz,dh$. On each normalized branch, write
$dh=a_i u_i$ for a primitive plane kernel generator,
and put $v_i=\operatorname{ord}a_i$.
The inclusion of each of these two free summands into
its normalized module has length $\delta(C,p)$.
The second summand has the additional vanishing orders
$v_i$. Hence
\[
 \tau_p=2\delta(C,p)+\sum_i v_i.
\]
We verify $\sum_i(v_i-r_i)=2\delta(C,p)$ directly by
curve duality. Let $A$ be the reduced plane-curve ring,
$B$ its normalization, and $I=\operatorname{Hom}_A(B,A)$
its conductor. Applying $\operatorname{Hom}_A(-,A)$ to
$0\to A\to B\to B/A\to0$ gives
$A/I\simeq\operatorname{Ext}^1_A(B/A,A)$.
Here $\operatorname{Ext}^1_A(B,A)=0$ because $B$ is
maximal Cohen--Macaulay over the one-dimensional
Gorenstein ring $A$. The dual of a finite-length module
has the same length, so $\operatorname{length}(A/I)=\delta$,
and $\operatorname{length}(B/I)=2\delta$.
If the conductor exponent on branch $i$ is $c_i$,
this says $\sum_i c_i=2\delta$.
Choose a generic plane coordinate $x$ with minimal
valuation on every normalized branch. The dualizing
generator $dx/h_y$ then pulls back with
order $r_i-v_i$. The identity
$\omega_B=I\omega_A$ therefore gives
$c_i=v_i-r_i$. Combining these equalities proves
$\tau_p-r_p=4\delta$ and the stated planar formula.

\noindent\textbf{Step 5.} Now consider a nonplanar unibranch curve of multiplicity
$m$. Its embedding dimension is three. As in P4.30,
all gradient entries have valuation at least $m$, so
$\tau_p\ge2m$, while $r_p=m-1$ and
$\delta\ge m-1$. A nonplanar unibranch curve has
$m\ge3$. To see the last assertion, if $x=t^2$, write
each other coordinate as $a(x)+t b(x)$. After subtracting
the even parts, choose the nonzero $b$ of least order.
The other odd coordinate is an analytic multiple of it
over $\mathbb C\{x\}$. This is a smooth ambient equation
containing the curve, so its embedding dimension is at
most two. The case where both odd parts vanish contradicts
the birational normalization parametrization.

\noindent\textbf{Step 6.} The estimates give $d_p\ge2m\ge6$. Equality at six
would require $m=3$, $\delta=2$, and $\tau_p=6$.
The valuation semigroup of a unibranch curve has exactly
$\delta$ gaps, as follows by filtering its local ring
and its normalization by order. Thus the two gaps here
are one and two, and we can choose generators with
initial orders three, four, and five. They form an
ambient coordinate system: their classes generate the
three-dimensional cotangent space, since the square
of the maximal ideal has order at least six.
In a defining relation, the quadratic monomials
$x^2$ and $xy$ have respectively the unique least
orders six and seven, so their coefficients vanish.
Every gradient with valuation three consequently has
its initial vector in the single $z$ direction
(from a possible $xz$ term). Two such initial
gradient vectors are dependent. Every two-by-two
minor therefore has valuation at least seven,
contradicting $\tau_p=6$. Hence $d_p\ge7$ in this case.

\noindent\textbf{Step 7.} Finally consider a nonplanar point with two branches
of multiplicities $m_1,m_2$. Two smooth curve branches
in a smooth threefold are contained in a smooth
surface. Indeed choose a common immersive coordinate
$x$ and write them as graphs
$(y,z)=(a_i(x),b_i(x))$. Divide their difference pair
by its common power of $x$; the equation
\[
 \frac{b_1-b_2}{x^k}(y-a_1)
 -\frac{a_1-a_2}{x^k}(z-b_1)=0
\]
has a nonzero linear derivative and contains both.
Thus in the nonplanar case $m_1+m_2\ge3$.
The absence of a linear defining equation gives,
on the two branches separately,
$\tau_p\ge2(m_1+m_2)$.
The reduced-curve normalization sequence gives
$\delta(C,p)=\delta_1+\delta_2+
\operatorname{length}\mathcal O_Y/(P_1+P_2)
\ge m_1+m_2-1$.
Also $r_p=m_1+m_2-2$.
It follows that
$d_p\ge2(m_1+m_2)+1\ge7$.
This proves the gap in all cases.

\noindent\textbf{Step 8.} At the unique two-branch point, a planar singularity
has defect five if its delta invariant is one, and
at least ten otherwise. A nonplanar point has defect
at least seven. Every other singular point is
unibranch and contributes at least five if planar,
and at least seven if nonplanar. Smooth points
contribute zero. Thus the sum of the defects can
be five or at least seven, but never six.
The required value $\chi(\mathcal O_{C_{\rm sch}})=-1$
would give exactly six. This contradiction excludes
$e(S)=1$.

\noindent\textbf{Step 9.} The final statement uses P4.26: the sum of all the
nonnegative fibre Euler numbers is at most two.
The rational-resolution case is strictly positive,
and we have just excluded one, leaving exactly two.
\end{proof}


\subsection{The conductor graph excludes Euler number one in the rational case}

\begin{proposition}[P4.35]\label{night:P4.35}\label{stmt:P4.35}
\StatementRecord{night:P4.35}{proved}{Historical events: GPT-6 Astra (OpenAI), ChatGPT Pro mode, effort Pro, S6\_06\_BLIND, 16--18 November 2026: targeted proof check. Also checked in this writing. Pre-continuation writing events: GPT-6 Astra (OpenAI), ChatGPT Pro mode, effort Pro, S6\_10\_PAPERS, 18 November 2026: targeted proof check. Continuation metadata events: GPT-6 Astra (OpenAI), ChatGPT Pro mode, effort Pro, S6\_10\_PAPERS continuation, 19 November 2026: used and rechecked in this run; status-vocabulary reconciliation; metadata reconciliation only; no new mathematical check. \emph{Pre-continuation recorded proof scope.} The night-run re-derivation was by GPT-6 Astra (OpenAI, ChatGPT Pro mode, effort Pro), 16--18 November 2026. Also checked in this writing by GPT-6 Astra (OpenAI), ChatGPT Pro mode, effort Pro, 18 November 2026. The pairwise-intersection two-label argument forces a conductor graph cycle. Rational exceptional fibres were checked by the ideal sequence and coherent relative dimension one, and contraction of their simply connected trees preserves the conductor fundamental group. The Euler-one case is excluded without a transverse-type assumption}{\ref{night:P4.26}; \ref{night:P4.33}; \ref{night:P4.8}; \ref{night:P4R.6}}

No irreducible support with rationally resolved normalization
has Euler number one. Consequently, in a fibration whose
supports are all irreducible or normal crossing, every
irreducible singular support with rationally resolved
normalization has Euler number exactly two.

More generally, an irreducible support cannot have reduced
image conductor equal to a rational curve with exactly one
point-identification cycle when its reduced normalization
conductor has no graph cycles. This excludes that conductor
case for every ruled-resolution double-horizontal pole tree.
\end{proposition}
\begin{proof}
\noindent\textbf{Step 1.} We first prove the graph assertion independently of
the transverse singularity type and of ramification.
Let $D\subset N$ be the reduced normalization conductor,
and suppose its point-identification graph is a tree.
Let $c$ be its number of normalized compact components.
Suppose its image $C$ is irreducible with rational
normalization and exactly one point-identification
cycle. There is then a unique point $p\in C$ with
two analytic branches; all other points are unibranch.
Denote the two preimages of $p$ on $C^\nu$ by
$p_+,p_-$. Every normalized component of $D$ maps
finitely and surjectively to $C^\nu$.

\noindent\textbf{Step 2.} The finite set $\nu^{-1}(p)$ has $n$ points and
parametrizes the local irreducible hypersurface
branches of $S$ at $p$. Every pair of such branches
intersects along a curve: two nonunit local
hypersurface equations in a regular local ring of
dimension three generate an ideal of height at most
two. Their common curve lies in $C$. Consequently
the nonempty subsets of $\{+,-\}$ recording which
branches of $C$ are contained in each hypersurface
branch pairwise intersect. Such a family of subsets
of a two-element set has a common element. If a
singleton occurs its element belongs to every set;
if no singleton occurs every set is the whole set.
Interchanging the labels, every point of
$\nu^{-1}(p)$ therefore has at least one conductor
branch mapping to the plus branch of $C$.

\noindent\textbf{Step 3.} Let $r_+$ and $r_-$ count the points on $D^\nu$
over $p_+$ and $p_-$ respectively. The preceding
paragraph gives $r_+\ge n$. Surjectivity of the
map from each of the $c$ normalized components
gives $r_-\ge c$. The contribution of the points
over $p$ to the point-identification graph of $D$
is exactly
\[
 \sum_{q\in\nu^{-1}(p)}(r_q-1)
       =r_++r_--n\ge c.
\]
For a connected reduced compact curve with $c$
normalized components, the cycle rank of this graph
is $1-c+\sum_q(r_q-1)$. All other contributions
are nonnegative. Its rank is therefore at least
one, contradicting the tree hypothesis.
This proves the graph assertion.

\noindent\textbf{Step 4.} Now let the normalization have rational resolution
and suppose $e(S)=1$. The normal-crossing alternative
has Euler number two by P4.8. In the other alternative
P4.26 gives a connected simply connected resolved
pole support $A\subset T$, and
\[
 \pi_1(S)=\pi_1(C)=F_b,\qquad b\le1,
 \qquad e(C)=c_C+1-b.
\]
Its normalized components are rational.
Since $e(S_{\rm reg})\ge0$, Euler number one forces
$c_C=b=1$ and $e(S_{\rm reg})=0$.
It remains to verify the tree property for the
actual conductor $D\subset N$, allowing singular $N$.

\noindent\textbf{Step 5.} P4.33 makes every singularity of $N$ rational.
Every connected exceptional fibre of $T\to N$
is simply connected. Here is the precise curve
calculation. Over a small Stein neighbourhood,
$R^1\rho_{N*}\mathcal O_T=0$ gives
$H^1(\mathcal O_T)=0$ on its inverse image.
This proper map has every fibre of dimension at
most one over that Stein neighbourhood. Its coherent
cohomological dimension is consequently at most
one, so the second higher direct image of every
coherent sheaf vanishes. Applied to the ideal of the reduced
exceptional curve $E$, Leray and the Stein property
give $H^2(\mathcal O_T(-E))=0$ there. Its ideal
sequence consequently gives $H^1(E,\mathcal O_E)=0$.
The curve normalization exact sequence now implies
that all normalized components are rational and
that their point-identification graph is a tree.
Thus $E$ is simply connected.

\noindent\textbf{Step 6.} By P4R.6, $A$ is the full inverse image of $D$.
The proper map $A\to D$ is the quotient which
collapses these connected exceptional curves to
points and is an isomorphism elsewhere.
Compatible triangulation makes these curve
inclusions cofibrations. Collapsing a connected
simply connected subcomplex preserves fundamental
groups by van Kampen. Collapsing the finitely many
exceptional fibres therefore gives $\pi_1(D)=0$.
Each normalized component of $D$ is rational,
being birational to a strict pole component in
$T$. Hence its point-identification graph is a
tree. The already proved graph assertion contradicts
the description of $C$. This excludes $e(S)=1$.

\noindent\textbf{Step 7.} For the final more general application, the
resolved pole support in the ruled double-section
case is $H$ with rational trees attached.
Every exceptional fibre is connected and is a
union of its components. If it does not contain
$H$, it is a connected rational subtree. If it
contains $H$, all the remaining components form
rational trees attached to that component once.
Contracting such a connected subcurve cannot
create a cycle in the graph: contraction of a
connected subgraph of a tree again has a tree
as quotient. The only positive-genus vertex,
if contracted, is replaced by a point, and this
also introduces no graph edge. Thus the image
conductor $D$ has no graph cycles, and the same
contradiction applies.

\noindent\textbf{Step 8.} Finally, for a fibration in the two requested
support classes, P4.26 makes the total sum of
nonnegative fibre Euler numbers at most two.
The rational-resolution case is strictly positive;
excluding one leaves exactly two.
\end{proof}


\subsection{Eliminating all positive-Euler Hopf resolutions}

\begin{proposition}[P4.36]\label{night:P4.36}\label{stmt:P4.36}
\StatementRecord{night:P4.36}{proved}{Historical events: GPT-6 Astra (OpenAI), ChatGPT Pro mode, effort Pro, S6\_06\_BLIND, 16--18 November 2026: targeted proof check. Also checked in this writing. Pre-continuation writing events: GPT-6 Astra (OpenAI), ChatGPT Pro mode, effort Pro, S6\_10\_PAPERS, 18 November 2026: targeted proof check; used and rechecked in this run. Continuation metadata events: GPT-6 Astra (OpenAI), ChatGPT Pro mode, effort Pro, S6\_10\_PAPERS continuation, 19 November 2026: used and rechecked in this run; status-vocabulary reconciliation; metadata reconciliation only; no new mathematical check. \emph{Pre-continuation recorded proof scope.} The night-run re-derivation was by GPT-6 Astra (OpenAI, ChatGPT Pro mode, effort Pro), 16--18 November 2026. Also checked in this writing by GPT-6 Astra (OpenAI), ChatGPT Pro mode, effort Pro, 18 November 2026. The Hopf finite-cover reduction, saturated derivative line and rank-two image, Chern identity, integral quadratic blowup increment, and equality propagation through every exceptional curve were re-derived. Minimality forces the actual normalization to be smooth and minimal; the Hopf foliation normal-form input remains a separately cited dependency}{\ref{night:P4.24}; \ref{night:P4R.6}}

Every irreducible singular support in Part~4.2 whose
normalization has a Hopf minimal model has smooth minimal
Hopf normalization and Euler number zero. The assertion
does not assume initially that the normalization is smooth
or that its Euler number vanishes.
\end{proposition}
\begin{proof}
\noindent\textbf{Step 1.} Let $T\to N\to S\subset Y$ be the minimal normalization
resolution, and let $\pi:T\to T_0$ contract its $k$ blowups
to a minimal Hopf surface. We first take a finite unramified
primary Hopf cover $T'_0\to T_0$ and pull back the actual
sequence of point blowups. This gives a finite unramified
map $T'\to T$ and a blowup morphism $T'\to T'_0$.
Compose with $T\to Y$. All Chern calculations below take
place on this explicit smooth cover; no cover of the
singular normalization $N$ is asserted. To avoid a second
set of symbols we write $T,T_0,k,\rho$ for these covered
objects until the last paragraph.

\noindent\textbf{Step 2.} The primary Hopf normal form and proof of P4.24 give a
nowhere-zero holomorphic foliation on $T_0$ tangent to
every compact curve. Its saturated pullback to $T$ is a
line subbundle away from finitely many simple saddle
zeros. It is tangent to every strict conductor component
and every created exceptional curve. The actual conductor
centres lie on these invariant curves. Every blowup centre
is either a regular point of this foliation or one of its
saddles. Let $\varepsilon=0$ in the first case and
$\varepsilon=1$ in the second. The same explicit two-chart
calculation as in P4.24 gives
\[
 c_1(L)=-K_T+V,\qquad
 V'=\mu^*V+\varepsilon E,
 \qquad K_T^2=-k,\qquad K_TV-V^2=0.
\]
Topologically trivial line-bundle factors from $T_0$ do
not change these integral intersection calculations.

\noindent\textbf{Step 3.} Put $B=\rho^*TY$. Its first and second rational Chern
classes vanish: $b_2(Y)=0$ gives $c_1(TY)=0$ rationally,
and Poincar\'e duality gives $H^4(Y;\mathbb Q)=0$ and
hence $c_2(TY)=0$. The derivative $L\to B$ can now have
a divisorial zero, since $\rho$ contracts exceptional
curves. Denote that effective integral divisor by $Z$.
Dividing by it gives the saturated line injection in
codimension one
\[
 L'=L(Z)\longrightarrow B,
 \qquad c_1(L')=-K_T+V+Z.
\]
The quotient $Q=B/L'$ is torsionfree of rank two.
Write $z_0\ge0$ for the length of $Q^{**}/Q$.
The double dual is locally free on the smooth surface,
and the Chern sequence gives
\[
 c_2(Q^{**})=(-K_T+V+Z)^2-z_0.
\]

\noindent\textbf{Step 4.} Let $E_B\subset B$ be the saturated rank-two image
of $d\rho:TT\to B$. Its determinant is
$\det E_B=\omega_T^{-1}(R)$, where the effective
integral divisor $R$ is the divisorial zero of the
determinant map. Since $E_B$ is saturated in $B$,
the divided foliation line $L'$ is contained in it.
The image of $E_B/L'$ in $Q^{**}$ is saturated in
codimension one. Its rank-one double dual
$J=(E_B/L')^{**}$ is a line bundle. The map into
$Q^{**}$ extends across the finitely many omitted
points by Hartogs, since both sheaves are locally
free there. Its saturation is therefore a line bundle
of first Chern class $R-V-Z$; the rank-one quotient
is an ideal sheaf of a zero-dimensional scheme
times a line bundle of first Chern class $K_T-R$.
Let its length be $z_1\ge0$. Thus
\[
 (-K_T+V+Z)^2-z_0
       =(R-V-Z)(K_T-R)+z_1.
\]
Using $K_T^2=-k$ and $K_TV-V^2=0$, this is precisely
\[
 \mathcal B(R,Z)
 :=R^2-K_TR-VR-ZR-K_TZ+2VZ+Z^2
 =k+z_0+z_1.                                      \tag{P4.36.1}
\]

\noindent\textbf{Step 5.} We now evaluate the left side through the actual
point blowups, retaining all integral coefficients.
Push $R$ and $Z$ down along the chosen sequence.
At one reversed blowup write
\[
 R'=\mu^*R+cE,\quad Z'=\mu^*Z+dE,
 \quad K'=\mu^*K+E,\quad V'=\mu^*V+\varepsilon E.
\]
The integers $c,d$ are not assumed positive.
Since $E^2=-1$ and a pulled-back divisor has zero
intersection with $E$, direct expansion gives
\[
 \mathcal B(R',Z')-\mathcal B(R,Z)
 =-c^2+(1+\varepsilon)c+cd+(1-2\varepsilon)d-d^2.
\]
Set $a=c-1$ and $b=d-(1-\varepsilon)$. For each
$\varepsilon\in\{0,1\}$ this is exactly
\[
 1-(a^2+b^2-ab)\le1.                            \tag{P4.36.2}
\]
The quadratic form is
$\frac12(a^2+b^2+(a-b)^2)$; it is nonnegative
and is zero exactly when $a=b=0$.
All divisor intersections on the initial Hopf
surface are zero because $b_2(T_0)=0$.
Summing (P4.36.2) therefore gives
$\mathcal B(R,Z)\le k$.
Equation (P4.36.1) forces equality, $z_0=z_1=0$,
and at every single blowup
\[
 c=1,\qquad d=1-\varepsilon.                    \tag{P4.36.3}
\]

\noindent\textbf{Step 6.} Every created exceptional curve has positive
coefficient in $Z$. At a regular centre, (P4.36.3)
adds one to its nonnegative pulled-back coefficient.
At a saddle centre it adds zero, but at least one
previously created exceptional curve passes through
that saddle. The original Hopf foliation has no
zero; the two-chart blowup calculation places each
new saddle on the new exceptional curve. Induction
therefore gives a positive pulled-back coefficient
at every later saddle blowup as well.

\noindent\textbf{Step 7.} The foliation is tangent to each exceptional curve.
If its image under $\rho$ were a curve, the derivative
along its tangent line would be nonzero at a generic
point: a nonconstant holomorphic map of smooth
normalized curves in characteristic zero is
generically immersive. Positive divisorial order
of $d\rho|_L$ contradicts that fact. Thus every
created exceptional curve is contracted by $\rho$,
and hence by the normalization resolution, since
the normalization map is finite.

\noindent\textbf{Step 8.} Return to the original, possibly secondary Hopf
surface. Every exceptional curve downstairs has
an exceptional lift in the finite unramified cover;
the preceding argument makes its image a point.
Thus all exceptional curves of the original
$\pi:T\to T_0$ would be contracted by the minimal
resolution $T\to N$. If $k>0$, the exceptional curve
of its last blowup is a smooth rational curve of
self-intersection $-1$, which a minimal resolution
cannot contract. Hence $k=0$.
There is no further exceptional contraction from
a minimal Hopf surface: every compact curve has
self-intersection zero by $b_2=0$, whereas a
nonempty exceptional configuration has negative
definite intersection matrix, as proved in P4R.6.
Therefore $T=N=T_0$ is a smooth minimal Hopf
surface. P4.24 gives $e(S)=0$ and its complete
elliptic-conductor description.
\end{proof}


\subsection{Excluding the remaining class-VII minimal models}

\begin{proposition}[P4.37]\label{night:P4.37}\label{stmt:P4.37}
\StatementRecord{night:P4.37}{proved}{Historical events: GPT-6 Astra (OpenAI), ChatGPT Pro mode, effort Pro, S6\_06\_BLIND, 16--18 November 2026: targeted proof check. Also checked in this writing. Pre-continuation writing events: GPT-6 Astra (OpenAI), ChatGPT Pro mode, effort Pro, S6\_10\_PAPERS, 18 November 2026: targeted proof check; used and rechecked in this run. Continuation metadata events: GPT-6 Astra (OpenAI), ChatGPT Pro mode, effort Pro, S6\_10\_PAPERS continuation, 19 November 2026: used and rechecked in this run; status-vocabulary reconciliation; metadata reconciliation only; no new mathematical check. \emph{Pre-continuation recorded proof scope.} The night-run re-derivation was by GPT-6 Astra (OpenAI, ChatGPT Pro mode, effort Pro), 16--18 November 2026. Also checked in this writing by GPT-6 Astra (OpenAI), ChatGPT Pro mode, effort Pro, 18 November 2026. The class-VII Chern identity, negative-definite initial quadratic form, and at-most-one contribution of every blowup were re-derived. The resulting strict inequality contradicts the ambient derivative identity. The exact Favre and Dloussky-Oeljeklaus source loci for the initial foliation and curve basis are tracked separately}{\ref{night:P4.23}; \ref{night:P4.26}; \ref{night:P4.36}}

No irreducible singular support in Part~4.2 has normalization
whose minimal resolution has a class-VII minimal model with
positive second Betti number. Consequently every class-VII
normalization occurring here is a smooth minimal Hopf surface
and has Euler number zero.
\end{proposition}
\begin{proof}
\noindent\textbf{Step 1.} The numerically anticanonical classification used in P4.23
and P4.26 reduces this question to an intermediate Kato
minimal model $T_0$: the Inoue--Hirzebruch and parabolic
alternatives were excluded there by their reduced
normal-crossing conductor and the Chern identity.
Write $n=b_2(T_0)>0$, let $A_0$ be its maximal reduced
curve divisor, and let $D_0$ be the effective integral
pole divisor pushed down from the normalization resolution.
Thus $K_{T_0}+D_0$ is numerically zero. The support of
$D_0$ is $A_0$.

\noindent\textbf{Step 2.} We state the precise additional primary input.
Favre, \emph{Dynamique des applications rationnelles},
Theorem~1.2.24(1)--(2), printed p.~35, says that the
$n$ rational curve classes span $H^2(T_0;\mathbb Q)$
and that their intersection form is negative definite.
Theorem~1.2.30, printed p.~42, and its proof on
pp.~44 and~46 give a holomorphic foliation with conormal
line
\[
 N_{\mathcal F}^*=\mathcal O_{T_0}(-A_0)
                       \otimes\text{a flat line}.
\]
Theorem~1.2.32 and its proof on p.~46 show that every
component of $A_0$ is invariant and that its singularities
are precisely the intersection points, with local
first integrals $z^\alpha w^\beta$, $\alpha,\beta>0$.
The invariant-curve and order-one assertions are also
Dloussky--Oeljeklaus, Theorem~3.1, pp.~1528--1529.
The actual source URLs and exact theorem locators are
recorded in the bibliography and literature log.

\noindent\textbf{Step 3.} It follows from the determinant of tangent and normal
lines that the foliation tangent line has numerical
first Chern class $-K_{T_0}-A_0$. The local generator
at an intersection is
$\beta z\partial_z-\alpha w\partial_w$.
Blowing up a regular point or one of these saddles
again gives only regular points and saddles with
opposite-sign nonzero eigenvalues. Thus every actual
blowup centre has the same two possibilities
$\varepsilon=0,1$ as in P4.36, with the same formulas.

\noindent\textbf{Step 4.} Let $T\to T_0$ contain $k$ blowups. Write
$c_1(L)=-K_T+V$ for the saturated foliation line;
initially $V_0=-A_0$, and
$V'=\mu^*V+\varepsilon E$ at a blowup.
The reduced divisor $A_0$ is a rational cycle with
rational trees attached. Its normalization has
$n$ rational components and $n$ nodes, including
self-nodes, so $\chi(\mathcal O_{A_0})=0$.
Adjunction, or the divisor exact sequence on $T_0$,
therefore gives
\[
 (K_{T_0}+A_0)A_0=0.
\]
Hence $K_{T_0}V_0-V_0^2=0$.
Under each blowup this quantity changes by
$\varepsilon^2-\varepsilon=0$.
Since $K_{T_0}^2=-n$, we obtain
\[
 K_TV-V^2=0,\qquad K_T^2=-n-k.
\]

\noindent\textbf{Step 5.} Define $B=\rho^*TY$, the divisorial determinant zero
$R$, the divisorial foliation zero $Z$, and the
nonnegative lengths $z_0,z_1$ exactly as in P4.36.
Its saturated sheaf calculation, including the
Hartogs extension of the rank-one double dual,
now gives
\[
 \mathcal B(R,Z)
 =R^2-K_TR-VR-ZR-K_TZ+2VZ+Z^2
 =n+k+z_0+z_1.                                  \tag{P4.37.1}
\]
Every reversed blowup changes $\mathcal B$ by at
most one, by the unrestricted integer identity
(P4.36.2). It remains to evaluate the initial term.

\noindent\textbf{Step 6.} Push $R,Z$ down to effective divisors $R_0,Z_0$
on $T_0$. Their supports are contained in $A_0$:
the derivative is generically injective away from
the resolved conductor, and in any event all
compact curves of the intermediate surface are
components of its maximal divisor. Put $D=D_0$,
$A=A_0$. Using $K=-D$ numerically and $V=-A$ gives
\[
 \mathcal B_0
 =R_0^2+DR_0+AR_0-Z_0R_0+DZ_0-2AZ_0+Z_0^2.
\]
For divisor classes define
$q(X,Y)=X^2-XY+Y^2$.
Adjunction above gives $DA=A^2$, while $D^2=-n$.
Expansion consequently gives the exact completion
\[
 \mathcal B_0
       =n+q(R_0+D,Z_0+D-A).                     \tag{P4.37.2}
\]
The form on divisor classes is negative definite.
Moreover
\[
 q(X,Y)=\tfrac12\bigl(X^2+Y^2+(X-Y)^2\bigr),
\]
so it is nonpositive, and it vanishes only for
$X=Y=0$. The first argument $R_0+D$ is nonzero:
it is a nonzero effective combination of the
$n$ independent curve classes. Thus
$\mathcal B_0<n$.
After the $k$ blowups we still have
$\mathcal B(R,Z)<n+k$, contradicting (P4.37.1).
This excludes the intermediate minimal model.

\noindent\textbf{Step 7.} All class-VII cases with positive second Betti
number have now been excluded. The zero-Betti
Inoue case has no conductor curves and was
excluded in P4.26. P4.36 treats the remaining
Hopf case and proves the final assertion.
\end{proof}


\subsection{Two exact local exclusions for ordinary-double curves}

\begin{proposition}[P4.38]\label{night:P4.38}\label{stmt:P4.38}
\StatementRecord{night:P4.38}{proved}{Historical events: GPT-6 Astra (OpenAI), ChatGPT Pro mode, effort Pro, S6\_06\_BLIND, 16--18 November 2026: targeted proof check. Also checked in this writing. Pre-continuation writing events: GPT-6 Astra (OpenAI), ChatGPT Pro mode, effort Pro, S6\_10\_PAPERS, 18 November 2026 05:10:56.410177+00:00: computation or formal certificate; targeted proof check. Continuation metadata events: GPT-6 Astra (OpenAI), ChatGPT Pro mode, effort Pro, S6\_10\_PAPERS continuation, 19 November 2026: used and rechecked in this run; status-vocabulary reconciliation; metadata reconciliation only; no new mathematical check. \emph{Pre-continuation recorded proof scope.} The night-run re-derivation was by GPT-6 Astra (OpenAI, ChatGPT Pro mode, effort Pro), 16--18 November 2026. Also checked in this writing by GPT-6 Astra (OpenAI), ChatGPT Pro mode, effort Pro, 18 November 2026. proved. Full planar-singular and cusp/tacnode-plus-axis ordinary-double exclusion, including the actual boundary comparison. Re-derived P\textasciicircum{}(2)=P\textsuperscript{2} for the reduced complete-intersection plane curve and the exact Morse splitting f=Z\textsuperscript{2}+h\textsuperscript{2}u. A morsification and nonzero critical-value shift give Hessian block 2h\_s Hess(h\_s), hence a Morse sphere. For P=(xz,yz,y\textsuperscript{2}−x\textasciicircum{}k), k=3,4, directly intersected (z,h)\textsuperscript{2} with (x,y)\textsuperscript{2}, obtaining the five displayed generators. The axis quadratic determinant after h\ensuremath{\to}h−sx\textsuperscript{2} has divided factor F=(-sA+Bz)(A+Dz)−C\textsuperscript{2}z\textsuperscript{2}/4, with F\_z=AB and F\_s=−A\textsuperscript{2}. Its simple nonzero root gives the actual coordinate Y\textsuperscript{2}+x\textsuperscript{2}Z. Followed the compact sphere of square −4 in P4.16 and both nested proper boundary transports; the contradiction uses the null-homotopy of every sphere in the actual nearby T4. The unchanged critical\_curve\_node.py replay passed. Review provenance: New derivation after reading the supplied proof; not a blind review and not an independent human check}{\ref{night:P4.16}; \ref{night:P4.20}}

An actual irreducible support in Part~4.2 cannot have a singular
planar reduced critical curve at a point at which every generic
transverse germ is an ordinary double point. Nor can its reduced
critical curve there be the transverse union of a plane cusp
and a smooth branch, again with ordinary-double transverse germs
on both components. The latter assertion also holds with the
cusp replaced by a plane tacnode.
These exclusions allow singular normalization.
\end{proposition}
\begin{proof}
\noindent\textbf{Step 1.} We retain the boundary transport of P4.20 throughout. At a fixed
compact boundary, ordinary-double transverse germs form an open
condition. A sufficiently small deformation retaining that
condition gives the same nearby nonzero-level boundary
submersion. The finite-critical-value argument in that proof
then transports any compact vanishing sphere from the deformed
germ into the preceding nearby fibre. When a coordinate change
or a unit requires a smaller ball, we perform this construction
on that ball first, and then transport through the earlier
larger ball. Thus every sphere constructed below reaches the
original nearby four-torus.

\noindent\textbf{Step 2.} First suppose that the reduced critical curve has ideal
$P=(z,h(x,y))$, where $h$ is reduced and singular at the
origin. The equation $f$ lies in $P^{(2)}$: its generic
transverse ordinary double point vanishes to order two on
each component. The complete-intersection square $P^2$ is
unmixed, so $P^{(2)}=P^2$. We therefore have the exact expression
\[
 f=A z^2+Bzh+Ch^2.
\]
Adding a small constant multiple of $z^2$ preserves the critical
curve and the ordinary-double boundary condition. We choose it
so that the $z^2$ coefficient at the origin is nonzero. The
Morse lemma in $z$, with parameters $x,y$, then gives
\[
 f=Z^2+h(x,y)^2u(x,y).
\]
Indeed the critical $z$ coordinate vanishes on $h=0$ and
is divisible by $h$, so its substitution has the displayed
square factor. Add a small constant to $u$ and shrink the
ball so that $u$ is a unit. Absorbing its square root in $h$
gives $Z^2+H(x,y)^2$, with the same singular reduced zero
curve. A reduced singular plane curve has an isolated critical
point. A generic small morsification $H_s$ has a Morse critical
point, and a small additive constant makes all these critical
values nonzero. At any such point $Z^2+H_s^2$ has a
nondegenerate isolated critical point: its plane Hessian is
$2H_s\operatorname{Hess}H_s$, and its $Z$ Hessian is $2$.
Its vanishing sphere has self-intersection $-2$.
The moving curve $H_s=0$ remains transverse to the chosen
boundary, with nondegenerate normal quadratic form there.
The preceding boundary transport therefore contradicts
$\pi_2(T^4)=0$.

\noindent\textbf{Step 3.} For the second assertion write
\[
 h=y^2-x^k,\qquad P=(xz,yz,h),\qquad k=3\text{ or }4.
\]
This is the union of the plane curve $z=h=0$ and the
transverse line $x=y=0$. Its symbolic square is
\[
 P^{(2)}=(z,h)^2\cap(x,y)^2
  =(x^2z^2,xyz^2,y^2z^2,zh,h^2).
\]
For $k=4$, the first equality uses the two smooth branches of
$h=0$ separately; their complete-intersection square is unmixed
and hence is still $(z,h)^2$. To verify the intersection
formula directly, write an element of $(z,h)^2$ as
$az^2+bzh+ch^2$. The last two terms belong to $(x,y)^2$,
and $z$ is a nonzero divisor modulo this ideal. The first term
belongs to it exactly when $a\in(x,y)^2$.
Consequently
\[
 f=Azh+z^2(Bx^2+Cxy+Dy^2)+Gh^2.
\]
Small constant changes in $A,B$ retain the boundary condition
and make $A(0)B(0)\ne0$. Replace $h$ in this expression by
$h_s=y^2-x^k-sx^2$. The new critical curve contains
$z=h_s=0$ and $x=y=0$, and its boundary ordinary-double
germs persist. On the latter axis the transverse quadratic
matrix has coefficients
\[
 a=-sAz+Bz^2,\qquad b=Cz^2,\qquad c=Az+Dz^2,
\]
with off-diagonal entry $b/2$. Here the four coefficient
functions are restricted to that axis. Its determinant
divided by $z^2$ is
\[
 F(s,z)=(-sA+Bz)(A+Dz)-C^2z^2/4.
\]
At $(s,z)=(0,0)$ we have $F_z=AB\ne0$ and
$F_s=-A^2$. The implicit-function theorem gives a simple root
$z(s)=sA(0)/B(0)+O(s^2)$, nonzero for small nonzero $s$.
There the $y^2$ coefficient is nonzero, the Hessian has rank
one, and its determinant has a simple zero in the axis
direction. The parametrized Morse lemma gives
$Y^2+x^2H(x,z)$, with $H_z\ne0$. The exact coordinate
change $Z=H(x,z)$ yields $Y^2+x^2Z$.
P4.16 gives its compact spherical class of square $-4$.
Nested boundary transport once more places this class in
the original nearby four-torus, a contradiction.
\end{proof}

The exact quadratic matrix, determinant, and the two nonzero
derivatives in this proof are verified by
\texttt{code/part4/critical\_curve\_node.py}; its saved output
reports PASS.


\subsection{A coefficient-two exceptional curve over an ordinary-triple germ}

\begin{proposition}[P4.39]\label{night:P4.39}\label{stmt:P4.39}
\StatementRecord{night:P4.39}{proved}{Historical events: GPT-6 Astra (OpenAI), ChatGPT Pro mode, effort Pro, S6\_06\_BLIND, 16--18 November 2026: targeted proof check. Also checked in this writing. Pre-continuation writing events: GPT-6 Astra (OpenAI), ChatGPT Pro mode, effort Pro, S6\_10\_PAPERS, 18 November 2026 05:10:56.410177+00:00: computation or formal certificate; targeted proof check. Continuation metadata events: GPT-6 Astra (OpenAI), ChatGPT Pro mode, effort Pro, S6\_10\_PAPERS continuation, 19 November 2026: used and rechecked in this run; status-vocabulary reconciliation; metadata reconciliation only; no new mathematical check. \emph{Pre-continuation recorded proof scope.} The night-run re-derivation was by GPT-6 Astra (OpenAI, ChatGPT Pro mode, effort Pro), 16--18 November 2026. Also checked in this writing by GPT-6 Astra (OpenAI), ChatGPT Pro mode, effort Pro, 18 November 2026. proved. Full coefficient-two exceptional-curve exclusion for uniform ordinary-triple transverse germs; arithmetic-genus wording repair recorded. The ambient blowup residue comparison gives D1+(d−2)e and 2=(d−2)n+a. Re-derived positivity using the negative exceptional intersection matrix and arithmetic adjunction. The two received expressions 2g(H)−2 must be 2p\_a(H)−2 until smoothness is deduced: positive geometric genus gives p\_a\ensuremath{\geq}1, and coefficient one forces p\_a=1, hence g=1 and zero delta, so H is smooth elliptic. For d=4 the quartic Hessian vanishes along an actual critical arc, giving multiplicity at least3 and excluding every positive-genus reduced tangent component. For d=3, depressed Weierstrass preparation and P\textasciicircum{}(3)=P\textsuperscript{3} give u\textsuperscript{3}+h\textsuperscript{2}Bu+h\textsuperscript{3}C. Re-derived both the double-line quadratic normalization, whose discriminant has degree\ensuremath{\leq}2, and the isolated-point chart v\textasciicircum{}e+xG. At a simple zero the actual coordinate W=G gives v\textasciicircum{}e+xW; its cyclic-quotient resolution is rational. The three pole cases and all first-blowup degree bounds passed the unchanged exact programs. This proves the statement with the explicitly supplied adjunction correction. Review provenance: New derivation after reading the supplied proof; not a blind review and not an independent human check. The receiving adjunction calculation now uses arithmetic genus until its equality proves that the exceptional curve is smooth elliptic}{\ref{night:P4R.6}}

Let $S=(f=0)$ be a reduced hypersurface germ in a smooth
threefold, with one-dimensional reduced critical locus $C$.
Suppose that every generic transverse germ along $C$ is an
ordinary triple point. On the minimal resolution $T$ of its
normalization, let $D$ be the integral total pole divisor
defined by the residue of an ambient volume form divided by
$f$. No exceptional curve of positive genus has coefficient
two in $D$.

In particular, the contracted horizontal section of a
positive-genus ruled resolution in Part~4.2 cannot lie over
a point where all local critical-curve branches have generic
type $D_4$. This assertion does not replace that local
hypothesis by a mixture of transverse types.
\end{proposition}
\begin{proof}
\noindent\textbf{Step 1.} We blow up the ambient origin, write $S_1$ for the strict
transform and $N_1$ for its normalization, and take the
minimal resolution $T_1$ of $N_1$. The proper map $N_1\to S$
lifts to the normalization $N$ by its universal property.
Thus $T_1\to N$ is a resolution and dominates $T$ by a
sequence of point blowups. The strict transform of a proposed
positive-genus exceptional curve $H$ remains present and has
the same pole coefficient two.

\noindent\textbf{Step 2.} Put $d=\operatorname{mult}_0f$, and denote the ambient
exceptional divisor by $E$. Actual adjunction gives
\[
 \omega_{S_1}=\pi^*\omega_S\otimes
        \mathcal O((2-d)E)|_{S_1}.
\]
Let $D_1$ be the total conductor divisor on $T_1$ relative
to $\omega_{S_1}$, and let $e$ be the Cartier pullback of
$E$. With the specified residue comparison, the total pole
of the original germ on $T_1$ is
\[
                    D_1+(d-2)e.                       \tag{4.39a}
\]
The divisor $D_1$ is effective. Its strict conductor
coefficients are positive integers. For exceptional curves
of the minimal resolution of $N_1$, the negative-matrix
argument of P4R.6 applies locally: $K_{T_1}F\ge0$, since
there is no smooth rational exceptional $(-1)$-curve;
adding the strict conductor retains nonnegative intersection.
The inverse exceptional intersection matrix has nonpositive
entries, so all exceptional coefficients are nonnegative.

\noindent\textbf{Step 3.} If $H$ is exceptional over $N_1$, its coefficient in $D_1$
is strictly positive. A zero coefficient would give
\[
 0=(K_{T_1}+D_1)H
   =2p_a(H)-2-H^2+D_1H>0.
\]
If that coefficient is one, adjunction gives
\[
 0=2p_a(H)-2+(D_1-H)H.
\]
The normalization of $H$ has positive genus, so $p_a(H)\ge g(H^\nu)\ge1$. Both terms in the last equality are nonnegative, hence $p_a(H)=1$ and $(D_1-H)H=0$. The normalization formula $p_a(H)=g(H^\nu)+\sum\delta(H,p)$ now gives $g(H^\nu)=1$ and zero delta invariant at every point. Thus $H$ is smooth elliptic and is disjoint from every other positive component of $D_1$. A zero-coefficient exceptional
curve cannot meet a positive one, by its own intersection
equation and $K_{T_1}F\ge0$. The connected exceptional
fibre is consequently $H$ alone and meets no strict conductor.
The entire germ of $S_1$ at its image is normal with an
isolated simple elliptic singularity. Indeed, any additional
hypersurface branch would intersect this branch in a curve
and would give strict conductor meeting $H$.

\noindent\textbf{Step 4.} Set $n=\operatorname{ord}_H(e)>0$ and
$a=\operatorname{coeff}_H(D_1)\ge0$. Formula (4.39a)
gives the exact identity
\[
                       2=(d-2)n+a.                    \tag{4.39b}
\]
Thus $d\ge5$ is impossible. For $d=4$ we have
$n=1,a=0$, and $H$ is not exceptional over $N_1$.
At its generic point the normalization is an isomorphism,
so $H$ is birational to a reduced component of the projective
quartic tangent cone. Every second derivative of $f$ vanishes
on $C$. Its leading term along any analytic arc of $C$
therefore makes that quartic have a point of multiplicity
at least three. An irreducible quartic has arithmetic genus
three and loses at least three there. A smooth elliptic cubic
component together with the remaining line could have
multiplicity at most two at the point. All components of
degree at most two are rational. Thus every reduced tangent
component is rational, excluding $H$ for $d=4$.

\noindent\textbf{Step 5.} We are left with $d=3$ and
\[
                   (n,a)=(2,0)\quad\hbox{or}\quad(1,1).
\]
Choose a coordinate $u$ with nonzero pure cubic initial
coefficient, apply Weierstrass preparation, and depress the
monic cubic. Multiplication by the preparation unit does not
affect Hessian vanishing on $C$, because $f=df=0$ there.
We obtain
\[
 f=\text{unit}\cdot\bigl(u^3+b(x,z)u+c(x,z)\bigr).
\]
The second $u$ derivative is $6u$, so $C$ lies in the smooth
surface $u=0$. Write its reduced ideal as $(u,h(x,z))$.
The complete-intersection cube is unmixed and equals its
symbolic cube. Generic order three along $C$ and coefficient
comparison therefore give the exact expression
\[
 f=\text{unit}\cdot
       \bigl(u^3+h^2B(x,z)u+h^3C(x,z)\bigr).           \tag{4.39c}
\]
If $C$ is singular, $\operatorname{ord}h\ge2$ and the
tangent cone is the triple plane $u^3$. The degree-four
term restricts to zero on this plane. After the ambient
blowup the entire reduced exceptional line is critical,
so none of its points is an isolated normal singularity.
If $C$ is smooth, take $x=h$. Every homogeneous term $f_j$
then has degree at least three in $u,x$. The cubic tangent
cone is a binary cubic in these two variables, hence three
lines, counted with multiplicities, through the tangent
direction of $C$.

\noindent\textbf{Step 6.} First suppose that $H$ dominates a tangent-cone line.
For $(n,a)=(2,0)$ the normalization is generically an
isomorphism, and $H$ is birational to that line. For
$(n,a)=(1,1)$, a transverse normalization branch has conductor
exponent one. It is smooth and meets exactly one other smooth
branch transversely: the intrinsic conductor exponent of a
singular plane branch is at least two, and each other branch
contributes a positive intersection multiplicity. Thus the
generic transverse germ of $S_1$ is a node. A multiplicity-one
tangent line is generically smooth and cannot have $a=1$.
For a multiplicity-three line, the two geometric branches of
that node have exceptional orders whose weighted sum is three.
A connected double cover would give equal orders and an even
sum. Thus the order-one branch has degree one over the line
and is rational.

\noindent\textbf{Step 7.} The remaining possibility is a double line together with a
distinct simple line. Here $C$ is smooth by (4.39c). Make a
linear change so that $f_3=xu^2$. On the $x$ chart of the
blowup put $u=xv,z=xw$. The strict equation is
\[
 v^2+x f_4(1,v,w)+x^2f_5(1,v,w)+\cdots.
\]
An ordinary double curve on $x=v=0$ requires
$f_4(1,0,w)=0$ identically. Its transverse quadratic form is
\[
 v^2+xvB(w)+x^2C(w),
 \qquad \deg B\le1,\quad\deg C\le2.                  \tag{4.39d}
\]
The degree bounds follow from the degree-at-least-three
condition in $u,x$ on every $f_j$: here
$B=\partial_v f_4(1,0,w)$ and $C=f_5(1,0,w)$.
The normalization over the exceptional line adjoins a root
of $q^2+B(w)q+C(w)$. Its discriminant has degree at most
two. The resulting quadratic extension of $\mathbb C(w)$
is branched at at most two points, including infinity.
Riemann--Hurwitz makes its compact normalization rational;
the components of a split cover are rational as well.

\noindent\textbf{Step 8.} Finally suppose that $H$ is exceptional over a point of $N_1$.
We must have $n=a=1$, and the first part of the proof would
make $S_1$ normal with an isolated simple elliptic singularity
there. This cannot happen. For singular $C$ the whole
exceptional line is critical, as already proved. For smooth
$C$ its strict transform passes through the common point of
the tangent lines, so that point is not an isolated normal
singularity. Away from it, at a line of multiplicity
$e=2$ or $3$, divide by the nonzero factor from any other
line. The strict equation has the form
\[
                         v^e+xG(v,w,x),
\]
where $G(0,w,0)$ is affine linear in $w$, by the same
normal-degree bound. If this function is nonzero at a point,
the surface is smooth. If it is identically zero, the entire
line is critical. At a simple zero, $G_w$ is a unit and the
exact coordinate change $W=G(v,w,x)$ gives $v^e+xW$, a
rational double point of type $A_{e-1}$. A multiplicity-one
line is smooth away from intersections. There is therefore
no isolated simple elliptic singularity. All possibilities
for $H$ have been excluded.
\end{proof}

The positive-integer solutions of (4.39b), the monomial degree
bounds in (4.39d), and the forced fourth-order factor of the
homogeneous quadratic discriminant are checked in
\texttt{code/part4/d4\_pole.py}. Its saved output reports PASS.


\subsection{The excess-one curve calculation and the ordinary-double sector}

\begin{proposition}[P4.40]\label{night:P4.40}\label{stmt:P4.40}
\StatementRecord{night:P4.40}{proved}{Historical events: GPT-6 Astra (OpenAI), ChatGPT Pro mode, effort Pro, S6\_06\_BLIND, 16--18 November 2026: targeted proof check. Also checked in this writing. Pre-continuation writing events: GPT-6 Astra (OpenAI), ChatGPT Pro mode, effort Pro, S6\_10\_PAPERS, 18 November 2026 05:10:56.410177+00:00: computation or formal certificate; targeted proof check. Continuation metadata events: GPT-6 Astra (OpenAI), ChatGPT Pro mode, effort Pro, S6\_10\_PAPERS continuation, 19 November 2026: used and rechecked in this run; status-vocabulary reconciliation; metadata reconciliation only; no new mathematical check. \emph{Pre-continuation recorded proof scope.} The night-run re-derivation was by GPT-6 Astra (OpenAI, ChatGPT Pro mode, effort Pro), 16--18 November 2026. Also checked in this writing by GPT-6 Astra (OpenAI), ChatGPT Pro mode, effort Pro, 18 November 2026. proved. Full excess-one analytic curve-ring classification and its global ordinary-double application. For A\ensuremath{\subset}R=Cdiag+\ensuremath{\oplus}t\_i C\{t\_i\} with length(R/A)=1, the quotient is killed by m\_A. Writing a missing element v with v\textsuperscript{2}=a+cv, c\ensuremath{\neq}0 would let the implicit-function theorem solve a root in A and force v into A; therefore m\_R\textsuperscript{2}\ensuremath{\subset}A. Classified the exact first-jet hyperplanes: one nonzero coefficient gives a cusp plus independent lines; two give the actual tacnode generators (t,t),(t\textsuperscript{2},−t\textsuperscript{2}); three give a plane ordinary triple because the second-jet product determinant is nonzero. Minimal generator lifts give the asserted ambient coordinate isomorphisms. In the global pole-tree case H must contract because its coefficient2 cannot dominate a generic conductor of exponent1. The current conductor duality gives \ensuremath{\chi}OC=−1; the two rational normalized components and one graph cycle give total delta3, identification length2 and excess1. Every resulting embedded singular germ is one of the actual exclusions in P4.38. The unchanged curve\_excess.py replay passed. Review provenance: New derivation after reading the supplied proof; not a blind review and not an independent human check}{\ref{night:P4.33}; \ref{night:P4.35}; \ref{night:P4.38}; \ref{night:P4R.4}}

In the scope of Part~4.2 where every fibre support is
irreducible or normal crossing, no irreducible support with
positive-genus ruled resolution has generic transverse type
$A_1$ on every component of its critical curve.
\end{proposition}
\begin{proof}
\noindent\textbf{Step 1.} We first give the exact local curve calculation needed in
the proof. Let $A$ be a reduced complex curve local ring with
$r$ branches, let $B=\bigoplus_{i=1}^r\mathbb C\{t_i\}$
be its normalization, and let $R\subset B$ be the ring whose
constant terms agree. Then
\[
 \operatorname{length}(R/A)=\delta(A)-(r-1).
\]
Suppose this length is one. The simple $A$-module $R/A$
is annihilated by $\mathfrak m_A$, so
$\mathfrak m_A R\subset A$. Choose
$v\in\mathfrak m_R$ spanning $R/A$ and write
$v^2=a+cv$, where $a\in\mathfrak m_A$ and $c\in\mathbb C$.
If $c\ne0$, the convergent square-root series solves
$b^2-cb-a=0$ with $b\in\mathfrak m_A$. Then
$(v-b)(v+b-c)=0$ and the second factor is a unit, forcing
$v=b$, a contradiction. Thus $v^2\in A$ and
\[
                         \mathfrak m_R^2\subset A.
\]
Consequently $A$ is exactly the inverse image of a hyperplane
\[
                         \sum_i a_i f_i'(0)=0
\]
in $\mathfrak m_R/\mathfrak m_R^2=\mathbb C^r$.
Rescaling the parameters with nonzero $a_i$ gives actual
analytic ring isomorphisms with the following models.

\noindent\textbf{Step 2.} With one nonzero coefficient, that branch has ring
$\mathbb C\{t^2,t^3\}$ and the others are independent smooth
branches joined at their constant terms. The embedding
dimension is $r+1$: for $r=1$ this is a plane cusp, for
$r=2$ a cusp and transverse line, and for $r=3$ it does
not embed in a smooth threefold. With two nonzero
coefficients, the pair has equal first derivatives. Its
ring is generated by $(t,t)$ and $(t^2,-t^2)$, giving the
plane tacnode $y^2=x^4$. The other $r-2$ branches are
independent transverse lines. For $r=3$ and three nonzero
coefficients, the first jets are the restrictions of linear
forms to three distinct lines in a plane. Their products
span all three second-jet coordinates, so the ring is
exactly the ordinary plane triple point. Equivalently,
the latter ring has precisely the same constant and
first-derivative conditions and colength three in $B$.
These cases exhaust $r\le3$. The analytic ring
isomorphisms lift to ambient coordinate changes: choose
minimal algebra generators, lift their independent linear
parts, and apply the inverse-function theorem. In embedding
dimension two we first straighten the smooth containing
surface. Every realizable model is therefore one of the
local classes excluded in P4.38 for ordinary-double
transverse type.

\noindent\textbf{Step 3.} Now suppose the stated ruled support exists. Its horizontal
pole section has coefficient two. It cannot map
nonconstantly to the conductor, since an ordinary double
point has conductor exponent one on each normalization
branch. The section is contracted. P4.33 gives
$\chi(\mathcal O_{C_{\rm sch}})=-1$.
The conductor image is Cohen--Macaulay and generically
reduced in type $A_1$. It is therefore reduced, so
$C_{\rm sch}=C$ and $\chi(\mathcal O_C)=-1$.
All normalized components of $C$ are rational.

\noindent\textbf{Step 4.} The already proved group and Euler calculations leave
exactly two normalized rational components and one
point-identification cycle. Indeed the rational-tree case
has Euler number at least three by P4R.4; the
one-component cycle is excluded by P4.35; and the global
nonnegative Euler sum is at most two. The normalization
sequence and the graph formula consequently give
\[
 \sum_p\delta_p=3,\qquad
 \sum_p(r_p-1)=2,
 \qquad
 \sum_p\bigl(\delta_p-(r_p-1)\bigr)=1.
\]
Every summand in the last sum is nonnegative. There is
exactly one excess-one point and every point has at most
three local branches. The local calculation classifies that
point by an actual analytic isomorphism, and P4.38 excludes
each possible model. This is the required contradiction.
\end{proof}

The first-jet cases, the second-jet product determinant, and
the global excess computation are verified in
\texttt{code/part4/curve\_excess.py}; its saved output reports
PASS. The proof retains the ordinary-double boundary
condition and makes no substitution of transverse type
$A_2$ for $A_1$.


\subsection{The ordinary-node defect in transverse type \texorpdfstring{$A_3$}{A3}}

\begin{proposition}[P4.41]\label{night:P4.41}\label{stmt:P4.41}
\StatementRecord{night:P4.41}{proved}{Historical events: GPT-6 Astra (OpenAI), ChatGPT Pro mode, effort Pro, S6\_06\_BLIND, 16--18 November 2026: targeted proof check. Also checked in this writing. Pre-continuation writing events: GPT-6 Astra (OpenAI), ChatGPT Pro mode, effort Pro, S6\_10\_PAPERS, 18 November 2026: used and rechecked in this run. GPT-6 Astra (OpenAI), ChatGPT Pro mode, effort Pro, S6\_10\_PAPERS, 18 November 2026 05:10:56.410177+00:00: computation or formal certificate; targeted proof check; used and rechecked in this run. Continuation metadata events: GPT-6 Astra (OpenAI), ChatGPT Pro mode, effort Pro, S6\_10\_PAPERS continuation, 19 November 2026: used and rechecked in this run; status-vocabulary reconciliation; metadata reconciliation only; no new mathematical check. \emph{Pre-continuation recorded proof scope.} The night-run re-derivation was by GPT-6 Astra (OpenAI, ChatGPT Pro mode, effort Pro), 16--18 November 2026. Also checked in this writing by GPT-6 Astra (OpenAI), ChatGPT Pro mode, effort Pro, 18 November 2026. proved. Full ordinary-node conductor-module calculation, exact ring-condition matrices, defect bound and equality normalization. Classified the normalized rank-one module at A=C\{x,y\}/(xy) as the diagonal line or the full two-coordinate space, giving J\ensuremath{\cong}A,q=1 or J\ensuremath{\cong}B,q=0 after the recorded branch factors. Opened and re-derived de Jong–van Straten Theorem1.12 with its codimension-two Cohen–Macaulay and Gorenstein hypotheses. In the free case I=(g,k\textsuperscript{2}), its adjugate forces f\ensuremath{\in}I\textsuperscript{2}. A unit g\textsuperscript{2} coefficient yields a boundary-preserving nonzero-value Morse point, so its vanishing gives defect\ensuremath{\geq}5. In the nonfree case, solving xy=\ensuremath{\alpha}a+\ensuremath{\beta}b identifies the exact pole term V(0)\ensuremath{\alpha}(x,0)/x in the null cubic; its vanishing forces the displayed branch orders. Equality is d=(4,1), b=(0,1), q=0. The invertible coordinate Y=y−ell(x)a gives I=(z\textsuperscript{2},Yz,x\textsuperscript{2}Y). Computed every new 2×2 minor: they force a=yA,c=zC,b=−x\textsuperscript{2}C+zB modulo I, and f\ensuremath{\in}(yz\textsuperscript{2})+I\textsuperscript{2}. The smooth-factor substitution y=z\textsuperscript{2}v and the Morse centre divisible by x\textsuperscript{2}y produce Z\textsuperscript{2}+x4Y. Its actual finite birational normalization (x,u)\ensuremath{\to}(x,−u\textsuperscript{2},x\textsuperscript{2}u) is smooth. The unchanged a3\_node.py replay passed. Review provenance: New derivation after reading the supplied proof; not a blind review and not an independent human check}{\ref{night:P4.28}; \ref{night:P4R.10}; \ref{night:P4R.14}}

At a point of an actual support in Part~4.2, suppose the
entire reduced critical curve is the ordinary node
$C=\{z=xy=0\}$ and both branches have generic transverse
type $A_3$. Let $d_i,b_i,q$ be the Hessian zero orders,
projected conormal colengths, and normalization descent
defect of P4R.10. Then
\[
 \delta(C)+\sum_{i=1}^2\left(\frac{d_i}{2}+b_i\right)+q
                         \ge\frac92.                 \tag{4.41a}
\]
Here $R=0$ and $\delta(C)=1$. Equality makes the normalization
of the hypersurface smooth over the point.
\end{proposition}
\begin{proof}
\noindent\textbf{Step 1.} Put $P=(z,xy)$, and let $I$ be the ambient ideal of the
scheme-theoretic conductor image. We have $P^2=P^{(2)}\subset I$,
and $J=P/I$ is torsionfree of rank one on both branches.
After dividing its normalized image by the two branchwise
vanishing orders, it is an $A=\mathcal O_C$ submodule of
$B=\mathbb C\{x\}\oplus\mathbb C\{y\}$ generating both
summands. The conductor of $A$ in $B$ is $\mathfrak m_A$,
so $\mathfrak m_AB\subset J$. Thus $J$ is determined by
a subspace of $\mathbb C^2$ projecting onto both coordinates.
If that subspace is a line, rescaling gives $J\simeq A$ and
$q=1$. If it is the whole space, $J\simeq B$ and $q=0$.
These are actual module isomorphisms. The divided orders are
the $b_i$ already in the formula.

\noindent\textbf{Step 2.} We use the ring condition of the actual conductor:
\[
 \operatorname{Hom}_{\mathcal O_S}(I,I)=
 \operatorname{Hom}_{\mathcal O_S}(I,\mathcal O_S).
                                                               \tag{4.41b}
\]
The right side is the normalization algebra by finite duality,
and $I$ is an ideal of that algebra. For the determinant
presentation of a Cohen--Macaulay codimension-two conductor
inside a Gorenstein hypersurface, this condition is equivalent
to membership of the adjugate entries in $I$. This is precisely
de Jong--van Straten, Theorem~1.12, printed p.~532
\cite{DeJongVanStraten1990Normalization}. Its hypotheses hold
here: the ambient space is regular, the reduced support is a
hypersurface, and its conductor quotient is Cohen--Macaulay.
We give all the matrices used below.

\noindent\textbf{Step 3.} Suppose first that $q=1$. The surjection
$P/P^2\to J$ is a surjection $A^2\to A$, with primitive
kernel generator $g$. Taking a complementary generator $k$
of $P$ gives $I=(g,k^2)$. A determinant presentation
$f=ag+bk^2$ has two-by-two adjugate containing $a,b,g,k^2$.
Condition (4.41b) says $a,b\in I$, so the exact equation is
\[
                   f=Ag^2+Bgk^2+Ck^4.                \tag{4.41c}
\]
We claim that $A(0)$ cannot be a unit. If $g_z(0)\ne0$,
straighten $g$, take $k$ to be a nodal equation in $g=0$,
and apply the parameter Morse lemma in $g$. The residual
is $k^4$ times a unit after a small change in its fourth-order
coefficient. A morsification of the node and a shift of its
critical value away from zero give an isolated Morse point
of nonzero critical value.

\noindent\textbf{Step 4.} If $g_z(0)=0$, its $xy$ coefficient is a unit and we take
$k=z$. Completing the square in (4.41c), including the
square root of $A$ in its first factor, gives exactly
\[
                         f=G^2+z^4U,
\]
with $G(x,y,0)=xy$ times a unit and $G_z(0)=0$.
Add a small constant to $U$, then take the actual coordinate
$Z=zU^{1/4}$. The parameter Morse lemma for $G$ in $x,y$
gives $G=XY+a(Z)$, where $a\in(Z^2)$. A small change
$a\mapsto a+\eta Z^2$ makes $a''(0)\ne0$. Replacing
$G$ by $G-\epsilon$ produces at the origin a critical point
of value $\epsilon^2$. The $XY$ Hessian entry is
$-2\epsilon$, and the $ZZ$ entry is
$-2\epsilon a''(0)$, so the point is Morse.
These deformations retain type $A_3$ along the moving
critical curve at each fixed boundary. The nested boundary
transport of P4.28 and P4R.14 places its square-$-2$ sphere
in the actual nearby torus, a contradiction. Hence $A(0)=0$.

\noindent\textbf{Step 5.} If $g_z(0)$ is a unit, both saturated Hessian lines approach
$dz$, so $b_1=b_2=1$ and $d_i\ge1$. If $g_z(0)=0$,
the differential of $g$ on each branch has exactly the
common factor $t_i$ and a remaining nonzero coefficient
in the other normal direction. Thus $b_i=0,d_i\ge3$.
In either case the left side of (4.41a) is at least five,
including $q=1$.

\noindent\textbf{Step 6.} Now let $q=0$. Choose actual generators $a,b$ of $P$ mapping
to $(1,0),(0,1)$ in $J\simeq B$. Their coefficient matrix
is invertible modulo the maximal ideal, hence invertible
locally. We obtain
\[
 I/P^2=(ya,xb),\qquad
 f=xU(x)b^2+yV(y)a^2\pmod {P^3}.                      \tag{4.41d}
\]
The second identity records the Hessian kernel $b$ on the
$x$ branch and $a$ on the $y$ branch. Write the exact
relation $xy=\alpha a+\beta b$. Near a general point of
the $x$ axis, solve this for $y$ using coordinates $x,a,b$.
The null cubic coefficient contains
$V(0)\alpha(x,0)/x$; the remaining contribution from
$P^3$ is holomorphic at the node. Its identical vanishing
in type $A_3$ forces $V(0)\alpha(0)=0$. The other branch
gives $U(0)\beta(0)=0$.

\noindent\textbf{Step 7.} If both $a_z(0),b_z(0)$ are nonzero, the relation for $xy$
has $\alpha(0)\beta(0)\ne0$. Hence $U(0)=V(0)=0$,
both $b_i=1$, and both $d_i\ge2$. The defect is at least
five. Otherwise interchange the branches so that $a_z(0)$
is a unit and $b_z(0)=0$. Then $\alpha(0)=0$ and
$\beta(0)$ is a unit, forcing $U(0)=0$. On the $x$ axis,
$db$ has exactly the common factor $x$ and remaining
$dy$ coefficient a unit. On the $y$ axis, $da$ is primitive
with a unit $dz$ coefficient. Thus
\[
 b_x=0,\quad d_x=3+\operatorname{ord}_xU\ge4,
 \qquad
 b_y=1,\quad d_y=1+\operatorname{ord}_yV\ge1.
\]
This proves (4.41a). Equality requires
$\operatorname{ord}_xU=1$ and $V(0)\ne0$.

\noindent\textbf{Step 8.} We identify the normalization in this exact equality case.
Use $a$ as the ambient coordinate $z$. The generator relation
with $\beta$ a unit gives
\[
                    I=(a^2,ya,x^2y-x\alpha a).
\]
Modulo its first two generators, replace $\alpha$ by
$\alpha(x,0,0)=x\ell(x)$. The actual coordinate
$Y=y-\ell(x)a$ puts the ideal into
$I=(z^2,Yz,x^2Y)$. Renaming $Y$ as $y$, its Hilbert--Burch
matrix and determinant extension are
\[
 \begin{pmatrix}y&0\\-z&x^2\\0&-z\end{pmatrix},
 \qquad
 M=\begin{pmatrix}y&0&a\\-z&x^2&b\\0&-z&c\end{pmatrix},
 \qquad
 \det M=az^2+byz+cx^2y.
\]
The new two-by-two minors are
$yb+za,yc,-zc,-x^2a,-za,x^2c+zb$, up to signs.
Their membership in $I$ gives exactly
\[
             a=yA,\qquad c=zC,\qquad b=-x^2C+zB
                          \pmod I.
\]
Therefore $f\in(yz^2)+I^2$, so
\[
           f=Uyz^2+Vz^4+Wx^2y^2z+Tx^4y^2.            \tag{4.41e}
\]
The equality orders $d_y=1,d_x=4$ say precisely
$U(0)T(0)\ne0$.

\noindent\textbf{Step 9.} Substitute $y=z^2v$ and divide (4.41e) by $z^4$.
At the finite root its $v$ derivative is $U(0)$.
The implicit-function theorem supplies the actual smooth
factor $g_1=y-z^2v(x,z)$. Division by it gives
\[
                  g_2=z^2A+x^2yzB+x^4yC,
                  \qquad A(0)C(0)\ne0.
\]
For clarity, this division retains the factors $x^2,x^4$:
subtracting the expression at $y=z^2v$ supplies one factor
$y-z^2v$; every extra substituted $y$ supplies $z^2$ and
is absorbed into the first term. The parameter Morse
centre in $z$ is divisible by $x^2y$. Substitution leaves
$x^4y$ times a unit. An actual coordinate change therefore
gives $g_2=Z^2+x^4Y$. Its finite birational normalization is
\[
                 (x,u)\longmapsto(x,-u^2,x^2u),
\]
with target coordinates $(x,Y,Z)$. It is smooth; finiteness
follows from $u^2=-Y$ and birationality from $u=Z/x^2$
away from $x=0$. The other branch $g_1$ is smooth already.
Thus the full normalization is smooth over the node.
\end{proof}

All determinant minors, the four minimum-defect cases,
the factor substitution, and the normalization equation
are verified by \texttt{code/part4/a3\_node.py}.
The output reports PASS. Its initial failed assertion
compared an expanded polynomial with a factored expression;
that output is retained, and the corrected assertion tests
their exact zero difference.

\begin{proposition}[P4.42]\label{night:P4.42}\label{stmt:P4.42}
\StatementRecord{night:P4.42}{proved}{Historical events: GPT-6 Astra (OpenAI), ChatGPT Pro mode, effort Pro, S6\_06\_BLIND, 16--18 November 2026: targeted proof check. Also checked in this writing. Pre-continuation writing events: GPT-6 Astra (OpenAI), ChatGPT Pro mode, effort Pro, S6\_10\_PAPERS, 18 November 2026 05:10:56.410177+00:00: computation or formal certificate; targeted proof check. Continuation metadata events: GPT-6 Astra (OpenAI), ChatGPT Pro mode, effort Pro, S6\_10\_PAPERS continuation, 19 November 2026: used and rechecked in this run; status-vocabulary reconciliation; metadata reconciliation only; no new mathematical check. \emph{Pre-continuation recorded proof scope.} The night-run re-derivation was by GPT-6 Astra (OpenAI, ChatGPT Pro mode, effort Pro), 16--18 November 2026. Also checked in this writing by GPT-6 Astra (OpenAI), ChatGPT Pro mode, effort Pro, 18 November 2026. proved. Full two-two-branch cycle exclusion, including planar different, nonplanar small-delta model and equality rigidity. For E=t\textasciicircum{}a e1\ensuremath{\oplus}t\textasciicircum{}b e2 and primitive ell=u e1+v e2, independently derived d/2\ensuremath{\geq}max(a−ord u,b−ord v) and projected length min(a+ord v,b+ord u); the sum is at least a+b. The exact plane-curve dualizing pullback dx/h\_y gives sum(a+b)=2delta+degR and local defect\ensuremath{\geq}3delta+q. Nonplanar branches have gradient and Hessian order at least their multiplicity, giving the stated positive contribution after ramification. The sole small delta2 nonplanar two-branch model is cusp plus transverse line; its symbolic-square extra class zh has maximal-ideal annihilator, and rank-one Hessian forces its constant coefficient to vanish, so f\ensuremath{\in}P\textsuperscript{2}. Its conormal orders are (2,3),(1,1), giving defect\ensuremath{\geq}8. Thus each two-branch junction contributes\ensuremath{\geq}9/2, with equality only at the ordinary node of P4.41 and smooth normalization. Re-derived \ensuremath{\chi}OCsch=8−sum(defects) on two normalized P1s. The uncontracted case contradicts \ensuremath{\chi}=g−1\ensuremath{\geq}0; the contracted case forces equality everywhere, and P4R.14 plus absence of any additional positive defect makes the normalization smooth, contradicting a contracted positive-genus section. Review provenance: New derivation after reading the supplied proof; not a blind review and not an independent human check}{\ref{night:P4.33}; \ref{night:P4.34}; \ref{night:P4.40}; \ref{night:P4.41}; \ref{night:P4R.10}; \ref{night:P4R.13}; \ref{night:P4R.14}; \ref{night:P4R.6}}

In the requested global support classes, a positive-genus
ruled-resolution support whose critical curve has generic
type $A_3$ everywhere cannot have two two-branch points as
its point-identification cycle. The two points may be
nonordinary. This includes additional
unibranch singular points of the critical curve.
\end{proposition}
\begin{proof}
\noindent\textbf{Step 1.} After P4R.13, the remaining positive ruled case has two
rational normalized conductor components and one graph
cycle. In the stated instance the cycle is provided by
two two-branch points. We first extend the local lower
bound of P4.41 to these points.

\noindent\textbf{Step 2.} At a planar curve point, write its ideal as $(z,h)$.
Its symbolic square equals its ordinary square, so the
rank-one Hessian belongs to $\operatorname{Sym}^2 E_0$,
where $E_0$ is the normalized conormal image. On a branch
put $E_0=\langle t^a e_1,t^b e_2\rangle$ and write the
primitive Hessian direction as $\ell=u e_1+v e_2$.
One of $u,v$ is a unit. Membership of
$t^d\ell^2$ in $\operatorname{Sym}^2E_0$ and projection
onto the quotient by $\ell$ give
\[
 \frac d2\ge\max\{a-\operatorname{ord}u,
                         b-\operatorname{ord}v\},\qquad
 b_{\rm proj}=\min\{a+\operatorname{ord}v,
                         b+\operatorname{ord}u\}.
\]
The sum of the maximum and minimum on the right is
$a+b$. Thus $d/2+b_{\rm proj}\ge a+b$.
The plane different formula proved in P4.34 is
$\sum(a+b)=2\delta+\deg R$. The whole local defect
is consequently at least $3\delta+q$.
A nonordinary planar two-branch point has $\delta\ge2$
and contributes at least six.

\noindent\textbf{Step 3.} For a nonplanar point, the ambient Hessian vanishes: a
nonzero rank-one Hessian would give $f=X^2+g(Y,Z)$ by
the parameter Morse lemma and would put the entire
critical curve in the smooth plane $X=0$.
On a branch of multiplicity $m$, all conormal generators
have order at least $m$, and every Hessian entry has
order at least $m$. Thus $b_{\rm proj}\ge m$, $d\ge m$,
and, since $r=m-1$,
\[
                  d/2+b_{\rm proj}-r\ge1+m/2\ge1.
\]
When $\delta\ge3$, a nonplanar two-branch point
therefore contributes at least five.
The only remaining possibility has $\delta=2$.
The exact excess-one classification in P4.40 identifies
it with a cusp and transverse line,
\[
                  P=(xz,yz,h),\qquad h=y^2-x^3.
\]
Here $P^{(2)}/P^2$ is generated by $zh$ and is
annihilated by the maximal ideal. Write $f=Azh+P^2$.
Along the cusp, the normal quadratic determinant is
$4KG-A^2$, with $K(0)=0$. Its rank is one generically,
so $A(0)=0$ and in fact $f\in P^2$.
The conormal lattices on the two normalized branches
are respectively $\operatorname{diag}(t^2,t^3)$ and
$\operatorname{diag}(t,t)$. Their total colength is seven,
while $\deg R=1$ and $\delta=2$.
Applying the same symmetric-square inequality gives
defect at least $2+7-1=8$.
Every two-branch point thus contributes at least $9/2$,
and equality is possible only at the ordinary node of
P4.41, whose normalization is smooth.

\noindent\textbf{Step 4.} Apply the rank-one Hessian calculation
of P4R.10 to the two normalized projective lines. On each
line the rank-two conormal kernel has degree $2+\deg R$.
The rank-one quotient has Euler characteristic
$3+\deg R-\deg L$. Combining both lines with the
normalization sequence of $C$ gives the exact formula
\[
 \chi(\mathcal O_{C_{\rm sch}})
   =8-\sum_p\left[
       \delta_p+\sum_{i\text{ over }p}
          \left(\frac{d_i}{2}+b_i-r_i\right)+q_p
                 \right].                            \tag{4.42a}
\]
Every extra unibranch singular point has a strictly positive
bracket by the local conormal inequalities of P4R.10.
At a smooth point the bracket is nonnegative. The two
two-branch points contribute at least nine together by the
local calculation just given.

\noindent\textbf{Step 5.} If the horizontal section is not contracted, P4.33 gives
$\chi(\mathcal O_{C_{\rm sch}})=g-1\ge0$, contradicting
(4.42a), whose right side is at most $-1$.
If the section is contracted, that Euler characteristic is
$-1$. Equality must then hold at both two-branch points and every
other bracket must vanish. There are no additional
unibranch singular points. Both cycle points are ordinary
nodes and the normalization is smooth over them by P4.41.
At every smooth point of the
critical curve, P4R.14 gives either the product tacnode
or the displayed $A_1$ normalization degeneration. The
latter has Hessian zero order one and contributes $1/2$,
so it is excluded by the equality in (4.42a). Thus the
normalization is smooth along all of the conductor, and
it is smooth away from the conductor by P4R.6. A contracted
positive-genus horizontal section in its minimal resolution
is impossible. This proves the assertion.
\end{proof}


\subsection{The uniform cuspidal sector}

\begin{proposition}[P4.43]\label{night:P4.43}\label{stmt:P4.43}
\StatementRecord{night:P4.43}{proved}{Historical events: GPT-6 Astra (OpenAI), ChatGPT Pro mode, effort Pro, S6\_06\_BLIND, 16--18 November 2026: targeted proof check. Also checked in this writing. Pre-continuation writing events: GPT-6 Astra (OpenAI), ChatGPT Pro mode, effort Pro, S6\_10\_PAPERS, 18 November 2026 05:10:56.410177+00:00: computation or formal certificate; targeted proof check. Continuation metadata events: GPT-6 Astra (OpenAI), ChatGPT Pro mode, effort Pro, S6\_10\_PAPERS continuation, 19 November 2026: used and rechecked in this run; status-vocabulary reconciliation; metadata reconciliation only; no new mathematical check. \emph{Pre-continuation recorded proof scope.} The night-run re-derivation was by GPT-6 Astra (OpenAI, ChatGPT Pro mode, effort Pro), 16--18 November 2026. Also checked in this writing by GPT-6 Astra (OpenAI), ChatGPT Pro mode, effort Pro, 18 November 2026. proved. Full uniform-cusp positive-Euler exclusion via the actual normalization and conductor pushout maps. Two local hypersurface branches intersect in a curve in the regular threefold. That would make the generic transverse germ reducible, so uniform A2 makes the finite normalization bijective; properness gives an actual homeomorphism. The resolution contracts vertical rational trees without changing \ensuremath{\pi}1, and contracting the horizontal section kills the entire ruling surface group because its inclusion induces an isomorphism. Thus \ensuremath{\pi}1(N) is either a surface group of genus g\ensuremath{\geq}1 or trivial. The explicit conductor pushout and the positive-Euler graph reduction give \ensuremath{\pi}1(S)=\ensuremath{\pi}1(C)=Z for two rational components and one cycle. Abelianization has rank2g or0 for N, so the homeomorphism cannot induce the required group isomorphism. Followed the exact roles of R7/R9/R13, including exclusion of an elliptic conductor component. Review provenance: New derivation after reading the supplied proof; not a blind review and not an independent human check}{\ref{night:P4R.13}; \ref{night:P4R.7}}

In the requested full-fibration classes, a support with
positive-genus ruled resolution cannot have positive
Euler number and generic transverse type $A_2$ along
every component of its critical curve.
\end{proposition}
\begin{proof}
\noindent\textbf{Step 1.} We use the exact normalization map
$\nu:N\longrightarrow S$ and its minimal resolution
$\rho:T\longrightarrow N$. A germ of $S$ cannot have
two distinct analytic hypersurface branches. Indeed,
their intersection has dimension one in the smooth
ambient threefold. A general point of an irreducible
component of that intersection has at least two local
hypersurface branches, so its transverse curve germ is
reducible. The transverse cusp $A_2$ is irreducible.
This contradicts the stated uniform transverse type.
The points of the finite normalization over a germ are
its analytic branches. Consequently $\nu$ is bijective.
It is proper, so its continuous map from a compact space
to the Hausdorff space $S$ is a homeomorphism.

\noindent\textbf{Step 2.} Let $B$ be the base curve of the minimal ruled surface,
of genus $g\ge1$, and let $H$ be the horizontal section
of coefficient two in the total pole divisor. Point
blowups preserve the fundamental group, and the ruling
gives $\pi_1(T)=\pi_1(B)$. By the fibre-tree calculation
in P4R.7, every vertical exceptional configuration is a
tree of rational curves. Collapsing any such connected
tree preserves the fundamental group: apply van Kampen
to its tubular neighbourhood and the complement, or
equivalently to the cofibration quotient by that
simply connected tree. These quotients are the actual
topological contractions of the resolution map.

\noindent\textbf{Step 3.} If $H$ is not contracted by $\rho$, all resolution
exceptional configurations are of this vertical kind.
It follows that $\pi_1(N)=\pi_1(B)$. If $H$ is
contracted, its connected exceptional configuration
contains $H$. The inclusion $H\hookrightarrow T$
induces an isomorphism on fundamental groups, since
the ruling restricts to the identity section and all
point blowups preserve that isomorphism. The quotient
which collapses this configuration therefore kills
the whole group. The other contracted configurations
are rational trees, giving $\pi_1(N)=1$.

\noindent\textbf{Step 4.} On the other hand, the proved positive ruled reduction,
including P4R.13, gives a conductor curve $C$ with two
rational normalized components and one graph cycle.
The conductor cofibration pushout in P4R.7 gives
$\pi_1(S)=\pi_1(C)=\mathbb Z$. This contradicts the
homeomorphism $N\simeq S$: neither the trivial group
nor the genus-$g$ surface group is $\mathbb Z$, as their
abelianizations have ranks zero and $2g$, respectively.
\end{proof}


\subsection{Further uniform transverse sectors}

\begin{proposition}[P4.44]\label{night:P4.44}\label{stmt:P4.44}
\StatementRecord{night:P4.44}{proved}{Historical events: GPT-6 Astra (OpenAI), ChatGPT Pro mode, effort Pro, S6\_06\_BLIND, 16--18 November 2026: targeted proof check. Also checked in this writing. Pre-continuation writing events: GPT-6 Astra (OpenAI), ChatGPT Pro mode, effort Pro, S6\_10\_PAPERS, 18 November 2026 05:10:56.410177+00:00: computation or formal certificate; targeted proof check. Continuation metadata events: GPT-6 Astra (OpenAI), ChatGPT Pro mode, effort Pro, S6\_10\_PAPERS continuation, 19 November 2026: used and rechecked in this run; status-vocabulary reconciliation; metadata reconciliation only; no new mathematical check. \emph{Pre-continuation recorded proof scope.} The night-run re-derivation was by GPT-6 Astra (OpenAI, ChatGPT Pro mode, effort Pro), 16--18 November 2026. Also checked in this writing by GPT-6 Astra (OpenAI), ChatGPT Pro mode, effort Pro, 18 November 2026. proved. Full contraction of the horizontal section in the remaining uniform-A3 three-branch case. After the two-node-cycle exclusion, the conductor has two rational normalized components and a single three-branch junction. If H survives, its map to one normalized P1 has degree2, using both generic tacnode preimages. The explicit P4R.14 normalizations each have two distinct points over a smooth critical point, so this double cover is unramified there. The exact Euler identity bounds the total defect by9−g\ensuremath{\leq}8. A nonplanar triple consumes\ensuremath{\geq}5, a planar triple\ensuremath{\geq}9, and any extra unibranch singularity\ensuremath{\geq}2. Hence H has at most two possible branch values over the triple plus one elsewhere. A degree2 map from genus g\ensuremath{\geq}1 requires2g+2\ensuremath{\geq}4 distinct branch values. The mismatch proves contraction using the actual finite map, not only the incidence graph. Review provenance: New derivation after reading the supplied proof; not a blind review and not an independent human check}{\ref{night:P4.33}; \ref{night:P4.42}; \ref{night:P4R.14}; \ref{night:P4R.6}}

In the remaining positive ruled case with generic transverse
type $A_3$ on every conductor component, the horizontal
section of the pole divisor is contracted by the resolution
map. The point-identification cycle occurs at a single
three-branch point of the critical curve.
\end{proposition}
\begin{proof}
\noindent\textbf{Step 1.} P4.42 excludes a cycle at two two-branch points. The
remaining graph has two rational normalized components
and one cycle, so its cycle occurs at one three-branch
point $p$. Suppose that the horizontal section $H$ is
not contracted. It maps finitely to one normalized
conductor component $C_H^\nu=\mathbb P^1$.
The generic tacnode has two conductor preimages.
The degree of $H\to\mathbb P^1$ is greater than one,
since $H$ has genus $g\ge1$, and hence is exactly two.
No other, vertical conductor component can map onto
$C_H$: the horizontal map already supplies both generic
preimages.

\noindent\textbf{Step 2.} At a point where the entire critical curve is smooth,
P4R.14 gives either the product tacnode or the germ
$X(Y^2+XZ)$. Each has two distinct normalization points
over the critical-curve point. Each point lies on the
strict conductor, by P4R.6. Thus the two preimages must
both lie on $H$, and the degree-two map is unramified
there. All ramification is over singular points of $C$.

\noindent\textbf{Step 3.} By (4.42a) and P4.33 the sum of the local defects is
$9-g\le8$. A nonplanar three-branch point has
$\delta\ge2$, and the branchwise estimate in P4.42
adds at least one per branch; its defect is at least
five. A planar three-branch point has $\delta\ge3$
and contributes at least $3\delta\ge9$, which is
already impossible. Every additional unibranch singular
point contributes at least two: in the planar case
the bound is $3\delta$, and in the nonplanar case it is
at least $\delta+1$. There is consequently at most one
additional unibranch singular point.

\noindent\textbf{Step 4.} The normalization $C_H^\nu$ has at most two points
above $p$, since the other global conductor component
supplies at least one of its three branches. Including
the possible additional unibranch point leaves at most
three possible branch values for $H\to\mathbb P^1$.
Every ramification point of a degree-two map has index
two, and its branch value has only that ramified
preimage. Riemann--Hurwitz requires $2g+2\ge4$ distinct
branch values. This contradiction proves that $H$ is
contracted.
\end{proof}

\begin{proposition}[P4.45]\label{night:P4.45}\label{stmt:P4.45}
\StatementRecord{night:P4.45}{proved}{Historical events: GPT-6 Astra (OpenAI), ChatGPT Pro mode, effort Pro, S6\_06\_BLIND, 16--18 November 2026: targeted proof check. Also checked in this writing. Pre-continuation writing events: GPT-6 Astra (OpenAI), ChatGPT Pro mode, effort Pro, S6\_10\_PAPERS, 18 November 2026 05:10:56.410177+00:00: computation or formal certificate; targeted proof check. Continuation metadata events: GPT-6 Astra (OpenAI), ChatGPT Pro mode, effort Pro, S6\_10\_PAPERS continuation, 19 November 2026: used and rechecked in this run; status-vocabulary reconciliation; metadata reconciliation only; no new mathematical check. \emph{Pre-continuation recorded proof scope.} The night-run re-derivation was by GPT-6 Astra (OpenAI, ChatGPT Pro mode, effort Pro), 16--18 November 2026. Also checked in this writing by GPT-6 Astra (OpenAI), ChatGPT Pro mode, effort Pro, 18 November 2026. proved. Full planar-singular ordinary-triple exclusion with explicit critical-value deformation. The complete-intersection symbolic cube gives f=Au\textsuperscript{3}+Bu\textsuperscript{2}h+Cuh\textsuperscript{2}+Dh\textsuperscript{3}. Replaced h by an independent H in the actual analytic F and made arbitrarily small coefficient changes so its binary cubic has a simple critical root with nonzero critical value. The implicit-function root u=Hv gives F=H\textsuperscript{3}beta with beta a unit. Inverted the analytic map H\ensuremath{\to}H beta\textasciicircum{}(1/3), then chose H\_s so the critical-value function is exactly a morsification hhat\_s\textsuperscript{3}. At a nonzero-value plane Morse point the Schur complement is3hhat\_s\textsuperscript{2} Hess(hhat\_s), while F\_uu is nonzero. The resulting full Morse point retains ordinary-triple transverse boundary germs. The exact P4.20/P4.28 boundary transport embeds its square−2 sphere in the actual nearby torus, giving the exclusion. Review provenance: New derivation after reading the supplied proof; not a blind review and not an independent human check}{\ref{night:P4.28}}

In an actual torus degeneration, a point of a planar
singular critical curve cannot have generic transverse
type $D_4$ along all its local branches.
\end{proposition}
\begin{proof}
\noindent\textbf{Step 1.} Write the full reduced critical-curve ideal as
$P=(u,h(x,y))$, where $h$ is reduced and singular.
Generic multiplicity three gives $f\in P^{(3)}$.
Powers of this complete-intersection prime-by-components
ideal are unmixed, so $P^{(3)}=P^3$. Thus there are
holomorphic coefficients such that
\[
 f=A u^3+B u^2h+Cuh^2+Dh^3.
\]
Introduce an independent variable $H$ and the actual
holomorphic function
\[
 F(u,x,y,H)=A(u,x,y)u^3+B(u,x,y)u^2H
               +C(u,x,y)uH^2+D(u,x,y)H^3.
\]
Then $f=F(u,x,y,h)$. Arbitrarily small constant changes
of the four coefficients preserve the $D_4$ transverse
germs on a fixed boundary and make the binary cubic
at the origin have distinct roots and a simple critical
point as a polynomial in $u/H$. Its critical value is
nonzero: a zero critical value would be a repeated root.
We also choose its leading coefficient nonzero.

\noindent\textbf{Step 2.} After substituting $u=Hv$, the quotient $F_u/H^2$
is holomorphic. The implicit-function theorem gives
a holomorphic critical root $v=v(x,y,H)$ and
\[
          F(Hv(x,y,H),x,y,H)=H^3\beta(x,y,H),
                       \qquad \beta(0)\ne0.
\]
Choose the local cube root of the unit $\beta$.
The map $H\mapsto H\beta(x,y,H)^{1/3}$ has nonzero
$H$ derivative at the origin and has a holomorphic
inverse with parameters $x,y$. Put
\[
             \widehat h(x,y)=h(x,y)
                           \beta(x,y,h(x,y))^{1/3}.
\]
It defines the same singular reduced plane curve as $h$.
Choose a small morsification $\widehat h_s$ with a
Morse critical point of nonzero critical value. This
exists by the positive local degree of the gradient
of the singular reduced plane function; a small constant
shift avoids zero critical values. Use the displayed
inverse to solve exactly
\[
         H_s\beta(x,y,H_s)^{1/3}=\widehat h_s,
             \qquad f_s=F(u,x,y,H_s(x,y)).
\]
On its $u$-critical branch the critical-value function
is exactly $\widehat h_s^3$. At the selected plane
Morse point, $H_s\ne0$ and $F_{uu}\ne0$. The Schur
complement of the full Hessian is
$3\widehat h_s^2\operatorname{Hess}(\widehat h_s)$.
Consequently $f_s$ has a nondegenerate three-variable
critical point with nonzero critical value.

\noindent\textbf{Step 3.} All deformations above can be chosen successively small
on the fixed boundary. There the moving curve $H_s=0$
has the same ordinary-triple transverse germ: its binary
cubic has three distinct linear factors. The boundary
transport used in P4.28 therefore identifies the local
nearby fibres through nested representatives. The Morse
vanishing sphere embeds in a nearby global torus fibre
and has self-intersection $-2$, contradicting its
asphericity. This proves the assertion.
\end{proof}

\begin{proposition}[P4.46]\label{night:P4.46}\label{stmt:P4.46}
\StatementRecord{night:P4.46}{proved}{Historical events: GPT-6 Astra (OpenAI), ChatGPT Pro mode, effort Pro, S6\_06\_BLIND, 16--18 November 2026: targeted proof check. Also checked in this writing. Pre-continuation writing events: GPT-6 Astra (OpenAI), ChatGPT Pro mode, effort Pro, S6\_10\_PAPERS, 18 November 2026 05:10:56.410177+00:00: computation or formal certificate; targeted proof check. Continuation metadata events: GPT-6 Astra (OpenAI), ChatGPT Pro mode, effort Pro, S6\_10\_PAPERS continuation, 19 November 2026: used and rechecked in this run; status-vocabulary reconciliation; metadata reconciliation only; no new mathematical check. \emph{Pre-continuation recorded proof scope.} The night-run re-derivation was by GPT-6 Astra (OpenAI, ChatGPT Pro mode, effort Pro), 16--18 November 2026. Also checked in this writing by GPT-6 Astra (OpenAI), ChatGPT Pro mode, effort Pro, 18 November 2026. proved. Full uniform-D4 positive-Euler exclusion using the complete conormal/descent defect. The generic conductor is the actual symbolic double ideal P\textasciicircum{}(2), by the conductor square and depth-one argument; Riemann–Roch on the two normalized P1s gives \ensuremath{\chi}=10−sum(delta+tau−r), so the uncontracted total is11−g\ensuremath{\leq}10. At a nonplanar triple of smooth branches, a generic linear section has length3 and embedding dimension2, hence ideal(v,w)\textsuperscript{2}. The nonzerodivisor Tor sequence and Nakayama lift exactly3 defining generators. The normalized conormal needs6 generator values and its branch lattices already have colength\ensuremath{\geq}6; the descent therefore adds\ensuremath{\geq}3, producing defect\ensuremath{\geq}11. A singular branch gives total multiplicity\ensuremath{\geq}4; the only apparent small case m=4,delta=3 is the cusp plus two independent lines, of embedding dimension4, excluded by the explicit P4.40 ring model. The remaining two-junction alternative forces two planar ordinary nodes and g=1; each node is excluded by P4.45. Contracted H is separately excluded by the corrected P4.39 proof. Review provenance: New derivation after reading the supplied proof; not a blind review and not an independent human check}{\ref{night:P4.31}; \ref{night:P4.33}; \ref{night:P4.34}; \ref{night:P4.39}; \ref{night:P4.40}; \ref{night:P4.45}}

In the requested full-fibration classes, a support with
positive-genus ruled resolution cannot have positive
Euler number and generic transverse type $D_4$ along
every component of its critical curve.
\end{proposition}
\begin{proof}
\noindent\textbf{Step 1.} If the horizontal section is contracted, P4.39 already
gives a contradiction. Suppose it is not contracted.
The remaining critical curve has two rational normalized
components and one graph cycle. For the actual
conductor scheme, ordinary-triple transverse type gives
the ideal $P^{(2)}$, as in P4.31. Let $Q=P/P^{(2)}$,
and let $\tau_p$ be the full colength of its differentiated
image in the direct image of the normalized rank-two
kernel $K$. This includes both the branchwise lattice
colength and the further descent colength. Riemann--Roch
on the two normalized projective lines gives
\[
 \chi(\mathcal O_{C_{\rm sch}})
       =10-\sum_p(\delta_p+\tau_p-r_p).
\]
P4.33 identifies the left side with $g-1$, so the total
defect is $11-g\le10$.

\noindent\textbf{Step 2.} We prove that a three-branch point contributes at least
eleven. A planar point contributes $5\delta\ge15$ by
the exact different calculation in P4.34. Consider a
nonplanar point. If its three branches are smooth, the
curve has multiplicity three and embedding dimension
three. A generic linear parameter has an Artinian
section of length three and embedding dimension two,
with ideal $(v,w)^2$. The defining curve ideal therefore
has exactly three minimal generators. To justify the
lifting, the parameter is a nonzerodivisor on the
Cohen--Macaulay curve ring, so the Tor term vanishes in
the ideal sequence modulo that parameter; Nakayama
then lifts precisely the three generators.
Each branchwise conormal lattice has rank two and
colength at least two, since every defining gradient
has positive order. Their sum has colength at least
six. Modulo the three branch maximal ideals its
normalization has six generator values. The image of
$Q$, generated by three elements over the local curve
ring, supplies a vector space of dimension at most
three. The additional descent colength is therefore
at least three. Since $\delta\ge2$ and $R=0$, the
full defect is at least $2+6+3=11$.

\noindent\textbf{Step 3.} If one of the three branches is singular, let $m\ge4$
be total multiplicity. The branchwise gradient bound
gives $\tau\ge2m$, while $\deg R=m-3$.
The general inequality $\delta\ge m-1$ leaves only
$m=4,\delta=3$ as a possible value below the required
bound. That value has excess $\delta-(3-1)=1$.
The exact classification in P4.40 shows that its only
singular-branch model is a cusp and two independent
lines, of embedding dimension four. It does not embed
in the ambient smooth threefold. Thus $\delta\ge4$
in the present case, and the full defect is at least
$\delta+2m-(m-3)=\delta+m+3\ge11$.
A cycle at one three-branch point is impossible.

\noindent\textbf{Step 4.} The other possible cycle has two two-branch points.
Each contributes at least five by P4.34; equality is
possible only for a planar ordinary node. A nonplanar
two-branch point contributes at least seven, and a
nonordinary planar point at least ten. The total bound
of ten therefore forces two ordinary planar nodes
and $g=1$. Each such node is forbidden by P4.45.
This proves the assertion.
\end{proof}

The local Hessian identities, rank-one lattice identity,
and exact integer defect bounds are checked in
\texttt{code/part4/uniform\_sectors.py}; the saved output
reports PASS.


\subsection{The full quadratic conductor at three axes}

\begin{proposition}[P4.47]\label{night:P4.47}\label{stmt:P4.47}
\StatementRecord{night:P4.47}{proved}{Historical events: GPT-6 Astra (OpenAI), ChatGPT Pro mode, effort Pro, S6\_06\_BLIND, 16--18 November 2026: targeted proof check. Also checked in this writing. Pre-continuation writing events: GPT-6 Astra (OpenAI), ChatGPT Pro mode, effort Pro, S6\_10\_PAPERS, 18 November 2026 05:10:56.410177+00:00: computation or formal certificate; targeted proof check; used and rechecked in this run. Continuation metadata events: GPT-6 Astra (OpenAI), ChatGPT Pro mode, effort Pro, S6\_10\_PAPERS continuation, 19 November 2026: used and rechecked in this run; status-vocabulary reconciliation; metadata reconciliation only; no new mathematical check. \emph{Pre-continuation recorded proof scope.} The night-run re-derivation was by GPT-6 Astra (OpenAI, ChatGPT Pro mode, effort Pro), 16--18 November 2026. Also checked in this writing by GPT-6 Astra (OpenAI), ChatGPT Pro mode, effort Pro, 18 November 2026. proved. Full all-three-coefficient ordinary-axis equality normalization and boundary-preserving cone-polar exclusion. The maximal-ideal annihilator of xyz modulo P\textsuperscript{2} and the leading rank-one Hessian force f\ensuremath{\in}P\textsuperscript{2}. Defect9 gives descent length1; the rank-two quotient-value matrix and actual branch-coordinate multiplications lift its kernel to g\ensuremath{\in}I with quadratic part axy+bxz+cyz. The complete intersection(g,xyz) contains P\textsuperscript{2}, agrees generically with I, and its Cohen–Macaulay torsionfree quotient rules out a finite kernel, proving exact ideal equality. The adjugate condition gives f\ensuremath{\in}I\textsuperscript{2}, hence f=G\textsuperscript{2}+(xyz)\textsuperscript{2}U with a nondegenerate quadratic G. A unit U gives two ordinary quadratic surface branches. For U(0)=0, the fixed-axis scaling yields the explicit tangential function r6 H(t)\textsuperscript{2}(r ell(t)−1). Independently recomputed its derivative and all3 nonzero-value, nondegenerate critical points, including their common rational value and Hessian determinant. The implicit-function theorem lifts these to actual nonzero-value Morse points with critical values of order epsilon7. Boundary A3 is preserved because only U changes and its outer transverse coefficients stay units. The unchanged a3\_three\_axes.py replay passed. Review provenance: New derivation after reading the supplied proof; not a blind review and not an independent human check}{\ref{night:P4.41}}

Consider an actual torus degeneration at a point where the
entire reduced critical curve is the three coordinate axes,
with generic transverse type $A_3$ on each axis. Suppose
the local Hessian--conductor defect is nine, the three
Hessian orders are two, and the quartic equation is the
square of $a xy+b xz+c yz$ with $abc\ne0$.
Then the normalization has only rational singularities
over that point. In particular it cannot contain the
contraction of a positive-genus horizontal section.
\end{proposition}
\begin{proof}
\noindent\textbf{Step 1.} Put $P=(xy,xz,yz)$ and $h=xyz$. The symbolic square is
$P^{(2)}=P^2+(h)$, and the actual conductor image ideal
$I$ contains $P^{(2)}$. Also $f\in P^{(2)}$.
Write $f=Ah+P^2$. On the $x$-axis, a nonzero constant
term of $A$ would give a leading rank-two transverse
Hessian, since the quadratic term is then a nonzero
multiple of $xyz$. Generic type $A_3$ has rank one.
Thus $A(0)=0$. Since $\mathfrak m h\subset P^2$,
we have $f\in P^2$.

\noindent\textbf{Step 2.} The normalized conormal lattice is $tK$ on each axis,
so its projected colength is one. The curve has
$\delta=2$ and $R=0$. The stated defect therefore
forces the remaining normalization descent colength
to be one. The normalized quotient values form a
two-dimensional subspace of $\mathbb C^3$ with one
relation involving all three coordinates. Multiplying
a source generator that is a unit on one branch by
that branch's axis coordinate generates its positive
orders alone. Hence all positive-order branch functions
occur in the image, and the leading kernel vector
lifts to an actual element $g\in I$ with quadratic
part $a xy+b xz+c yz$.

\noindent\textbf{Step 3.} We prove the exact ideal equality $I=(g,h)$.
On each coordinate plane, $g$ is the product of its
two coordinate functions times a unit. Thus
$I_0=(g,h)$ has radical $P$. Furthermore $P^2\subset I_0$:
modulo $h$, multiplication of $g$ by $xy$ gives a unit
times $(xy)^2$, and the other two squares follow in
the same way; all mixed products are multiples of $h$.
At the generic point of every axis, $I_0$ is the
tacnode conductor ideal, generated by a normal coordinate
and the square of the null coordinate. It agrees
generically with $I$. The ideal $I_0$ is a
codimension-two complete intersection, so
$\mathcal O/I_0$ is Cohen--Macaulay and $P/I_0$ is
torsionfree on the reduced curve. The map
$P/I_0\to P/I$ has finite-length kernel, which must
vanish. This proves $I=I_0$.

\noindent\textbf{Step 4.} The actual conductor ring condition, with the hypotheses
verified in P4.41, gives $f\in I^2$ by the two-by-two
adjugate calculation of de Jong--van Straten,
Theorem~1.12. Write $f=A g^2+Bgh+Ch^2$.
Its quartic part makes $A(0)\ne0$. Completing the square
and absorbing $\sqrt A$ into the first factor gives
the exact equation
\[
                         f=G^2+h^2U.                 \tag{4.47a}
\]
The quadratic part of $G$ is nondegenerate, and
$(G,h)=(g,h)$. If $U(0)\ne0$, the two factors
$G\pm h\sqrt{-U}$ have nondegenerate quadratic
Hessians. Each is an isolated $A_1$ surface singularity,
so the full normalization has rational singularities.

\noindent\textbf{Step 5.} We exclude $U(0)=0$ by a deformation that keeps the
three axes fixed. Add an arbitrarily small generic
linear form to $U$, and then replace $U$ by
$U-\epsilon$. On a fixed boundary the transverse
$A_3$ coefficient is a unit, so these changes retain
the boundary transverse type. A diagonal linear change
puts the quadratic part of $G$ in the form
$G_2=xy+xz+yz$, changing $h$ only by a nonzero constant.
Let $L$ be the resulting linear part of $U$.
Under $(x,y,z)=\epsilon(X,Y,Z)$, equation (4.47a) becomes
\[
 \frac{f_\epsilon}{\epsilon^4}
   =G_\epsilon(X,Y,Z)^2+\epsilon^3(XYZ)^2
                          \bigl(L(X,Y,Z)-1+O(\epsilon)\bigr),
 \qquad G_\epsilon=G_2+O(\epsilon).
\]
Near a nonzero point of the smooth cone $G_2=0$,
use $G_\epsilon$ as the normal coordinate. The
parameter Morse centre is $O(\epsilon^3)$.
Substitution leaves the tangential function
$\epsilon^3(\Phi+O(\epsilon))$, where
\[
                 \Phi=(XYZ)^2(L-1)|_{G_2=0}.
\]

\noindent\textbf{Step 6.} We verify a nonempty open set of choices of $L$ for
which $\Phi$ has a nondegenerate critical point of
nonzero value. Parametrize the projective conic by
\[
              [X:Y:Z]=[t(t+1):t+1:-t].
\]
For $L=X+Y+Z$, set
$H(t)=-t^2(t+1)^2$ and $\ell(t)=t^2+t+1$.
In radial coordinate $r$,
$\Phi=r^6H(t)^2(r\ell(t)-1)$.
Radial criticality gives $r=6/(7\ell(t))$, and
\[
 \frac{d}{dt}\left(\frac{H}{\ell^3}\right)
  =\frac{t(t-1)(t+1)(t+2)(2t+1)}{(t^2+t+1)^4}.
\]
The three critical points $t=1,-2,-1/2$ are simple
and avoid zeros and poles. Their radial coordinates
are $2/7,2/7,8/7$. At each point,
\[
       \Phi=-\frac{1024}{823543},\qquad
       \det\operatorname{Hess}_{r,t}\Phi
                     =\frac{4718592}{1977326743}.
\]
These numbers are nonzero. Simplicity and avoidance
of zeros and poles are nonvanishing algebraic conditions
on the coefficients of $L$. The explicit example proves
that they hold on a nonempty Zariski open set, so an
arbitrarily small generic linear perturbation of any
initial linear part is sufficient.

\noindent\textbf{Step 7.} After eliminating the normal Morse coordinate, divide
the two tangential critical equations by $\epsilon^3$.
They are analytic in $\epsilon$ and have invertible
Jacobian at the displayed points. The implicit-function
theorem lifts them to nondegenerate three-variable
critical points of $f_\epsilon$. Their critical values
have leading term $\epsilon^7$ times the displayed
nonzero value. They tend to the original point.
Since the deformation retains boundary $A_3$, the
nested boundary transport embeds a Morse vanishing
sphere of square $-2$ into a nearby global torus fibre.
This contradicts asphericity. Therefore $U(0)=0$
cannot occur, proving the assertion.
\end{proof}

The conic derivative, all three exact critical values
and Hessian determinants, and the related conductor
matrices are executed in
\texttt{code/part4/a3\_three\_axes.py}.
The output reports PASS. The working record explicitly
rejects the earlier proposed change $G\mapsto G-\epsilon$:
with $h$ fixed that change replaces the boundary tacnode
by two ordinary double points and does not supply the
required boundary transport.


\subsection{The three-branch tacnode defect budget}

\begin{proposition}[P4.48]\label{night:P4.48}\label{stmt:P4.48}
\StatementRecord{night:P4.48}{proved}{Historical events: GPT-6 Astra (OpenAI), ChatGPT Pro mode, effort Pro, S6\_06\_BLIND, 16--18 November 2026: targeted proof check. Also checked in this writing. Pre-continuation writing events: GPT-6 Astra (OpenAI), ChatGPT Pro mode, effort Pro, S6\_10\_PAPERS, 18 November 2026 05:10:56.410177+00:00: computation or formal certificate; targeted proof check. Continuation metadata events: GPT-6 Astra (OpenAI), ChatGPT Pro mode, effort Pro, S6\_10\_PAPERS continuation, 19 November 2026: used and rechecked in this run; status-vocabulary reconciliation; metadata reconciliation only; no new mathematical check. \emph{Pre-continuation recorded proof scope.} The night-run re-derivation was by GPT-6 Astra (OpenAI, ChatGPT Pro mode, effort Pro), 16--18 November 2026. Also checked in this writing by GPT-6 Astra (OpenAI), ChatGPT Pro mode, effort Pro, 18 November 2026. proved. Full planar/ordinary-axis triple defect bound and exact quartic-square classification. A planar conductor quotient is generated by2 elements while its three normalized branches need3 values, so q\ensuremath{\geq}1 and 3delta+q\ensuremath{\geq}10. At the three axes, P\textasciicircum{}(2)=P\textsuperscript{2}+(xyz) and the leading transverse determinant force f\ensuremath{\in}P\textsuperscript{2}; all normalized conormal lattices are exactly tK. Re-derived the three principal-minor equations for the quartic tensor. When all diagonal coefficients are nonzero, substituting the quadratic null direction on the x-axis makes its cubic coefficient linear in x proportional to F−BD/A; higher total-degree terms cannot change that coefficient. Hence the full quartic is a square. A zero square gives defect\ensuremath{\geq}19/2; a single nonzero coefficient leaves one Hessian order\ensuremath{\geq}3 and a quotient-value matrix of rank\ensuremath{\leq}2, again giving\ensuremath{\geq}19/2. At least two coefficients give Hessian orders(2,2,2), projected orders(1,1,1), and the displayed rank-two three-row matrix forces q\ensuremath{\geq}1. Equality is exactly the retained data. The finite matrix and rational-cone checks passed in the unchanged replay. Review provenance: New derivation after reading the supplied proof; not a blind review and not an independent human check}{\ref{night:P4.42}; \ref{night:P4.47}}

At a planar three-branch critical-curve point with uniform
generic transverse type $A_3$, the local conductor defect
is at least ten. At an ordinary nonplanar three-branch
point, the defect is at least nine. Equality in the
second assertion forces Hessian orders $(2,2,2)$,
projected colengths $(1,1,1)$, descent colength one,
and a quartic square $(a xy+b xz+c yz)^2$ with at
least two nonzero coefficients. The case $abc\ne0$
has rational normalization by P4.47.
\end{proposition}
\begin{proof}
\noindent\textbf{Step 1.} For a planar critical curve the ideal is $P=(u,h)$.
The rank-one nilradical $J=P/I$ of the conductor image
is generated by at most two elements over $\mathcal O_C$.
Its normalized image has at least one generator value
on each of the three branches. Modulo the three branch
maximal ideals, the source supplies a vector space of
dimension at most two, while the normalized image
supplies dimension three. Hence its descent colength
$q$ is at least one. The planar symmetric-square and
different calculation in P4.42 gives local defect
$\ge3\delta+q$. Since a planar three-branch point has
$\delta\ge3$, its defect is at least ten.

\noindent\textbf{Step 2.} For an ordinary nonplanar point use actual coordinates
with $P=(xy,xz,yz)$. The argument in P4.47 gives
$f\in P^2$. On each axis, its normalized conormal
lattice is exactly $tK$, so the projected colength
is exactly one. Every Hessian order is at least two.
Write the quartic form in the basis $(xy,xz,yz)$ as
\[
 A(xy)^2+C(xz)^2+E(yz)^2
      +2B(xy)(xz)+2D(xy)(yz)+2F(xz)(yz).
\]
The rank-one Hessians give
$AC=B^2$, $AE=D^2$, $CE=F^2$.
If a diagonal coefficient vanishes, its adjacent
off-diagonal coefficients vanish and the remaining
two-by-two block has rank at most one. If all three
diagonal coefficients are nonzero, the vanishing null
cubic on the $x$-axis gives $F=BD/A$. Higher-degree
terms cannot alter the coefficient linear in $x$
of that normal cubic. Thus the whole quartic is a
square $(a xy+b xz+c yz)^2$.

\noindent\textbf{Step 3.} If the quartic is zero, all three Hessian orders are
at least three, and the local defect is at least
$2+3+9/2=19/2$. If the square has just one nonzero
coefficient, say $a xy$, the $z$-axis Hessian order
is at least three. In the source-generator basis
$(xy,xz,yz)$, the three leading quotient-value rows
have the forms
\[
                      (0,1,0),\qquad(0,0,1),
                      \qquad(0,*,*).
\]
Their rank is at most two, so $q\ge1$ and again
the defect is at least $19/2$.

\noindent\textbf{Step 4.} If at least two coefficients are nonzero, all three
Hessian orders are two. The leading quotient-value
matrix has rank two. For example, up to rescaling
its rows it is
\[
 \begin{pmatrix}-b&a&0\\-c&0&a\\0&-c&b\end{pmatrix}.
\]
It has zero determinant and rank two whenever at
least two of $a,b,c$ are nonzero. Therefore $q\ge1$,
and the local defect is at least $2+3+3+1=9$.
Equality forces $q=1$ and the displayed Hessian and
projection orders. P4.47 applies when all three
quadratic coefficients are nonzero.
\end{proof}


\subsection{The remaining ordinary-axis conductor}

\begin{proposition}[P4.49]\label{night:P4.49}\label{stmt:P4.49}
\StatementRecord{night:P4.49}{proved}{Historical events: GPT-6 Astra (OpenAI), ChatGPT Pro mode, effort Pro, S6\_06\_BLIND, 16--18 November 2026: targeted proof check. Also checked in this writing. Pre-continuation writing events: GPT-6 Astra (OpenAI), ChatGPT Pro mode, effort Pro, S6\_10\_PAPERS, 18 November 2026 05:10:56.410177+00:00: computation or formal certificate; targeted proof check; used and rechecked in this run. Continuation metadata events: GPT-6 Astra (OpenAI), ChatGPT Pro mode, effort Pro, S6\_10\_PAPERS continuation, 19 November 2026: used and rechecked in this run; status-vocabulary reconciliation; metadata reconciliation only; no new mathematical check. \emph{Pre-continuation recorded proof scope.} The night-run re-derivation was by GPT-6 Astra (OpenAI, ChatGPT Pro mode, effort Pro), 16--18 November 2026. Also checked in this writing by GPT-6 Astra (OpenAI), ChatGPT Pro mode, effort Pro, 18 November 2026. proved. Full two-form conductor classification, both normalization algebras, all polar deformations and exact rational-singularity deductions. Lifted the missing-coefficient quadratic kernel to g=x(y+z)−y z c. The Hilbert–Burch minors and the torsionfree generic comparison give I=(g,xyz,(yz)\textsuperscript{2}). Analytic division modulo(x,y+z,yz), followed by the invertible map x\ensuremath{\to}x−yzB, gives exactly the two lambda=0,1 ideals. For lambda0 independently solved every adjugate condition using I:h,I:x,I:k, obtaining f=A x\textsuperscript{2}k\textsuperscript{2}+B x\textsuperscript{2}kh+Cxh\textsuperscript{2}+Dh4 and the stated square/coordinate completion. A unit C gives a smooth factor and the finite birational normal quadric u(y+z)=yz with X=−u\textsuperscript{2}. C(0)=0 gives the explicit weighted polar; recomputed its reduced critical polynomial, nonzero critical values and Hessian resultants. For lambda1 independently followed the three colon ideals, the coefficient module and the exact identities reducing to Ag\textsuperscript{2}+Bgxh+Cgh\textsuperscript{2}+D(xh)\textsuperscript{2}+kappa p. Re-derived the four multiplication relations from the actual adjugate fractions v=−g/h and w=−(ak+bh)/h\textsuperscript{2}. For kappa a unit, the last relation eliminates k, and the displayed invertible linear map identifies the tangent-cone ideal with the prime rational cubic cone ideal; same dimension prevents an additional initial equation. Its blowup is smooth along the degree3 rational exceptional curve, whose O(3n) thickening layers have H1=0, giving rationality by formal functions. For kappa(0)=0, completed the normal square and used XH\textsuperscript{3}=−Yp,H4=−Kp on G=0 to obtain the stated tangential polar. Its exact nonzero value and Hessian at(t,r)=(1/2,−7/6) give actual Morse points of order epsilon8. Both deformations preserve the fixed-axis A3 boundary data. The unchanged a3\_missing\_coefficient.py passed all colon, multiplication, resultant and tangent-cone checks. Review provenance: New derivation after reading the supplied proof; not a blind review and not an independent human check}{\ref{night:P4.41}; \ref{night:P4.47}; \ref{night:P4.48}}

At an ordinary three-axis critical-curve point with uniform
generic transverse type $A_3$ and local defect nine, the
normalization has rational singularities. This includes the
quartic square with exactly one missing coefficient left
by P4.47--P4.48.
\end{proposition}
\begin{proof}
\noindent\textbf{Step 1.} Only the one-missing-coefficient case remains. Rescale
its leading conductor form to $x(y+z)$.
As in P4.47, the descent colength one lifts the leading
kernel vector to $g\in I$ with this quadratic part.
Write $g=x(\alpha y+\beta z)+yz\gamma$, where
$\alpha,\beta$ are units and $\gamma(0)=0$.
Replacing $y,z$ by $\alpha y,\beta z$ is an analytic
coordinate change preserving the axes. Put $k=y+z$,
$h=yz$. Then $g=xk-hc$, with $c(0)=0$.
The ideal $I_0=(g,xh,h^2)$ lies in the actual conductor
$I$ and has Hilbert--Burch matrix
\[
                 \begin{pmatrix}h&0\\-k&h\\c&-x\end{pmatrix}.
\]
It has radical $(xy,xz,yz)$, is Cohen--Macaulay, and
agrees with $I$ generically on each axis. The
torsionfree finite-kernel comparison used in P4.47
therefore gives the exact equality $I=I_0$.

\noindent\textbf{Step 2.} There are exactly two required ideal forms. Since
$\mathcal O/(x,k,h)=\mathbb C\{y\}/(y^2)$, analytic
division writes $c=\lambda y+xA+kB+hC$ for a scalar
$\lambda$. Replace $g$ by $g+xhA+h^2C$, and then
set $x'=x-hB$. These generator and coordinate changes
give
\[
                    I=(x'k-\lambda yh,x'h,h^2).
\]
If $\lambda=0$ this is $(xk,xh,h^2)$. If
$\lambda\ne0$, rescaling $x'=\lambda X$ gives
$(Xk-yh,Xh,h^2)$. Every change is invertible near
the origin. We treat the two forms separately.

\paragraph{The ideal $(xk,xh,h^2)$.}
Appending $(a,b,c)$ to its Hilbert--Burch matrix gives
determinant $axk+bxh+ch^2$. The adjugate ring condition
of de Jong--van Straten, Theorem~1.12, applies to the
actual conductor as in P4.41. The colon ideals
\[
 I:h=(x,h),\qquad I:x=(k,h),\qquad I:k=(x,h^2)
\]
give exactly
$f=A x^2k^2+B x^2kh+Cxh^2+Dh^4$, with $A(0)\ne0$.
Complete the square in $k$, absorb $\sqrt A$ into $x$,
and replace $y$ by $y(1+\gamma z)$ to make
$k+\gamma h$ the new $y+z$. These are analytic
coordinate maps; $h$ changes only by a unit.
Absorbing that unit in the coefficients gives
\[
                       f=x^2k^2+xh^2C+h^4D.          \tag{4.49a}
\]
If $C(0)\ne0$, substitute $x=h^2v$ and divide by
$h^4$. The finite $v$-root is simple, giving a smooth
factor $x-h^2v(y,z)$. Its quotient is $xk^2+h^2U$
with $U$ a unit. The coordinate $X=x/U$ turns it
into $Xk^2+h^2$. Its finite birational normalization is
\[
                 u=h/k,\qquad X=-u^2,
                 \qquad u(y+z)=yz.
\]
The last equation is an $A_1$ quadric, so the
normalization has rational singularities.

\noindent\textbf{Step 3.} Suppose $C(0)=0$. Perturb $D$ by a small constant
and the linear part of $C$ generically, retaining
$C(0)=0$, and then replace $C$ by $C-\epsilon$.
On a fixed boundary the transverse tacnode coefficients
are units, so these changes retain boundary $A_3$.
Under $x=\epsilon^3X$, $y=\epsilon Y$,
$z=\epsilon Z$, the leading scaled equation is
\[
 \frac f{\epsilon^8}=F+O(\epsilon),\qquad
 F=X^2K^2+XH^2(L-1)+D_0H^4,
 \quad K=Y+Z,\ H=YZ.
\]
For $L=K$, $D_0=1$, put $W=Y-Z$.
At $W=0$ the $X$-critical coordinate is
$X=-K^2(K-1)/32$, and the remaining derivative is
$K^5(12K^2+7K-3)/512$.
At either root of $12K^2+7K-3$, the full critical
value and Hessian determinant are respectively
\[
 \frac{K^6(K+1)(3K-1)}{1024},\qquad
 -\frac{3K^{10}(K+1)(3K-1)(28K^2+14K-5)}{32768}.
\]
They are nonzero: their numerator resultants with the
critical polynomial are $8748$ and $4923269424$.
Nondegeneracy, nonzero value and avoidance of the axes
hold on a nonempty algebraic open set of coefficients.
An arbitrarily small generic perturbation is therefore
available. The implicit-function theorem lifts these
Morse points through the higher weighted terms, with
nonzero critical values of order $\epsilon^8$.
The boundary comparison embeds their square-$-2$
vanishing spheres in a nearby torus fibre, a
contradiction. Thus $C(0)=0$ is excluded.

\paragraph{The ideal $(g,xh,h^2)$, where $g=xk-yh$.}
Append $(a,b,d)$ to
\[
                         \begin{pmatrix}h&0\\-k&h\\y&-x\end{pmatrix}.
\]
Its determinant is $ag+bxh+dh^2$. The new adjugate
entries give the conditions
\[
       ax,\ ah,\ bx+dh,\ kd+by,\ dh-ay,\ ak+bh\in I.
\]
The exact colon identities are
\[
 I:h=(x,h),\qquad I:x=(xk,h,z^2),\qquad I:x^2=(k,h).
\]
The first two conditions give $a=hA$ modulo $I$;
then $dh-ay$ gives $d=Ay+xD+hE$.
Initially $bx+dh$ permits
$b=-Ak+hB+z^2B_2$ modulo $I$.
The condition $ak+bh$ forces $z^2B_2\in(x,h)$,
so $B_2\in(x,y)$. Its $x$-part lies in $I$, since
$xz^2=h^2-xh$ modulo $I$, and its $y$-part is
absorbed in $hB$. Thus $b=-Ak+hB$ modulo $I$.
The remaining condition is $y(D+B)+kE\in(x,h)$.

\noindent\textbf{Step 4.} Write $T=D+B$. Modulo $(x,h,k)$, the equality
$yT=0$ in $\mathbb C\{y\}/(y^2)$ forces $T(0)=0$.
Write $T=xW+yU+kV$. Since $k$ is a nonzerodivisor
modulo $(x,h)$, the general solution is
$E=-yU-yV+xE_1+hE_2$. Substitution in the determinant
gives $f\in(p)+I^2$, where
\[
                            p=yh^2(x-h).
\]
The exact identities
$xh^3=ygh^2-yp$ and $h^4=ygh^2-kp$ then put the
equation in the form
\[
 f=A g^2+B gxh+Cgh^2+D(xh)^2+\kappa p,
                       \qquad A(0)\ne0.              \tag{4.49b}
\]

\noindent\textbf{Step 5.} We first treat $\kappa(0)\ne0$.
Use the appended coefficients
$a=Ag+Bxh+Ch^2$, $b=Dxh+\kappa yh$,
$d=-\kappa yh$. The actual conductor identifies
$\operatorname{Hom}(I,\mathcal O_S)$ with the
normalization, generated by $1,n_1,n_2$, where the
adjugate fractions are
\[
          n_1=-a/h,\qquad n_2=-(ak+bh)/h^2.
\]
Since $A$ is a unit, the same actual normalization
algebra is generated by $1,v,w$, where
$v=(n_1+Bx+Ch)/A=-g/h$ and $w=n_2$.
Direct multiplication gives the exact relations
\begin{align*}
 Av^2-Bxv-Chv+Dx^2+\kappa xy-\kappa yh&=0,\\
 xw&=Ayv-Bxy-(\kappa+C)yh,\\
 vw&=\kappa h-Dxy,\\
 w^2+[(\kappa+B)y+Ck]w-A\kappa k+ADy^2
             -D(\kappa+C)yh-B\kappa h&=0.
\end{align*}
At the origin they give $v=w=0$, so there is one
normalization point. The last equation eliminates $k$
analytically, with $k$ in the maximal-ideal square.
Thus $h=y(k-y)$ has quadratic initial term $-y^2$.
Write $a=A_0$, $b=B_0$, $d=D_0$, $u=\kappa_0$.
The three quadratic equations become
\[
 av^2-bxv+dx^2+uxy=0,\qquad
 xw-ayv+bxy=0,\qquad vw+uy^2+dxy=0.
\]
Set $V=av-(b/2)x$, $\delta=ad-b^2/4$,
$Y=\delta x+auy$, $W_0=w+(b/2)y$, and
$W=auW_0+\delta V$. This invertible linear change
turns them into
\[
                       V^2+xY=0,\qquad
                       VW+Y^2=0,\qquad xW-YV=0.
\]
They define the irreducible rational normal cubic cone.
The full tangent-cone ideal contains this prime ideal
and has the same dimension two, so it equals it.
Blowing up the normalization point is smooth along its
exceptional rational normal cubic: the two independent
initial equations have independent derivatives along
that smooth projective curve in every blowup chart.
The exceptional curve is $\mathbb P^1$ with normal
bundle $\mathcal O(-3)$. The exact sequences of its
thickenings have successive terms $\mathcal O(3n)$,
all with $H^1=0$. Formal functions gives
$R^1\rho_*\mathcal O=0$. The singularity is rational.

\noindent\textbf{Step 6.} Finally suppose $\kappa(0)=0$. Perturb $D$ and the
linear part of $\kappa$ generically, and replace
$\kappa$ by $\kappa-\epsilon$. The fixed axes retain
boundary $A_3$: equation (4.49b) stays in the same
ring-condition module, and its boundary transverse
quartic coefficients remain units.
Scale $x=\epsilon^2X$, $y=\epsilon Y$, $z=\epsilon Z$.
Put $K=Y+Z$, $H=YZ$, $G=XK-YH$.
After division by $\epsilon^6$, the first terms are
\[
 A_0G^2+\epsilon(B_0GXH+C_0GH^2)
   +\epsilon^2\bigl[D_0(XH)^2+p(L-1)\bigr].
\]
Near a nonzero smooth point of $G=0$, eliminate the
normal Morse coordinate. On that cone the identities
$XH^3=-Yp$, $H^4=-Kp$ absorb the square-completion
corrections into the coefficients, leaving the leading
tangential function
\[
                    \Phi=D'(XH)^2+p(L'-1)|_{G=0}.
\]
Choose $D'=1$, $L'=-Z$ to exhibit a nonempty open
set of suitable coefficients. The parametrization
\[
 K=r,\quad Y=tr,\quad Z=(1-t)r,\quad
 X=t^2(1-t)r^2
\]
gives $\Phi=t^4(1-t)^4r^7[1+r(t^2-t+1)]$.
At $t=1/2$, $r=-7/6$, its value and Hessian are
\[
 -\frac{823543}{573308928},\qquad
 \operatorname{diag}\left(\frac{117649}{1990656},
                          \frac{15647317}{214990848}\right).
\]
The Hessian determinant is
$1840891197733/427972821516288$, nonzero.
The angular critical numerator factors as
$(2t-1)(3t^2-3t-4)$, so this point is simple and
avoids the axes and $K=0$. The generic coefficient
condition is algebraic and nonempty. After normal
elimination, divide the tangential critical equations
by $\epsilon^2$ and apply the implicit-function theorem.
The resulting full Morse points have nonzero critical
values of order $\epsilon^8$ and tend to the origin.
Boundary transport gives the forbidden square-$-2$
vanishing sphere in a nearby torus fibre. Thus
$\kappa(0)=0$ is excluded, completing both ideal forms.
\end{proof}

All adjugates, colon ideals, normalization relations,
linear coordinate identities, and both polar Hessian
calculations are executed in
\texttt{code/part4/a3\_missing\_coefficient.py}.
The saved output reports PASS.

\begin{proposition}[P4.50]\label{night:P4.50}\label{stmt:P4.50}
\StatementRecord{night:P4.50}{proved}{Historical events: GPT-6 Astra (OpenAI), ChatGPT Pro mode, effort Pro, S6\_06\_BLIND, 16--18 November 2026: targeted proof check. Also checked in this writing. Pre-continuation writing events: GPT-6 Astra (OpenAI), ChatGPT Pro mode, effort Pro, S6\_10\_PAPERS, 18 November 2026 05:10:56.410177+00:00: computation or formal certificate; targeted proof check. Continuation metadata events: GPT-6 Astra (OpenAI), ChatGPT Pro mode, effort Pro, S6\_10\_PAPERS continuation, 19 November 2026: used and rechecked in this run; status-vocabulary reconciliation; metadata reconciliation only; no new mathematical check. \emph{Pre-continuation recorded proof scope.} The night-run re-derivation was by GPT-6 Astra (OpenAI, ChatGPT Pro mode, effort Pro), 16--18 November 2026. Also checked in this writing by GPT-6 Astra (OpenAI), ChatGPT Pro mode, effort Pro, 18 November 2026. proved. Full final uniform-A3 positive-Euler exclusion, with every local budget alternative followed. The verified P4.42/P4.44 maps leave a contracted positive-genus section, two rational normalized conductor components, one triple point and exact total defect9. P4.48 excludes planar triples; P4R.16 excludes the tacnode-plus-line excess-one model; P4R.17 exhausts delta4 nonplanar models and proves each bound>9. The uniform branch estimate delta+3+m/2 excludes delta\ensuremath{\geq}5. At delta2 the exact ring is the three axes, which consumes the full budget, and P4.47/P4.49 make its normalization rational. Any additional singular critical-curve point contributes positively and is absent. At each smooth critical point the actual forms in P4R.14 give smooth or A1 normalization, and R6 excludes exceptional blocks away from the conductor. Rational normalization singularities cannot contract H: the ruled restriction H1(T,O)\ensuremath{\to}H1(H,O) is an isomorphism, whereas a contraction would kill that map, as quantified in P4.33. This gives the stated contradiction. Review provenance: New derivation after reading the supplied proof; not a blind review and not an independent human check}{\ref{night:P4.33}; \ref{night:P4.42}; \ref{night:P4.44}; \ref{night:P4.47}; \ref{night:P4.48}; \ref{night:P4.49}; \ref{night:P4R.14}; \ref{night:P4R.16}; \ref{night:P4R.17}; \ref{night:P4R.6}}

In the requested full-fibration classes, no support
with positive-genus ruled resolution has positive
Euler number and generic transverse type $A_3$
on every critical-curve component.
\end{proposition}
\begin{proof}
\noindent\textbf{Step 1.} P4.42 and P4.44 leave only a contracted horizontal
positive-genus section and a cycle at one three-branch
point of a conductor with two rational normalized
components. The total defect is nine by P4.33 and
(4.42a). P4.48 excludes planar triples; P4R.16
excludes the nonplanar excess-one tacnode-plus-line
case; the exact excess-two classification and conormal
calculation in P4R.17 exclude
every delta-four nonplanar triple. The branchwise bound
$\delta+3+m/2$ excludes all larger delta invariants,
where $m\ge3$ is total multiplicity.
The only remaining triple is the ordinary three-axis
point. P4.48 gives local defect at least nine, so it
uses the entire budget, and P4.47 and P4.49 give
rational normalization there.

\noindent\textbf{Step 2.} Every additional singular critical-curve point has
strictly positive defect and therefore is absent.
At a smooth point of the entire critical curve,
P4R.14 gives the product tacnode or its $A_1$
normalization degeneration; both have rational
normalization. Away from the conductor, P4R.6
excludes exceptional configurations. The normalization
therefore has rational singularities everywhere.
Its minimal resolution cannot contract the horizontal
positive-genus section. This contradiction proves
the assertion.
\end{proof}


\subsection{Local intersection forms and the quartic tangent obstruction}

\begin{proposition}[P4.51]\label{night:P4.51}\label{stmt:P4.51}
\StatementRecord{night:P4.51}{proved}{Historical events: GPT-6 Astra (OpenAI), ChatGPT Pro mode, effort Pro, S6\_06\_BLIND, 16--18 November 2026: targeted proof check. Also checked in this writing. Pre-continuation writing events: GPT-6 Astra (OpenAI), ChatGPT Pro mode, effort Pro, S6\_10\_PAPERS, 18 November 2026: targeted proof check. Continuation metadata events: GPT-6 Astra (OpenAI), ChatGPT Pro mode, effort Pro, S6\_10\_PAPERS continuation, 19 November 2026: used and rechecked in this run; status-vocabulary reconciliation; metadata reconciliation only; no new mathematical check. \emph{Pre-continuation recorded proof scope.} The night-run re-derivation was by GPT-6 Astra (OpenAI, ChatGPT Pro mode, effort Pro), 16--18 November 2026. Also checked in this writing by GPT-6 Astra (OpenAI), ChatGPT Pro mode, effort Pro, 18 November 2026. The local Milnor-cycle lift to the torus cover associated with the specialization kernel W was re-derived; dim W<=3 implies the pairing on Lambda\textasciicircum{}2 W is zero}{\ref{night:P4R.4}}

Let $S$ be an irreducible reduced support in the torus
fibration of Part~4.2 with $b_1(S)>0$. Every local
Milnor fibre of a reduced local equation of $S$ has
identically zero rational intersection form. This applies
in particular to the remaining positive ruled support,
whose fundamental group is $\mathbb Z$.
\end{proposition}
\begin{proof}
\noindent\textbf{Step 1.} Let $U$ be a sufficiently small full tube and $F_t$ a
nearby compact smooth torus fibre. P4R.4 identifies the
homomorphism
\[
 \pi_1(F_t)=\mathbb Z^4\longrightarrow
 \pi_1(U)=\pi_1(S)
\]
and gives a finite cyclic cokernel, of order equal to
the fibre multiplicity. Its map on rational first
homology is therefore onto. Write
\[
 W=\ker\bigl(H_1(F_t;\mathbb Q)\longrightarrow
                  H_1(S;\mathbb Q)\bigr).
\]
Then $\dim W=4-b_1(S)\le3$.

\noindent\textbf{Step 2.} Choose a coordinate ball about the point, with its
closed ball contained in $U$. The local full fibre
equation is $u g^m$, with $g$ reduced and $u$ a
holomorphic unit. We take a holomorphic $m$-th root of
$u$ on the ball and absorb it into $g$. The actual
nearby fibre in the ball is consequently the disjoint
union of the local Milnor fibres
$g=\zeta t^{1/m}$. Each such fibre $M$ is an oriented
codimension-zero submanifold with boundary of $F_t$.
Its inclusion into $U$ factors through the contractible
ball. Thus every loop of $M$ maps to zero in
$H_1(U;\mathbb Q)$, and the image subgroup
\[
 H=\operatorname{im}\bigl(\pi_1(M)\longrightarrow
                                 \mathbb Z^4\bigr)
\]
satisfies $H\otimes\mathbb Q\subset W$.

\noindent\textbf{Step 3.} The covering criterion lifts $M\to F_t$ to the
connected cover $F_H\to F_t$ corresponding to $H$.
This cover is a $K(H,1)$, with $H$ free abelian. The
map on second rational homology is therefore exactly
\[
 \bigwedge^2(H\otimes\mathbb Q)\longrightarrow
              \bigwedge^2 H_1(F_t;\mathbb Q).
\]
Every class from $H_2(M;\mathbb Q)$ consequently maps
into $\bigwedge^2W$. The torus intersection pairing is
exterior multiplication followed by the orientation
isomorphism $\bigwedge^4H_1(F_t;\mathbb Q)=\mathbb Q$.
It vanishes on $\bigwedge^2W$ because $\dim W\le3$.

\noindent\textbf{Step 4.} This ambient pairing is the local intersection
pairing on $M$. We represent finitely many cycles
inside its interior and make them transverse there;
the orientation-preserving inclusion into $F_t$
retains all intersection points and signs. Thus the
local pairing is zero. The symbolic exterior-square
matrix and its zero Gram matrix are recorded in
\texttt{code/part4/local\_intersection.py} and its
executed output.
\end{proof}

\begin{proposition}[P4.52]\label{night:P4.52}\label{stmt:P4.52}
\StatementRecord{night:P4.52}{proved}{Historical events: GPT-6 Astra (OpenAI), ChatGPT Pro mode, effort Pro, S6\_06\_BLIND, 16--18 November 2026: targeted proof check. Also checked in this writing. Pre-continuation writing events: GPT-6 Astra (OpenAI), ChatGPT Pro mode, effort Pro, S6\_10\_PAPERS, 18 November 2026: targeted proof check; used and rechecked in this run. Continuation metadata events: GPT-6 Astra (OpenAI), ChatGPT Pro mode, effort Pro, S6\_10\_PAPERS continuation, 19 November 2026: used and rechecked in this run; status-vocabulary reconciliation; metadata reconciliation only; no new mathematical check. \emph{Pre-continuation recorded proof scope.} The night-run re-derivation was by GPT-6 Astra (OpenAI, ChatGPT Pro mode, effort Pro), 16--18 November 2026. Also checked in this writing by GPT-6 Astra (OpenAI), ChatGPT Pro mode, effort Pro, 18 November 2026. The compact-cycle, conductor-pole, and multiplicity reduction were followed. Siersma 1991 section 7.1, printed page 462, gives (j2)\_*(T-I)=VAR'' j; sections 6.2-6.5, printed pages 460-461, preserve the required ordinary degree-two groups. The boundary Wang sequence yields the claimed rank bound; no degree-two variation isomorphism was assumed}{\ref{night:P4.20}; \ref{night:P4.39}; \ref{night:P4.51}; \ref{night:P4R.15}; \ref{night:P4R.6}}

In the remaining positive ruled support of Part~4.2,
the image of a contracted horizontal pole section has
ambient multiplicity exactly three. This restriction
imposes no generic-transverse-type hypothesis on
any component of the critical curve.
\end{proposition}
\begin{proof}
\noindent\textbf{Step 1.} We first give the comparison with a leading
homogeneous polynomial. Write the reduced equation
as $g=g_d+g_{d+1}+\cdots$, and suppose that $g_d$ is
reduced. Every finite collection of cycles in
$M_d=\{g_d=1\}$, together with its intersections,
transports into a local Milnor fibre of $g$.
Indeed, choose a compact neighbourhood of the cycle
representatives. Euler's identity shows that
$dg_d\ne0$ on $M_d$. The functions
\[
 G_\epsilon(x)=\epsilon^{-d}g(\epsilon x)
\]
converge with their derivatives to $g_d$ on that
neighbourhood. We solve
\[
 dG_\epsilon(V_\epsilon)
                     =-\partial_\epsilon G_\epsilon
\]
using the Hermitian gradient divided by its squared
norm. A cutoff supported in a slightly larger
compact neighbourhood, equal to one along the
trajectories under consideration, gives a flow for
small positive real $\epsilon$. It transports those
cycles to $G_\epsilon=1$. Holomorphic scaling by
$\epsilon$ then embeds them in $g=\epsilon^d$ inside
the fixed original ball, preserving orientations
and intersections. This compares smooth nonzero
level sets on compact neighbourhoods; no assertion
about equisingularity of the zero set is needed.
By P4.51 the pairing on $M_d$ is consequently zero.

\noindent\textbf{Step 2.} For a reduced homogeneous polynomial its global
fibre and a sufficiently large compact Milnor fibre
have the same homology and scalar monodromy. One
can see this directly from the squared radius on
the affine fibre. This is a proper semialgebraic
function with finitely many critical values; above
the largest, its gradient supplies a product end.
A large invariant ball therefore gives a compact
core. Rescaling identifies that core with a small
local Milnor fibre. The Milnor-fibre CW dimension
bound gives zero homology above degree two, and
reducedness gives connectedness.

\noindent\textbf{Step 3.} We need one exact variation identity. For such a
homogeneous polynomial, let $\mu_i$ be the
transverse plane-curve Milnor numbers at the
singular points of its reduced projective curve,
and put $M=\sum_i\mu_i$. Write $F$ for its compact
Milnor fibre and $\partial_2F$ for the boundary
part over the singular-link circles. Let
\[
 \begin{split}
 j&:H_2(F;\mathbb Q)\longrightarrow
                         H_2(F,\partial F;\mathbb Q),\\
 j_2&:H_2(F;\mathbb Q)\longrightarrow
                         H_2(F,\partial_2F;\mathbb Q).
 \end{split}
\]
Siersma's chain-level variation identity
\cite[\S3.6, pp.~451--452; \S7.1, p.~462]
{Siersma1991Variation} is
\[
                 j_2(T-I)=\operatorname{VAR}^{II}j.
                                                            \tag{4.52a}
\]
Here $T$ is the Milnor monodromy on $H_2(F;\mathbb Q)$.
The identity does not require the variation map to
be an isomorphism in dimension two. The dimension-two
Wang sequence is explicitly treated in that source,
\S6.3. Poincar\'e--Lefschetz duality identifies $j$
with the intersection form. Its vanishing makes
$\operatorname{im}(T-I)$ lie in $\ker j_2$.
The exact sequence of $(F,\partial_2F)$ identifies
that kernel with the image of $H_2(\partial_2F)$.
Each boundary piece is the mapping torus of a
transverse Milnor fibre. Its Wang sequence gives
\[
 H_2(\partial_2F;\mathbb Q)
       =\bigoplus_i\ker(A_i-I)\subset
                        \bigoplus_iH_1(F_i;\mathbb Q),
\]
where $A_i$ is the vertical monodromy. Thus zero
intersection form implies
\[
                 \operatorname{rank}(T-I)\le M.             \tag{4.52b}
\]
An empty $\partial_2F$ gives the same assertion with
$M=0$.

\noindent\textbf{Step 4.} Suppose now that the reduced projective quartic
$D=\{g_4=0\}$ has a normalization component of
positive genus. By degree, either $D$ is an
irreducible quartic with normalization genus
$g\ge1$, or it is a smooth elliptic cubic and a
line. In the first case
$\sum\delta_i=3-g\le2$, so
$M=\sum(2\delta_i-r_i+1)\le4$.
In the second case B\'ezout gives total intersection
length three and at least one intersection point.
Hence $\sum\delta_i=3$ and
$\sum(r_i-1)\ge1$, giving $M\le5$.

\noindent\textbf{Step 5.} The map $\{g_4=1\}\to\mathbb P^2\setminus D$ is
an unramified fourfold cover. The normalization
sequence and plane-curve genus formula give
\[
              \chi(F)=4\chi(\mathbb P^2\setminus D)
                      =4(7-M).
\]
The three nonidentity scalar monodromies act freely
on the compact invariant core, so their Lefschetz
numbers are zero. It follows by the character
formula for the cyclic group of order four that
\[
 H_0(F;\mathbb C)-H_1(F;\mathbb C)+H_2(F;\mathbb C)
       =(7-M)\,\mathbb C[\mathbb Z/4]
\]
as virtual representations. Every nontrivial
character therefore has multiplicity at least
$7-M$ in $H_2$. The finite-order monodromy is
semisimple, and hence
\[
             \operatorname{rank}(T-I)
                   \ge3(7-M)\ge6>M.                       \tag{4.52c}
\]
This contradicts (4.52b).

\noindent\textbf{Step 6.} We apply this obstruction through the actual pole
comparison. Blow up the ambient origin, normalize
the strict transform, and take its minimal
resolution $T_1$. As in P4.39, this resolution maps
to the original normalization and dominates its
minimal resolution by point blowups. The strict
transform of the positive-genus horizontal section
$H$ retains pole coefficient two. If $E$ is the
ambient exceptional divisor, put
$n=\operatorname{ord}_H(E)>0$, and let $a\ge0$
be the coefficient of $H$ in the total conductor
divisor relative to the strict-transform dualizing
sheaf. Adjunction gives exactly
\[
                          2=(d-2)n+a.                     \tag{4.52d}
\]
Effectiveness follows from the exceptional
intersection matrix argument of P4R.6. An
exceptional positive-genus curve over the new
normalization has strictly positive coefficient,
by its adjunction equation. Thus $d\ge5$ is
impossible. If $d=4$, then $n=1,a=0$; the curve
$H$ is neither exceptional over the new
normalization nor over its conductor. The
normalization is an isomorphism at its generic
point. Its image is a multiplicity-one component
of the quartic exceptional curve, birational to
$H$. It has positive genus. Such a component has
degree three or four, leaving at most one linear
factor; the entire quartic polynomial is reduced.
The preceding contradiction excludes this case.

\noindent\textbf{Step 7.} For clarity, an isolated critical point cannot
hide a contracted $H$. A generic small linear
perturbation of an isolated hypersurface critical
point has a positive finite number of Morse
critical points, the local degree of its gradient.
A small constant shift makes a chosen critical
value nonzero. On a fixed sufficiently small ball
the original zero-level boundary is smooth and
transverse; compactness retains the boundary
submersion over a uniform punctured value disc.
The finite-critical-value and path-transport
argument of P4.20 then carries the Morse
vanishing sphere, of square $-2$, into the
original local Milnor fibre. P4.51 excludes it.
This argument does not assume smoothness of the
normalization.

\noindent\textbf{Step 8.} At a nonisolated point with $d=2$, the Hessian is
nonzero. P4R.15 gives the exact product equation
$X^2+Y^m$. For even $m$ its two analytic branches
are smooth. For odd $m$ its normalization is the
explicit smooth parametrization
$(Y,X,Z)=(u^2,\mathrm i u^m,Z)$.
Neither case has an exceptional curve in its
minimal normalization resolution. A point with
$d=1$ is smooth as well. Therefore the image of
$H$ has multiplicity exactly three.

\noindent\textbf{Step 9.} The cyclic character calculation and all integer
bounds are independently recomputed in
\texttt{code/part4/quartic\_tangent.py}; no numerical
approximation enters the argument.
\end{proof}


\subsection{The cubic double line and its weighted leading fibre}

\begin{proposition}[P4.53]\label{night:P4.53}\label{stmt:P4.53}
\StatementRecord{night:P4.53}{proved}{Historical events: GPT-6 Astra (OpenAI), ChatGPT Pro mode, effort Pro, S6\_06\_BLIND, 16--18 November 2026: targeted proof check. Also checked in this writing. Pre-continuation writing events: GPT-6 Astra (OpenAI), ChatGPT Pro mode, effort Pro, S6\_10\_PAPERS, 18 November 2026: targeted proof check. Continuation metadata events: GPT-6 Astra (OpenAI), ChatGPT Pro mode, effort Pro, S6\_10\_PAPERS continuation, 19 November 2026: used and rechecked in this run; status-vocabulary reconciliation; metadata reconciliation only; no new mathematical check. \emph{Pre-continuation recorded proof scope.} The night-run re-derivation was by GPT-6 Astra (OpenAI, ChatGPT Pro mode, effort Pro), 16--18 November 2026. Also checked in this writing by GPT-6 Astra (OpenAI), ChatGPT Pro mode, effort Pro, 18 November 2026. The pole identity, weighted degree-five discriminant, the connected double-cover branch bound, both coprime and common-factor cases, and the variation-rank contradiction were re-derived. The exceptional horizontal case was checked by adjunction and the minimal-resolution intersection inequality}{\ref{night:P4.51}; \ref{night:P4.52}}

In the remaining positive ruled support, blow up the
multiplicity-three image of the contracted horizontal
section $H$, normalize the strict transform $S_1$, and
take its minimal resolution $T_1$. The strict transform
of $H$ is exceptional over a point of $S_1$. The entire
germ of $S_1$ there is normal and has an isolated simple
elliptic singularity. In particular $H$ cannot dominate
a curve in the ambient exceptional plane. No generic
transverse type is imposed on the original critical curve.
\end{proposition}
\begin{proof}
\noindent\textbf{Step 1.} Use the comparison maps of P4.52, and let $D_1$ be the
effective total conductor relative to $\omega_{S_1}$.
Write $n=\operatorname{ord}_H(E)>0$ and
$a=\operatorname{coeff}_H(D_1)$. The pole identity is
$2=n+a$, so $(n,a)=(2,0)$ or $(1,1)$.

\noindent\textbf{Step 2.} Suppose $H$ dominates a reduced component $L$ of the
projective cubic tangent curve. At its generic point
a transverse Weierstrass presentation has degree
$m$, the multiplicity of that component in the cubic.
Factoring over the completed valuation field of the
exceptional coordinate $r$ gives the exact degree
identity
\[
 \sum_{H'}\operatorname{ord}_{H'}(r)
                 [\mathbb C(H'):\mathbb C(L)]=m.             \tag{4.53a}
\]
The factors group the roots by valuation branches;
the two factors in each summand are the ramification
index and residue-field degree. Residue extensions
are separable in characteristic zero.

\noindent\textbf{Step 3.} A positive-genus component of a cubic is a smooth
elliptic cubic of multiplicity one. The strict
transform is smooth along it, with $n=1,a=0$,
contrary to $2=n+a$. Every other reduced component
is rational. The finite map from $H$ then has degree
$h\ge2$. Equation (4.53a) excludes $n=2$, since
$m\le3$. Thus $n=a=1$.

\noindent\textbf{Step 4.} At a generic point of $L$, the $h$ points on $H$
give distinct branches of the transverse plane germ,
each with conductor exponent one. The branch formula
\[
             c_i=2\delta_i+\sum_{j\ne i}I_{ij}
\]
forces $h=2$ and exactly two smooth branches meeting
once, with no additional branch. Each has $r$-order
one, so their total Weierstrass degree is two. Hence
$m=2$, and the tangent polynomial is $y^2\ell$,
where $y=0$ is the double line and $\ell$ a different
linear form.

\noindent\textbf{Step 5.} Choose complementary coordinates $x,z$. The
restriction $\ell(x,z)$ is nonzero; any $y$ term in
$\ell$ has higher weight for the weights used below.
In a generic blowup chart along the double line,
write $v$ for the transverse coordinate. The
strict-transform equation begins
\[
                 \ell v^2+r(f_4|_L)+\cdots.
\]
A nonzero $f_4|_L$ would give ramification order two
and residue degree one. Therefore $f_4$ is divisible
by $y$. The two branches meet once, so after
$v=ru$ their distinct leading slopes solve
\[
                         \ell u^2+G_3u+G_5=0,
\]
where $G_3,G_5$ are forms of degrees three and five
in $x,z$. The discriminant
$\Delta=G_3^2-4\ell G_5$ is nonzero and nonsquare
in $\mathbb C(\mathbb P^1)$. The actual residue-field
extension of $H\to L$ is
$\mathbb C(\mathbb P^1)(\sqrt\Delta)$.

\noindent\textbf{Step 6.} Consequently the original equation has weighted
leading polynomial
\[
          P_5=\ell(x,z)y^2+G_3(x,z)y+G_5(x,z),               \tag{4.53b}
\]
of degree five for weights $(1,2,1)$ on $(x,y,z)$,
with no smaller-weight terms. This identifies the
normalized projective component birationally with
$H$ through its function field.

\noindent\textbf{Step 7.} Let $D=\{P_5=0\}\subset\mathbb P(1,2,1)$.
The quadratic is irreducible over $\mathbb C(x/z)$,
since the degree-two residue cover $H$ is connected.
The only possible further factor is a common
vertical factor of $\ell,G_3,G_5$. Since $\ell$
is linear, either $P_5$ is irreducible or
$P_5=\ell Q_4$, with $Q_4$ monic quadratic in $y$.
Both alternatives are reduced.

\noindent\textbf{Step 8.} In the irreducible case, let $k_i$ be the
multiplicities of the distinct zeros of the
homogeneous degree-six discriminant. The normalized
double cover has genus
$(\#\{i:k_i\text{ odd}\}-2)/2$. Positive genus
therefore requires at least four odd multiplicities,
and gives
\[
              \sum_i(k_i-1)=6-\#\{i\}\le2.                  \tag{4.53c}
\]
Away from $\ell=0$ and the weighted vertex,
completing the quadratic identifies the projective
curve germ with $w^2-s^{k_i}$, up to a unit and
holomorphic coordinates. Its transverse Milnor
number is $k_i-1$. At $\ell=0$, if $G_3\ne0$ the
finite root is simple and the other point is at
the weighted vertex. If $G_3=0$, absence of a
common factor gives $G_5\ne0$, and $\Delta$ has
a simple zero. There is no additional singularity.
In the chart $y=1$ the vertex is the quotient of
$\mathbb C^2$ by $(x,z)\mapsto(-x,-z)$; the
lifted equation $\ell+G_3+G_5$ has nonzero
differential. The original polynomial also has
nonzero derivative in the $\ell$ direction on
the $y$-axis, so that axis is not critical.
The total transverse Milnor-number sum is thus
$M\le2$.

\noindent\textbf{Step 9.} In the reducible case write
$G_3=\ell A_2$, $G_5=\ell A_4$. Then
$Q_4=y^2+A_2y+A_4$ has discriminant of degree four.
A connected positive-genus double cover requires
four distinct simple roots. Its normalization is
elliptic and $Q_4$ is smooth on the weighted plane;
it avoids the weighted vertex. The rational
component $\ell=0$ meets it in total length two,
read directly from the restricted monic quadratic.
Either there are two transverse intersections,
with total Milnor number two, or one tangency of
order two, with Milnor number three. There are no
other critical rays. Hence in all cases
\[
                              M\le3.                       \tag{4.53d}
\]

\noindent\textbf{Step 10.} The weighted plane has Euler number three: its
strata $x\ne0$, $x=0,z\ne0$, and $x=z=0$ are
$\mathbb C^2$, $\mathbb C$, and a point.
In the irreducible case $e(D)\le e(H)\le0$.
In the reducible case
$e(D)=0+2-\#(Q_4\cap\{\ell=0\})\le1$.
Therefore
\[
               \chi\bigl(\mathbb P(1,2,1)\setminus D\bigr)
                         \ge2.                            \tag{4.53e}
\]

\noindent\textbf{Step 11.} The weighted scalar action of a primitive fifth
root, with weights $(1,2,1)$, acts freely on
$F=\{P_5=1\}$. Its quotient is
$\mathbb P(1,2,1)\setminus D$: the only possible
nontrivial stabilizer of the weighted
$\mathbb C^*$ action is at the weighted vertex,
which belongs to $D$. This is an unramified
fivefold cover. The compact-core and character
calculation of P4.52 gives
\[
 \operatorname{rank}(T-I)
 \ge4\chi\bigl(\mathbb P(1,2,1)\setminus D\bigr)
 \ge8.
\]
The variation identity (4.52a) applies to every
one-dimensional hypersurface critical locus. Its
boundary mapping-torus estimate would give
$\operatorname{rank}(T-I)\le M\le3$ if the
intersection form were zero.

\noindent\textbf{Step 12.} It is zero here by the actual anisotropic scaling
comparison. Set
\[
 D_\epsilon(x,y,z)=(\epsilon x,\epsilon^2y,\epsilon z).
\]
Then $\epsilon^{-5}f(D_\epsilon(x,y,z))$ converges
with derivatives to $P_5$. The compact-cycle flow
of P4.52 transports its pairings into the actual
local Milnor fibre, whose pairing vanishes by
P4.51. This contradiction excludes every dominating
$H$.

\noindent\textbf{Step 13.} Thus $H$ is exceptional over the new normalization.
Its coefficient $a$ is positive, so $n=a=1$.
Adjunction gives
\[
                    0=2g(H)-2+(D_1-H)\cdot H.
\]
It follows that $H$ is elliptic and meets no other
positive pole component. A zero-coefficient
exceptional curve cannot meet a positive one:
its own intersection equation has nonnegative
canonical term on a minimal resolution and
nonnegative remaining terms, with a strictly
positive term at such a meeting. The connected
exceptional fibre is therefore $H$ alone and is
disjoint from the strict conductor. The
normalization is an isomorphism at its image;
another analytic hypersurface branch would
intersect this branch in a curve and create
conductor meeting $H$. Hence the entire
strict-transform germ is normal, with the single
smooth elliptic exceptional curve $H$. This is
precisely an isolated simple elliptic singularity.

\noindent\textbf{Step 14.} The discriminant multiplicity partitions and the
exact rank bounds are recomputed in
\texttt{code/part4/cubic\_weighted.py} and its
executed output.
\end{proof}


\subsection{Exclusion of the contracted horizontal section}

\begin{proposition}[P4.54]\label{night:P4.54}\label{stmt:P4.54}
\StatementRecord{night:P4.54}{proved}{Historical events: GPT-6 Astra (OpenAI), ChatGPT Pro mode, effort Pro, S6\_06\_BLIND, 16--18 November 2026: targeted proof check. Also checked in this writing. Pre-continuation writing events: GPT-6 Astra (OpenAI), ChatGPT Pro mode, effort Pro, S6\_10\_PAPERS, 18 November 2026: targeted proof check; used and rechecked in this run. Continuation metadata events: GPT-6 Astra (OpenAI), ChatGPT Pro mode, effort Pro, S6\_10\_PAPERS continuation, 19 November 2026: used and rechecked in this run; status-vocabulary reconciliation; metadata reconciliation only; no new mathematical check. \emph{Pre-continuation recorded proof scope.} The night-run re-derivation was by GPT-6 Astra (OpenAI, ChatGPT Pro mode, effort Pro), 16--18 November 2026. Also checked in this writing by GPT-6 Astra (OpenAI), ChatGPT Pro mode, effort Pro, 18 November 2026. The cofinal ideal filtrations, exact associated graded section ring, analytic coordinate changes preserving r, original-coordinate blowdown map, and all three weighted Milnor-fibre contradictions were re-derived. Grauert 1960 section 6 Hauptsatze I and IIa, printed page 59, was opened. The prime and order-four quotient-curve arguments were checked separately}{\ref{night:P4.51}; \ref{night:P4.52}; \ref{night:P4.53}; \ref{night:P4R.18}}

In the remaining positive ruled support of Part~4.2,
the horizontal pole section is not contracted by the
normalization resolution. It consequently maps with
degree two or three onto a rational normalized conductor
component whose generic transverse germ is $D_4$.
No transverse-type restriction is imposed on the other
conductor component.
\end{proposition}
\begin{proof}
\noindent\textbf{Step 1.} By P4.53, a contracted section would be the sole smooth
elliptic exceptional curve $H$ of an isolated normal
hypersurface singularity on the first ambient strict
transform. Its self-intersection is $-d<0$, and the
actual exceptional coordinate $r$ has order one along
$H$. Write $\pi:\widetilde V\to(V,0)$ for this local
resolution. We establish the following three coordinate
weights and leading degrees, retaining the actual $r$:
\[
 (1,1,1;3),\qquad(1,1,2;4),\qquad(1,2,3;6).
                                                        \tag{4.54w.1}
\]

\paragraph{1. Analytic formal functions and the exact associated graded ring}

\noindent\textbf{Step 2.} Put \(L=\mathcal O_H(-H)\), of degree \(d\), and
\[
 I_m=(\pi_*\mathcal O_{\widetilde V}(-mH))_0
     =\{a\in\mathcal O_{V,0}:\operatorname{ord}_H(a)\ge m\}.
\]
For \(m\ge1\), the restrictions of
\(\mathcal O(-mH)\) to successive infinitesimal neighbourhoods
of \(H\) have graded factors \(L^{m+j}\), \(j\ge0\).
Each has zero first cohomology, by Serre duality on the elliptic
curve and its strictly positive degree. Their extensions also
have zero first cohomology.

\noindent\textbf{Step 3.} The analytic formal-functions theorem is Grauert1960,
Section6, Hauptsatz IIa, printed p.59; Hauptsatz I on the same
page supplies coherence \cite[\S6, Hauptsätze I and IIa]{Grauert60}. The maximal-ideal pullback and the
ideal of \(H\) have the same radical. On a finite cover of
the compact exceptional curve, Noetherianity gives one uniform
power of the latter contained in the former. Their powers are
therefore cofinal. Applying the theorem with
\(\mathcal O(-mH)\) gives a zero completion for
\(R^1\pi_*\mathcal O(-mH)\). This coherent module is
supported at the origin, so it is zero. Consequently
\[
 I_m/I_{m+1}\simeq H^0(H,L^m)\qquad(m\ge1),             \tag{4.54w.2}
\]
through the actual restriction maps in
\[
 0\longrightarrow\mathcal O(-(m+1)H)
 \longrightarrow\mathcal O(-mH)
 \longrightarrow L^m\longrightarrow0.
\]
These maps respect multiplication. Since normality gives
\(\pi_*\mathcal O_{\widetilde V}=\mathcal O_V\), and a
function vanishes at the contracted point exactly when its
restriction to \(H\) vanishes, \(I_1=\mathfrak m\).
The degree-zero piece is \(\mathbb C\). Thus
\[
 \operatorname{gr}_{I}\mathcal O_{V,0}
   =\bigoplus_{m\ge0}H^0(H,L^m)                         \tag{4.54w.3}
\]
as graded algebras, not only as vector spaces.

\noindent\textbf{Step 4.} The same formal-functions theorem in degree zero identifies
\[
 \widehat{\mathcal O}_{V,0}
      \simeq\varprojlim_m H^0(mH,\mathcal O_{mH})
      \simeq\varprojlim_m\mathcal O_{V,0}/I_m.          \tag{4.54w.4}
\]
The second equality uses the just-proved first-direct-image
vanishing. This supplies the exact completeness and separation
needed in the leading-equation argument below; no unproved equivalence between a
divisorial and an ordinary degree filtration is being used.

\noindent\textbf{Step 5.} Because \(\mathfrak m^2\subset I_2\), (4.54w.2) gives a
surjection
\(\mathfrak m/\mathfrak m^2\to H^0(H,L)\).
The source has dimension three and the target dimension \(d\).
Hence \(d\le3\).

\paragraph{2. The three section rings and their actual elliptic curves}

\noindent\textbf{Step 6.} For \(d=3\), Riemann--Roch shows that \(L\) separates
points and tangent vectors: subtracting any length-two
subscheme leaves a degree-one line bundle with one section.
The complete linear series therefore embeds \(H\) as a smooth
plane cubic. The restriction sequence on \(\mathbb P^2\),
using \(H^1(\mathbb P^2,\mathcal O(k))=0\), gives its whole
section ring as
\[
 \mathbb C[x,y,z]/(F_3),\qquad \deg(x,y,z)=(1,1,1),
\]
where \(F_3=0\) is that smooth cubic.

\noindent\textbf{Step 7.} For \(d=2\), the complete pencil is base-point free and gives
a degree-two map \(p:H\to\mathbb P^1\), with
\(L=p^*\mathcal O(1)\). Trace splits
\(p_*\mathcal O_H=\mathcal O\oplus B^{-1}\).
Its Euler characteristic is zero, so \(\deg B=2\).
The algebra structure is therefore a double-cover equation
\(z^2=F_4(x,y)\), where the binary quartic has four distinct
zeros because the double cover is smooth. Projection formula
gives the entire section ring with generator weights
\((1,1,2)\) and that single degree-four relation.

\noindent\textbf{Step 8.} For \(d=1\), write \(L=\mathcal O_H(p)\).
Riemann--Roch gives a degree-one section \(x\), one new
degree-two section \(y\), and one new degree-three section
\(z\). Dividing by powers of \(x\) gives functions of pole
orders two and three at \(p\). Their Weierstrass relation
has weighted degree six. Equivalently, after graded changes
only for describing this section ring, it is
\[
 z^2=y^3+A x^4y+B x^6,
 \qquad 4A^3+27B^2\ne0,
 \qquad\deg(x,y,z)=(1,2,3).
\]
The pole-order count shows that the monomials with at most one
factor \(z\), reduced by this relation, give every space
\(H^0(H,L^m)\). Thus this is the whole section ring. Its
weighted projective curve is the actual smooth \(H\).
The coordinate choice below does not require completing the
square or applying any of these optional graded simplifications.

\paragraph{3. Retaining the given coordinate using only permitted changes}

\noindent\textbf{Step 9.} Start with analytic ambient coordinates \((r,u,v)\), retaining
the given \(r\). Their images form a basis of
\(\mathfrak m/\mathfrak m^2\). The map to
\(I_1/I_2=H^0(L)\) has the rank just computed.

\noindent\textbf{Step 10.} When \(d=3\), it is an isomorphism, so any linear complement
to \(r\) supplies the other two valuation-one coordinates.

\noindent\textbf{Step 11.} When \(d=2\), choose a linear coordinate \(s\) whose
leading section complements that of \(r\), and choose the
remaining independent linear coordinate \(v_0\) in the
kernel of the degree-one restriction. Thus
\(\operatorname{ord}_H(v_0)\ge2\). Its degree-two leading
section must be independent of \(r^2,rs,s^2\). Otherwise
all degree-two initial sections of analytic functions would
lie in that three-dimensional span: a Taylor expansion shows
that every term of ordinary order at least three has valuation
at least three, and the linear terms of valuation two are
accounted for by \(v_0\). This contradicts
\(h^0(L^2)=4\). Therefore \(v_0\) has valuation exactly
two and is the required indecomposable degree-two generator.
The coordinates \((r,s,v_0)\) were obtained by an invertible
linear change.

\noindent\textbf{Step 12.} When \(d=1\), first choose two independent linear coordinates
\(u_0,v_0\) in the kernel of the degree-one restriction.
At least one linear combination, denoted \(s\), has a
degree-two initial section independent of \(r^2\), since
\(h^0(L^2)=2\). Choose the remaining independent linear
coordinate \(v_1\), and subtract a scalar multiple of \(s\)
so that its degree-two initial term is \(\alpha r^2\).
Set
\[
                         v_2=v_1-\alpha r^2.
\]
This is a triangular analytic coordinate change, and
\(\operatorname{ord}_H(v_2)\ge3\). Its degree-three
initial section must be independent of \(r^3,rs\).
Indeed, if it were absent or in that span, any Taylor series
of valuation three would have degree-three initial section
in that same span: the possible terms are precisely the
linear \(v_2\) term and the weighted-degree-three monomials
\(r^3,rs\); its valuation-one and valuation-two terms
must first vanish. This contradicts \(h^0(L^3)=3\).
Thus \(v_2\) has valuation exactly three and supplies the
indecomposable degree-three generator. The three ambient
coordinates remain analytically invertible, because only
an invertible linear change and the displayed quadratic
translation were performed.

\noindent\textbf{Step 13.} This argument proves directly that the chosen analytic
coordinates have the prescribed orders and generate the
graded ring. It avoids a circular appeal to filtered
Nakayama or to a coordinate normalization that might move
the fixed \(r\).

\paragraph{4. The actual weighted leading equation}

\noindent\textbf{Step 14.} Write the defining analytic equation in the coordinates just
constructed, with the weights in (4.54w.1). Its weighted initial
form is the single defining relation of the graded section
ring, up to a nonzero scalar. Here is an explicit proof of that
last assertion.

\noindent\textbf{Step 15.} Let \(F_D\) be the graded relation from the section-ring calculation above, expressed
in the actual chosen initial generators. The analytic function
\(F_D(r,s,v)\) has valuation at least \(D+1\).
At each successive valuation, (4.54w.3) writes its initial
section as a weighted homogeneous polynomial in the three
generators. Subtract that polynomial. This produces a formal
series \(G=F_D+F_{D+1}+\cdots\) whose image in every
\(\mathcal O_{V,0}/I_m\) is zero. Since all generator
weights are at most three, the correction polynomials tend
to zero in ordinary power-series degree; their sum is a
well-defined formal power series. By (4.54w.4), its image in
\(\widehat{\mathcal O}_{V,0}\) is zero.

\noindent\textbf{Step 16.} The completed hypersurface ideal is principal, generated by
the actual analytic equation \(f\). Thus \(G=h f\) for
a formal series \(h\). Taking weighted initial forms gives
\(F_D=\operatorname{in}(h)\operatorname{in}(f)\).
The polynomial \(F_D\) is irreducible because its weighted
projective curve is the smooth connected elliptic curve
\(H\). Also \(f\) is nonconstant and vanishes at the
origin. Hence \(\operatorname{in}(h)\) is a nonzero
constant, and \(\operatorname{in}(f)\) is a scalar multiple
of \(F_D\). This proves (4.54w.1) for the actual defining
equation, with all higher terms retained.

\paragraph{5. The exact descent to the original blowdown coordinates}

\noindent\textbf{Step 17.} In the chart with original coordinates
\(x=r u\), \(y=r v\), every change above has the form
\[
 s=a u+b v+c r,\qquad
 v'=a'u+b'v+c'r-\alpha r^2,
 \qquad ab'-a'b\ne0.
\]
It descends to the actual holomorphic coordinate map
\[
 \begin{aligned}
 x'&=r s=a x+b y+c r^2,\\
 y'&=r v'=a'x+b'y+c'r^2-\alpha r^3,\\
 r'&=r.
 \end{aligned}                                        \tag{4.54w.5}
\]
Its Jacobian determinant at the origin is \(ab'-a'b\),
so the inverse-function theorem supplies an analytic inverse.
This is the required morphism preserving the original
parameter and blowdown structure.

\paragraph{Leading polynomials in the original coordinates.}
Apply this calculation to the simple elliptic point supplied
by P4.53. The original equation is $r^3f_1$, and its original
coordinates are $(x,y,z)=(rs,rv,r)$. The three pairs of
original weights and degrees are
\[
 (2,2,1;6),\qquad (2,3,1;7),\qquad (3,4,1;9).
\]
For a strict-transform monomial $r^as^bv^c$, multiplication
by $r^3$ gives $x^by^cz^{3+a-b-c}$. Its original weighted
degree is three plus its strict-transform weighted degree.
Analyticity of the original equation makes every
negative-power coefficient zero. Thus the actual weighted
leading term is a polynomial of the indicated degree.
On $z\ne0$, the ratios
\[
     \left(x/z^{1+w_s},\ y/z^{1+w_v}\right)
\]
identify its projective component with the strict-transform
elliptic relation on $r\ne0$. This is a birational map
of the actual residue-field curves; their smooth
normalizations are the same elliptic curve $H$.

\noindent\textbf{Step 18.} Anisotropic scaling with these weights and the compact-cycle
flow in P4.52 transports every pairing of the leading
Milnor fibre into the actual local torus fibre. By P4.51,
all these pairings must vanish. We now contradict this
in each of the three cases.

\paragraph{Degree six.}
The leading polynomial is exactly
\[
                    P_6(x,y,z)=G_3(x,y,z^2),
\]
where $G_3(X,Y,Z)$ defines a smooth projective cubic.
The map $(x,y,z)\mapsto(x,y,z^2)$ gives a finite double
cover from $F=\{P_6=1\}$ to $\{G_3=1\}$, branched over
\[
                       B=\{G_3(x,y,0)=1\}.
\]
The homogeneous cubic fibre has Euler number nine: it
is a threefold unramified cover of the complement of
an elliptic cubic in $\mathbb P^2$, whose Euler number
is three. Let $r\in\{1,2,3\}$ be the number of distinct
zeros of the nonzero binary cubic $G_3(x,y,0)$ on
$\mathbb P^1$. The curve $B$ is a threefold cover of
that punctured projective line. Hence
\[
          \chi(B)=3(2-r),\qquad \chi(F)=18-\chi(B)=12+3r.
\]

\noindent\textbf{Step 19.} Weighted monodromy has order six. Its only
nonidentity element with a fixed set is its cube,
which negates $z$ and fixes $x,y$; its fixed set is
$B$. The other four nonidentity elements have no
fixed point on $F$. We choose an invariant
semialgebraic compact core large enough for $F$
and $B$. The finite-order Lefschetz character of
$H_0-H_1+H_2$ is therefore $\chi(F)$ at the identity,
$\chi(B)$ at the cube, and zero elsewhere. Its
trivial-character multiplicity is
\[
                       \frac{\chi(F)+\chi(B)}6=3.
\]
Since $H_0$ has only the trivial character, the
nontrivial-character dimension in $H_2$ is at least
$\chi(F)-3=9+3r\ge12$.

\noindent\textbf{Step 20.} Every nonzero critical orbit of $P_6$ lies on $z=0$
at a tangency of the cubic with $Z=0$. If the
intersection orders are $e_i$, then $\sum e_i=3$.
At a tangency, smoothness of $G_3$ makes its
$Z$-derivative a unit, and the exact transverse
germ is $u^{e_i}+z^2$, after holomorphic coordinates
and a unit. Consequently
\[
                       M=\sum(e_i-1)=3-r\le2.
\]
This also shows reducedness of $P_6$: a repeated
factor would give a two-dimensional critical set,
whereas the derivative calculation gives only
these critical orbits. The variation inequality
(4.52b) contradicts the lower bound twelve for
the nontrivial monodromy rank.

\paragraph{Weighted curve cohomology.}
We record the two curve facts needed for the last
cases, with their quotient comparisons. For
$\mathbb P(a,b,c)=\operatorname{Proj}\mathbb C[u,v,w]$
and a reduced equation of weighted degree $e$,
the exact graded sequence sheafifies to
\[
 0\longrightarrow\mathcal O_{\mathbb P}(-e)
 \longrightarrow\mathcal O_{\mathbb P}
 \longrightarrow\mathcal O_D\longrightarrow0.
\]
The three standard affine charts give its \v Cech
complex. The middle cohomology of every
$\mathcal O_{\mathbb P}(k)$ is zero. The top
cohomology of $\mathcal O_{\mathbb P}(-e)$ has
basis the Laurent monomials $u^{-i}v^{-j}w^{-k}$,
$i,j,k\ge1$, of weighted degree $-e$. Thus
\[
 h^1(\mathcal O_D)=
 \#\{\hbox{nonnegative weighted monomials of degree }
                         e-(a+b+c)\},
 \qquad h^0(\mathcal O_D)=1.                         \tag{4.54a}
\]
In particular the arithmetic genus is zero below
degree $a+b+c$, and is one in degrees $a+b+c$
and $a+b+c+1$ when there is exactly one variable
of weight one.

\noindent\textbf{Step 21.} At a cyclic quotient vertex of prime order two
or three, with no invariant linear ambient
coordinate, a smooth coarse curve has a smooth
plane pullback. To prove this, its normalized
branches form one orbit. A single singular
branch of multiplicity $m\ge2$ has a faithful
normalization action of order $p$; every invariant
ambient function has order at least $2m>p$,
whereas a regular quotient requires a function
of normalization order $p$. For $p$ transitively
permuted branches the stabilizer is trivial,
and every invariant function has order at least
two, whereas a regular quotient requires order
one. Both alternatives are impossible. The
remaining single branch is smooth, and its
reduced equation has the character of a nonzero
linear ambient coordinate. These arguments
use the actual analytic invariant rings:
averaging makes the invariants functor exact,
so every quotient-curve function lifts to an
invariant ambient function.

\paragraph{Degree seven.}
Consider weights $(1,2,3)$ and degree seven.
Formula (4.54a) gives arithmetic genus zero for
curves of degree at most five and arithmetic
genus one in degrees six and seven. Thus a
positive-genus component has degree six or
seven, genus one, and multiplicity one. There
is at most one further degree-one factor;
the polynomial is reduced.

\noindent\textbf{Step 22.} In degree seven the whole curve $D$ is
irreducible with arithmetic and normalization
genus one. Its total delta invariant is zero,
so it is a smooth coarse elliptic curve.
At the quotient vertices of orders two and
three, the preceding argument makes its plane
pullback smooth. Elsewhere the ambient plane
is smooth. Thus the affine cone has no
nonzero critical orbit and $M=0$.

\noindent\textbf{Step 23.} In degree six write $D=E\cup\{z=0\}$, with
$z$ the weight-one coordinate. The curve $E$
is smooth coarse elliptic. It cannot contain
either quotient vertex: its degree-six
equation has trivial character there, while
both local-coordinate characters are
nontrivial, so a smooth lift would require
an impossible invariant linear equation.
The genus formula for $D$ gives total delta
one. Its two smooth components therefore
meet once transversely in the smooth part
of the ambient weighted plane. There is
one ordinary-double critical orbit and
$M=1$.

\noindent\textbf{Step 24.} The weighted plane has Euler number three,
from its affine-cell stratification. In the
two cases $e(D)=0$ or one, so its complement
has Euler number three or two. Both quotient
vertices belong to $D$, since degree seven
is divisible by neither two nor three.
The weighted $\mathbb C^*$ action is free
off $D$, so $\{P_7=1\}$ is an unramified
sevenfold cover of that complement. The
free order-seven monodromy gives nontrivial
$H_2$ rank at least eighteen or twelve,
respectively. This contradicts (4.52b),
whose upper bound is zero or one.

\paragraph{Degree nine.}
For weights $(1,3,4)$, (4.54a) gives
arithmetic genus zero below degree eight
and one in degrees eight and nine.
A positive-genus component must have one
of these two degrees, and multiplicity one.

\noindent\textbf{Step 25.} An elliptic degree-eight component is
impossible. It necessarily contains the
weight-three vertex, since eight is not
divisible by three. The local plane
characters are $1,1$ modulo three, whereas
its equation has character two. The
prime-order quotient argument would give
a smooth lift, whose equation must have
character one. This is a contradiction.

\noindent\textbf{Step 26.} Thus the entire degree-nine curve $D$ is
irreducible and smooth coarse elliptic.
It avoids the weight-three vertex: its
equation there has character zero,
incompatible with a smooth plane lift.
At the weight-four vertex the plane
characters are $1,3$ and its equation
has character one. We check the lift
without assuming that four is prime.
The transitive orbit of normalized
branches has size one, two, or four.

\noindent\textbf{Step 27.} For four branches the stabilizer is
trivial and every invariant maximal-ideal
function has order at least two on a
branch, contradicting regularity of the
quotient. For two branches the stabilizer
has order two. A singular branch gives
invariant orders at least four, exceeding
two. Two smooth branches would have a
quadratic initial equation, but the
quadratic characters of the ambient
action are only zero and two, never one.
For one branch the normalization action
is faithful of order four. Both ambient
coordinate orders are odd, since their
characters are one and three. A singular
branch therefore has multiplicity at
least three, and all invariant functions
have order at least six, exceeding four.
These contradictions leave a single
smooth branch.

\noindent\textbf{Step 28.} The degree-nine cone is therefore smooth
away from its vertex. It is singular
at the vertex, because no variable has
weight nine. Its hypersurface critical
point is isolated, so the isolated-point
Morse transport of P4.52 gives a spherical
class of square $-2$ in its Milnor fibre.
This again contradicts the zero pairing
transported from the actual torus.

\noindent\textbf{Step 29.} All three possibilities for the sole
simple elliptic exceptional curve have
now been excluded. Together with P4.53
this proves that the horizontal pole
section cannot be contracted. P4R.18
then gives the asserted generic $D_4$
image and degree two or three. The
other conductor component retains its
unrestricted transverse type.

\noindent\textbf{Step 30.} The exact weighted cohomology counts,
order-six characters, quotient quadratic
characters and an explicit smooth cubic
are recomputed in
\path{code/part4/}\allowbreak\path{simple_elliptic_weighted.py}
and its executed outputs.
\end{proof}


\subsection{The conductor graph and the horizontal spectral cover}

\begin{proposition}[P4.55]\label{night:P4.55}\label{stmt:P4.55}
\StatementRecord{night:P4.55}{proved}{Historical events: GPT-6 Astra (OpenAI), ChatGPT Pro mode, effort Pro, S6\_06\_BLIND, 16--18 November 2026: targeted proof check. Also checked in this writing. Pre-continuation writing events: GPT-6 Astra (OpenAI), ChatGPT Pro mode, effort Pro, S6\_10\_PAPERS, 18 November 2026: targeted proof check. Continuation metadata events: GPT-6 Astra (OpenAI), ChatGPT Pro mode, effort Pro, S6\_10\_PAPERS continuation, 19 November 2026: used and rechecked in this run; status-vocabulary reconciliation; metadata reconciliation only; no new mathematical check. \emph{Pre-continuation recorded proof scope.} The night-run re-derivation was by GPT-6 Astra (OpenAI, ChatGPT Pro mode, effort Pro), 16--18 November 2026. Also checked in this writing by GPT-6 Astra (OpenAI), ChatGPT Pro mode, effort Pro, 18 November 2026. The two-sheet graph budgets, symbolic double-conductor defect bound, ruled cubic spectral genus formula, degree-three exclusion, and existence of the unique remaining vertical sheet were re-derived, including both junction and cube exclusions}{\ref{night:P4.31}; \ref{night:P4.34}; \ref{night:P4.54}; \ref{night:P4R.13}; \ref{night:P4R.18}; \ref{night:P4R.19}}

In the remaining positive ruled case, the image $A$ of the
horizontal pole section is a smooth rational curve. The map
$H\to A$ has degree exactly two. The third generic
normalization preimage along $A$ is supplied by one rational
vertical conductor component, mapping isomorphically to $A$.
The other image component retains its unrestricted transverse
type. A separated self-identification point on either image
component is impossible.
\end{proposition}
\begin{proof}
\noindent\textbf{Step 1.} We write $T\to N\to S$ for the minimal normalization
resolution. By P4.54 and P4R.18, the pole section $H$ is
noncontracted, has genus $g\ge1$, and maps with degree two
or three to $A^\nu=\mathbb P^1$. The full generic transverse
germ along $A$ is $D_4$. The reduced image conductor is
$C=A\cup B$, with two rational normalized components and
one graph cycle.

\noindent\textbf{Step 2.} The reduced normalization conductor $D$ has a tree as its
point-identification graph. Its positive-genus component
$H$ is smoothly embedded: the ruling descends to $N$, and
restriction of that ruling is an inverse to $H\to N$.
In particular no two points of $H$ are identified by the
normalization resolution. Every other normalized conductor
component is rational. These are the actual conductor
curves; vertical components and their identifications are
retained throughout the argument.

\noindent\textbf{Step 3.} At any point of $C$, points of $N$ over it correspond to
local irreducible hypersurface branches of $S$. Two such
branches intersect along a curve, by the height bound for
two equations in a regular local ring of dimension three.
Their nonempty supports among the local analytic branches
of $C$ therefore pairwise intersect.

\paragraph{Separated self-identifications.}
Suppose $p\in B\setminus A$ has two analytic branches.
Let $D_B$ be the union of the conductor components mapping
to $B$, with $c_B$ normalized components and $s_B\ge1$
connected components. Its graph is a forest. Every component
maps finitely and surjectively to $B^\nu$, so the numbers
$r_+,r_-$ of its normalized points over the two preimages
of $p$ satisfy $r_+,r_-\ge c_B$.

\noindent\textbf{Step 4.} There are $n$ normalization points of $S$ over $p$.
Their nonempty supports in $\{+,-\}$ pairwise intersect.
They have a common element: a singleton forces its element
into every support, and without a singleton all supports
are the entire two-element set. After interchanging signs,
$r_+\ge n$. The graph contribution at $p$ is consequently
\[
                  r_++r_--n\ge c_B.
\]
The whole forest has contribution $c_B-s_B<c_B$, a
contradiction. The same proof applies to a self-identification
of $A$ away from $B$, using all components mapping to $A$.
Thus the image cycle consists either of two distinct
intersections of $A$ and $B$, or of one three-branch point
belonging to both components.

\paragraph{The actual primary conductor and smoothness of $A$.}
Let $C_{A,\mathrm{sch}}$ be the primary component of the
actual image conductor supported on $A$. The primary
restriction maps of P4R.13, used in P4R.18, give
\[
                    g-1\le\chi(\mathcal O_{C_{A,\mathrm{sch}}}).
                                                            \tag{4.55a}
\]
The generic ordinary-triple conductor is the transverse
maximal ideal squared. Purity of the actual conductor
quotient identifies this primary ideal with $P_A^{(2)}$,
as in P4.31. The exact normalized conormal calculation of
P4.34 therefore gives
\[
 \chi(\mathcal O_{C_{A,\mathrm{sch}}})
       =5-\sum_p\bigl(\delta(A,p)+\tau_p-r_p\bigr).
                                                            \tag{4.55b}
\]
A planar point contributes exactly $5\delta(A,p)$, and a
nonplanar singular point with at most two branches contributes
at least seven. There cannot be three branches of $A$ at one
point: the identity
$\sum_p(r_p(C)-1)=2$ would then leave no identification for
the necessary connection with $B$; placing that connection
at the same point gives at least four branches and also
exceeds the identity.

\noindent\textbf{Step 5.} Since $g\ge1$, (4.55a)--(4.55b) bound the total defect by
five. If $A$ is singular, it is planar everywhere, has total
delta invariant one, and $g=1$. It is then a locally complete
intersection curve in $Y$. Its conormal
$Q=P_A/P_A^2$ is locally free of rank two. Let
$j:\mathbb P^1\to A$ be its normalization and put $E_0=j^*Q$.
Adjunction gives
\[
             \omega_A=(\omega_Y|_A)\otimes(\det Q)^{-1}.
\]
The degree of $j^*\omega_Y$ is zero because $b_2(Y)=0$.
Also $p_a(A)=1$, so $\deg j^*\omega_A=0$. The degree of a
line bundle is preserved by finite birational normalization:
the normalization sequences for that line bundle and for
$\mathcal O_A$ have the same finite-length correction.
Consequently
\[
                              \deg E_0=0.                 \tag{4.55c}
\]

\noindent\textbf{Step 6.} The actual reduced equation has generic order three along
$A$. Powers of a complete-intersection curve ideal are
unmixed, so $P_A^{(3)}=P_A^3$. Its third normal term is an
actual global section
\[
 q_3\in H^0\!\left(A,
          \operatorname{Sym}^3 Q\otimes\mathcal O_Y(S)|_A\right).
\]
The line twist has degree zero and becomes trivial on
$\mathbb P^1$. On the quotient projective bundle
$\mathbb P(E_0)$, with tautological class $\xi$ and fibre
class $F$, the pulled-back cubic has zero divisor of class
$3\xi$. Remove its common zero fibres, of total degree
$b\ge0$. The residual spectral curve $\Gamma$ has class
$3\xi-bF$. Its horizontal components are reduced, because
the generic $D_4$ cubic has three distinct roots.

\noindent\textbf{Step 7.} We specify the spectral morphisms. At a generic point of
$H$, the first normal derivative of $T\to S\subset Y$
gives the direction of its smooth transverse branch and
hence a point of $\mathbb P(E_0)$. On the smooth compact
curve $H$, its two local coordinates extend after cancelling
their common power. This gives a holomorphic map over the
actual map $H\to\mathbb P^1$. Distinct generic preimages
on $H$ give distinct cubic roots. Hence its map onto the
spectral component $\Gamma_H$ is finite and birational,
and $H$ is that component's normalization. In degree three,
$\Gamma=\Gamma_H$. In degree two, precisely one remaining
rational vertical conductor component supplies the third
root, with degree one over $A$. Its direction curve
$\Gamma_R$ is a section and is isomorphic to $\mathbb P^1$.
Thus $\Gamma=\Gamma_H\cup\Gamma_R$. These exhaust the
three generic roots.

\noindent\textbf{Step 8.} For a rank-two bundle of degree $e$ on $\mathbb P^1$,
\[
 \xi^2=e,\qquad \xi F=1,\qquad F^2=0,\qquad
 K_{\mathbb P(E_0)}=-2\xi+(e-2)F.
\]
Adjunction gives the exact arithmetic genus
\[
                    p_a(3\xi-bF)=3e-2-2b.               \tag{4.55d}
\]
With (4.55c) this is $-2-2b$. In degree three it is at
least $g\ge1$. In degree two the exact union sequence gives
\[
 p_a(\Gamma)=p_a(\Gamma_H)+p_a(\Gamma_R)
           +\operatorname{length}(\Gamma_H\cap\Gamma_R)-1
       \ge g-1\ge0.
\]
Both alternatives are impossible. Therefore $A$ is smooth.

\paragraph{The three-branch image point in degree three.}
Suppose now that $H\to A$ has degree three. Then $H$ is
the only conductor component mapping onto $A$. Since $A$
is smooth, a single three-branch image point consists of
one $A$ branch and two $B$ branches, denoted $+$ and $-$.

\noindent\textbf{Step 9.} Let $c_B$ be the number of normalized conductor components
mapping to $B$. The entire conductor has $1+c_B$
normalized components and a tree graph, so its total
point-identification contribution is $c_B$. At the proposed
point write $r_A\ge1$, $r_+,r_-\ge c_B$ for the numbers
of conductor branches, and $n$ for the number of normalization
points. Each normalization point has at most one $A$
conductor branch, because $H$ is embedded.

\noindent\textbf{Step 10.} A singleton support $\{A\}$ forces every point to carry
$A$. Then $n=r_A$, and the graph contribution is
$r_++r_-\ge2c_B>c_B$. A singleton $\{+\}$ forces
$r_+\ge n$, giving contribution at least
$r_A+r_-\ge1+c_B$; the same applies to $\{-\}$.
Thus every support has at least two elements. Consequently
$r_A+r_++r_-\ge2n$, and its graph contribution is at least
\[
       \frac{r_A+r_++r_-}{2}
                         \ge c_B+\frac{r_A}{2}>c_B.
\]
This is again impossible. Only two distinct $A$--$B$
junctions remain in degree three.

\paragraph{Both degree-three junctions are common cubic zeros.}
Let these junctions be $p,q$. If a local normalization
branch over $p$ has conductor only over $A$, pairwise
hypersurface intersections force every normalization point
over $p$ to carry $A$. Thus $n=r_A$, and the graph
contribution at $p$ is $r_B\ge c_B$, using the entire
budget. At $q$ some normalization point must carry both
$A$ and $B$: otherwise an $A$-only point and a $B$-only
point have disjoint supports. That point contributes
positively to the graph, contradicting the budget. There
is therefore no $A$-only local branch at either junction.

\noindent\textbf{Step 11.} We retain all factors in the following local calculation.
Write the smooth curve as $A=\{x=y=0\}$. If its nonzero
cubic normal form has a simple root, use the normal blowup
chart $x=yu$. The implicit-function theorem on the strict
equation gives that root as $u=\alpha(y,z)$. The actual
function has the analytic factor
\[
                         f_1=x-y\alpha(y,z).
\]
This branch is smooth. Exact division writes $f=f_1f_2$,
and restriction of $f_2$ to this branch is $y^2$ times a
unit: the other cubic directions differ from the simple
root at the point. Its intersections with all other
branches are therefore supported only on $A$. It is the
forbidden $A$-only normalization branch.

\noindent\textbf{Step 12.} Every nonzero binary cubic has a simple root unless it is
a cube. P4R.19 excludes a nonzero cube by its exact finite
weighted maps and compact intersection class. Hence the
whole cubic tensor vanishes at both $p$ and $q$.

\paragraph{The degree-three genus contradiction.}
For the now smooth curve $A=\mathbb P^1$, adjunction gives
$\deg N_{A/Y}=-2$, so $E=N_{A/Y}^*$ has degree two.
The actual cubic has at least two common zeros. Its residual
spectral divisor is $3\xi-bF$ with $b\ge2$, and is
irreducible with normalization $H$. Formula (4.55d) gives
\[
                    p_a(\Gamma)=4-2b\le0,
\]
contradicting $p_a(\Gamma)\ge g\ge1$. Degree three is
excluded. The degree-two alternative and its unique rational
third sheet are precisely the possibilities already identified
by the spectral morphism above.

\noindent\textbf{Step 13.} The exact intersection polynomial and the graph inequalities
are recomputed in
\path{code/part4/d4_graph_spectral.py}; its saved output
reports PASS. The inequalities in this proof establish the
unrestricted graph counts.
\end{proof}


\subsection{The other image curve with at most two generic branches}

\begin{proposition}[P4.56]\label{night:P4.56}\label{stmt:P4.56}
\StatementRecord{night:P4.56}{proved}{Historical events: GPT-6 Astra (OpenAI), ChatGPT Pro mode, effort Pro, S6\_06\_BLIND, 16--18 November 2026: targeted proof check. Also checked in this writing. Pre-continuation writing events: GPT-6 Astra (OpenAI), ChatGPT Pro mode, effort Pro, S6\_10\_PAPERS, 18 November 2026: targeted proof check. Continuation metadata events: GPT-6 Astra (OpenAI), ChatGPT Pro mode, effort Pro, S6\_10\_PAPERS continuation, 19 November 2026: used and rechecked in this run; status-vocabulary reconciliation; metadata reconciliation only; no new mathematical check. \emph{Pre-continuation recorded proof scope.} The night-run re-derivation was by GPT-6 Astra (OpenAI, ChatGPT Pro mode, effort Pro), 16--18 November 2026. Also checked in this writing by GPT-6 Astra (OpenAI), ChatGPT Pro mode, effort Pro, 18 November 2026. The rank-two finite-flat ribbon calculation, the graph alternatives, exact completed-square normalization, and the equality 2b=m versus b>=m were re-derived. The SNC conductor support used here follows by tracking the minimal ruled pole divisor through each point blowup in P4R.7 and P4R.9}{\ref{night:P4.55}; \ref{night:P4R.7}; \ref{night:P4R.9}; \ref{night:P4R.19}}

In the remaining case of P4.55, the other conductor image
component $B$ has at least three generic normalization
branches. This conclusion imposes no preliminary bound on
its transverse multiplicity, conductor exponent, or tangency
order.
\end{proposition}
\begin{proof}
\noindent\textbf{Step 1.} We retain the smooth curve $A$, its degree-two horizontal
cover $H$, and its degree-one vertical cover $R$. All other
normalized conductor components map to $B$. Let $r_B$ be
the generic normalization branch number along $B$, and
$c_B$ the number of global conductor components over it.
Their finite degrees sum to $r_B$. We exclude $r_B=1,2$.
The normalized conductor has $2+c_B$ components and a tree
graph, with total identification contribution $1+c_B$.
Distinct points on one normalized component cannot be
identified, since that would create a graph loop.

\paragraph{A length-two intersection calculation.}
Let $F,G$ be two local hypersurface branches whose
scheme-theoretic intersection is supported on the smooth
curve $A$ and has generic length $m\in\{1,2\}$ along it.
The ring
\[
                 Q=\mathbb C\{x,y,z\}/(F,G)
\]
is a one-dimensional complete intersection. Its only
associated prime is the ideal of $A$. Choose $z$ as a
coordinate on $A$. Then $Q$ is finite and torsion-free
over $\mathbb C\{z\}$, hence free of rank $m$. The
reduction map to $\mathbb C\{z\}$ splits through the
chosen coordinate.

\noindent\textbf{Step 2.} For $m=1$, the ring is $\mathbb C\{z\}$, and the linear
parts of $F,G$ have rank two. Both branches are smooth.
For $m=2$, the nilradical is a free rank-one module over
$\mathbb C\{z\}$. Its square vanishes generically and
therefore vanishes, by torsion-freeness. A generator
$\epsilon$ gives the exact analytic algebra isomorphism
\[
                    Q\simeq\mathbb C\{z,\epsilon\}/(\epsilon^2).
                                                            \tag{4.56a}
\]
Thus $Q$ has embedding dimension two, so one of $dF,dG$
is nonzero. If $G$ has generic normal order two along $A$,
it belongs to the ordinary square of its smooth ideal and
$dG$ vanishes on all of $A$. Then $dF\ne0$ and $F$ is
smooth. We use this calculation on the original intersection
scheme, retaining its multiplicity.

\paragraph{Two distinct image junctions.}
At a two-branch image point, let $r_A,r'_B$ count the
normalized conductor points over its two branches, and
let $n$ count normalization points of $S$. Pairwise
intersection of the hypersurface branches gives a common
element in their supports in $\{A,B\}$. Hence
$n\le\max(r_A,r'_B)$, and the graph contribution is at
least $\min(r_A,r'_B)$. Here $r_A=2$ or three, from the
one or two $H$ points and the unique $R$ point.

\noindent\textbf{Step 3.} If $c_B=2$, both degrees are one, and $r'_B=2$ at either
junction. The two contributions total at least four,
exceeding the budget three. If $c_B=1$ and its degree is
two, an unramified junction costs at least two and the
other at least one, exceeding the budget two. Thus one
component $V$ remains, with exactly one normalized point
over each junction.

\noindent\textbf{Step 4.} If $B$ belongs to every branch support at a junction,
there is only one normalization point there and its
contribution is $r_A\ge2$. The other junction makes this
impossible. Therefore $A$ belongs to every support. The
budget forces contribution exactly one at each junction:
every normalization point has one $A$ branch, and the
sole $V$ point is paired with one of them. Connectedness
forces the graph path $H$--$V$--$R$, with its edges at
the two different junctions.

\noindent\textbf{Step 5.} At the $R$--$V$ junction, let $F$ be the corresponding
hypersurface branch. The other branches come from $H$
and carry only $A$. If its two preimages are distinct,
intersection with either such branch has support $A$ and
generic length one, because the full generic $A$ germ
is an ordinary triple point. The preceding calculation
makes $F$ smooth. If the two $H$ preimages coalesce,
their one hypersurface branch $G$ has generic normal
order two along $A$, and the intersection with $F$ has
length two. Formula (4.56a) again makes $F$ smooth.

\noindent\textbf{Step 6.} This is impossible. The branch $F$ carries all the nearby
generic $B$ preimages: either the single branch of a
singular unibranch transverse germ, or both branches
when $V\to B^\nu$ has degree two and ramifies at this
point. In either case $F$ is intrinsically nonnormal
along the nearby $B$ germ. It cannot be smooth in a
neighbourhood of the junction.

\paragraph{A three-branch image point with one $B$ component.}
The image branches are $A,B_+,B_-$, since $A$ is smooth.
If $c_B=1$, the normalized points on $V$ over the two
$B$ signs are distinct and cannot be identified in $N$.
No normalization point can therefore carry both signs.
Pairwise hypersurface intersections force every point
carrying either sign to carry $A$ as well.

\noindent\textbf{Step 7.} The graph budget two now forces exactly one point over
each sign, each paired with exactly one $A$ branch, and
no other identification. The pairs are $H$--$V$ and
$R$--$V$. An extra $H$ preimage is $A$-only. At the
$R$--$V$ point, the other $H$ branch or branches intersect
its hypersurface branch only along $A$, because their
$B$ supports have the opposite sign or are empty. The
same length-one or length-two calculation makes this
branch smooth, contradicting its intrinsic nonnormality
along the $B$ germ whose generic preimages it contains.

\paragraph{A three-branch point with two $B$ components.}
Write the two degree-one components as $V,W$. If $h=1$
or two is the number of $H$ points over the $A$ value,
the total normalized branch count is
\[
                         T=(h+1)+2+2=h+5.
\]
The graph budget three gives $n\ge h+2$ normalization
points. No singleton image support is possible: a
singleton $\{A\}$ forces $n\le h+1$, and a singleton
$\{B_+\}$ or $\{B_-\}$ forces $n\le2$. Hence every
point has at least two image labels and $n\le T/2$.
This excludes $h=2$. For $h=1$ there are precisely
three points, each with exactly two conductor branches
of different image labels. Up to signs and interchange
of $V,W$, their supports are
\[
 \begin{array}{c|c}
 \text{hypersurface branch}&\text{conductor branches}\\ \hline
 F_R&R,\ B_-\text{ on }V\\
 F_H&H,\ B_+\text{ on }W\\
 F_0&B_+\text{ on }V,\ B_-\text{ on }W.
 \end{array}
                                                            \tag{4.56b}
\]
The tree is the path $H$--$W$--$V$--$R$, and $H\to A$
is ramified with local degree two at its point.

\noindent\textbf{Step 8.} The intersection of $F_R$ and $F_H$ is supported on $A$
and has generic length two. The equation of $F_H$ has
normal order two along $A$, so (4.56a) makes $F_R$
smooth. Choose its equation as $X$, a coordinate $Y$
on that surface cutting $A$, and a coordinate $z$ on
$A$. Preparation and depression in $Y$, preserving the
ideal $(X,Y)$, give the actual equation
\[
                       F_H=\text{unit}\,[Y^2+X^2k(X,z)].
                                                            \tag{4.56c}
\]
The second $X$ follows from normal multiplicity two
along every nearby point of $A$.

\noindent\textbf{Step 9.} The monic finite partial normalization is
\[
      N'=\{u^2+k(X,z)=0\},\qquad Y=Xu.
                                                            \tag{4.56d}
\]
It is an isomorphism away from $A$. The branch $F_H$
is generically smooth along $B_+$: the two generic
$B$ sheets belong separately to $F_H$ and $F_0$.
Thus $N'$ has at most an isolated singularity here.
Its hypersurface Cohen--Macaulay property and regularity
in codimension one make it normal, so it is the actual
normalization.

\noindent\textbf{Step 10.} The reduced divisor $X=0$ is supported on $H$ and has
equation $u^2+k(0,z)=0$. The actual $H$ is smoothly
embedded and has degree two with ramification over $A$.
An even zero of $k(0,z)$ gives two branches; an odd
zero of order at least three gives a singular curve.
Therefore its zero is simple, $k_z(0,0)\ne0$, and
$N'$ is smooth. The minimal normalization resolution
is the identity in this neighbourhood.

\noindent\textbf{Step 11.} Let $m\ge1$ be the generic intersection multiplicity of
the two branches along $B$. It is the same on both
local signs, because they belong to the same irreducible
image component. Let $G$ define $F_0$, and set
\[
                      b=\operatorname{ord}_{p,A}(G|_A)>0.
\]
On the smooth branch $F_R$, its divisor is $mB_-$.
Restricting to $A$ gives
\[
                 b=m I_{F_R}(A,B_-)\ge m.                \tag{4.56e}
\]
On the smooth normalization (4.56d), the divisor of $G$
is $mW$. The actual pole support is normal crossing,
so $H$ and $W$ meet transversely. Its restriction to
$H$ consequently has order $m$. But it is the pullback
of the same ambient restriction $G|_A$ under the local
degree-two map. Hence
\[
                              2b=m,
\]
contradicting (4.56e). Every case with $r_B\le2$ is
excluded.

\noindent\textbf{Step 12.} The exact incidence partitions and the normalization
substitution are checked in
\texttt{code/part4/d4\_other\_two\_branches.py}.
Its corrected run reports PASS; the initial structural
symbolic-equality assertion failure is preserved separately.
\end{proof}


\subsection{Exclusion of the two-junction spectral configuration}

\begin{proposition}[P4.57]\label{night:P4.57}\label{stmt:P4.57}
\StatementRecord{night:P4.57}{proved}{Historical events: GPT-6 Astra (OpenAI), ChatGPT Pro mode, effort Pro, S6\_06\_BLIND, 16--18 November 2026: targeted proof check. Also checked in this writing. Pre-continuation writing events: GPT-6 Astra (OpenAI), ChatGPT Pro mode, effort Pro, S6\_10\_PAPERS, 18 November 2026: targeted proof check. Continuation metadata events: GPT-6 Astra (OpenAI), ChatGPT Pro mode, effort Pro, S6\_10\_PAPERS continuation, 19 November 2026: used and rechecked in this run; status-vocabulary reconciliation; metadata reconciliation only; no new mathematical check. \emph{Pre-continuation recorded proof scope.} The night-run re-derivation was by GPT-6 Astra (OpenAI, ChatGPT Pro mode, effort Pro), 16--18 November 2026. Also checked in this writing by GPT-6 Astra (OpenAI), ChatGPT Pro mode, effort Pro, 18 November 2026. The Hessian-zero consequence of at least three branches, cubic spectral identity, two-junction ruling, and every allocation of the simple common cubic zero between linear and quadratic factors were re-derived. Each case contradicts the required number of branches}{\ref{night:P4.55}; \ref{night:P4.56}; \ref{night:P4R.19}}

The remaining degree-two configuration of P4.55--P4.56
cannot have its image-conductor graph cycle realized at
two distinct intersections of its two components.
\end{proposition}
\begin{proof}
\noindent\textbf{Step 1.} The component $B$ has at least three generic normalization
branches, hence transverse multiplicity at least three.
Along the smooth $A$, choose coordinates $A=\{x=y=0\}$,
with $z$ on $A$. The actual equation belongs to $(x,y)^3$;
write $q_3(x,y;z)$ for its normal cubic.

\noindent\textbf{Step 2.} A local branch of $B$ has first nonzero normal displacement
$(a,b)t^q$, with $(a,b)\ne0$. All ambient Hessian entries
of the actual equation vanish identically along $B$, by
its generic multiplicity. The first normal-order term of
its normal Hessian is
\[
       t^q\operatorname{Hess}_{x,y}
                    (q_3(\,\cdot\,,\,\cdot\,;0))(a,b).
\]
Terms of higher normal order and terms from varying $z$
have larger order in $t$. This matrix is therefore zero.
For a homogeneous binary cubic $Q$, a nonzero vector
annihilating this Hessian makes the directional derivative
of $Q$ the zero quadratic. In coordinates with that vector
as first axis, $Q$ is a cube of the other variable, or zero.
P4R.19 excludes the nonzero cube. Thus the cubic vanishes
at every $A$--$B$ junction, including tangent junctions.

\paragraph{The global spectral equality.}
Let $p,q$ be the two distinct junctions. The common cubic
zero divisor has degree $b\ge2$. On $\mathbb P(N_{A/Y}^*)$
its residual spectral curve is $\Gamma_H\cup\Gamma_R$,
of class $3\xi-bF$. The bundle has degree two, $\Gamma_R$
is a section, and $H$ normalizes $\Gamma_H$. Adjunction
and the exact union sequence give
\[
       g(H)+\delta(\Gamma_H)+\Gamma_H\cdot\Gamma_R=5-2b.
\]
More precisely, the displayed formula is an equality of
integers with $\delta(\Gamma_H)$ the total normalization
length. Since $g(H)\ge1$, it forces
\[
 b=2,\qquad g(H)=1,\qquad \delta(\Gamma_H)=0,
                  \qquad\Gamma_H\cdot\Gamma_R=0.          \tag{4.57a}
\]
The two common cubic zeros are simple, the double spectral
cover is smooth, and it is disjoint from the third section.

\paragraph{The graph support at a junction.}
Let $c_B$ count the global conductor components over $B$.
The graph budget is $c_B+1$. At either junction the
contribution is at least $\min(r_A,r'_B)$, where
$r_A\ge2$ and $r'_B\ge c_B$ count normalized points
over its two image branches.

\noindent\textbf{Step 3.} If $c_B=2$, the two contributions total at least four,
exceeding the budget three. If $c_B=1$, a junction with
at least two $B$ preimages costs at least two, and the
other at least one, exceeding the budget two. Thus both
have one $B$ preimage. The argument of P4.56 gives the
path $H$--$V$--$R$, with one edge at each junction.
At the $R$--$V$ junction the local branch carries all
nearby generic $B$ sheets of $V$, regardless of its
finite degree. Its intersection with the $H$ branch or
branches is a length-one or length-two thickening of
smooth $A$. P4.56 makes it smooth, contradicting its
intrinsic nonnormality from the at least three $B$ sheets.

\noindent\textbf{Step 4.} Therefore $c_B\ge3$. If $A$ belongs to every branch
support at a junction, that point costs at least
$r'_B\ge c_B$, and the other costs at least two.
This exceeds $c_B+1$. The two-element support argument
consequently forces $B$ to belong to every local
hypersurface branch at each junction. Branches carrying
only $B$ have been retained in this conclusion.

\paragraph{Choosing distinct normalization points.}
The ruling descends to $r_N:N\to H$. Its restriction
to $H$ is the identity, and $R$ maps to a fixed point
$b_0\in H$. If a point of $H$ and a point of $R$ have
the same image in $N$, that $H$ point is $b_0$.
This includes identifications through a contracted
exceptional chain. All such coincidences therefore lie
over the single value $a_0\in A$ which is the image of
$b_0$. Choose a junction different from $a_0$, and call
it $p$. Its $H$ preimages and its $R$ preimage are
distinct normalization points. Different $H$ points
are already distinct because $H$ is embedded.

\noindent\textbf{Step 5.} The actual local equation factors into the corresponding
irreducible hypersurface factors and any factors carrying
only $B$. For an $A$-containing factor, let $b_i\ge0$
be the common $z$-order of its first nonzero normal term.
For a factor not containing $A$, let $c_j\ge1$ be the
order of its nonzero restriction to $A$ at $p$.
Gauss's lemma over $\mathbb C\{z\}$ gives
\[
                         \sum_i b_i+\sum_j c_j=1.         \tag{4.57b}
\]
Indeed primitive homogeneous polynomials have nonzero
product modulo $z$, so their common coefficient orders
add exactly. The right side is the simple cubic zero
of (4.57a).

\noindent\textbf{Step 6.} If $H$ is unramified here, there are three
$A$-containing factors of normal order one. At least
two have nonzero linear normal terms, by (4.57b).
Their directions differ by (4.57a). The implicit-function
theorem makes their intersection exactly $A$, whereas
both must contain $B$. This is a contradiction.
Thus $H$ is ramified at $p$.

\noindent\textbf{Step 7.} There are now two $A$-containing factors, denoted by
$F_H$ and $F_R$, of normal orders two and one. Their
residual spectral directions differ. Equality (4.57b)
has precisely the following three possibilities.

\paragraph{A factor carrying only $B$ uses the zero.}
Here $b_H=b_R=0$ and one such factor restricts to $A$
with order one. The branch $F_R$ is smooth. In coordinates
$X=F_R$, $A=\{X=Y=0\}$, restriction of $F_H$ to $X=0$
is $Y^2$ times a unit, since its residual direction
differs from the $R$ direction. Thus these two branches
intersect only along $A$. They both contain $B$, which
is impossible.

\paragraph{The quadratic factor uses the zero.}
Here $b_H=1$, $b_R=0$, and there is no other factor.
The branch $F_R$ is again smooth, and in its coordinates
exact division gives
\[
                         F_H|_{X=0}=Y^2h(Y,z).
\]
The value $h(0,z)$ has a simple zero: the residual
quadratic is nonzero on the $R$ direction, so its
common scalar factor survives with order one on this
restriction. Thus $h_z(0,0)\ne0$. Its zero curve is the
extra image branch $B$, with intersection multiplicity
one between $F_H$ and $F_R$.

\noindent\textbf{Step 8.} The smooth $F_R$ branch supplies one generic $B$
normalization branch. Each branch supplied by $F_H$
has positive intersection multiplicity with it. There
can therefore be at most one such branch. The full
$B$ germ would have at most two branches, contradicting
P4.56.

\paragraph{The linear factor uses the zero.}
Finally $b_R=1$, $b_H=0$, with no other factor.
The nonzero quadratic form of $F_H$ has a double root.
Preparation and depression preserving $A$ give
\[
                   F_H=\text{unit}\,[Y^2+X^2k(X,z)].
\]
Its spectral curve is $u^2+k(0,z)=0$. By (4.57a) this
is smooth and its double cover is ramified at $p$.
Hence $k_z(0,0)\ne0$. The exact finite map
\[
 \{u^2+k(X,z)=0\}\longrightarrow\{F_H=0\},
                         \qquad Y=Xu
\]
has smooth source, is birational, and is an isomorphism
off $A$. It is the normalization. Thus $F_H$ supplies
one generic $B$ branch.

\noindent\textbf{Step 9.} For $F_R$, the simple common zero of its linear normal
term makes its quadratic Taylor part
\[
                       zL(X,Y)+Q_2(X,Y),\qquad L\ne0.
\]
There is no $z^2$ term because $F_R$ vanishes on $A$.
Its Hessian has rank at least two. Rank three gives an
isolated critical point. In rank two the holomorphic
Morse lemma with a parameter gives the exact equation
$U^2+V^2+h(W)$. A nonisolated critical curve forces
$h'\equiv0$, hence $h=0$ because the factor vanishes
at the point. The resulting factorization
$(U+iV)(U-iV)$ contradicts local irreducibility.
Therefore $F_R$ has at most isolated singularities and
also supplies one generic $B$ branch. Again the full
branch number is two, contradicting P4.56.

\noindent\textbf{Step 10.} All allocations in (4.57b) have been excluded, proving
the assertion. The exact genus budget and nonzero
Hessian minor are recomputed in
\texttt{code/part4/d4\_two\_junction.py}; its saved
output reports PASS.
\end{proof}


\subsection{The completed irreducible classification and sign theorem}

\begin{proposition}[P4.58]\label{night:P4.58}\label{stmt:P4.58}
\StatementRecord{night:P4.58}{proved}{Historical events: GPT-6 Astra (OpenAI), ChatGPT Pro mode, effort Pro, S6\_06\_BLIND, 16--18 November 2026: targeted proof check. Also checked in this writing. Pre-continuation writing events: GPT-6 Astra (OpenAI), ChatGPT Pro mode, effort Pro, S6\_10\_PAPERS, 18 November 2026: targeted proof check. Continuation metadata events: GPT-6 Astra (OpenAI), ChatGPT Pro mode, effort Pro, S6\_10\_PAPERS continuation, 19 November 2026: used and rechecked in this run; status-vocabulary reconciliation; metadata reconciliation only; no new mathematical check. \emph{Pre-continuation recorded proof scope.} The night-run re-derivation was by GPT-6 Astra (OpenAI, ChatGPT Pro mode, effort Pro), 16--18 November 2026. Also checked in this writing by GPT-6 Astra (OpenAI), ChatGPT Pro mode, effort Pro, 18 November 2026. Every displayed assembly step was checked against its stated local input. The positive ruled branch uses the actual P4.54 contraction exclusion, P4R.18 reduction, P4.55-57 graph and spectral exclusions, and P4R.20 endpoint calculation. The theorem retains the hypothesis that every support is irreducible or NC; individual dependency checks remain separately recorded}{\ref{night:P4.10}; \ref{night:P4.11}; \ref{night:P4.26}; \ref{night:P4.32}; \ref{night:P4.35}; \ref{night:P4.36}; \ref{night:P4.37}; \ref{night:P4.54}; \ref{night:P4.55}; \ref{night:P4.56}; \ref{night:P4.57}; \ref{night:P4.6}; \ref{night:P4.7}; \ref{night:P4.8}; \ref{night:P4R.13}; \ref{night:P4R.18}; \ref{night:P4R.20}}

Let $f:Y\to B$ be the fibration of Part~4.2, with
$b_1(Y)=b_2(Y)=0$, connected fibres and general fibre a
complex two-torus. Suppose every reduced fibre support
is irreducible or has normal crossings. Then
\[
 B=\mathbb P^1,\qquad c_3(Y)\in\{0,2\},
                       \qquad b_3(Y)=2-c_3(Y).
\]
Every reduced support has nonnegative Euler number.
An irreducible singular support whose normalization has
a rational resolution has Euler number two. Every other
irreducible support has Euler number zero. Such a support
is either smooth with a finite unramified torus cover,
or singular with smooth minimal elliptic ruled
normalization, or singular with smooth minimal Hopf
normalization. In particular the class-VII normalization
alternative is realized only by a smooth minimal Hopf
surface.

For a normal-crossing support the Euler number equals
its number of triple points, and is even. Its zero case
is exactly P4.8: in the reducible case the normalization
is a disjoint union of smooth minimal Hopf or non-K\"ahler
elliptic surfaces, with smooth elliptic conductor curves
and finite unramified branch maps. The equality
$c_3(Y)=0$ holds exactly when every support has one of
the corresponding zero-Euler forms.
\end{proposition}
\begin{proof}
\noindent\textbf{Step 1.} P4.26 proves the sign for every irreducible support,
including singular normalizations, and P4.6--P4.8 give
the full normal-crossing calculation. P4.35 proves the
exact value two for the rationally resolved singular
case. P4.36--P4.37 show that every class-VII case has
smooth minimal Hopf normalization and Euler number
zero. P4.32 gives the complete irreducible zero-Euler
classification and its converses.

\noindent\textbf{Step 2.} We finish the positive-genus ruled alternative without
a transverse-type assumption on its vertical conductor
branches. P4R.13 and the preceding reductions show that
a positive-Euler ruled support would have a doubled
horizontal pole section and an image conductor with two
rational normalized components and one graph cycle.
P4.54 excludes contraction of that section, using the
actual weighted normalization and blowdown maps.
P4R.18 reduces its generic image to $D_4$ with degree
two or three. P4.55 makes that image smooth and excludes
degree three, leaving the degree-two section and one
rational vertical third sheet. P4.56 forces at least
three generic branches on the other image component.

\noindent\textbf{Step 3.} The graph alternatives are exhaustive by P4.55. P4.57
excludes two distinct junctions. P4R.20 excludes a sole
three-branch image point, with the exact conductor-path
and endpoint normalization calculation. Thus no
positive-Euler ruled support remains. This proves the
asserted value zero for every nonrational irreducible
support, and P4.32 supplies its actual normalization.

\noindent\textbf{Step 4.} Finally Euler additivity gives $c_3(Y)$ as the sum of
these nonnegative support Euler numbers. The base,
Poincar\'e-duality calculation $e(Y)=2-b_3(Y)$, and
parity are proved in P4.26. They yield the displayed
global values. A sum of nonnegative integers vanishes
exactly when every summand vanishes, giving the equality
classification through P4.32 and P4.8.

\noindent\textbf{Step 5.} The hypothesis checks are the explicit constructions
P4.10 and P4.11. Removing $b_2=0$ permits the proved
Euler value $24$. Removing the torus-fibre hypothesis
permits the Hopf-fibre example, which fails the stated
normalization equality classification. The numerical
normal-crossing sign itself follows from the stronger
calculation in P4.7 and does not require that hypothesis.
\end{proof}


\subsection{The specialization group and the unibranch case}

\begin{proposition}[P4R.1]\label{night:P4R.1}\label{stmt:P4R.1}
\StatementRecord{night:P4R.1}{proved}{Historical events: GPT-6 Astra (OpenAI), ChatGPT Pro mode, effort Pro, S6\_06\_BLIND, 16--18 November 2026: targeted proof check. Also checked in this writing. Pre-continuation writing events: GPT-6 Astra (OpenAI), ChatGPT Pro mode, effort Pro, S6\_10\_PAPERS, 18 November 2026: targeted proof check; used and rechecked in this run. Continuation metadata events: GPT-6 Astra (OpenAI), ChatGPT Pro mode, effort Pro, S6\_10\_PAPERS continuation, 19 November 2026: used and rechecked in this run; status-vocabulary reconciliation; metadata reconciliation only; no new mathematical check. \emph{Pre-continuation recorded proof scope.} The night-run re-derivation was by GPT-6 Astra (OpenAI, ChatGPT Pro mode, effort Pro), 16--18 November 2026. Also checked in this writing by GPT-6 Astra (OpenAI), ChatGPT Pro mode, effort Pro, 18 November 2026. Compatible triangulation is applied in the smooth real analytic total space. Its open-star retraction and a proper small tube give an actual retraction to the irreducible support. General-position discs identify the complement-group kernel with the normal closure of one positive meridian, whose base winding is the full-fibre multiplicity. The image of the torus lattice is abelian and normal with finite cyclic quotient, and the retraction surjects this virtually abelian group onto pi1 of the support}{the hypotheses and explicit local calculation stated here}

For a proper holomorphic degeneration of complex two-tori with
smooth total threefold and central fibre $mS$, where $S$ is
irreducible and reduced, the group $\pi_1(S)$ has an abelian
normal subgroup with cyclic quotient of order dividing $m$.
\end{proposition}
\begin{proof}
We give the neighbourhood and group comparison explicitly.
The compatible triangulation of Lojasiewicz, Theorem~2,
\noindent\textbf{Step 1. Triangulation of semi-analytic sets}, p.~464 \cite{Lojasiewicz1964}, applies to
the smooth real analytic total space and its closed analytic
subset $S$. Its hypotheses hold because the total space is
countable and finite dimensional, and the collection consisting
of $S$ is locally finite and semianalytic. After barycentric
subdivision $S$ is a full subcomplex. In an open star of this
subcomplex, erase the barycentric coordinates at vertices
outside $S$ and normalize the remaining coordinates to sum
to one. Linear interpolation gives a strong deformation
retraction $r:V\to S$. The star can be chosen inside a
relatively compact neighbourhood of $S$.

Properness gives a small disc $\Delta_\epsilon$ with
$U=f^{-1}(\Delta_\epsilon)\subset V$: on the compact
complement of $V$ in a larger closed tube, $|f|$ has a
positive minimum. Thus the inclusion $i:S\to U$ and the
continuous restriction $r:U\to S$ satisfy $r\circ i=\id$.
This retraction comparison also follows from the statement
immediately after Definition~1 of Fernandez de Bobadilla--Kollar,
\noindent\textbf{Step 2. Homotopically trivial deformations}, p.~85 \cite{BobadillaKollar2012}; the preceding
construction supplies the particular comparison needed here.

Write $U^*=U\setminus S$. The proper smooth fibration
$U^*\to\Delta_\epsilon^*$ has fibre $T^4$. Its homotopy
sequence and $\pi_2(\Delta_\epsilon^*)=0$ give
\[
 \pi_1(U^*)=\Lambda\rtimes\langle t\rangle,
 \qquad \Lambda=\mathbb Z^4.
\]
The extension splits because its quotient is infinite cyclic.
The inclusion $U^*\to U$ induces a surjection on fundamental
groups. A loop can be perturbed off the real-codimension-two
smooth stratum of $S$ and off its higher-codimension singular
strata. The same general-position argument for a null-homotopy
disc avoids the singular strata and meets $S_{\rm reg}$ in
finitely many transverse points. Deleting little discs about
these intersections expresses the kernel as the normal closure
of the corresponding meridians. The smooth locus of the
irreducible $S$ is connected, so those meridians are conjugate,
with their positive orientation fixed by a complex transverse
coordinate.

At a smooth point of $S$, the defining map is a unit times
$z^m$. A positive meridian therefore maps to $m$ under
$\pi_1(U^*)\to\mathbb Z$. Let $A$ be the image of $\Lambda$
in $\pi_1(U)$. It is abelian and normal, and
$\pi_1(U)/A$ is cyclic with order dividing $m$.
The homomorphism $r_*:\pi_1(U)\to\pi_1(S)$ is surjective
because $r_*i_*=\id$. Its image $r_*(A)$ is abelian and
normal, and the induced quotient remains cyclic of order
dividing $m$. This proves the assertion for the actual support,
without replacing it by its normalization.
\end{proof}



\begin{proposition}[P4R.2]\label{night:P4R.2}\label{stmt:P4R.2}
\StatementRecord{night:P4R.2}{proved}{Historical events: GPT-6 Astra (OpenAI), ChatGPT Pro mode, effort Pro, S6\_06\_BLIND, 16--18 November 2026: targeted proof check. Also checked in this writing. Pre-continuation writing events: GPT-6 Astra (OpenAI), ChatGPT Pro mode, effort Pro, S6\_10\_PAPERS, 18 November 2026: targeted proof check; used and rechecked in this run. Continuation metadata events: GPT-6 Astra (OpenAI), ChatGPT Pro mode, effort Pro, S6\_10\_PAPERS continuation, 19 November 2026: used and rechecked in this run; status-vocabulary reconciliation; metadata reconciliation only; no new mathematical check. \emph{Pre-continuation recorded proof scope.} The night-run re-derivation was by GPT-6 Astra (OpenAI, ChatGPT Pro mode, effort Pro), 16--18 November 2026. Also checked in this writing by GPT-6 Astra (OpenAI), ChatGPT Pro mode, effort Pro, 18 November 2026. The smooth bijective normalization is a homeomorphism by compactness and the Hausdorff property, so its Euler number and fundamental group are the actual support invariants. Torsion adjunction with a nonzero conductor kills plurigenera. A ruled base of genus at least two has no virtually abelian finite-index surface group. The resulting rational, elliptic-ruled and classVII Chern formulas prove nonnegativity and the precise smooth minimal zero cases; the nonempty conductor excludes the Inoue alternative}{\ref{night:P4.12}; \ref{night:P4.8}; \ref{night:P4R.1}}

In the standing torus-fibration setting of this section, every irreducible singular support
whose normalization is smooth and whose normalization morphism
is bijective has nonnegative topological Euler number.
Its Euler number is zero precisely when that normalization is
a minimal ruled surface over an elliptic curve or a Hopf surface.
\end{proposition}
\begin{proof}
\noindent\textbf{Step 1.} Let $\nu:T\to S$ denote this finite normalization morphism.
A continuous bijection from the compact $T$ to the Hausdorff
$S$ is a homeomorphism. It therefore induces an isomorphism
of fundamental groups and equality of Euler numbers. P4R.1
makes $\pi_1(T)$ virtually abelian.

\noindent\textbf{Step 2.} The effective-conductor and canonical-torsion argument of
P4.12 applies to this same $T$ and proves $\kappa(T)=-\infty$.
The opened classification theorem used there leaves rational,
ruled, and class-VII minimal models, followed by point blowups.
Blowing up a smooth point preserves the fundamental group:
van Kampen applies to the common punctured four-ball with
simply connected boundary $S^3$. A ruled surface over a
genus-$g$ curve has the same fundamental group as that curve,
by the homotopy sequence of its $\mathbb P^1$-bundle.
For $g\ge2$ this group is not virtually abelian. Indeed a
finite-index subgroup corresponds to an unramified cover
of genus $1+d(g-1)>1$, and a genus-$h$ surface group maps
onto the nonabelian free group on $h$ generators by sending
the $b$ generators to the identity in its usual presentation.
Consequently $g\le1$.

\noindent\textbf{Step 3.} A rational surface has positive Euler number. A ruled surface
over an elliptic curve, blown up $k$ times, has Euler number
$k$. A class-VII surface has Euler number $b_2\ge0$, and
each blowup increases it by one. These formulas give
$e(S)=e(T)\ge0$. Equality in the ruled case is exactly
$k=0$. Equality in class VII gives a minimal surface with
$b_2=0$. The conductor is nonempty because $S$ is singular
and nonnormal; P4.12 gives an effective conductor curve on
$T$. The precise class-VII classification and absence of
curves on Inoue surfaces verified in P4.8 therefore force
$T$ to be Hopf. The converse follows from the same Euler
formulas and the homeomorphism $\nu$.
\end{proof}

\subsection{The conductor pushout for ruled normalizations of higher genus}

\begin{proposition}[P4R.3]\label{night:P4R.3}\label{stmt:P4R.3}
\StatementRecord{night:P4R.3}{proved}{Historical events: GPT-6 Astra (OpenAI), ChatGPT Pro mode, effort Pro, S6\_06\_BLIND, 16--18 November 2026: targeted proof check. Also checked in this writing. Pre-continuation writing events: GPT-6 Astra (OpenAI), ChatGPT Pro mode, effort Pro, S6\_10\_PAPERS, 18 November 2026: targeted proof check. Continuation metadata events: GPT-6 Astra (OpenAI), ChatGPT Pro mode, effort Pro, S6\_10\_PAPERS continuation, 19 November 2026: used and rechecked in this run; status-vocabulary reconciliation; metadata reconciliation only; no new mathematical check. \emph{Pre-continuation recorded proof scope.} The night-run re-derivation was by GPT-6 Astra (OpenAI, ChatGPT Pro mode, effort Pro), 16--18 November 2026. Also checked in this writing by GPT-6 Astra (OpenAI), ChatGPT Pro mode, effort Pro, 18 November 2026. Adjunction and Riemann-Hurwitz force the minimal ruled pole divisor over genus at least two to have horizontal part twice a section. The blowup coefficient rule gives only attached rational trees and no coefficient-one exceptional component. The actual finite normalization quotient is the conductor pushout; the conductor inclusion induces the isomorphism of fundamental groups needed for the van Kampen calculation. The curve graph free-product presentation and virtual abelianness bound its Euler number. In the only zero case all rational poles disappear, forcing the original minimal section and contradicting the independently proved nonnegative open-support Euler number}{\ref{night:P4.12}; \ref{night:P4.13}; \ref{night:P4R.1}; \ref{night:P4R.2}}

In the standing torus-fibration setting of this section, let an irreducible singular fibre
support $S$ have smooth normalization $\nu:T\to S$ whose
minimal model is ruled over a curve of genus $g\ge2$.
Write $D$ for the conductor divisor on $T$ and
$C=\operatorname{Sing}(S)_{\rm red}$.
The inclusion $D_{\rm red}\to T$ induces an isomorphism on
fundamental groups, and the normalization pushout induces
$\pi_1(S)\simeq\pi_1(C)$. Moreover,
\[
 e(C)\ge0,\qquad e(S)>0.
\]
The assertion includes conductor components on different
ruling fibres that map onto the same curve in $C$.
\end{proposition}
\begin{proof}
\noindent\textbf{Step 1.} The canonical comparison of P4.12 gives
$K_T+D\equiv0$, with $D$ effective. We first determine the
reduced conductor as an embedded topological space; this
step does not require the ruling to descend to $S$.
Let $\pi:T\to T_0$ be the sequence of point blowdowns to
the minimal ruled surface $r_0:T_0\to B$, and put
$D_0=\pi_*D$. Then $K_{T_0}+D_0\equiv0$ and
$D_0\cdot F=2$ for a ruling fibre $F$.

\noindent\textbf{Step 2.} There are three possibilities for the horizontal part of
$D_0$: a reduced bisection, two distinct reduced sections,
or twice one section. The first two are excluded directly.
For an irreducible horizontal component $H$ of coefficient
one, adjunction and the effectiveness of the other
components give
\[
 2p_a(H)-2=(K_{T_0}+H)\cdot H
       =-(D_0-H)\cdot H\le0.
\]
On the normalization of $H$, Riemann--Hurwitz for its
finite map of degree $d\ge1$ to $B$ gives
$2g(\widetilde H)-2\ge d(2g-2)>0$.
Since $p_a(H)\ge g(\widetilde H)$, these inequalities
are incompatible. Thus
\[
 D_0=2H+\sum_{j=1}^{v}a_jF_j,\qquad a_j\ge1,
\]
where $H$ is a section and the $F_j$ are distinct ruling
fibres. A horizontal curve of degree one is a section:
its normalization maps isomorphically to the smooth base,
and the resulting map to the projective bundle is a
section with closed image. In particular $H$ is smooth.

\noindent\textbf{Step 3.} At one step $\beta:T'\to T''$ of the blowup sequence,
let $\mu$ be the multiplicity of the current conductor
$D''$ at the centre and $E$ the exceptional curve.
Write $D'=D''_{\rm strict}+aE$. The canonical blowup
formula gives
\[
 K_{T'}+D'
   =\beta^*(K_{T''}+D'')+(1+a-\mu)E.
\]
Intersection with $E$ and numerical triviality give
$a=\mu-1$. Effectiveness therefore places every centre
on the conductor. At a smooth conductor point of
coefficient one, the exceptional curve is absent from
the new conductor. At a smooth point of coefficient at
least two, it is a rational leaf attached at one point.
At an intersection of two conductor components it
subdivides their edge. Starting with $H$ and the
transverse fibres $F_j$, these operations preserve
normal crossings and make $D_{\rm red}$ a copy of $H$
with rational trees attached. No cycle is introduced.

\noindent\textbf{Step 4.} The projection of $H$ to $B$ induces an isomorphism on
fundamental groups, and every attached rational tree
is simply connected. A point blowup leaves
$\pi_1(T)$ unchanged, as shown in P4R.2. The preceding
description also proves, at every blowup, that the
inclusion of the reduced conductor induces the same
isomorphism to the fundamental group of the surface.
Hence $\pi_1(D_{\rm red})\to\pi_1(T)$ is an
isomorphism.

\noindent\textbf{Step 5.} We next verify the pushout for the actual singular
support. The morphism $\nu$ is finite, is an
isomorphism away from the conductor, and has smooth
source. Its nonisomorphism locus equals $C$: a normal
point has a neighbourhood identified with its
normalization and is consequently smooth. Conversely
every smooth point is normal. Thus the natural map
\[
 T\underset{D_{\rm red}}{\cup}C\longrightarrow S
\]
is a continuous bijection. Both source and target are
compact Hausdorff spaces, so it is a homeomorphism.
The compatible triangulation used in P4R.1 makes
$D_{\rm red}\subset T$ a subcomplex and hence a
cofibration. Replacing the map $D_{\rm red}\to C$ by
its mapping cylinder computes the same homotopy
pushout: the homotopy extension property for this
cofibration contracts the inserted cylinder without
altering the homotopy type. Van Kampen on that double
mapping cylinder therefore gives
\[
 \pi_1(S)=\pi_1(T)
       *_{\pi_1(D_{\rm red})}\pi_1(C)
       \simeq\pi_1(C).
\]
All three spaces in the amalgam are connected. This
argument retains every identification made by the
finite conductor map, including identifications
between horizontal and vertical branches.

\noindent\textbf{Step 6.} P4R.1 makes this group virtually abelian. Let $c$ be
the number of irreducible components of $C$, let
$g_i$ be the genera of their smooth normalizations,
and let $b$ be the first Betti number of the incidence
graph of the finite point identifications. Replacing
each identification by a small tree joining its
preimages gives a homotopy model consisting of these
closed oriented surfaces joined along a connected
graph. Choosing a spanning tree gives
\[
 \pi_1(C)\simeq
   \Bigl(*_{i=1}^{c}\pi_1(C_i^{\rm norm})\Bigr)*F_b,
 \qquad
 e(C)=c+1-2\sum_i g_i-b.
\]
The formula also covers unibranch singular points:
their normalization is a local topological
homeomorphism and they contribute no graph edge.

\noindent\textbf{Step 7.} A surface factor of genus at least two cannot occur
in a virtually abelian group, by the finite-cover
argument of P4R.2. Two nontrivial surface factors,
two free generators, or one elliptic factor together
with a free generator contain a nonabelian free
subgroup. Such a subgroup cannot lie in a virtually
abelian group: its intersection with a finite-index
abelian subgroup would have finite index and would
be abelian, whereas a finite-index subgroup of a
nonabelian free group has rank at least two. The
possibilities are therefore all $g_i=0$ with
$b\le1$, or exactly one $g_i=1$ with $b=0$.
The displayed Euler formula gives $e(C)\ge0$.
Additivity of compactly supported Euler numbers and
P4.13 now give
\[
 e(S)=e(S_{\rm reg})+e(C)\ge0.
\]

\noindent\textbf{Step 8.} Finally equality is impossible here. It would force
$e(C)=0$, so the same formula gives $c=1$, $g_1=1$
and $b=0$. Every irreducible component of
$D_{\rm red}$ maps finitely and nonconstantly onto
this curve. A rational vertical component would
therefore give a nonconstant map from $\mathbb P^1$
to its genus-one normalization, contrary to
Riemann--Hurwitz. There are no vertical conductor
components. The initial conductor is consequently
$2H$, and there can have been no blowup: the first
centre would lie on $2H$ and would create a vertical
conductor component of coefficient one. Thus
$T=T_0$ and $D_{\rm red}=H$. It follows that
\[
 e(S_{\rm reg})=e(T\setminus H)
       =(4-4g)-(2-2g)=2-2g<0,
\]
contradicting P4.13. This proves strict positivity.
\end{proof}

P4R.3 proves strict positivity for its stated subclass.
The later classification P4.58 excludes every positive-Euler
ruled case, including this subclass, by the complete local
and global conductor arguments P4.54--P4.57 and P4R.18--P4R.20.
The earlier live status is retained in
\path{records/part4_root_superseded_status.md}.

\begin{proposition}[P4R.4]\label{night:P4R.4}\label{stmt:P4R.4}
\StatementRecord{night:P4R.4}{proved}{Historical events: GPT-6 Astra (OpenAI), ChatGPT Pro mode, effort Pro, S6\_06\_BLIND, 16--18 November 2026: targeted proof check. Also checked in this writing. Pre-continuation writing events: GPT-6 Astra (OpenAI), ChatGPT Pro mode, effort Pro, S6\_10\_PAPERS, 18 November 2026: targeted proof check. Continuation metadata events: GPT-6 Astra (OpenAI), ChatGPT Pro mode, effort Pro, S6\_10\_PAPERS continuation, 19 November 2026: used and rechecked in this run; status-vocabulary reconciliation; metadata reconciliation only; no new mathematical check. \emph{Pre-continuation recorded proof scope.} The night-run re-derivation was by GPT-6 Astra (OpenAI, ChatGPT Pro mode, effort Pro), 16--18 November 2026. Also checked in this writing by GPT-6 Astra (OpenAI), ChatGPT Pro mode, effort Pro, 18 November 2026. The nested open-star and tube maps, punctured radial bundle equivalence and unchanged meridian normal closures give the explicit pi1 isomorphism between support and tube. The complement is Lambda semidirect\_A Z and its filling relation is lambda t\textasciicircum{}m. Rational abelianization has dimension ker(A-I); quotienting by Lambda has order m. A rational-tree conductor makes the support simply connected, hence m=1 and det(I-A) is a positive integer. The trace formula and the tree Euler count then give e(S) at least three}{\ref{night:P4.13}; \ref{night:P4R.1}; \ref{night:P4R.3}}

In the setting of P4R.1, choose a positive small-loop
monodromy $A$ on $\Lambda=H_1(T^4;\mathbb Z)$.
For a suitable $\lambda\in\Lambda$, a positive
meridian of the reduced support is represented by
$\mu=\lambda t^m$ in $\Lambda\rtimes_A\langle t\rangle$.
The natural small-tube comparison gives
\[
 \pi_1(S)\simeq
 (\Lambda\rtimes_A\langle t\rangle)/\langle\!\langle\mu\rangle\!\rangle,
 \qquad
 \pi_1(S)/\operatorname{im}\Lambda\simeq\mathbb Z/m,
 \qquad
 b_1(S)=\dim_{\mathbb Q}\ker(A-I).
\]
For the support in P4R.3, a singular-locus incidence
graph that is a tree of rational normalization
components gives $e(S)\ge3$.
\end{proposition}
\begin{proof}
\noindent\textbf{Step 1.} We strengthen the neighbourhood comparison of P4R.1
without assuming that a preassigned tube is itself
one of the chosen open stars. Choose an open-star
neighbourhood $V_0$ retracting onto $S$, a small tube
$U_0\subset V_0$, a second open-star neighbourhood
$V_1$ with $S\subset V_1\subset U_0$, and a smaller
tube $U_1\subset V_1$. Properness supplies both tube
inclusions. The punctured-disc fibrations give an
isomorphism
\[
 \pi_1(U_1\setminus S)\longrightarrow
 \pi_1(U_0\setminus S):
\]
the fibres can be identified by radial parallel
transport and both base fundamental groups are
generated by the same positive loop. As in P4R.1,
the kernels of the maps to $\pi_1(U_i)$ are the
normal closures of the positive smooth-support
meridian. These meridians correspond under the
displayed isomorphism. Hence
$\pi_1(U_1)\to\pi_1(U_0)$ is an isomorphism.

\noindent\textbf{Step 2.} This isomorphism factors through $\pi_1(V_1)$,
which is identified with $\pi_1(S)$ by its
retraction. Thus $\pi_1(S)\to\pi_1(U_0)$ is
surjective. The restriction to $U_0$ of the
retraction of $V_0$ proves injectivity. This
identifies $\pi_1(S)$ with the exact meridian
quotient, not just with a further quotient.

\noindent\textbf{Step 3.} Projection to the cyclic factor sends $\mu$ to
$m$. After dividing by the image of $\Lambda$,
the remaining presentation is therefore exactly
$\langle t\mid t^m=1\rangle$. In the rational
abelianization the presentation becomes
\[
 \frac{(\Lambda/(A-I)\Lambda)\otimes\mathbb Q
       \ \oplus\ \mathbb Qt}
      {\mathbb Q(\overline\lambda+mt)}.
\]
The single displayed relation has nonzero
$t$ coefficient and reduces the dimension by
one. This proves the asserted formula for $b_1$.

\noindent\textbf{Step 4.} In the rational-tree case of P4R.3 the same
pushout calculation gives $\pi_1(S)=1$.
The cyclic quotient then forces $m=1$, and the
$b_1$ formula gives $\det(I-A)\ne0$.
P4.13 gives
$e(S_{\rm reg})=\det(I-A)\ge1$.
A nonempty connected tree of $c\ge1$ rational
normalization components has Euler number
$c+1\ge2$, including any unibranch singular
points. Additivity gives $e(S)\ge3$.
\end{proof}

\subsection{Hopf normalizations and the conductor graph}

\begin{proposition}[P4R.5]\label{night:P4R.5}\label{stmt:P4R.5}
\StatementRecord{night:P4R.5}{proved}{Historical events: GPT-6 Astra (OpenAI), ChatGPT Pro mode, effort Pro, S6\_06\_BLIND, 16--18 November 2026: targeted proof check. Also checked in this writing. Pre-continuation writing events: GPT-6 Astra (OpenAI), ChatGPT Pro mode, effort Pro, S6\_10\_PAPERS, 18 November 2026: targeted proof check. Continuation metadata events: GPT-6 Astra (OpenAI), ChatGPT Pro mode, effort Pro, S6\_10\_PAPERS continuation, 19 November 2026: used and rechecked in this run; status-vocabulary reconciliation; metadata reconciliation only; no new mathematical check. \emph{Pre-continuation recorded proof scope.} The night-run re-derivation was by GPT-6 Astra (OpenAI, ChatGPT Pro mode, effort Pro), 16--18 November 2026. Also checked in this writing by GPT-6 Astra (OpenAI), ChatGPT Pro mode, effort Pro, 18 November 2026. Adjunction on a minimal Hopf surface and the lift of any rational normalization to C2 minus zero show that every compact curve is smooth elliptic and that distinct curves are disjoint. Tracking the complete pole divisor through blowups gives e(Sreg)=k-v. The disconnected conductor pushout has free quotient of rank b+r-s. Its virtually abelian constraint and the exact curve Euler formula leave only one putative negative configuration; the two-branch hypersurface intersection argument makes that image locally unibranch, which forces b=0 and contradicts the negative count}{\ref{night:P4.12}; \ref{night:P4R.1}; \ref{night:P4R.3}}

In the standing torus-fibration setting of this section, an irreducible singular
support whose normalization is smooth and has a Hopf
surface as its minimal model has nonnegative Euler
number. More precisely, for its reduced conductor
image $C$ one has $e(C)\ge0$. If the normalization is
obtained by $k$ point blowups and $v$ of the exceptional
curves belong to the conductor, then
\[
 e(S_{\rm reg})=k-v,\qquad
 e(S)=k-v+e(C)\ge0.
\]
\end{proposition}
\begin{proof}
\noindent\textbf{Step 1.} Let $T_0$ be the minimal Hopf surface. Every
irreducible curve $E$ on $T_0$ is smooth elliptic.
Indeed $b_2(T_0)=0$ gives $E^2=K_{T_0}\cdot E=0$,
so adjunction gives $p_a(E)=1$. A rational
normalization would define a nonconstant map
$\mathbb P^1\to T_0$. It lifts to the universal
cover $\mathbb C^2\setminus\{0\}$, where its two
coordinate functions are constant. This is
impossible. Thus the normalization genus is one
and the sum of the local delta invariants is zero.
Distinct curves are disjoint, since their
intersection number is zero and every local
intersection of distinct complex curves is
positive.

\noindent\textbf{Step 2.} Push the effective conductor $D$ to $D_0$ on
$T_0$. The canonical comparison of P4.12 makes
$K_{T_0}+D_0$ torsion. The divisor $D_0$ is
nonempty: otherwise a positive power of $K_{T_0}$
would be trivial, contradicting the Kodaira
dimension $-\infty$ of a Hopf surface. Its
reduced support consists of $r\ge1$ disjoint
smooth elliptic curves. The blowup coefficient
$\mu-1$ calculation of P4R.3 applies without a
ruling. At a smooth conductor point it adds a
rational leaf precisely when $\mu\ge2$; at a
node it subdivides the corresponding edge.
Consequently $D_{\rm red}$ has $r$ connected
components, each a smooth elliptic root with a
rational tree attached. With $v$ rational
components it has exactly $v$ double points,
and hence $e(D_{\rm red})=v$. Since
$e(T)=k$, normalization away from the conductor
gives $e(S_{\rm reg})=k-v$. Each rational
component was created by one of the $k$ blowups,
so $k-v\ge0$.

\noindent\textbf{Step 3.} Write $s$ for the number of connected components
of $C$, $c$ for its number of irreducible
components, and $g$ for the number of its
normalizations of genus one. Every normalized
component has genus zero or one: it receives a
finite nonconstant map from a normalized
conductor component. Each elliptic component
must receive an elliptic root, and distinct
image components require distinct roots.
Thus $c\ge g$ and $r\ge g$. Let $b$ be the sum
of the cycle ranks of the point-identification
graphs of the $s$ connected components. The
normalization formula gives
\[
 e(C)=c+s-2g-b.
\]

\noindent\textbf{Step 4.} We retain the disconnected conductor in the
fundamental-group comparison. Its pushout has
one vertex for $T$ and one for each connected
component of $C$, with $r$ edges supplied by
the connected components of $D_{\rm red}$.
This graph is connected, because the finite
map $D_{\rm red}\to C$ is surjective.
The van Kampen presentation has $r-s$ stable
letters, in addition to the vertex groups.
The group of each conductor component is the
group of its elliptic root. Its map into $C$
factors, on that root, through the smooth
normalization of one image component; the
attached rational tree adds no fundamental
group generator. Kill the vertex group
$\pi_1(T)$ and the normal subgroups generated
by all normalized component groups of $C$.
Every conductor-group relation then disappears.
The resulting quotient is exactly
\[
 \pi_1(S)\twoheadrightarrow F_{b+r-s}.
\]
This is also the group of the homotopy-pushout
graph obtained by making the same quotients in
the normalization graph model of $C$. P4R.1
therefore gives $b+r-s\le1$. Substitution yields
\[
 e(C)\ge c+r-2g-1\ge-1.
\]

\noindent\textbf{Step 5.} We exclude the only negative possibility by
the actual local hypersurface geometry.
The value $e(C)<0$ in these inequalities
forces $c=r=g$ and $b+r-s=1$. Every image
component then has elliptic normalization.
No rational conductor component can map
finitely onto one of them, by
Riemann--Hurwitz. Thus $D_{\rm red}$ consists
of disjoint smooth elliptic curves. This
conclusion permits blowups at coefficient-one
points; their exceptional curves are absent
from the conductor and do not affect the
following argument.

\noindent\textbf{Step 6.} At a point $p\in C$, suppose there were two
distinct reduced curve germs $C_1,C_2$ of
$C$. Choose normalization preimages over $p$
whose conductor germs map onto $C_1,C_2$.
They must be distinct preimages, since the
reduced conductor has just one smooth germ
at each of its points. They correspond to
distinct irreducible analytic surface
branches $S_1,S_2$ of $S$ at $p$.
The support $S$ is a hypersurface in the
smooth total threefold. Its distinct
surface branches have a one-dimensional
intersection germ through $p$: two local
nonassociate irreducible equations in the
regular three-dimensional local ring have
height two. That intersection lies in the
conductor image of both branches. The
unique conductor germ on each chosen
normalization branch therefore has the
same image germ, giving $C_1=C_2$, a
contradiction.

\noindent\textbf{Step 7.} Hence $C$ is locally unibranch and its
irreducible components are disjoint. Its
point-identification graphs have $b=0$
and $s=c=g$. The Euler formula gives
$e(C)=0$, contrary to the negative
possibility. We have proved $e(C)\ge0$,
and the asserted formulas finish the proof.
\end{proof}

\subsection{The integral pole divisor on a normalization resolution}

\begin{proposition}[P4R.6]\label{night:P4R.6}\label{stmt:P4R.6}
\StatementRecord{night:P4R.6}{proved}{Historical events: GPT-6 Astra (OpenAI), ChatGPT Pro mode, effort Pro, S6\_06\_BLIND, 16--18 November 2026: targeted proof check. Also checked in this writing. Pre-continuation writing events: GPT-6 Astra (OpenAI), ChatGPT Pro mode, effort Pro, S6\_10\_PAPERS, 18 November 2026: targeted proof check; used and rechecked in this run. Continuation metadata events: GPT-6 Astra (OpenAI), ChatGPT Pro mode, effort Pro, S6\_10\_PAPERS continuation, 19 November 2026: used and rechecked in this run; status-vocabulary reconciliation; metadata reconciliation only; no new mathematical check. \emph{Pre-continuation recorded proof scope.} The night-run re-derivation was by GPT-6 Astra (OpenAI, ChatGPT Pro mode, effort Pro), 16--18 November 2026. Also checked in this writing by GPT-6 Astra (OpenAI), ChatGPT Pro mode, effort Pro, 18 November 2026. Finite Gorenstein duality produces the effective horizontal conductor in the pulled meromorphic canonical divisor. The actual principal-function exceptional intersection equation proves negative definiteness and strict positivity of minus the inverse matrix on each block. Minimal-resolution adjunction gives nonnegative right-hand coefficients, so all exceptional pole coefficients are nonnegative. A zero block is a disjoint rational minus-two configuration outside the conductor. Its isolated hypersurface morsification supplies an embedded vanishing sphere of square minus two in a nearby torus via the normalized root cover, contradicting asphericity. Hence every exceptional block meets the conductor and has positive coefficients}{\ref{night:P4.12}}

Let $S$ be an irreducible singular support in
Part~4.2, let $\nu:N\to S$ be its normalization,
and let $\rho:T\to N\to S$ use the minimal
resolution of $N$. There is an effective integral
divisor $D$ on $T$ such that
\[
 \omega_T(D)\simeq\rho^*\omega_S,
 \qquad \operatorname{Supp}D
       =\rho^{-1}(\operatorname{Sing}S)_{\rm red}.
\]
In particular $K_T+D$ is torsion, by P4.12.
The strict-transform conductor has its usual
positive integral coefficients. Every connected
exceptional configuration over a singular point
of $N$ has strictly positive coefficients in $D$.
Every such configuration meets the strict-transform
conductor.
\end{proposition}
\begin{proof}
\noindent\textbf{Step 1.} The reduced support is a Cartier hypersurface in
the smooth threefold, so $\omega_S$ is invertible.
Choose a local generator $\eta$ of this line
bundle. On $S_{\rm reg}$ it is an ordinary
holomorphic two-form; its pullback is a
meromorphic two-form on $T$. Its divisorial
orders are integers. Two generators differ by
a holomorphic unit on $S$, whose pullback is
still a unit. Consequently their negative
divisors define one integral divisor $D$ and
the displayed line-bundle isomorphism.
Finite duality, in the form
\[
 \nu_*\omega_N
  =\mathcal Hom_{\mathcal O_S}
       (\nu_*\mathcal O_N,\omega_S),
\]
identifies the generic strict-transform pole
ideal with the conductor. This is the same
normalization adjunction used in P4.12. Thus
the coefficients away from the exceptional
curves are the positive integral conductor
exponents; away from conductor and exceptional
curves the pullback has no zero or pole.
No assertion that $K_N$ alone is
$\mathbb Q$-Cartier is needed.

\noindent\textbf{Step 2.} We prove positivity along an exceptional
configuration. Its intersection matrix $M$
is negative definite, and we give the local
comparison establishing the matrix property
used here. Choose a nonzero local function
on the normal surface $N$ vanishing at the
image point. Its pullback divisor is
$H+\sum_i n_iE_i$, with $n_i>0$ and $H$
effective without exceptional components.
The strict transform meets the exceptional
configuration, because a curve through the
image point has closure meeting its proper
inverse image. The divisor is principal
near the configuration, so
\[
 Mn=-h,\qquad h_i=H\cdot E_i\ge0,
 \qquad h\ne0.
\]
The off-diagonal entries of $M$ are
nonnegative. On a connected configuration,
writing a real vector as $v_i=n_iw_i$
gives the identity
\[
 v^tMv
 =-\sum_{i<j}M_{ij}n_in_j(w_i-w_j)^2
   -\sum_i n_ih_iw_i^2.
\]
It is strictly negative for every nonzero
$v$: equality forces all $w_i$ equal along
the connected graph and then forces that
common value to be zero. The same maximum
principle, applied to $M(n_iw_i)=-b$ with
$b_i\ge0$, shows that $-M^{-1}$ has
nonnegative entries, strictly positive on
each connected block. This also follows
directly by taking a negative minimum of
the $w_i$ in the displayed weighted
Laplacian identity.

\noindent\textbf{Step 3.} Write $D=D_h+\sum_i d_iE_i$, with $D_h$
the strict conductor. Minimality of the
resolution gives
\[
 K_T\cdot E_i=2p_a(E_i)-2-E_i^2\ge0.
\]
Indeed the only exception for a negative
curve of arithmetic genus zero would be
a smooth rational exceptional curve of
self-intersection $-1$, which is absent
from a minimal resolution. Also
$D_h\cdot E_i\ge0$. Intersecting
$K_T+D=\rho^*K_S$ with the contracted
curves gives
\[
 Md=-b,
 \qquad b_i=(K_T+D_h)\cdot E_i\ge0.
\]
Hence $d_i\ge0$. A zero coefficient on
a connected block forces $b=0$ there.
Every curve in that block is then a
smooth rational $(-2)$-curve, and the
block is disjoint from the strict
conductor.

\noindent\textbf{Step 4.} In fact, every exceptional block disjoint from
the strict conductor is forbidden, independently
of its coefficients. Finite duality identifies
the conductor on the normal $N$, up to the
invertible pullback of $\omega_S$, with the
rank-one reflexive sheaf $\omega_N$. Thus its
support has no isolated point. Any conductor
curve through the image point has strict
transform meeting the exceptional block:
properness makes the closure of its punctured
inverse image meet the fibre. Therefore the
image of a block disjoint from the strict
conductor is away from the conductor, and
$\nu$ is an isomorphism near it. The
corresponding point of $S$ is a normal isolated
hypersurface singularity. The local
Milnor-fibre argument of P4.12 applies
at this point itself. A local reduced
equation $g$ has an isolated critical
point and positive Milnor number. A
Morse perturbation supplies a vanishing
two-sphere of square $-2$ in its Milnor
fibre. The original map is a unit times
$g^m$; taking a local root of the unit,
each of its nearby local fibres is a
disjoint union of copies of a Milnor
fibre of $g$. Such a copy embeds into
the nearby four-torus. Every map
$S^2\to T^4$ is null-homotopic, so that
embedded sphere would have zero
homology class and square zero, a
contradiction. Thus every exceptional
block has positive pole coefficients.
This proves both effectiveness and the
asserted support equality.
\end{proof}

\begin{proposition}[P4R.7]\label{night:P4R.7}\label{stmt:P4R.7}
\StatementRecord{night:P4R.7}{proved}{Historical events: GPT-6 Astra (OpenAI), ChatGPT Pro mode, effort Pro, S6\_06\_BLIND, 16--18 November 2026: targeted proof check. Also checked in this writing. Pre-continuation writing events: GPT-6 Astra (OpenAI), ChatGPT Pro mode, effort Pro, S6\_10\_PAPERS, 18 November 2026: targeted proof check. Continuation metadata events: GPT-6 Astra (OpenAI), ChatGPT Pro mode, effort Pro, S6\_10\_PAPERS continuation, 19 November 2026: used and rechecked in this run; status-vocabulary reconciliation; metadata reconciliation only; no new mathematical check. \emph{Pre-continuation recorded proof scope.} The night-run re-derivation was by GPT-6 Astra (OpenAI, ChatGPT Pro mode, effort Pro), 16--18 November 2026. Also checked in this writing by GPT-6 Astra (OpenAI), ChatGPT Pro mode, effort Pro, 18 November 2026. The degree-two minimal ruled pole divisor, preservation of its SNC tree under blowup, the pushout fundamental group, and the normalization-trace Euler comparison were re-derived}{\ref{night:P4.13}; \ref{night:P4R.1}; \ref{night:P4R.3}; \ref{night:P4R.6}}

In Part~4.2, every irreducible singular
support whose normalization has a
resolution birationally ruled over a
curve of genus at least two has
nonnegative Euler number. The
normalization may be singular.
\end{proposition}
\begin{proof}
\noindent\textbf{Step 1.} Use the minimal resolution and effective
integral divisor of P4R.6. Push this
divisor to the minimal ruled model
$T_0$. Its intersection with a ruling
fibre is two. The adjunction and
Riemann--Hurwitz calculation in P4R.3
therefore gives exactly
\[
 D_0=2H+\sum_j a_jF_j,
 \qquad a_j\in\mathbb Z_{>0},
\]
with $H$ a smooth section. The reverse
blowup coefficient is the integral
number $\mu-1$. At a conductor crossing
one has $\mu\ge2$, so its exceptional
curve remains in the support and
subdivides the edge. At a smooth
conductor point the operation either
adds a rational leaf or leaves the
reduced curve unchanged. Thus
$A=\operatorname{Supp}D$ is connected,
is $H$ with rational trees attached,
and $\pi_1(A)\to\pi_1(T)$ is an
isomorphism.

\noindent\textbf{Step 2.} Put $C=\operatorname{Sing}(S)_{\rm red}$.
The map $\rho$ is proper and is an
isomorphism outside $A$, and P4R.6
identifies $A$ with its full inverse
image of $C$. The continuous map
$T\cup_A C\to S$ is therefore a
compact-to-Hausdorff bijection and a
homeomorphism. The same cofibration
and van Kampen argument as P4R.3 gives
$\pi_1(S)\simeq\pi_1(C)$. This retains
the exceptional components that map
to points, as well as the strict
conductor components that map to
curves. P4R.1 and the normalization
graph calculation in P4R.3 give
$e(C)\ge0$. Finally P4.13 gives
$e(S_{\rm reg})\ge0$, so
$e(S)=e(S_{\rm reg})+e(C)\ge0$.
\end{proof}

\begin{proposition}[P4R.8]\label{night:P4R.8}\label{stmt:P4R.8}
\StatementRecord{night:P4R.8}{proved}{Historical events: GPT-6 Astra (OpenAI), ChatGPT Pro mode, effort Pro, S6\_06\_BLIND, 16--18 November 2026: targeted proof check. Also checked in this writing. Pre-continuation writing events: GPT-6 Astra (OpenAI), ChatGPT Pro mode, effort Pro, S6\_10\_PAPERS, 18 November 2026: targeted proof check. Continuation metadata events: GPT-6 Astra (OpenAI), ChatGPT Pro mode, effort Pro, S6\_10\_PAPERS continuation, 19 November 2026: used and rechecked in this run; status-vocabulary reconciliation; metadata reconciliation only; no new mathematical check. \emph{Pre-continuation recorded proof scope.} The night-run re-derivation was by GPT-6 Astra (OpenAI, ChatGPT Pro mode, effort Pro), 16--18 November 2026. Also checked in this writing by GPT-6 Astra (OpenAI), ChatGPT Pro mode, effort Pro, 18 November 2026. The Hopf rooted-tree count, disconnected conductor pushout, free quotient rank b+r-s, and the equality case with contracted rational subtrees were re-derived. The hypersurface-branch intersection argument forces the image curve to be unibranch in the only potentially negative case}{\ref{night:P4R.5}; \ref{night:P4R.6}; \ref{night:P4R.7}}

In Part~4.2, every irreducible singular support
whose normalization has a resolution with Hopf
minimal model has nonnegative Euler number.
The normalization may be singular. For the
minimal normalization resolution, let $k$ be
the number of blowups from its Hopf minimal
model and $v$ the number of rational components
of the pole divisor. Then
\[
 e(S_{\rm reg})=k-v\ge0,\qquad
 e(\operatorname{Sing}S)\ge0.
\]
\end{proposition}
\begin{proof}
\noindent\textbf{Step 1.} Use $\rho:T\to N\to S$ and the effective
integral divisor $D$ of P4R.6. Put
$A=\operatorname{Supp}D$ and
$C=\operatorname{Sing}(S)_{\rm red}$.
The Hopf-curve and blowup calculations of
P4R.5 use only integrality, effectiveness,
and $K_T+D\equiv0$. They therefore apply
to this divisor: $A$ has $r\ge1$ connected
components, each a smooth elliptic root
with rational trees attached. Each of its
$v$ rational components came from one of
the $k$ blowups; $e(A)=v$, $e(T)=k$, and
$\rho$ is an isomorphism off $A$. This
proves the displayed formula for the
regular locus.

\noindent\textbf{Step 2.} Every one-dimensional component of $C$
receives a nonconstant proper map from a
component of $A$, so its normalization
has genus zero or one. There are no
isolated points of $C$: an isolated
singular point of the hypersurface is
normal by the hypersurface $S_2$ property
and regularity in codimension one, and
P4R.6's vanishing-sphere argument excludes
it. Let $s,c,g,b$ respectively count the
connected components of $C$, its
irreducible components, its genus-one
normalizations, and the cycle ranks of
its point-identification graphs.
Then $e(C)=c+s-2g-b$, $r\ge g$, and
$c\ge g$.

\noindent\textbf{Step 3.} The proper-map pushout $T\cup_A C=S$
and its cofibration justification are
exactly those of P4R.7. A root which
maps to a point gives the zero
homomorphism on its fundamental group;
a noncontracted root maps through the
normalization of one component of $C$.
Every rational tree is simply connected.
Consequently killing $\pi_1(T)$ and
the normalized component groups of $C$
in the disconnected van Kampen
presentation leaves the free quotient
$F_{b+r-s}$, as in P4R.5. Virtual
abelianness gives $b+r-s\le1$, hence
\[
 e(C)\ge c+r-2g-1\ge-1.
\]

\noindent\textbf{Step 4.} A negative value forces $c=r=g$.
Every root is then noncontracted, maps
onto an elliptic normalization, and
supplies one of the $c$ image components.
Every rational component of $A$ is
contracted: a nonconstant map from its
normalization $\mathbb P^1$ to an
elliptic curve contradicts
Riemann--Hurwitz. Thus every exceptional
block is a rational subtree of one
rooted component of $A$. It meets that
root in exactly one point. At least one
intersection is forced by the final
assertion of P4R.6, since the roots are
the only strict conductor components;
two intersections would make a cycle
in the rooted tree. The exceptional
fibre is connected. One can see this
directly by Stein factorization of
$T\to N$: its finite birational factor
is an isomorphism because $N$ is normal.

\noindent\textbf{Step 5.} The conductor on $N$ consequently has
one local analytic branch at each of
its points. At a singular point of $N$,
each local branch has a strict transform
meeting the exceptional fibre, whereas
there is exactly one such intersection
with the smooth root. Away from those
fibres the assertion follows from the
smooth disjoint roots themselves.

\noindent\textbf{Step 6.} The local hypersurface argument of
P4R.5 now applies with $N$ in place of
its smooth resolution. Two different
curve germs of $C$ would require two
different normalization points, because
each point has just one conductor germ.
Their corresponding hypersurface
branches intersect in a curve germ.
That germ belongs to both conductor
images and forces the two chosen germs
to coincide. Thus $C$ is locally
unibranch with disjoint irreducible
components. This gives $b=0$, $s=c=g$,
and $e(C)=0$, a contradiction.
Therefore $e(C)\ge0$, proving the
assertion.
\end{proof}

\begin{proposition}[P4R.9]\label{night:P4R.9}\label{stmt:P4R.9}
\StatementRecord{night:P4R.9}{proved}{Historical events: GPT-6 Astra (OpenAI), ChatGPT Pro mode, effort Pro, S6\_06\_BLIND, 16--18 November 2026: targeted proof check. Also checked in this writing. Pre-continuation writing events: GPT-6 Astra (OpenAI), ChatGPT Pro mode, effort Pro, S6\_10\_PAPERS, 18 November 2026: targeted proof check. Continuation metadata events: GPT-6 Astra (OpenAI), ChatGPT Pro mode, effort Pro, S6\_10\_PAPERS continuation, 19 November 2026: used and rechecked in this run; status-vocabulary reconciliation; metadata reconciliation only; no new mathematical check. \emph{Pre-continuation recorded proof scope.} The night-run re-derivation was by GPT-6 Astra (OpenAI, ChatGPT Pro mode, effort Pro), 16--18 November 2026. Also checked in this writing by GPT-6 Astra (OpenAI), ChatGPT Pro mode, effort Pro, 18 November 2026. Adjunction on the reduced horizontal pole divisor forces equality in every degree inequality; the etale bisection or two disjoint sections, absence of vertical poles, and absence of resolution or nonminimal exceptional curves were re-derived}{\ref{night:P4.13}; \ref{night:P4.22}; \ref{night:P4R.6}; \ref{night:P4R.7}}

In Part~4.2, every irreducible singular
support whose normalization has a
resolution birationally ruled over an
elliptic curve has nonnegative Euler
number. When the horizontal part of
the pole divisor on the minimal ruled
model is reduced, the normalization
is smooth and minimal and the support
has Euler number zero.
\end{proposition}
\begin{proof}
\noindent\textbf{Step 1.} Push the integral divisor of P4R.6 to
the minimal ruled surface $T_0\to B$.
Its horizontal degree is two. A reduced
horizontal component $H$ satisfies
\[
 0\le 2g(\widetilde H)-2
 \le 2p_a(H)-2
 =-(D_0-H)H\le0.
\]
The first inequality is
Riemann--Hurwitz over the elliptic
base. All inequalities are equalities:
every such component is smooth elliptic,
unramified over $B$, and disjoint from
the other components of $D_0$. Every
vertical fibre meets the horizontal
part, so no vertical component occurs.
The reduced-horizontal possibilities
are therefore one smooth elliptic
bisection or two disjoint elliptic
sections, each with coefficient one.

\noindent\textbf{Step 2.} The reverse blowup formula $\mu-1$
then keeps $A=\operatorname{Supp}D$
a disjoint union of their smooth
strict transforms; each new exceptional
curve is absent from $A$. Every
exceptional curve of the normalization
resolution must belong to $A$, by
P4R.6. An exceptional component would
therefore be one of these disjoint
elliptic curves and would be disjoint
from the strict conductor, contrary
to P4R.6. Thus the normalization is
smooth. P4.22, with its hypotheses now
verified, proves minimality and
$e(S)=0$.

\noindent\textbf{Step 3.} In the remaining case
$D_0=2H+\sum a_jF_j$ with $H$ a
smooth section. The same reverse
blowups make $A$ the elliptic root $H$
with rational trees attached. Its
inclusion in $T$ induces an
isomorphism on fundamental groups,
by the ruling and invariance under
point blowups. The resolution pushout
gives $\pi_1(S)=\pi_1(C)$ as in
P4R.7. The virtually abelian curve
group calculation gives $e(C)\ge0$,
and P4.13 gives $e(S_{\rm reg})\ge0$.
Their sum proves $e(S)\ge0$ also
in this remaining case.
\end{proof}
\subsection{The Hessian line and the tacnodal equality case}

\begin{proposition}[P4R.10]\label{night:P4R.10}\label{stmt:P4R.10}
\StatementRecord{night:P4R.10}{proved}{Historical events: GPT-6 Astra (OpenAI), ChatGPT Pro mode, effort Pro, S6\_06\_BLIND, 16--18 November 2026: targeted proof check. Also checked in this writing. Pre-continuation writing events: GPT-6 Astra (OpenAI), ChatGPT Pro mode, effort Pro, S6\_10\_PAPERS, 18 November 2026 05:10:56.410177+00:00: computation or formal certificate; targeted proof check. Continuation metadata events: GPT-6 Astra (OpenAI), ChatGPT Pro mode, effort Pro, S6\_10\_PAPERS continuation, 19 November 2026: used and rechecked in this run; status-vocabulary reconciliation; metadata reconciliation only; no new mathematical check. \emph{Pre-continuation recorded proof scope.} The night-run re-derivation was by GPT-6 Astra (OpenAI, ChatGPT Pro mode, effort Pro), 16--18 November 2026. Also checked in this writing by GPT-6 Astra (OpenAI), ChatGPT Pro mode, effort Pro, 18 November 2026. proved. Full elliptic-primary A3 conductor Euler inequality, exact differentiated morphisms and equality case. Finite duality makes the conductor curve Cohen–Macaulay; generic containment and absence of a zero-dimensional submodule give P\textasciicircum{}(2)\ensuremath{\subset}I. Its nilradical J=P/I is torsionfree rank1. Differentiation embeds P/P\textasciicircum{}(2) into the normalized kernel K of j*OmegaY\ensuremath{\to}OmegaE, with degK=degR. The actual Hessian is a section of L\textsuperscript{2}\ensuremath{\otimes}j*O(S); the latter factor has degree0, so 2degL=sum d. Projection to K/L has kernel exactly I/P\textasciicircum{}(2), determined generically and extended by torsionfreeness. Its branchwise and descent colengths give the displayed exact Euler formula. For a nonplanar branch, all gradient minors have order\ensuremath{\geq}2m; writing a+b=tau and extracting a unit column of t\textasciicircum{}d ell\textsuperscript{2} gives d\ensuremath{\geq}a, hence d/2+b\ensuremath{\geq}m>m−1. For a planar branch, the exact generators dz,dh give E=<dz,t\textasciicircum{}a u>, a\ensuremath{\geq}m; the symmetric-square condition yields d/2+b\ensuremath{\geq}a. Equality of the global Euler number forces delta=0, then R=0 and all remaining nonnegative lengths vanish, so the curve is smooth and the Hessian nowhere zero. The bounded tacnode\_hessian replay is retained only as a supplement to this uniform proof. Review provenance: New derivation after reading the supplied proof; not a blind review and not an independent human check}{\ref{night:P4.29}; \ref{night:P4.30}}

Let $S$ be an irreducible reduced fibre support in Part~4.2.
Suppose that its reduced singular curve $C$ is irreducible and
unibranch at every point, that its normalization $j:E\to C\hookrightarrow Y$
is elliptic, and that its generic transverse germ is a tacnode
$x^2-y^4=0$. Let $C_{\mathrm{sch}}$ be the image conductor scheme.
Then
\[
 \chi(\mathcal O_{C_{\mathrm{sch}}})\leq0.
\]
Equality makes $C$ smooth and makes the transverse Hessian of
the reduced defining equation have rank one at every point of $C$.
\end{proposition}
\begin{proof}
\noindent\textbf{Step 1.} Let $P\subset\mathcal O_Y$ be the reduced ideal of $C$, and let
$I\subset\mathcal O_Y$ be the inverse image of the conductor ideal
of $S$. The conductor depth argument of P4.29 gives a
Cohen--Macaulay quotient $\mathcal O_Y/I$ of pure dimension one.
At the generic tacnode the normalization has two smooth branches;
its conductor is $(x,y^2)$ in the displayed coordinates.
Consequently $I$ contains $P^{(2)}$. Indeed a section of $P^{(2)}$
has zero image in $\mathcal O_Y/I$ at the generic point, and any
nonzero image would generate a finite-length submodule, contrary
to depth one. We obtain
\[
 0\longrightarrow J:=P/I\longrightarrow
 \mathcal O_{C_{\mathrm{sch}}}\longrightarrow\mathcal O_C
 \longrightarrow0.
\]
Here $J^2=0$, because $P^2\subset P^{(2)}\subset I$.
The module $J$ is therefore an $\mathcal O_C$-module. It is
torsionfree of rank one: a finite-support torsion submodule would
again be a finite-length submodule of $\mathcal O_Y/I$.

\noindent\textbf{Step 2.} On $E$, put
\[
 K=\ker\bigl(j^*\Omega_Y\longrightarrow\Omega_E\bigr),
 \qquad\operatorname{im}(dj)=\Omega_E(-R).
\]
These are a rank-two bundle and a line bundle respectively.
The coefficient of $R$ at a point with local parametrization
$x_i=x_i(t)$ is $r=\min_i\operatorname{ord}_t x_i'(t)$.
We have $\deg K=\deg R$, since $\deg j^*\Omega_Y=0$
and $E$ is elliptic. Differentiation identifies
$G=P/P^{(2)}$ with a subsheaf of $j_*K$, as in P4.30.
Write $E_0\subset K$ for the image of its pullback to $E$.
This is a rank-two locally free subsheaf with finite-length
quotient.

\noindent\textbf{Step 3.} The Hessian of a local reduced equation $f$ of $S$ annihilates
the tangent derivative of $C$, since all first derivatives of
$f$ vanish on $C$. It thus defines a quadratic tensor in
$\operatorname{Sym}^2K$. Under a change $f\mapsto u f$ its
restriction to $C$ multiplies by $u$, so globally it is a section
\[
 q\in H^0\bigl(E,\operatorname{Sym}^2K\otimes j^*\mathcal O_Y(S)\bigr).
\]
It has generic rank one. Let $L\subset K$ be the saturation of
its generic image line. On the smooth curve $E$ this is a line
subbundle. The tensor lies in $L^{\otimes2}\otimes
j^*\mathcal O_Y(S)$ everywhere, since that is a subbundle and
the assertion holds on a dense open set. Write its zero orders
as $d_p\geq0$. The normal line of $S$ has degree zero on $E$:
locally over the base disc the divisor $mS$ is principal.
It follows that
\[
 2\deg L=\sum_p d_p.
\]

\noindent\textbf{Step 4.} Generically $I/P^{(2)}$ is exactly the line of the linear
conductor equation $x$, which is the Hessian line. Projection
to $K/L$ therefore gives an injection
\[
 J\longrightarrow j_*(K/L).
\]
To check the kernel at the exceptional points, a section whose
projection vanishes generically has torsion image in $J$; the
torsionfreeness just proved makes that image zero. Let $b_p$ be
the colength of the projected lattice $E_0\to K/L$ at $p$.
Let $q_p$ be the additional nonnegative length between $J$ and
the direct image of this normalized image lattice. Riemann--Roch
on $E$ now gives the exact formula
\[
 \chi(J)=\deg R-\frac12\sum_p d_p-\sum_p b_p-\sum_p q_p.
\]
We prove locally that $d_p/2+b_p\geq r_p$.

\noindent\textbf{Step 5.} At a smooth point of $C$, $r_p=0$, so the claim follows from
nonnegativity. At a singular point put
$m=\min_i\operatorname{ord}_t x_i(t)$, so $r_p=m-1$.
First suppose the embedding dimension is three. Every defining
equation of $C$ has no linear term. Its gradient along the
parametrization has all entries of order at least $m$.
Every two-by-two minor of the conormal image consequently has
order at least $2m$. Since $K$ is saturated in the ambient free
rank-three module, a generator of $\bigwedge^2K$ has a unit
Pl\"ucker coordinate. Thus
\[
 \tau_p:=\operatorname{length}(K/E_0)_p\geq2m.
\]
Let $a_p=\operatorname{length}(L/(E_0\cap L))_p$.
The exact lattice sequence gives $\tau_p=a_p+b_p$.
Choose a primitive local generator $\ell$ of $L$ and write
$q=h\ell^2$ after trivializing the normal line, so
$\operatorname{ord}_t h=d_p$.
Every column of the Hessian is the differential of one of
the first derivatives of $f$. Those first derivatives belong
to $P$, so the columns lie in $E_0$. At least one ambient
coordinate of the primitive vector $\ell$ is a unit; its
corresponding column shows that $h\ell\in E_0$.
Therefore $d_p\geq a_p$, and
\[
 \frac{d_p}{2}+b_p\geq\frac{a_p}{2}+b_p
 \geq\frac{\tau_p}{2}\geq m>m-1=r_p.
\]

\noindent\textbf{Step 6.} Now suppose the embedding dimension is two. Choose coordinates
with $P=(z,H(x,y))$, where $H$ is an irreducible reduced plane
equation with no linear term. Here $P^{(2)}=P^2$.
For an explicit verification, $z,H$ is a regular sequence, and
$P/P^2$ is free over the domain $\mathcal O_Y/P$.
If $s\notin P$ and $sv\in P^2$, reduction modulo $P$ gives
$v\in P$; reduction in this free conormal module then gives
$v\in P^2$. Thus $P^2$ is $P$-primary and equals its symbolic
square. The equation $f$ belongs to this square, so
\[
 f=A z^2+B zH+C H^2.
\]
Let $u$ be a primitive generator of the plane kernel of the
derivative row. We have $dH=w u$, with
$a=\operatorname{ord}_t w\geq m$, because both plane entries
of $dH$ have order at least $m$. In the basis $(dz,u)$ of $K$
we consequently have
\[
 E_0=\langle dz,w u\rangle,
 \qquad q=2A\,dz^2+2Bw\,dz\,u+2Cw^2u^2.
\]
Write the primitive generator of $L$ as
$\ell=\ell_z dz+\ell_u u$, and again $q=h\ell^2$ with
$\operatorname{ord}_t h=d_p$.
When $\ell_z$ is a unit, the projected conormal lattice has
colength $b_p=\min(a,\operatorname{ord}_t\ell_u)$.
The coefficient of $u^2$ in the displayed tensor gives
\[
 d_p+2\operatorname{ord}_t\ell_u\geq2a,
 \qquad d_p/2+b_p\geq a\geq m.
\]
The same inequality holds when $\ell_u=0$, since then $b_p=a$.
When $\ell_z$ is not a unit, $\ell_u$ is a unit. Projection of
$dz$ generates $K/L$, so $b_p=0$; the $u^2$ coefficient gives
$d_p\geq2a$, proving the same bound. These cases cover all
singular curve points and establish the claimed local inequality.

\noindent\textbf{Step 7.} We conclude that $\chi(J)\leq0$. Since normalization gives
$\chi(\mathcal O_C)=-\sum\delta(C,p)$, the conductor exact
sequence yields
\[
 \chi(\mathcal O_{C_{\mathrm{sch}}})
 =-\sum_p\delta(C,p)+\chi(J)\leq0.
\]
Equality forces every delta invariant to vanish, so $C=E$ is
smooth. Then $R=0$, $E_0=K$, and all normalization descent
defects vanish. The exact formula reduces to
$\chi(J)=-\tfrac12\sum d_p$, so equality also forces
$d_p=0$ everywhere. The Hessian has rank one at every point.
The program \texttt{code/part4/tacnode\_hessian.py} verifies the
two lattice inequalities for the stated bounds, the exact
rank-one Hessian on the singular parametrization
$(t^2,t^3,0)$, and the simple Hessian zero of the permitted
local model $zx^2+xy^2$. Its JSON and text outputs retain
the bounds; the preceding inequalities prove the unrestricted
statement.
\end{proof}

\begin{proposition}[P4R.11]\label{night:P4R.11}\label{stmt:P4R.11}
\StatementRecord{night:P4R.11}{proved}{Historical events: GPT-6 Astra (OpenAI), ChatGPT Pro mode, effort Pro, S6\_06\_BLIND, 16--18 November 2026: targeted proof check. Also checked in this writing. Pre-continuation writing events: GPT-6 Astra (OpenAI), ChatGPT Pro mode, effort Pro, S6\_10\_PAPERS, 18 November 2026 05:10:56.410177+00:00: computation or formal certificate; targeted proof check. Continuation metadata events: GPT-6 Astra (OpenAI), ChatGPT Pro mode, effort Pro, S6\_10\_PAPERS continuation, 19 November 2026: used and rechecked in this run; status-vocabulary reconciliation; metadata reconciliation only; no new mathematical check. \emph{Pre-continuation recorded proof scope.} The night-run re-derivation was by GPT-6 Astra (OpenAI, ChatGPT Pro mode, effort Pro), 16--18 November 2026. Also checked in this writing by GPT-6 Astra (OpenAI), ChatGPT Pro mode, effort Pro, 18 November 2026. proved. Full zero-Euler double-horizontal A3 normalization, degree-two conductor and product-tacnode conclusion. Followed the pole-tree pushout to identify the reduced conductor normalization as one elliptic curve. Every rational pole component contracts, so H is its only strict preimage and has degree2. The actual ruling descends through the normalized graph because exceptional fibres are vertical; its Leray sequence gives R1rho O=0. Finite conductor duality and torsion-line Riemann–Roch give \ensuremath{\chi}OCsch=g−1. The independently rederived R10 bound forces g=1, smooth conductor and a nowhere-zero rank-one Hessian; Riemann–Hurwitz makes H\ensuremath{\to}C an unramified double cover. Distinct H points cannot merge under a vertical exceptional fibre. Thus the two local hypersurface factors are smooth graphs in the nondegenerate Morse variable. Their difference is a unit times the square of a coordinate defining the smooth conductor, since the generic intersection order is2. The exact unit-root map gives X\textsuperscript{2}−Y4; its two finite normalization maps are smooth. The remaining normalization and minimality inputs retain their stated R6/P4.22 scopes. Review provenance: New derivation after reading the supplied proof; not a blind review and not an independent human check}{\ref{night:P4.12}; \ref{night:P4.22}; \ref{night:P4.29}; \ref{night:P4R.10}; \ref{night:P4R.6}; \ref{night:P4R.9}}

Let $S$ be an irreducible singular support in Part~4.2 whose
normalization has ruled resolution, and suppose $e(S)=0$.
In the double-horizontal-divisor case with generic transverse
tacnode, the normalization is a smooth minimal elliptic ruled
surface. The conductor image is a smooth elliptic curve, its
reduced preimage is an unramified covering of degree two,
and every local germ is a product tacnode.
\end{proposition}
\begin{proof}
\noindent\textbf{Step 1.} Use $T\to N\to S$, the horizontal section $H$, and
$C=\operatorname{Sing}(S)_{\mathrm{red}}$ as in P4R.9.
The fundamental-group and Euler argument of P4.29 gives an
irreducible curve $C$ with elliptic normalization, unibranch
at every point. All rational curves in the resolved conductor
are contracted, since none maps nonconstantly to an elliptic
curve. Thus $H$ is the only strict conductor component.
Its finite map to $C$ has degree two, and its genus $g$ is at
least one by Riemann--Hurwitz.

\noindent\textbf{Step 2.} Every exceptional curve of $T\to N$ is vertical for the ruling
$T\to H$. The normalized-graph construction of P4.29 descends
this ruling to $N\to H$. Its Leray calculation proves
$R^1\rho_{N*}\mathcal O_T=0$, and hence
$\chi(\mathcal O_N)=1-g$. Finite duality identifies the
conductor on $N$ with $\omega_N\otimes\nu^*\omega_S^{-1}$.
The latter line twist is torsion by P4.12. Serre duality,
the rational-singularity comparison and Riemann--Roch on $T$
therefore give $\chi(I)=1-g$ and
$\chi(\mathcal O_{D_{\mathrm{sch}}})=0$.
The exact conductor square, with $\chi(\mathcal O_S)=0$,
then gives
\[
 \chi(\mathcal O_{C_{\mathrm{sch}}})=g-1.
\]
P4R.10 makes this at most zero. Thus $g=1$, $C$ is smooth,
and the reduced defining equation has rank-one Hessian along
all of $C$. Riemann--Hurwitz makes $H\to C$ unramified.

\noindent\textbf{Step 3.} We verify the local analytic conclusion. Over any point of
$C$, its two preimages on $H$ cannot be identified in $N$:
each exceptional block is vertical for the descended ruling
and hence meets only one point of the section $H$.
There are therefore two local irreducible branches of $S$.
The holomorphic Morse lemma with parameters writes its reduced
equation as $X^2+g(y,z)$. Its two analytic factors are monic
linear polynomials in $X$, by Weierstrass factorization, and
their coefficients sum to zero. Each branch is a graph and
is smooth. Their intersection has support exactly the smooth
curve $C$. Consequently, in coordinates with $C=\{X=y=0\}$,
the factor coefficient is a unit times a power of $y$.
The generic tacnode fixes that power to two. Extracting a square
root of the unit in the $y$ coordinate gives
$X^2-y^4=0$ times the $z$ disc. This also verifies the actual
normalization and conductor morphisms locally.
The normalization was already smooth away from the conductor
by P4R.6. It is therefore smooth everywhere, and P4.22 makes
this elliptic ruled surface minimal.
\end{proof}

\subsection{The components outside the pole divisor in a ruled fibre}

\begin{proposition}[P4R.12]\label{night:P4R.12}\label{stmt:P4R.12}
\StatementRecord{night:P4R.12}{proved}{Historical events: GPT-6 Astra (OpenAI), ChatGPT Pro mode, effort Pro, S6\_06\_BLIND, 16--18 November 2026: targeted proof check. Also checked in this writing. Pre-continuation writing events: GPT-6 Astra (OpenAI), ChatGPT Pro mode, effort Pro, S6\_10\_PAPERS, 18 November 2026 05:10:56.410177+00:00: computation or formal certificate; targeted proof check. Continuation metadata events: GPT-6 Astra (OpenAI), ChatGPT Pro mode, effort Pro, S6\_10\_PAPERS continuation, 19 November 2026: used and rechecked in this run; status-vocabulary reconciliation; metadata reconciliation only; no new mathematical check. \emph{Pre-continuation recorded proof scope.} The night-run re-derivation was by GPT-6 Astra (OpenAI, ChatGPT Pro mode, effort Pro), 16--18 November 2026. Also checked in this writing by GPT-6 Astra (OpenAI), ChatGPT Pro mode, effort Pro, 18 November 2026. proved. Full disjoint-leaf description of fibre components outside the pole support on the actual ruled blowup model. On the actual point-blown-up ruled model, every vertical component E has negative self-intersection by the full multiplicity equation m\_E E\textsuperscript{2}+sum m\_j(E\_j E)=0. If E lies outside D, then DE=2+E\textsuperscript{2}\ensuremath{\geq}0, so E\textsuperscript{2} is−1 or−2. The latter gives DE=0 and is disjoint from D; the resolution is an isomorphism there by R6. Removing the fibre multiplicity is unramified on E, and its finite etale pullback consists of P1s. Actual smooth-family transport places its normal square−2 sphere in a nearby T4, impossible. Thus E\textsuperscript{2}=−1 and DE=1, giving one transverse intersection with a coefficient-one pole component. Two outside components cannot be adjacent: the positive full-fibre multiplicity relations would give both m\_E>m\_Eprime and m\_Eprime>m\_E. Therefore all such components are pairwise disjoint leaves. No multiplicity-one assumption was used. Review provenance: New derivation after reading the supplied proof; not a blind review and not an independent human check}{\ref{night:P4R.3}; \ref{night:P4R.6}; \ref{night:P4R.9}}

Let $T\to N\to S$ be the minimal normalization resolution of an
irreducible support in Part~4.2 with ruled minimal model of positive
base genus. Let $r:T\to B_0$ be its ruling, and let $D$ be the
effective integral pole divisor of P4R.6. Every component $E$ of a
reducible ruling fibre which is absent from $\operatorname{Supp}D$
is a smooth rational curve of self-intersection $-1$. It meets
exactly one pole component, transversely once, and that component
has pole coefficient one. In the double-horizontal-section case,
$E$ is disjoint from the horizontal section. Distinct such
components are disjoint; each is a leaf of the reduced ruling fibre.
These assertions retain the multiplicities of the ruling fibre
and do not require the multiplicity of $E$ to be one.
\end{proposition}
\begin{proof}
\noindent\textbf{Step 1.} A reduced horizontal pole divisor occurs here only over an
elliptic base, by P4R.3, and P4R.9 then makes the normalization
smooth and minimal. In that case there is no reducible ruling
fibre. We may therefore work in the double-horizontal-section
case throughout the proof.
A ruling fibre on $T$ is obtained from a smooth $\mathbb P^1$
by the actual sequence of point blowups. Its reduced support is
a normal-crossing tree of smooth rational curves. Each component
of a reducible fibre has negative self-intersection. One can
verify the latter assertion without a contraction criterion:
if $F_b=\sum m_iE_i$ is the full fibre, then
$0=F_bE_i=m_iE_i^2+\sum_{j\ne i}m_jE_jE_i$;
connectedness and reducibility make the last sum positive.

\noindent\textbf{Step 2.} For $E$ outside the pole support, effectiveness and adjunction give
\[
 0\le DE=-K_TE=2+E^2.
\]
Thus $E^2$ is $-1$ or $-2$. In the second case $DE=0$,
so $E$ is disjoint from the full pole support. The map to $S$
then identifies a neighbourhood of this compact curve with a
neighbourhood in $S_{\rm reg}$, by P4R.6. The local cyclic
base change which removes the multiplicity of the support is
unramified over this regular locus. Its restriction to $E$ is
a disjoint union of copies of $\mathbb P^1$, because
$\pi_1(\mathbb P^1)=0$. A chosen copy has the same normal
degree $-2$. A smooth local trivialization over a sufficiently
small parameter disc transports it into a nearby four-torus,
still with self-intersection $-2$. Every map from $S^2$ to that
torus is null-homotopic, since its universal cover is
$\mathbb R^4$; its homology class and self-intersection must
therefore be zero. This contradiction leaves $E^2=-1$ and
$DE=1$.

\noindent\textbf{Step 3.} All pole coefficients and all intersection multiplicities are
positive integers at an intersection. Hence $DE=1$ says exactly
that there is one intersection, transverse and on a coefficient-one
component. A horizontal pole section of coefficient two cannot
be that component.

\noindent\textbf{Step 4.} It remains to exclude an intersection of two components $E,E'$
outside the pole divisor. Both have square $-1$, and both meet
the pole support. In the full ruling fibre, let their positive
multiplicities be $m,m'$. The two fibre-intersection equations
would give
\[
 m=m'+a,\qquad m'=m+b,
\]
where $a,b$ are positive sums of the multiplicities of their
other vertical neighbours; their respective pole neighbours make both
sums positive. Adding these equations gives $a+b=0$, a
contradiction. Thus no two outside components meet. Each has
exactly the one pole neighbour already found and is a leaf.
\end{proof}

\subsection{Restriction to an elliptic conductor component}

\begin{proposition}[P4R.13]\label{night:P4R.13}\label{stmt:P4R.13}
\StatementRecord{night:P4R.13}{proved}{Historical events: GPT-6 Astra (OpenAI), ChatGPT Pro mode, effort Pro, S6\_06\_BLIND, 16--18 November 2026: targeted proof check. Also checked in this writing. Pre-continuation writing events: GPT-6 Astra (OpenAI), ChatGPT Pro mode, effort Pro, S6\_10\_PAPERS, 18 November 2026: targeted proof check; used and rechecked in this run. Continuation metadata events: GPT-6 Astra (OpenAI), ChatGPT Pro mode, effort Pro, S6\_10\_PAPERS continuation, 19 November 2026: used and rechecked in this run; status-vocabulary reconciliation; metadata reconciliation only; no new mathematical check. \emph{Pre-continuation recorded proof scope.} The night-run re-derivation was by GPT-6 Astra (OpenAI, ChatGPT Pro mode, effort Pro), 16--18 November 2026. Also checked in this writing by GPT-6 Astra (OpenAI), ChatGPT Pro mode, effort Pro, 18 November 2026. The conductor restriction and canonical-trace maps, primary restriction Euler bound, degree-zero jet and discriminant sections, product analytic types, and impossibility of a first blowup were re-derived. Greb 2011 Remark 2.1, printed page 415, was opened for the surface rational-singularity trace criterion; the earlier local jet bounds remain separately recorded dependencies}{\ref{night:P4.30}; \ref{night:P4.31}; \ref{night:P4.33}; \ref{night:P4R.10}; \ref{night:P4R.6}; \ref{night:P4R.7}; \ref{night:P4R.9}}

In Part~4.2, let an irreducible singular support $S$ have a
normalization with ruled resolution. If its reduced image
conductor has a component with elliptic normalization, then
the normalization of $S$ is a smooth minimal elliptic ruled
surface, the image conductor is one smooth elliptic curve,
and $e(S)=0$. Every local germ along this curve is the product
of its generic transverse germ with a disc. In particular no
rational conductor component is attached to that curve.
\end{proposition}
\begin{proof}
\noindent\textbf{Step 1.} Let $\rho_N:T\to N$ be the minimal normalization resolution,
let $\nu:N\to S$ be the finite normalization, and put
$\rho=\nu\rho_N$. The reduced-horizontal case is already
proved by P4R.9. We therefore use the actual double-horizontal
pole divisor $D$ of P4R.6. Its support is a smooth section $H$
with rational trees attached. The pushout and curve-group
calculation of P4R.7--P4R.9 show that the reduced image
conductor $C$ is connected, that its point-identification
graph has no cycles, and that it has exactly one normalized
component of positive genus, of genus one. Denote the
corresponding reduced component by $C_E$ and its normalization
by $j:E\to C_E\hookrightarrow Y$. In particular every point
of $C_E$ is unibranch. Every other normalized component of
$C$ is rational.

\noindent\textbf{Step 2.} A rational component of $D$ cannot map nonconstantly to $E$.
Thus $H$ maps nonconstantly to $E$. Write $g=g(H)\ge1$.
At a general point of $C_E$, every normalization branch has
conductor exponent two, because the coefficient of $H$ in
$D$ is two. The intrinsic conductor exponent of a singular
plane branch is at least two, and each additional branch
increases it by its positive intersection multiplicity.
The possibilities are therefore a single ordinary cusp,
two smooth branches with intersection multiplicity two, or
three smooth branches with pairwise intersection multiplicity
one. These are $A_2,A_3,D_4$, respectively, and $H\to E$
has degree $1,2,3$ in the three cases.

\noindent\textbf{Step 3.} We first construct a map which detects global functions on
the entire conductor, including its rational components.
Since $H$ is not contracted, P4.33 proves that $N$ has
rational singularities. Put $L_N=\nu^*\omega_S$.
Finite duality identifies its conductor ideal with
\[
 I_N=\omega_N\otimes L_N^{-1}.
\]
The canonical trace is an isomorphism
$\rho_{N*}\omega_T=\omega_N$ for rational analytic
singularities. We use Greb, Remark~2.1, p.~415
\cite{Greb2011}. The hypotheses hold here: $N$ is normal
and Cohen--Macaulay, $\rho_N$ is a proper resolution,
P4.33 gives $R^1\rho_{N*}\mathcal O_T=0$, and its
one-dimensional fibres give vanishing of the higher
structure-sheaf direct images. The integral pole identity
of P4R.6 and projection formula consequently give
\[
 \rho_{N*}\mathcal O_T(-D)
 =\rho_{N*}(\omega_T\otimes\rho_N^*L_N^{-1})
 =I_N.
\]
These are equalities of ideal subsheaves of $\mathcal O_N$:
the pole identity uses the pullback of the same residue
section which defines finite duality, and the trace agrees
with it on the birational isomorphism locus. Equality there
determines the maps between these torsionfree sheaves.
The divisor ideal sequence therefore gives an injection
\[
 \mathcal O_N/I_N\longrightarrow\rho_{N*}\mathcal O_D.
\]
The finite conductor square gives an injection from the
image conductor scheme. Composing them yields the actual
morphism of coherent $\mathcal O_S$-algebras
\[
 \mathcal O_{C_{\rm sch}}\longrightarrow
 \nu_*(\mathcal O_N/I_N)\longrightarrow\rho_*\mathcal O_D,
 \qquad\hbox{both arrows injective}.                         \tag{P4R.13.1}
\]

\noindent\textbf{Step 4.} We show that a global function on $D$ is determined by
its value at the generic doubled horizontal section.
At every point blowup $\beta:T'\to T''$ in the actual
map to the minimal ruled surface, the pole divisors obey
$D'=\beta^*D''-E_\beta$. The elementary direct images are
\[
 \beta_*\mathcal O(E_\beta)=\mathcal O,
 \quad R^1\beta_*\mathcal O(E_\beta)=0,
 \quad\beta_*\mathcal O=\mathcal O,
 \quad R^1\beta_*\mathcal O=0.
\]
For the first equality, a meromorphic function allowed
a pole only on $E_\beta$ descends off the centre and
extends across that point by Hartogs. For the vanishing,
the two blowup charts, or their exceptional-curve
filtration, have successive terms
$\mathcal O_{\mathbb P^1}(j)$ with $j\ge-1$;
their first cohomology vanishes. Equivalently one applies
$0\to\mathcal O\to\mathcal O(E_\beta)
\to\mathcal O_{E_\beta}(-1)\to0$ to the usual
blowup structure-sheaf calculation. Projection formula,
the divisor ideal sequence and the absence of second
direct images then give
\[
 \beta_*\mathcal O_{D'}=\mathcal O_{D''},
 \qquad R^1\beta_*\mathcal O_{D'}=0.                         \tag{P4R.13.2}
\]
These particular maps identify $H^0(D,\mathcal O_D)$
with $H^0(D_0,\mathcal O_{D_0})$, where
\[
 D_0=2H_0+V_0,
 \qquad V_0=\sum_j a_jF_j
\]
on the minimal ruled surface. The kernel of restriction
from $D_0$ to $2H_0$ is $\mathcal O_{V_0}(-2H_0)$.
Its successive layers on each thick fibre $a_jF_j$ are
$\mathcal O_{\mathbb P^1}(-2)$, since $F_j^2=0$ and
$H_0F_j=1$. They have no global sections. Restriction
to $2H_0$ is therefore injective on global functions.
This double section is Cohen--Macaulay and has no
zero-dimensional associated point, so a section which
vanishes generically on it vanishes everywhere.

\noindent\textbf{Step 5.} Let $C_{E,\rm sch}$ be the primary component of
$C_{\rm sch}$ supported on $C_E$: locally its ideal
is obtained by localizing the conductor ideal at the
generic point of $C_E$ and contracting. Primary
decomposition of the pure one-dimensional conductor
quotient makes these ideals coherent and compatible
on overlaps. Restriction is a surjection
\[
 \mathcal O_{C_{\rm sch}}\longrightarrow
 \mathcal O_{C_{E,\rm sch}}.
\]
Its map on $H^0$ is injective. Indeed, a global element
of its kernel maps by (P4R.13.1) to a function on $D$
which vanishes generically on $2H$; (P4R.13.2) and the
preceding paragraph make it zero. Its map on $H^1$
is surjective, because the kernel is supported in
dimension at most one and has no second cohomology.
We have proved
\[
 \chi(\mathcal O_{C_{\rm sch}})
 \le\chi(\mathcal O_{C_{E,\rm sch}}).                       \tag{P4R.13.3}
\]

\noindent\textbf{Step 6.} We evaluate the right side using the actual primary
ideal, without removing the rational components from
the left side. In type $A_2$ the primary conductor
is reduced: a pure one-dimensional quotient which
is generically reduced has no nilpotent submodule
of finite support. Thus its Euler characteristic
is $-\sum_p\delta(C_E,p)\le0$. In type $D_4$ the
primary ideal is the symbolic square of the ideal
of $C_E$, by exactly the generic-conductor and
depth-one argument of P4.31. P4.30 gives the same
nonpositive bound, with equality only for smooth
$C_E$.

\noindent\textbf{Step 7.} In type $A_3$, the primary ideal contains that
symbolic square and its nilpotent module is
torsionfree of rank one on $C_E$. The proof of
P4R.10 applies to this primary quotient: its
differentiation map, saturated Hessian line and
projection to the quotient conormal line depend
only on this ideal and on the actual equation of
$S$. The presence of other conductor components
changes none of those maps. More explicitly its
Euler characteristic is
\[
 -\sum_p\delta(C_E,p)+\deg R
 -\frac12\sum_p d_p-\sum_p b_p-\sum_p q_p\le0,
\]
with the same ramification divisor $R$, Hessian
zero orders $d_p$, projected lattice colengths
$b_p$ and descent defects $q_p$ as in P4R.10.
The local inequalities there prove this bound
at every singular point of $C_E$. Equality
makes $C_E$ smooth and its Hessian nowhere zero.

\noindent\textbf{Step 8.} On the other hand P4.33 gives
$\chi(\mathcal O_{C_{\rm sch}})=g-1$. Thus
\[
 0\le g-1
 =\chi(\mathcal O_{C_{\rm sch}})
 \le\chi(\mathcal O_{C_{E,\rm sch}})\le0.                    \tag{P4R.13.4}
\]
All these quantities vanish. We obtain $g=1$,
$C_E=E$ smooth, and an unramified map $H\to E$.
In type $A_3$ its transverse Hessian has rank
one everywhere on $E$.

\noindent\textbf{Step 9.} It remains to exclude attachments to this smooth
elliptic component. Put
\[
 K=\mathcal I_E/\mathcal I_E^2,
 \qquad N_E=\mathcal O_Y(S)|_E.
\]
Both $\deg K$ and $\deg N_E$ are zero. The first
follows from the conormal exact sequence,
$c_1(TY)E=0$, and $\deg\Omega_E=0$; the second
follows because a positive multiple of $S$ is
principal in a tube of this fibre.

\noindent\textbf{Step 10.} In type $A_2$ let $L\subset K$ be the saturated
quadratic-image line and put $Q=K/L$. The nonzero
quadratic tensor is a section of
$L^2\otimes N_E$, so $2\deg L\ge0$.
The cubic jet on the null direction is an
intrinsic section of $Q^3\otimes N_E$.
Here is the coordinate check, including points
where the quadratic tensor vanishes. A change
of a second-order lift changes the cubic
coefficient by the polarization of the
quadratic tensor. That polarization vanishes
on the null direction. A change of tangential
lift contributes a tangential derivative of
the rank-one quadratic tensor; its evaluation
twice on the null direction is again zero.
Changing a reduced defining equation multiplies
the tensor by its transition unit. Thus the
stated line bundle and its holomorphic section
are defined on all of $E$. The section is
nonzero generically by type $A_2$, so
$-3\deg L\ge0$. Both degrees are zero,
and both sections are nowhere zero. The
parametrized Morse lemma and a cube root
of the residual unit now give the product
equation $X^2+Y^3=0$ at every point of $E$.

\noindent\textbf{Step 11.} In type $A_3$ (P4R.13.4) already makes the
quadratic tensor nowhere zero, whence
$\deg L=\deg Q=0$. Its cubic null jet is
identically zero. Removing the quadratic
direction gives an intrinsic residual fourth
jet in $Q^4\otimes N_E$. In local coordinates
its relevant terms and coefficient are
\[
 a x^2+bxy^2+cy^4,
 \qquad q_4=c-\frac{b^2}{4a},\qquad a\ne0.                 \tag{P4R.13.5}
\]
Substitution of the critical coordinate
$x=-by^2/(2a)+O(y^3)$ proves the formula.
The change $x\mapsto x+vy^2$ sends $b$ to
$b+2av$ and $c$ to $c+bv+av^2$, leaving
$q_4$ unchanged. Changes of normal frames
and of equation give the claimed
$Q^4\otimes N_E$ transition. Tangential
lift changes do not affect a function
whose residual order is four. This section
is nonzero generically and has degree zero;
it never vanishes. The residual function
is $y^4$ times a unit, giving a product
tacnode at every point of $E$.

\noindent\textbf{Step 12.} In type $D_4$, the equation belongs to
$\mathcal I_E^3$. This ideal is primary
because $E$ is smooth, so generic order
three proves membership everywhere.
The cubic normal jet belongs to
$\operatorname{Sym}^3K\otimes N_E$.
Its discriminant is a section of
$(\det K)^6\otimes N_E^4$, of degree
zero. It is nonzero generically, hence
nowhere zero. The three distinct normal
tangent directions give three smooth
holomorphic branches locally along $E$.
After straightening the first two to
$x=0,y=0$, the third is $ax+by=0$
with both coefficients units. Setting
$X=ax,Y=-by$ gives the product equation
$XY(X-Y)=0$.

\noindent\textbf{Step 13.} In all three cases the entire singular
locus near every point of $E$ is $E$
itself. A rational conductor tree could
not meet it. Connectedness of $C$ then
forces $C=E$. The displayed local
product forms make $N$ smooth along
its conductor, and P4R.6 already makes
it smooth elsewhere. Hence $T=N$.
A finite normalization cannot contract
a curve, and a rational curve cannot
map nonconstantly to $E$. Thus $D=2H$
has no vertical component. The first
blowup of $2H_0$ in a nonempty actual
blowup sequence would create a vertical
pole component of coefficient one;
later blowups never remove its strict
transform from the pole support.
Consequently no blowup occurs and $N$
is minimal elliptic ruled. Finally
\[
 e(S)=e(N)-e(H)+e(E)=0.
\]
The exact coefficient changes, discriminant
weights and degree identities are executed
in \texttt{code/part4/elliptic\_primary\_restriction.py};
its JSON and text outputs are retained.
\end{proof}

\subsection{The smooth critical curve of transverse type \texorpdfstring{$A_3$}{A3}}

\begin{proposition}[P4R.14]\label{night:P4R.14}\label{stmt:P4R.14}
\StatementRecord{night:P4R.14}{proved}{Historical events: GPT-6 Astra (OpenAI), ChatGPT Pro mode, effort Pro, S6\_06\_BLIND, 16--18 November 2026: targeted proof check. Also checked in this writing. Pre-continuation writing events: GPT-6 Astra (OpenAI), ChatGPT Pro mode, effort Pro, S6\_10\_PAPERS, 18 November 2026 05:10:56.410177+00:00: computation or formal certificate; targeted proof check. Continuation metadata events: GPT-6 Astra (OpenAI), ChatGPT Pro mode, effort Pro, S6\_10\_PAPERS continuation, 19 November 2026: used and rechecked in this run; status-vocabulary reconciliation; metadata reconciliation only; no new mathematical check. \emph{Pre-continuation recorded proof scope.} The night-run re-derivation was by GPT-6 Astra (OpenAI, ChatGPT Pro mode, effort Pro), 16--18 November 2026. Also checked in this writing by GPT-6 Astra (OpenAI), ChatGPT Pro mode, effort Pro, 18 November 2026. proved. Full smooth-critical-axis A3 classification, finite-degree polar argument and explicit simple-elliptic local test. Diagonalized the rank-one Hessian over C\{z\}; the weight-four tensor is hx\textsuperscript{2}+axy\textsuperscript{2}+by4 with discriminant4hb−a\textsuperscript{2}. Independent constant changes make any zero simple away from h=0, producing the actual binomial X\textsuperscript{2}+y4Z and its forbidden compact spherical class. Thus the discriminant is a unit. At h(0)=0, the simple finite root after x=y\textsuperscript{2}v yields f=X(Y\textsuperscript{2}+XH). The off-axis critical equations are H\_z=0 and2H+XH\_X=0; the entire-axis hypothesis makes their intersection isolated. Re-derived the finite positive-degree map Psi=(H\_z,(G−2zH\_z)/2), the measure-zero exclusions of X=0 and H\_X=0, and the exact nonzero critical value−X\textsuperscript{3}H\_X/2. The full Hessian determinant is2X4J, so regular parameter values give actual Morse points and the boundary transport forces H\_z to be a unit. The explicit coordinate Z=H then gives X(Y\textsuperscript{2}+XZ). Re-derived the retained double-plane counterexample: its blowup double cover of P1 has4 branch values, exceptional square−2, residue discrepancy−1 and extra order1 from x. The unchanged polar program verifies all symbolic identities and the6 simple exact test solutions. Review provenance: New derivation after reading the supplied proof; not a blind review and not an independent human check}{\ref{night:P4.16}; \ref{night:P4.28}}

Let $S$ be an actual reduced support in Part~4.2. At a point
where the entire reduced critical locus of its local defining
function is one smooth curve, with generic transverse type
$A_3$, holomorphic coordinates put that function in one of
the forms
\[
 X^2+Y^4,\qquad X(Y^2+XZ).                              \tag{P4R.14.1}
\]
The normalization germs are respectively two smooth branches,
or a smooth plane and an $A_1$ quadric. No positive-genus
exceptional curve occurs over such a point.
\end{proposition}
\begin{proof}
\noindent\textbf{Step 1.} Choose a coordinate $z$ on the critical curve and normal
coordinates $x,y$. Its quadratic transverse Hessian has generic
rank one. Diagonalization of the symmetric matrix over
$\mathbb C\{z\}$, as in P4.28, gives the quadratic term
$h(z)x^2$. The cubic coefficient on the null direction
vanishes identically, by the generic transverse type. With
weights two for $x$ and one for $y$, the weight-four terms are
\[
 h(z)x^2+a(z)xy^2+b(z)y^4,\qquad
 \Delta(z)=4h(z)b(z)-a(z)^2.                             \tag{P4R.14.2}
\]
Every remaining term has larger weight and $\Delta$ is
nonzero generically.

\noindent\textbf{Step 2.} We first prove that $\Delta$ is a unit. A zero persists,
counted with multiplicity, under small independent constant
perturbations of $h,a,b$, by the argument principle on a
small parameter disc. Generic such perturbations make the
zeros simple and disjoint from $h=0$. Indeed $h=\Delta=0$
also requires $a=0$, which the independent parameters avoid;
where $h\ne0$, the derivative of $\Delta$ in the $b$
parameter is $4h\ne0$. These perturbations preserve the
critical axis and its transverse $A_3$ germ near the fixed
boundary. At a simple zero with $h\ne0$, the parametrized
Morse lemma gives
\[
 X^2+y^4U(y,z),\qquad U(0,0)=0,\quad U_z(0,0)\ne0.
\]
The actual coordinate $Z=U(y,z)$ gives $X^2+y^4Z$.
P4.16 supplies a compact spherical class with nonzero square
in its smoothing. The boundary comparison of P4.28 applies
here: near the boundary critical curve, the parametrized
Morse lemma and a fourth root of a unit give $X^2+Y^4$;
the tangential boundary coordinate makes the nonzero levels
submersive. Proper boundary transport therefore carries the
sphere to the actual nearby torus. Its universal cover is
$\mathbb R^4$, so a sphere has zero homology class and zero
square, a contradiction. Thus $\Delta$ is a unit.

\noindent\textbf{Step 3.} For $h(0)\ne0$, Morse splitting and a fourth root of the
residual unit give the first form in (P4R.14.1). Suppose
$h(0)=0$. Then $a(0)\ne0$. The holomorphic function
\[
 \frac{f(y^2v,y,z)}{y^4}
\]
restricts at $y=0$ to $h(z)v^2+a(z)v+b(z)$. Its finite
root $-b(0)/a(0)$ at the origin is simple. The holomorphic
implicit-function theorem gives a root $v=v(y,z)$, and
division by the monic equation of its graph gives
\[
 X=x-y^2v(y,z),\qquad f=Xg.
\]
On $X=0$, the function $g$ is $y^2$ times a unit. Its
parameter Morse centre in $y$ is therefore zero there and
its critical value is zero. The parametrized Morse lemma,
with parameters $X,z$, gives the exact form
\[
 f=X\bigl(Y^2+XH(X,z)\bigr).                            \tag{P4R.14.3}
\]
The remaining critical value is divisible by $X$ by its
vanishing on $X=0$. The function $H(0,z)$ is not identically
zero, since the generic transverse Hessian has rank one;
at the point under consideration $H(0,0)=0$.

\noindent\textbf{Step 4.} Put $G=2H+XH_X$. Off $X=0$, the critical equations of
(P4R.14.3) are $Y=0,H_z=0,G=0$. We prove $H_z(0,0)\ne0$.
Otherwise $H_z$ is a nonzero plane germ not divisible by
$X$, because $H(0,z)$ has finite order at least two.
No component of $H_z=0$ is contained in $G=0$: it would
give an additional critical-curve branch off $X=0$, whereas
the entire critical locus is the stated smooth axis. Hence
$(H_z,G)$ is an isolated complete intersection of positive
intersection multiplicity.

\noindent\textbf{Step 5.} Replace $H$ by $H-\eta z-\varepsilon$. Its plane critical
equations are
\[
 H_z-\eta=0,\qquad G-2\eta z-2\varepsilon=0.              \tag{P4R.14.4}
\]
Their solutions are the fibres of the holomorphic map
\[
 \Psi(X,z)=\left(H_z,\frac{G-2zH_z}{2}\right).
\]
Its zero fibre has the positive local degree just proved.
Choose a small compact representative with no boundary zero.
That degree is unchanged on a sufficiently small open
parameter ball. The universal equations in (P4R.14.4) have
parameter Jacobian determinant two, so regular parameter
values give a nonempty finite set of transverse solutions.
We also choose such a value outside the images under $\Psi$
of $X=0$ and $H_X=0$. These exceptional images have real
four-dimensional measure zero: each source curve has
countably many smooth real two-dimensional strata, and the
smooth image of each has measure zero. The germ $H_X$ is
not identically zero. Indeed in that case $H$ would depend
only on $z$ and both $H_z$ and $G=2H$ would vanish on the
whole line $z=0$, contrary to the isolated intersection.
The two exceptional images cannot exhaust the open ball
on which the positive degree is retained. This argument
uses no analyticity assertion for the image of a nonproper
map.

\noindent\textbf{Step 6.} At each chosen solution, the critical value is
\[
 f_{\eta,\varepsilon}=-\frac{X^3H_X}{2}\ne0.              \tag{P4R.14.5}
\]
For the indicated order of the two equations, put
\[
 J=\det\frac{\partial(G-2\eta z-2\varepsilon,\ H_z-\eta)}
                  {\partial(X,z)}.
\]
The $Y$-Hessian entry is $2X$, and the remaining Hessian
block at a critical point is
\[
 \begin{pmatrix}
 X(3H_X+XH_{XX})&X^2H_{Xz}\\
 X^2H_{Xz}&X^2H_{zz}
 \end{pmatrix}.
\]
Its determinant is $X^3J$, and the full Hessian determinant
is $2X^4J\ne0$. Thus these are isolated holomorphic Morse
points with nonzero critical values. Near the fixed boundary,
$H(0,z)$ remains nonzero under the small perturbation, so the
critical axis retains its transverse $A_3$ type. The same
proper boundary comparison transports a Morse vanishing
sphere, of square $-2$, into an actual nearby torus. This is
impossible. We have proved $H_z(0,0)\ne0$.

\noindent\textbf{Step 7.} The holomorphic coordinate $Z=H(X,z)$, keeping $X,Y$, now
gives the second form of (P4R.14.1). Its normalization has
the plane $X=0$ and the normal quadric $Y^2+XZ=0$. The
quadric's minimal resolution has one rational $(-2)$-curve.
The first form factors as $(X+iY^2)(X-iY^2)$ and its two
normalization branches are smooth. This proves the stated
normalization conclusions.

\noindent\textbf{Step 8.} The distinction between a pole coefficient and this
torus-smoothing conclusion can be checked by an exact
retained local test. Take
\[
 f=x\bigl[y^2+xz(x-z)(x-2z)\bigr].
\]
It has only the smooth $A_3$ critical axis $x=y=0$. The
second normalization branch is the double plane branched
over four distinct lines. Blowing up their common point
and normalizing gives a double cover of the exceptional
$\mathbb P^1$ branched in four points, hence an elliptic
curve of square $-2$. Its discrepancy is $-1$ and the
order of $x$ on it is one, giving pole coefficient two
for the displayed total reduced hypersurface. It is excluded
here by the polar deformation, not by the coefficient alone.
For $H=z(x-z)(x-2z)$, the deformation $H\mapsto H-1$ has
the exact Groebner basis
\[
 3x-5z^4-14z,\qquad 25z^6+50z^3-2.
\]
It gives six simple solutions; adjoining $x$, $H_x$ or
$J$ to the critical ideal gives the unit ideal in each
case. For $H\mapsto H-\varepsilon$ all six scale by
$\varepsilon^{1/3}$ and converge to the origin. The
identities and these exact rational calculations are
executed in \texttt{code/part4/smooth\_a3\_polar.py},
with JSON and text output retained.
\end{proof}

\subsection{A nonzero Hessian and the entire critical curve}

\begin{proposition}[P4R.15]\label{night:P4R.15}\label{stmt:P4R.15}
\StatementRecord{night:P4R.15}{proved}{Historical events: GPT-6 Astra (OpenAI), ChatGPT Pro mode, effort Pro, S6\_06\_BLIND, 16--18 November 2026: targeted proof check. Also checked in this writing. Pre-continuation writing events: GPT-6 Astra (OpenAI), ChatGPT Pro mode, effort Pro, S6\_10\_PAPERS, 18 November 2026 05:10:56.410177+00:00: computation or formal certificate; targeted proof check. Continuation metadata events: GPT-6 Astra (OpenAI), ChatGPT Pro mode, effort Pro, S6\_10\_PAPERS continuation, 19 November 2026: used and rechecked in this run; status-vocabulary reconciliation; metadata reconciliation only; no new mathematical check. \emph{Pre-continuation recorded proof scope.} The night-run re-derivation was by GPT-6 Astra (OpenAI, ChatGPT Pro mode, effort Pro), 16--18 November 2026. Also checked in this writing by GPT-6 Astra (OpenAI), ChatGPT Pro mode, effort Pro, 18 November 2026. proved. Full nonzero-Hessian product theorem allowing initially different transverse branch types. A nonzero complex symmetric Hessian supplies an actual Morse coordinate X and f=X\textsuperscript{2}+g(Y,Z). Factored the nonzero analytic plane function into its distinct irreducible factors with multiplicities. Two distinct factors define a finite positive-degree map; small independent target values give a transverse intersection avoiding all other factors. The unit-root coordinate gives X\textsuperscript{2}+Y\textasciicircum{}a Z\textasciicircum{}b, and P4.18 supplies its nonzero-square compact spherical class. Verified the boundary identification for every multiplicity: exponent1 remains smooth, higher exponents retain the product A-type; proper transport contradicts the torus. With a single factor g=h\textasciicircum{}m, a singular reduced h has finite positive gradient degree; subtracting a small linear form and shifting its value yields a Morse point with nonzero h-value. The full Hessian is diag(2,m v\textasciicircum{}(m−1)Q), of determinant2m\textsuperscript{2}v\textasciicircum{}(2m−2)detQ, and the same boundary transport excludes it. Therefore h is smooth and the actual coordinate Y=h gives the product. The unchanged arbitrary-m symbolic program passed. Review provenance: New derivation after reading the supplied proof; not a blind review and not an independent human check}{\ref{night:P4.18}; \ref{night:P4.28}}

At every critical-curve point of an actual reduced support in
Part~4.2 with nonzero ambient Hessian, holomorphic coordinates
give the exact product germ
\[
                       X^2+Y^m,\qquad m\ge2.
\]
In particular its entire reduced critical curve is smooth.
At a singular point of the reduced critical curve, the ambient
Hessian vanishes. This assertion allows different generic
transverse types on the different branches before the exclusion.
\end{proposition}
\begin{proof}
\noindent\textbf{Step 1.} A nonzero complex symmetric Hessian has a vector on which its
quadratic form is nonzero. The parametrized holomorphic Morse
lemma in that direction gives the exact equation
\[
                        f=X^2+g(Y,Z).
\]
The plane function $g$ is nonzero because the support is
reduced. Factor it in the convergent power-series ring as
\[
             g=u\prod_{i=1}^{r}h_i^{m_i},
             \qquad u(0)\ne0,\quad m_i\ge1,
\]
with distinct reduced irreducible factors $h_i$ vanishing at
the point. At least one exponent is at least two: a reduced
plane function has an isolated critical point, while the
critical locus here contains a curve.

\noindent\textbf{Step 2.} We exclude $r\ge2$ by an actual deformation of these factors.
For two distinct factors $h_i,h_j$, the holomorphic map
\[
 (h_i,h_j):(\mathbb C^2,0)\longrightarrow(\mathbb C^2,0)
\]
has isolated zero fibre and positive finite local degree.
On a sufficiently small compact representative that degree
is retained for small target values. Regular such values
give a nonempty transverse intersection of
$h_i-s_i=0$ and $h_j-s_j=0$. Independently choose the
other constants $s_k$ so that at least one of these
intersection points avoids every other zero factor. This
is possible because the chosen intersection set is finite
and each forbidden $s_k$ belongs to a finite set of values.
All parameters can be arbitrarily small.

\noindent\textbf{Step 3.} At the selected intersection, the two transverse factors
are holomorphic coordinates and every other factor is a
unit. Taking an actual local $m_i$th root of their unit
product and multiplying the first coordinate by it gives
the germ
\[
                        X^2+Y^{m_i}Z^{m_j}.
\]
P4.18 excludes every such germ by its compact spherical
class of nonzero square. The comparison here is to the
original torus degeneration, as follows. Choose a small
Milnor ball before moving the factors. Near its boundary,
the distinct reduced factors are smooth and disjoint.
The moved factor with exponent $m_i\ge2$ retains exactly
the transverse product type $A_{m_i-1}$ there; a factor
of exponent one remains a regular point of the surface.
Parametrized Morse coordinates and the appropriate root
of the remaining unit identify these boundary germs.
Solving the tangential boundary coordinate makes the
nonzero levels submersive. The proper boundary transport
used in P4.28 therefore identifies the needed nearby
levels and transports the spherical class into the actual
nearby torus. Its square cannot be nonzero there. Hence
there is only one factor.

\noindent\textbf{Step 4.} Absorb the unit by writing $h=u^{1/m}h_1$, so that
$g=h^m$. The reduced plane function $h$ is irreducible.
For a singular $h$, its gradient has an isolated zero,
and the finite positive local degree of
$(h_Y,h_Z)$ gives nonempty finite regular fibres under
small constant changes of that gradient. Equivalently,
subtracting a sufficiently small generic linear form
from $h$ gives isolated Morse critical points. Add a
small constant to make their critical values nonzero.
Call the resulting function $h_s$, and choose one such
critical point with value $v\ne0$ and nonsingular Hessian
$Q$. For $f_s=X^2+h_s^m$, the critical value and Hessian
are exactly
\[
 f_s=v^m\ne0,\qquad
 \operatorname{Hess}(f_s)
   =\begin{pmatrix}2&0\\0&m v^{m-1}Q\end{pmatrix}.
\]
Thus the full Hessian determinant is
$2m^2v^{2m-2}\det Q\ne0$. The moving curve $h_s=0$
retains its transverse $A_{m-1}$ type on the fixed
boundary. The same proper comparison transports the
Morse vanishing sphere, of square $-2$, to a nearby
torus, a contradiction. Consequently $h$ is smooth,
and using it as the coordinate $Y$ gives the asserted
product equation.

\noindent\textbf{Step 5.} The Hessian identity is executed symbolically for the
arbitrary positive integer $m$ in
\texttt{code/part4/hessian\_product.py}. Its retained
output also checks the moving-node example
$X^2+(YZ-\varepsilon)^m$ at a nonzero critical value.
\end{proof}


\subsection{The tacnode and transverse-line critical curve}

\begin{proposition}[P4R.16]\label{night:P4R.16}\label{stmt:P4R.16}
\StatementRecord{night:P4R.16}{proved}{Historical events: GPT-6 Astra (OpenAI), ChatGPT Pro mode, effort Pro, S6\_06\_BLIND, 16--18 November 2026: targeted proof check. Also checked in this writing. Pre-continuation writing events: GPT-6 Astra (OpenAI), ChatGPT Pro mode, effort Pro, S6\_10\_PAPERS, 18 November 2026 05:10:56.410177+00:00: computation or formal certificate; targeted proof check. Continuation metadata events: GPT-6 Astra (OpenAI), ChatGPT Pro mode, effort Pro, S6\_10\_PAPERS continuation, 19 November 2026: used and rechecked in this run; status-vocabulary reconciliation; metadata reconciliation only; no new mathematical check. \emph{Pre-continuation recorded proof scope.} The night-run re-derivation was by GPT-6 Astra (OpenAI, ChatGPT Pro mode, effort Pro), 16--18 November 2026. Also checked in this writing by GPT-6 Astra (OpenAI), ChatGPT Pro mode, effort Pro, 18 November 2026. proved. Full tacnode-plus-transverse-line symbolic square, rank-one obstruction and defect11 bound. Used the exact embedded ideal(xz,yz,y\textsuperscript{2}−x4) and its three parametrizations. The normalization sequence gives delta3 and zero ramification. Direct intersection of (z,h)\textsuperscript{2} and(x,y)\textsuperscript{2} gives P\textasciicircum{}(2)=P\textsuperscript{2}+(zh); the three displayed product identities kill zh by the maximal ideal, while its cubic term proves it is nonzero modulo P\textsuperscript{2}. In f=kappa zh+P\textsuperscript{2} the transverse determinant on(t,t\textsuperscript{2},0) has constant term−kappa\textsuperscript{2}, so rank1 forces kappa0 and the exact inclusion f\ensuremath{\in}P\textsuperscript{2}. Recomputed all three primitive gradient matrices, with elementary orders(1,2),(1,2),(1,1). The symmetric-square lattice inequality then gives3+3+3+2=11 before adding the nonnegative descent length, contradicting the exact total9. The unchanged a3\_tacnode\_line.py certificate passed. Review provenance: New derivation after reading the supplied proof; not a blind review and not an independent human check}{\ref{night:P4.33}; \ref{night:P4.40}; \ref{night:P4.42}; \ref{night:P4.44}}

In the remaining positive ruled case with uniform transverse
type $A_3$, the three-branch point of the critical curve cannot
be the union of a plane tacnode and an independent transverse
line. Its local conductor defect is at least eleven, whereas
the total defect in the contracted-section case is nine.
\end{proposition}
\begin{proof}
\noindent\textbf{Step 1.} The analytic ring calculation in P4.40 supplies actual ambient
coordinates for this curve. Its reduced ideal is
\[
 P=(xz,yz,h),\qquad h=y^2-x^4.
\]
Its three parametrizations are $(t,t^2,0)$, $(t,-t^2,0)$,
and $(0,0,t)$. The two tangent branches have intersection
multiplicity two, and their union meets the independent line
with intersection length one. The normalization sequence
therefore gives $\delta=3$. Every branch is immersive, so
its ramification divisor has coefficient zero.

\noindent\textbf{Step 2.} We first determine the symbolic square and the actual
equation. The plane-tacnode ideal $(z,h)$ is a reduced
complete intersection. Its square is unmixed and agrees
with the square of each prime at its generic point; hence
its symbolic square is its ordinary square. Intersecting
with the square of the line ideal gives
\[
 P^{(2)}=(z,h)^2\cap(x,y)^2=P^2+(zh).
\]
Indeed an element of $(z,h)^2$ is $Az^2+Bzh+Ch^2$.
Its last two terms belong to $(x,y)^2$. The first belongs
to that ideal precisely when $A$ does, because $z$ is a
nonzerodivisor modulo $(x,y)^2$. The resulting three
generators $x^2z^2,xyz^2,y^2z^2$ are already in $P^2$.
Furthermore
\[
 xzh=(xz)h,\qquad yzh=(yz)h,\qquad
 z^2h=(yz)^2-x^2(xz)^2.
\]
The class of $zh$ modulo $P^2$ is therefore annihilated
by the ambient maximal ideal. It is nonzero, since it has
a nonzero cubic term and $P^2$ has order at least four.
Thus $P^{(2)}/P^2$ is exactly one copy of the residue field.

\noindent\textbf{Step 3.} Let $f$ be the reduced defining function of the support.
Its generic transverse multiplicity is two, so
$f\in P^{(2)}$ and
\[
                       f=\kappa zh+f_0,\quad f_0\in P^2,
                       \quad\kappa\in\mathbb C.
\]
At $(t,t^2,0)$ with $t\ne0$, use $z,h$ as transverse
coordinates. The quadratic transverse part of $f$ has the
form
\[
 t^2a(t)z^2+(\kappa+t b(t))zh+c(t)h^2,
\]
where $a,b,c$ are holomorphic at $t=0$. The transverse
Hessian has rank one, so its determinant vanishes for
all sufficiently small nonzero $t$. Its constant term at
zero is $-\kappa^2$ (up to the harmless common convention
for the quadratic-form matrix). Hence $\kappa=0$, and
we have proved the exact inclusion $f\in P^2$.

\noindent\textbf{Step 4.} Let $K_i$ be the saturated rank-two kernel of the derivative
of the $i$th parametrization, and let $E_i\subset K_i$ be
the span of the differentiated generators of $P$. On the
two tangent branches choose the bases
\[
                dz,\quad dy\mp2t\,dx.
\]
The generator matrices have columns
\[
 \begin{pmatrix}t&\pm t^2&0\\0&0&\pm2t^2\end{pmatrix}.
\]
Their elementary-divisor orders are $(1,2)$. On the
independent line use $dx,dy$; the matrix is
\[
                   \begin{pmatrix}t&0&0\\0&t&0\end{pmatrix},
\]
with orders $(1,1)$.

\noindent\textbf{Step 5.} For completeness, the rank-one quadratic estimate of
P4.42 applies to these exact lattices as follows. Write
$E=t^a\mathbb C\{t\}e_1\oplus
t^b\mathbb C\{t\}e_2$, and write the saturated Hessian
line as $\ell=u e_1+v e_2$, with $(u,v)$ primitive.
Since $f\in P^2$, the tensor $t^d\ell^2$ belongs to
$\operatorname{Sym}^2E$. Therefore
\[
 \frac d2\ge\max(a-\operatorname{ord}u,
                         b-\operatorname{ord}v).
\]
Projection of $E$ to $K/\ell$ has colength
$b_{\rm proj}=\min(a+\operatorname{ord}v,
b+\operatorname{ord}u)$. Adding these two expressions
gives $d/2+b_{\rm proj}\ge a+b$; the identity follows
by subtracting the same two candidate quantities, with
the zero-component cases read directly. The local defect
formula of P4.42 consequently gives
\[
 \delta+\sum_i(d_i/2+b_i-r_i)+q
       \ge3+(1+2)+(1+2)+(1+1)=11.
\]
Here $q\ge0$ is the further descent colength, retained
rather than discarded from the definition. P4.44 contracts
the horizontal section, and P4.33 and (4.42a) give total
defect nine. The displayed lower bound contradicts that
total. This excludes the stated analytic curve germ.
\end{proof}

The symbolic-square identities, all three normalized
conormal matrices, and the rank-one determinant obstruction
are checked exactly in
\texttt{code/part4/a3\_tacnode\_line.py}.


\subsection{The nonordinary three-branch tacnode sector}

\begin{proposition}[P4R.17]\label{night:P4R.17}\label{stmt:P4R.17}
\StatementRecord{night:P4R.17}{proved}{Historical events: GPT-6 Astra (OpenAI), ChatGPT Pro mode, effort Pro, S6\_06\_BLIND, 16--18 November 2026: targeted proof check. Also checked in this writing. Pre-continuation writing events: GPT-6 Astra (OpenAI), ChatGPT Pro mode, effort Pro, S6\_10\_PAPERS, 18 November 2026 05:10:56.410177+00:00: computation or formal certificate; targeted proof check. Continuation metadata events: GPT-6 Astra (OpenAI), ChatGPT Pro mode, effort Pro, S6\_10\_PAPERS continuation, 19 November 2026: used and rechecked in this run; status-vocabulary reconciliation; metadata reconciliation only; no new mathematical check. \emph{Pre-continuation recorded proof scope.} The night-run re-derivation was by GPT-6 Astra (OpenAI, ChatGPT Pro mode, effort Pro), 16--18 November 2026. Also checked in this writing by GPT-6 Astra (OpenAI), ChatGPT Pro mode, effort Pro, 18 November 2026. proved. Full unrestricted delta-four nonplanar triple classification, exact ambient morphisms and all local bounds. For any normalized conormal lattice E, the exact sequence along the saturated Hessian line gives a+b=e1+e2. A unit Hessian column gives d\ensuremath{\geq}a and E\ensuremath{\subset}t\textasciicircum{}e1K gives b\ensuremath{\geq}e1, hence d/2+b\ensuremath{\geq}e1+e2/2 without assuming f\ensuremath{\in}P\textsuperscript{2}. The nonplanar branch estimate gives delta+3+m/2. For three smooth branches and delta4, classified the entire finite-colength analytic ring from its first2 jets. One first-jet direction gives the exact common-tangent model with E=t\textsuperscript{2}K on each branch, defect13. Two first-jet directions force two proportional branch evaluations; after actual parameter changes, the residual coupling mu\ensuremath{\neq}0 gives two analytic generators and hence a planar ring. For mu=0 the exact contact-three-plus-line ideal has conormal orders(1,3),(1,3),(1,1), giving21/2. For multiplicity4, branch delta2 would leave the minimally glued ring of embedding dimension4; hence the singular branch is the ordinary cusp. The colength-one module makes m0\textsuperscript{2}\ensuremath{\subset}A, and embedding dimension3 forces mA\textsuperscript{2}=m0\textsuperscript{2}. Classified its cotangent hyperplane and product map. A nonzero cusp-cubic coefficient gives independent tangent axes; the only allowed initial cubic xyz has rank-two leading Hessian and therefore vanishes, giving f\ensuremath{\in}m4 and defect11. With zero cubic coefficient, the three other coefficients must all be nonzero; the actual model has ideal(z(x−y),xy(x−y),z\textsuperscript{2}−x\textsuperscript{2}y). Its Hilbert–Burch minors, generic reducedness and depth give exact equality with the curve ideal. Primitive conormal matrices have orders(3,3),(1,2),(1,2), so defect23/2. The delta3 and delta2 cases then follow from the explicitly rederived excess-one ring classification. The unchanged a3\_delta4.py verifies all9 lattices and the4 exact jet identities; its received historical checker wording is not counted as a new check. Review provenance: New derivation after reading the supplied proof; not a blind review and not an independent human check}{\ref{night:P4.40}; \ref{night:P4.42}; \ref{night:P4.48}; \ref{night:P4R.10}; \ref{night:P4R.15}; \ref{night:P4R.16}}

In the remaining positive ruled case with uniform generic transverse
type $A_3$, every nonordinary three-branch critical-curve point has
local conductor defect strictly greater than nine. In particular,
the delta-four nonplanar cases have the explicit lower bounds
$13$, $21/2$, $11$, and $23/2$ below. The only three-branch germ
compatible with the total defect budget nine is the ordinary
nonplanar union of three coordinate axes.
\end{proposition}
\begin{proof}
Use the notation of P4R.10 and P4.42: \[
 D(C,f)=\delta(C)+
        \sum_i\left(\frac{d_i}{2}+b_i-r_i\right)+q,\qquad q\ge0.
\] Here \(d_i\) is the order of the saturated rank-one Hessian, \(b_i\) is the projected conormal colength, and \(r_i=m_i-1\) for a branch of multiplicity \(m_i\).

\par\medskip\noindent\textbf{1. Bounds without an ordinary-square assumption}\par\label{bounds-without-an-ordinary-square-assumption}

Let \(E_i\subset K_i\) be the normalized conormal lattice, with elementary divisors \((t^{e_1},t^{e_2})\), \(e_1\le e_2\). For the saturated Hessian line \(L_i\), put \[
 a_i=\operatorname{length}\bigl(L_i/(L_i\cap E_i)\bigr).
\] The exact lattice sequence gives \(a_i+b_i=e_1+e_2\). Every first derivative of \(f\) belongs to the radical ideal \(P\) of the critical curve. Thus every Hessian column belongs to \(E_i\). Write the Hessian as \(ut^{d_i}\ell_i^2\), with \(u\) a unit and \(\ell_i\) primitive. One coordinate of \(\ell_i\) is a unit; its column gives \(t^{d_i}\ell_i\in E_i\), so \(d_i\ge a_i\). Also \(b_i\ge e_1\), since \(E_i\subset t^{e_1}K_i\). Consequently \[
                  \frac{d_i}{2}+b_i
                    \ge e_1+\frac{e_2}{2}.           \tag{A}
\] This remains valid when a symbolic-square term occurs in \(f\).

Since the curve is nonplanar, \(P\) has no linear equation. Every defining gradient therefore has order at least \(m_i\) on branch \(i\), and \(b_i\ge m_i\). P4R.15 gives zero ambient Hessian, so \(f\) has order at least three and \(d_i\ge m_i\). For total multiplicity \(m=\sum m_i\), \[
                  D(C,f)\ge\delta(C)+3+\frac m2.       \tag{B}
\] Thus \(m\ge5\) gives \(D\ge19/2>9\). Only multiplicities three and four need inspection.

\par\medskip\noindent\textbf{2. Three smooth branches: the complete excess-two classification}\par\label{three-smooth-branches-the-complete-excess-two-classification}

Put \[
 B=\bigoplus_{i=1}^3\mathbb C\{t_i\},\quad
 F^j=\bigoplus_{i=1}^3t_i^j\mathbb C\{t_i\},\quad
 A_0=\mathbb C_{\rm diag}+F^1.
\] The actual ring \(A\subset A_0\) has \(\operatorname{length}(A_0/A)=2\), since \(A_0\) has delta invariant two. Each branch projection of \(A\) is onto.

Let \(W_1\subset F^1/F^2=\mathbb C^3\) be its first-jet image. Each coordinate projection is nonzero. If \(W_1=\mathbb C^3\), Nakayama first gives \(\mathfrak m_AA_0=F^1\) and then, on the finite \(A\)-module \(A_0/A\), gives \(A=A_0\). This contradicts colength two. Thus \(\dim W_1=1\) or \(2\).

\par\medskip\noindent\textbf{2a. One-dimensional first jets}\par\label{a.-one-dimensional-first-jets}

Scale the branch parameters so \(W_1=\mathbb C(1,1,1)\). The two missing first jets use the entire colength, hence \(F^2\subset A\). The exact algebra is \[
 A=\mathbb C_{\rm diag}+\mathbb C(t_1,t_2,t_3)+F^2.
\] It is generated by \[
 x=(t_1,t_2,t_3),\quad y=(t_1^2,0,0),\quad z=(0,t_2^2,0).
\] Indeed \(y,z,x^2-y-z\) give all three quadratic branch ideals, and multiplication by \(x\) gives every higher order. The embedded model is exactly \[
       P=\bigl(y(y-x^2),yz,z(z-x^2)\bigr).             \tag{C}
\] The branches are \((t,t^2,0),(t,0,t^2),(t,0,0)\). Its differentiated generators give \(E_i=t_i^2K_i\) on all three branches. For example, on \((t,t^2,0)\) use the primitive basis \(dy-2t\,dx,dz\). Formula (A) gives \[
                         D\ge4+3\cdot3=13.            \tag{D}
\]

\par\medskip\noindent\textbf{2b. Two-dimensional first jets}\par\label{b.-two-dimensional-first-jets}

The three branch evaluations are nonzero linear forms on \(W_1\). If no two are proportional, their squares are independent: choose the first two as \(s,t\); the third is \(as+bt\), \(ab\ne0\), whose square supplies the nonzero \(st\) coefficient. Thus products give the whole second-jet space. Multiplication by one first jet nonzero on all branches then gives every later jet. This would leave only one missing positive-jet dimension, contrary to colength two.

Two evaluations are therefore proportional. Scale and label the parameters to obtain \(W_1=\{(a,a,b)\}\). Its products give the same two-dimensional second-jet space. The one remaining missing dimension occurs there: if the second-jet image were full, multiplication by a nonvanishing first jet would make every later one full. Hence \(F^3\subset A\), and the pure second-jet image is exactly \(\{(c,c,d)\}\).

A lift of \(W_1\) has quadratic difference \(\lambda a+\mu b\). The analytic changes \(t_i\mapsto t_i+c_it_i^2\), on the first two branches, remove \(\lambda a\). Thus the exact ring has common constants, with first and second coefficients satisfying \[
                     a_1=a_2,\qquad c_1-c_2=\mu a_3, \tag{E}
\] and arbitrary independent terms of order at least three.

For \(\mu\ne0\), the two elements \[
 x=(t_1,t_2,0),\qquad y=(\mu t_1^2,0,t_3)
\] generate this whole ring. The exact identities \[
 \begin{aligned}
 xy/\mu&=(t_1^3,0,0),\\
 x^3-xy/\mu&=(0,t_2^3,0),\\
 y^3-\mu(xy)^2&=(0,0,t_3^3),\\
 y^2-\mu x^2y&=(0,0,t_3^2)
 \end{aligned}
\] give the independent cubic ideals and every permitted quadratic jet. Multiplication by \(x\) or \(y\) gives all the higher branch series. This ring is planar, contrary to the present hypothesis.

Thus \(\mu=0\). The ring is generated by \((t_1,t_2,0),(0,0,t_3),(t_1^3,0,0)\). Naming the independent transverse coordinate \(z\) gives the exact embedded model \[
                       P=(xz,yz,y(y-x^3)).            \tag{F}
\] Its branches are \((t,0,0),(t,t^3,0),(0,0,t)\), and its conormal elementary divisors are \[
                         (1,3),\quad(1,3),\quad(1,1).
\] On the second branch use \(dy-3t^2dx,dz\); the gradients give \(t^3(dy-3t^2dx)\) and \(t\,dz\). Without restricting a symbolic-square coefficient of \(f\), (A) gives \[
                  D\ge4+\frac52+\frac52+\frac32
                       =\frac{21}{2}>9.               \tag{G}
\]

These are actual analytic ring isomorphisms. In each nonplanar case the three displayed generators form a basis of the curve cotangent space. Lift them to the ambient local ring. Their linear Jacobian is invertible, so the inverse-function theorem gives an ambient coordinate isomorphism to the stated model.

\par\medskip\noindent\textbf{3. Multiplicity four: one cusp and two smooth branches}\par\label{multiplicity-four-one-cusp-and-two-smooth-branches}

The singular branch has multiplicity two and delta at least one. Joining three branches costs at least two further delta units. If the branch delta is two, all remaining gluing is minimal, and the local ring is \[
 \mathbb C_{\rm diag}
 +\mathfrak m_{\rm singular}
 +t_2\mathbb C\{t_2\}+t_3\mathbb C\{t_3\}.
\] Its embedding dimension is at least \(2+1+1=4\), impossible in a smooth threefold. A larger branch delta is impossible. Thus this branch is the ordinary cusp \(\mathbb C\{t^2,t^3\}\).

Set \[
 A_0=\mathbb C_{\rm diag}
 +\mathfrak m_{\mathbb C\{t^2,t^3\}}
 +t_2\mathbb C\{t_2\}+t_3\mathbb C\{t_3\}.
\] Now \(A_0\) has delta three, and \(A\subset A_0\) has colength one. The one-dimensional \(A\)-module \(A_0/A\) is killed by \(\mathfrak m_A\), so \(\mathfrak m_AA_0\subset A\). Its elements have zero constant term, giving \(\mathfrak m_AA_0=\mathfrak m_A\). Hence \(\mathfrak m_A\) is an ideal of \(A_0\). The local algebra \(A_0/\mathfrak m_A\) has length two, so its maximal ideal squares to zero. Therefore \[
                         \mathfrak m_0^2\subset A.     \tag{H}
\]

The four cotangent coordinates of \(A_0\), denoted \(a,b,c,d\), are represented by the cusp functions \(t^2,t^3\) and the two smooth parameters. The three-dimensional subspace \(W=\mathfrak m_A/\mathfrak m_0^2\) is a hyperplane \[
                 \lambda a+\mu b+\nu c+\xi d=0.       \tag{I}
\] Its projection to \(a,b\) is onto, since the actual branch image is the full cusp ring, and the two smooth projections are nonzero. Ambient embedding dimension three gives \(\mathfrak m_A^2=\mathfrak m_0^2\). Thus products of \(W\) must generate \(\mathfrak m_0^2/\mathfrak m_0^3\), whose coordinates are \[
                              a^2,\quad ab,\quad c^2,\quad d^2.
                                                               \tag{J}
\]

\par\medskip\noindent\textbf{3a. A nonzero coefficient of the cusp cubic}\par\label{a.-a-nonzero-coefficient-of-the-cusp-cubic}

For \(\mu\ne0\), solve (I) for \(b\). Surjectivity of the cusp projection forces \((\nu,\xi)\ne(0,0)\). The three minimal ambient coordinate functions have parametrizations \[
 \begin{array}{c|c}
 \text{cusp}&(t^2-\lambda't^3,-\nu't^3,-\xi't^3)\\
 \text{second branch}&(0,t,0)\\
 \text{third branch}&(0,0,t).
 \end{array}
\] Their tangent directions are independent.

By P4R.15, \(f\in\mathfrak m^3\). Let \(G_3\) be its homogeneous cubic part. Since \(df\) vanishes on all three branches, \(dG_3\) vanishes on each independent tangent axis. Every cubic monomial other than \(xyz\) has a derivative nonzero on an axis. Hence \(G_3=cxyz\). On the \(x\)-axis its normal Hessian is \[
                         \begin{pmatrix}0&c\\c&0\end{pmatrix}.
\] A nonzero \(c\) gives a nonzero leading minor on the actual branch, contradicting generic rank one. Thus \(c=0\), and \(f\in\mathfrak m^4\).

The Hessian entries have order at least \(2m_i\) on branch \(i\), and \(b_i\ge m_i\). Its defect contribution is at least \(m_i+1\). Therefore \[
                            D\ge4+(4+3)=11>9.         \tag{K}
\]

\par\medskip\noindent\textbf{3b. A zero coefficient of the cusp cubic}\par\label{b.-a-zero-coefficient-of-the-cusp-cubic}

For \(\mu=0\), surjectivity in (J) forces \(\lambda,\nu,\xi\) all to be nonzero. Otherwise two of \(a,c,d\) are proportional on \(W\), so their squares are dependent, and (J) has rank at most three. That creates an extra minimal generator, contrary to embedding dimension three. Scale the branch parameters so that \(a+c+d=0\).

The pure hyperplane generators belong to \(A\) exactly by (H). They give the parametrizations \[
                       (t^2,t^2,t^3),\quad(t,0,0),\quad(0,t,0).
\] Their defining ideal is \[
                 P=(z(x-y),xy(x-y),z^2-x^2y).        \tag{L}
\] A Hilbert--Burch matrix is \[
 \begin{pmatrix}
 xy&z\\
 -z&-x\\
 0&-(x-y)
 \end{pmatrix}.
\] Its maximal minors give (L), with alternating signs. The ideal has codimension two and is Cohen--Macaulay. Its zero set is exactly the cusp and the two axes, and it is generically reduced on each branch. A nonzero nilpotent quotient supported at the origin would violate depth one, so it is the actual reduced intersection ideal.

On the cusp use the primitive conormal basis \[
                     k_1=dx-dy,\qquad k_2=3t\,dx-2dz.
\] The differentiated generators are \(t^3k_1,t^4k_1,t^3(tk_1-k_2)\). Thus its elementary divisors are \((3,3)\). On each smooth axis they are \((1,2)\). Here \(r=1\) on the cusp and zero on the axes. By (A), \[
                  D\ge4+\left(\frac92-1\right)+2+2
                       =\frac{23}{2}>9.               \tag{M}
\]

This exhausts multiplicity four. Together with Sections 1--2 it proves the strict local bound for every nonplanar three-branch curve with delta four. Only actual coordinate isomorphisms, conormal maps, and the rank-one Hessian were used; no boundary deformation has been introduced.

\par\medskip\noindent\textbf{4. Every nonordinary triple exceeds the budget}\par\label{every-nonordinary-triple-exceeds-the-budget}

For a planar three-branch point, P4.48 gives defect at least ten, including the ordinary plane triple point. At a nonplanar point with delta at least five, (B) and \(m\ge3\) give \[
                  D\ge5+3+\frac32=\frac{19}{2}>9.
\] Delta four was exhausted above.

At delta three the excess \(\delta-(r-1)\) is one. The actual analytic ring calculation in P4.40 gives precisely an ordinary plane triple, a plane tacnode with an independent transverse line, or a cusp with two independent lines. The first is excluded by P4.48. The second has defect at least eleven by P4R.16. The third has embedding dimension four and cannot occur in the ambient smooth threefold.

Finally a three-branch point with delta two has zero excess. Its ring equals \(\mathbb C_{\rm diag}+F^1\), generated by the independent branch parameters. It is exactly the three coordinate axes, with ideal \((xy,xz,yz)\). Thus every nonordinary three-branch point is excluded from the total defect budget nine. The only remaining local possibility is the ordinary nonplanar triple treated by P4.48 and the separate full-quartic conductor calculation.

\end{proof}

The analytic model identities, primitive conormal lattices,
Hilbert--Burch minors, cubic obstruction, and jet generators are
checked exactly in \path{code/part4/a3_delta4.py}.


\subsection{The horizontal image in the unrestricted ruled remainder}

\begin{proposition}[P4R.18]\label{night:P4R.18}\label{stmt:P4R.18}
\StatementRecord{night:P4R.18}{proved}{Historical events: GPT-6 Astra (OpenAI), ChatGPT Pro mode, effort Pro, S6\_06\_BLIND, 16--18 November 2026: targeted proof check. Also checked in this writing. Pre-continuation writing events: GPT-6 Astra (OpenAI), ChatGPT Pro mode, effort Pro, S6\_10\_PAPERS, 18 November 2026: targeted proof check; used and rechecked in this run. Continuation metadata events: GPT-6 Astra (OpenAI), ChatGPT Pro mode, effort Pro, S6\_10\_PAPERS continuation, 19 November 2026: used and rechecked in this run; status-vocabulary reconciliation; metadata reconciliation only; no new mathematical check. \emph{Pre-continuation recorded proof scope.} The night-run re-derivation was by GPT-6 Astra (OpenAI, ChatGPT Pro mode, effort Pro), 16--18 November 2026. Also checked in this writing by GPT-6 Astra (OpenAI), ChatGPT Pro mode, effort Pro, 18 November 2026. The A3/D4 conductor alternatives, primary defect budget, at-most-three branch-value graph bound, and degree-two Riemann-Hurwitz contradiction were re-derived, using the previously stated local product models and defect bounds}{\ref{night:P4.33}; \ref{night:P4.42}; \ref{night:P4R.10}; \ref{night:P4R.13}; \ref{night:P4R.14}; \ref{night:P4R.15}; \ref{night:P4R.6}}

In the remaining positive-Euler ruled case of Part~4.2, a
noncontracted horizontal pole section maps with degree two or
three to a rational normalized conductor component whose generic
transverse germ is $D_4$. This assertion imposes no transverse
type or conductor-exponent restriction on the other image component.
\end{proposition}
\begin{proof}
\noindent\textbf{Step 1.} Write $H$ for the horizontal section, $g=g(H)\ge1$, and
$C_H$ for its image in the reduced conductor $C$. The preceding
ruled reduction gives two rational normalized components of $C$
and one point-identification cycle. The nonconstant map
\[
                         H\longrightarrow C_H^\nu=\mathbb P^1
\]
has degree at least two: a degree-one finite map between smooth
curves is an isomorphism, whereas their genera differ. Each
generic point of $H$ has conductor exponent two, by the actual
pole divisor of P4R.6. Thus the generic transverse plane curve
along $C_H$ has at least two normalization branches of exponent two.

\noindent\textbf{Step 2.} For its branches $B_i$ the exact conductor formula is
\[
             c_i=2\delta(B_i)+\sum_{j\ne i}I(B_i,B_j).
\]
A singular branch with $c_i=2$ permits no other branch,
so the two specified branches are smooth. With two total
branches their intersection multiplicity is two, giving $A_3$.
With three total branches, the first two exponents force all
three pairwise intersection multiplicities to equal one;
the third branch is then smooth as well, giving $D_4$.
Four branches would give exponent at least three on every
smooth branch. These possibilities are exhaustive. In type
$A_3$, $H$ has degree exactly two and supplies both generic
normalization preimages. In type $D_4$ its degree is two or
three. We exclude the former transverse type next.

\noindent\textbf{Step 3.} Assume that the transverse germ along $C_H$ is $A_3$.
Let $C_{H,\rm sch}$ be the primary component of the actual
image conductor scheme supported on $C_H$. The primary
restriction construction of P4R.13 applies before its
elliptic-genus calculation, without change. More explicitly,
P4.33 gives rational singularities of the normalization when
$H$ is not contracted, so finite duality and the canonical
trace give the injections
\[
 \mathcal O_{C_{\rm sch}}\longrightarrow
 \nu_*(\mathcal O_N/I_N)\longrightarrow\rho_*\mathcal O_D.
\]
The exact blowup relation $D'=\beta^*D''-E$ identifies
global functions on $D$ with those on $2H_0+V_0$ of the
minimal ruled model. The kernel of restriction to $2H_0$
has successive vertical-fibre layers
$\mathcal O_{\mathbb P^1}(-2)$, with no global sections.
Hence a global function is determined at the generic doubled
horizontal section. It follows that restriction from
$C_{\rm sch}$ to $C_{H,\rm sch}$ is injective on $H^0$
and surjective on $H^1$: the first assertion follows from
those injections and generic determination, and the second
from the dimension-one kernel of the restriction sheaf map.
Therefore
\[
                 g-1=\chi(\mathcal O_{C_{\rm sch}})
                    \le\chi(\mathcal O_{C_{H,\rm sch}}).
                                                            \tag{P4R.18.1}
\]
Every arrow here is the primary restriction of the same
conductor algebra; no other conductor component is deleted
from the left side.

\noindent\textbf{Step 4.} We evaluate the right side on $C_H^\nu=\mathbb P^1$.
Use the actual differentiated conormal, the saturated Hessian
line and the rank-one nilradical of P4R.10. The kernel $K$
of the normalized derivative has degree $2+\deg R$, since
$c_1(Y)\cdot C_H=0$. The Hessian line has degree
$\tfrac12\sum d_p$. Riemann--Roch and the normalization
sequence give exactly
\[
 \chi(\mathcal O_{C_{H,\rm sch}})
 =4-\sum_p\left[\delta(C_H,p)
      +\sum_{i\mapsto p}(d_i/2+b_i-r_i)+q_p\right].
                                                            \tag{P4R.18.2}
\]
All terms are those of this primary curve and its primary
conductor. They remain defined at its intersections with
the other component, because the pure conductor quotient
extends its generic containment of the symbolic square.

\noindent\textbf{Step 5.} Every singular point of $C_H$ contributes at least three
to the sum in (P4R.18.2). At a planar point the ideal is
$(u,h)$, its symbolic square equals its ordinary square,
and the actual equation belongs to that square. The
rank-one symmetric-square estimate of P4.42 and the exact
plane-curve different calculation give a contribution at
least $3\delta(C_H,p)+q_p\ge3$.
At a nonplanar point every defining equation has order at
least two. On a branch of multiplicity $m_i$ its projected
conormal order is at least $m_i$. The point is singular
on the entire critical curve, so P4R.15 makes the ambient
Hessian zero and gives $d_i\ge m_i$. Since $r_i=m_i-1$,
the contribution is at least
\[
                      \delta(C_H,p)+\sum_i(1+m_i/2),
\]
which is at least three at a singular curve point. The
same arguments give nonnegative contributions everywhere;
at a smooth primary point its conormal is the full kernel
and its only possible contribution is the nonnegative
Hessian zero order divided by two.
Combining (P4R.18.1)--(P4R.18.2) gives total primary
defect at most $5-g\le4$. Thus $C_H$ has at most one
singular point, and two different singular points are
impossible even in the presence of arbitrary transverse
types on the other global component.

\noindent\textbf{Step 6.} We count the possible branch values of the degree-two map
$H\to\mathbb P^1$ with the identifications retained.
The normalization graph formula is
\[
                         \sum_{p\in C}(r_p-1)=2.
\]
Connectedness leaves either two two-branch points or one
three-branch point. In the two-point case, two intersections
of the global components give two points on $C_H^\nu$.
There can be at most one additional unibranch singularity
of $C_H$, so these supply at most three possible branch
values. The other connected possibility is a self-identification
on one component and one intersection of the two components.
A self-identification on $C_H$ supplies two normalization
points and already costs at least three in the primary
defect. It permits no additional singular point, so with
the intersection there are again at most three possible
branch values. A self-identification on the other component
supplies no additional point on $C_H^\nu$.

\noindent\textbf{Step 7.} At a three-branch point, the two global components must
supply two and one branches. Thus $C_H^\nu$ has at most
two points over it. Two points already make $C_H$ singular
there and exclude any further singular point by the defect
bound. One point permits at most one further unibranch
singularity. This case also has at most three possible
branch values.

\noindent\textbf{Step 8.} There is no ramification elsewhere. At a point where the
entire critical curve is smooth, P4R.14 gives precisely
the product tacnode or $X(Y^2+XZ)$. Both have two distinct
normalization points over that point of the conductor.
Their two generic preimages lie on $H$, so properness
places both special preimages on $H$ as well. The
degree-two map is therefore unramified there. We have
accounted for every possible ramification value.
Riemann--Hurwitz for a degree-two map from a genus-$g$
curve requires $2g+2\ge4$ distinct branch values; each
ramified fibre consists of its single point of index two.
This contradicts the count of at most three. Thus the
generic germ along $C_H$ must be $D_4$, with degree two
or three as proved at the start.
\end{proof}

The exact connected identification patterns and their primary
defect/branch-value counts are enumerated in
\texttt{code/part4/horizontal\_primary.py}.


\subsection{A triple spectral root with nonzero cubic tensor}

\begin{proposition}[P4R.19]\label{night:P4R.19}\label{stmt:P4R.19}
\StatementRecord{night:P4R.19}{proved}{Historical events: GPT-6 Astra (OpenAI), ChatGPT Pro mode, effort Pro, S6\_06\_BLIND, 16--18 November 2026: targeted proof check. Also checked in this writing. Pre-continuation writing events: GPT-6 Astra (OpenAI), ChatGPT Pro mode, effort Pro, S6\_10\_PAPERS, 18 November 2026: targeted proof check. Continuation metadata events: GPT-6 Astra (OpenAI), ChatGPT Pro mode, effort Pro, S6\_10\_PAPERS continuation, 19 November 2026: used and rechecked in this run; status-vocabulary reconciliation; metadata reconciliation only; no new mathematical check. \emph{Pre-continuation recorded proof scope.} The night-run re-derivation was by GPT-6 Astra (OpenAI, ChatGPT Pro mode, effort Pro), 16--18 November 2026. Also checked in this writing by GPT-6 Astra (OpenAI), ChatGPT Pro mode, effort Pro, 18 November 2026. The exact cubic preparation and all three Newton slopes were checked. The explicit finite proper maps to two-node conic fibrations carry compact cycle classes of nonzero square back with their degrees; the repeated-factor balanced cubic remains connected. Anisotropic scaling transfers these classes to an actual Milnor fibre, contradicting P4.51}{\ref{night:P4.51}; \ref{night:P4.52}}

Let $f$ be the actual reduced local equation of the remaining
positive ruled support, and let $C_0$ be a smooth local branch
of its critical curve with generic transverse type $D_4$.
At a point where the cubic normal tensor along $C_0$ is
nonzero, that tensor is not the cube of a linear form.
Other critical-curve branches at the point are retained,
without a restriction on their generic transverse types.
\end{proposition}
\begin{proof}
\noindent\textbf{Step 1.} Choose coordinates $(x,y,z)$ with $C_0=(x=y=0)$.
Generic normal multiplicity three gives
$f\in(x,y)^{(3)}$. Powers of this smooth
complete-intersection ideal are primary, so
$f\in(x,y)^3$ at the specified point.

\noindent\textbf{Step 2.} Suppose its nonzero cubic normal tensor is a cube.
A linear normal-coordinate change makes it $x^3$, up
to a nonzero scalar. Weierstrass preparation in $x$
gives the exact equality
\[
 f=U(x,y,z)\bigl[x^3+a_2(y,z)x^2+a_1(y,z)x+a_0(y,z)\bigr],
\]
where $U$ is a holomorphic unit. The bracket belongs to
$(x,y)^3$, so $a_2,a_1,a_0$ are divisible by $y,y^2,y^3$,
respectively. Depressing the monic cubic is a holomorphic
coordinate change preserving $(x,y)$ and yields
\[
 f=U(x,y,z)\bigl[x^3+y^2B(y,z)x+y^3C(y,z)\bigr],
 \qquad B(0,0)=C(0,0)=0.                         \tag{R19.1}
\]
The unit remains in this equality and in the scaling
comparison below.

\noindent\textbf{Step 3.} Set
\[
 k=\operatorname{ord}_z B(0,z),\qquad
 l=\operatorname{ord}_z C(0,z),\qquad
 \rho=\min(k/2,l/3),
\]
where the order of an identically zero series is infinity.
At least one order is finite: otherwise all generic normal
cubics would be $x^3$, contrary to generic $D_4$.
Thus $0<\rho<\infty$. Choose $L>3\rho$ and positive
rational weights
\[
       \operatorname{wt}(x)=L+\rho,\qquad
       \operatorname{wt}(y)=L,\qquad
       \operatorname{wt}(z)=1.
\]
Multiplication by a common denominator gives positive
integer weights. The weighted degree of the first face
is $3L+3\rho$. Terms of $B$ containing $y$ have weight
gap at least $L-2\rho>0$, and those of $C$ containing
$y$ have gap at least $L-3\rho>0$. Consequently the
leading polynomial, apart from the nonzero factor $U(0)$,
is exactly one of
\[
\begin{array}{ll}
 P=x(x^2+a y^2z^k),\quad a\ne0,
       & k/2<l/3,\\[2pt]
 P=x^3+b y^3z^l,\quad b\ne0,
       & l/3<k/2,\\[2pt]
 P=Q(x,yz^r),\quad
 Q(X,V)=X^3+aXV^2+bV^3,\quad a,b\ne0,
       & k=2r, l=3r, r\ge1.
\end{array}                                                   \tag{R19.2}
\]
In the last case $r$ is an integer. The depressed cubic
$Q$ is not a cube: a cube with $X^3$ coefficient one and
$X^2V$ coefficient zero would have $a=b=0$. A double
factor of $Q$ is allowed throughout the argument.

\par\medskip\noindent\textbf{The compact class used in the comparison.}
Let $B$ be a connected smooth complex curve and let
$h:B\to\mathbb C$ have two distinct simple zeros $p,q$.
Join them by a smooth embedded arc $\gamma$ whose
interior contains no zero. On the smooth surface
\[
             Z_h=\{(b,u,v)\in B\times\mathbb C^2:
                                      uv=h(b)\},
\]
consider the locus over $\gamma$ with $|u|=|v|$.
It is a circle over each interior point and collapses
to a point at each end. Near an endpoint choose the
arc so that $h$ is a fixed nonzero complex multiple
of a positive real parameter $s$. Then
$u=\sqrt{s}\,e^{\mathrm i\phi}$ and
$v=c\sqrt{s}\,e^{-\mathrm i\phi}$ are the smooth
polar coordinates of a real disc in the local $uv$
chart. The locus is thus a smoothly embedded sphere.

It is Lagrangian for a product K\"ahler form.
Writing $u=R e^{\mathrm i\phi}$,
$v=R e^{\mathrm i(\psi-\phi)}$, the sum of the two
fibre area forms restricts to
$d(R^2/2)\wedge d\psi=0$, since $R,\psi$ depend
only on the arc parameter. The base area form also
restricts to zero. Multiplication by $\mathrm i$
identifies tangent and normal bundles, with the
opposite normal orientation relative to the complex
ambient orientation. Its self-intersection is
therefore $-\chi(S^2)=-2$.
For the compactly supported Poincar\'e dual $\alpha$,
\[
                         \int_{Z_h}\alpha^2=-2.             \tag{R19.3}
\]

If $\pi:W\to Z_h$ is a finite proper holomorphic map
of degree $d>0$ between smooth complex surfaces, then
$\pi^*\alpha$ is compactly supported and
\[
              \int_W(\pi^*\alpha)^2
                       =d\int_{Z_h}\alpha^2=-2d.            \tag{R19.4}
\]
The equality follows from the proper oriented degree;
off the analytic branch locus the map has $d$ positive
sheets, and the branch locus has measure zero for
integration of compactly supported top forms.
Poincar\'e duality gives a compact homology two-class
with that square. Its representative need not itself
be a sphere.

\par\medskip\noindent\textbf{The exact finite maps.}
In the first case of (R19.2), $P=1$ forces $x\ne0$
and gives
\[
           y^2z^k=h(x),\qquad
           h(x)=\frac{1-x^3}{ax},\qquad x\in\mathbb C^*.
\]
The map
\[
 (x,y,z)\longmapsto(x,u=y^2,v=z^k)
\]
is finite proper of degree $2k$ onto $uv=h(x)$.
The target is smooth: at $u=v=0$ the zeros of $h$
are simple, and elsewhere a derivative in $u$ or
$v$ is nonzero. The source is smooth by the weighted
Euler identity on $P=1$. Join two of the three cube
roots of one by an arc in $\mathbb C^*$ avoiding the
third. Equations (R19.3)--(R19.4) give a compact
two-class of square $-4k$ on the leading fibre.

In the second case the exact finite proper map is
\[
 \begin{split}
 \{x^3+b y^3z^l=1\}&\longrightarrow
                 \{uv=(1-x^3)/b\},\\
 (x,y,z)&\longmapsto(x,u=y^3,v=z^l).
 \end{split}
\]
It has degree $3l$, and the base is $\mathbb C$.
The same three simple zeros supply a pulled-back
compact class of square $-6l$.

For the balanced case let
\[
                 B_Q=\{(X,V)\in\mathbb C^2:Q(X,V)=1\}.
\]
It is smooth by the ordinary Euler identity and is
connected even when $Q$ has a double factor. Its
projection to $\mathbb P^1$ minus the distinct zero
directions of $Q$ is an unramified cyclic cover of
degree three. Around a zero direction of multiplicity
one or two its monodromy generates $\mathbb Z/3$.
Such a direction exists because $Q$ is not a cube,
so the cover is connected.

The function $V$ on $B_Q$ has precisely the three
zeros $(\zeta,0)$ with $\zeta^3=1$, and they are
simple because $Q_X(\zeta,0)=3\zeta^2\ne0$.
The surface
\[
             Z_Q=\{(b,u,w):b\in B_Q,\ uw=V(b)\}
\]
is smooth. The exact map
\[
 \begin{split}
       \{Q(x,yz^r)=1\}&\longrightarrow Z_Q,\\
       (x,y,z)&\longmapsto
                    \bigl(b=(x,yz^r),u=y,w=z^r\bigr)
 \end{split}                                                   \tag{R19.5}
\]
is finite proper of degree $r$: $z$ satisfies the
monic equation $z^r=w$ and the other coordinates are
retained. An arc in $B_Q$ joining two $V$-zeros and
avoiding the third supplies a compact class of square
$-2r$ on the source. In particular this argument
retains $Q=(X-\lambda V)^2(X+2\lambda V)$,
$\lambda\ne0$, without a perturbation of its
repeated factor.

\par\medskip\noindent\textbf{Transport to the actual local fibre.}
Let $(w_x,w_y,w_z)$ be the positive integer weights
and let $d$ be the resulting degree. Put
\[
 D_\epsilon(x,y,z)=
       (\epsilon^{w_x}x,\epsilon^{w_y}y,\epsilon^{w_z}z),
 \qquad
 G_\epsilon=\epsilon^{-d}f\circ D_\epsilon.
\]
These actual functions converge with derivatives on
compact sets to $U(0)P$. The weighted Euler identity
makes $dP$ nonzero on $P=1$, including the repeated-factor
case on its zero set. Choose compact representatives
of the two-class above and of the transverse
perturbation computing its square. On a compact
neighbourhood of these chains, the gradient-flow
comparison of P4.52 solves
\[
       dG_\epsilon(V_\epsilon)
                         =-\partial_\epsilon G_\epsilon.
\]
The denominator is bounded away from zero on a
slightly larger neighbourhood; a cutoff gives a
flow for sufficiently small positive $\epsilon$.
It transports the chains and their intersections
to $G_\epsilon=U(0)$. Scaling by $D_\epsilon$ then
embeds them in the original fixed Milnor ball on
\[
                         f=U(0)\epsilon^d.
\]
The oriented intersection number is preserved and
equals $-4k$, $-6l$, or $-2r$. Each is nonzero,
contradicting P4.51. This proves the assertion.
\end{proof}

\par\medskip\noindent\textbf{The two-direction exception is retained.}
The assertion cannot be enlarged to every spectral
collision. The reduced equation
\[
                       f=y(x^2+y^2z^2)
\]
has generic $D_4$ along $x=y=0$ and generic $A_1$
along $x=z=0$. Only two spectral directions collide
at their intersection. Its non-plane normalization
sheets $x=\pm\mathrm i yz$ have pole divisor
$2(y=0)+(z=0)$, with simple normal crossings.
Its affine fibre $f=1$ is exactly $(\mathbb C^*)^2$:
\[
 \begin{split}
       a&=x+\mathrm i yz,\qquad b=x-\mathrm i yz,\\
       x&=(a+b)/2,\qquad y=(ab)^{-1},\qquad
       z=(a-b)ab/(2\mathrm i).
 \end{split}
\]
The displayed maps are inverse and substitution gives
$f=1$. Its intersection form is zero: the product
of two unit circles generates its second homology,
and a radial translation gives a disjoint homologous
copy. More generally the fibre of $y(x^2+y^2z^n)$
becomes $u^2+z^n=t^3$, $t\ne0$, through
$t=1/y$, $u=x/y$. Positive radial dilation gives
a product real direction, along which any compact
cycle can be displaced from another compact cycle.
Its intersection form is also zero. Thus the
other generic transverse type remains unrestricted.

The exact weighted gaps, all three finite-map
equations, their degree-square identities, and the
inverse maps of the two-direction exception are
executed in \path{code/part4/d4_spectral_triple.py};
its \path{.json} and \path{.out} outputs report PASS.


\subsection{The three-branch image point with a double horizontal cover}

\begin{proposition}[P4R.20]\label{night:P4R.20}\label{stmt:P4R.20}
\StatementRecord{night:P4R.20}{proved}{Historical events: GPT-6 Astra (OpenAI), ChatGPT Pro mode, effort Pro, S6\_06\_BLIND, 16--18 November 2026: targeted proof check. Also checked in this writing. Pre-continuation writing events: GPT-6 Astra (OpenAI), ChatGPT Pro mode, effort Pro, S6\_10\_PAPERS, 18 November 2026: targeted proof check. Continuation metadata events: GPT-6 Astra (OpenAI), ChatGPT Pro mode, effort Pro, S6\_10\_PAPERS continuation, 19 November 2026: used and rechecked in this run; status-vocabulary reconciliation; metadata reconciliation only; no new mathematical check. \emph{Pre-continuation recorded proof scope.} The night-run re-derivation was by GPT-6 Astra (OpenAI, ChatGPT Pro mode, effort Pro), 16--18 November 2026. Also checked in this writing by GPT-6 Astra (OpenAI), ChatGPT Pro mode, effort Pro, 18 November 2026. The two-component ruled spectral calculation, graph path, endpoint ribbon normalization, degree-one internal sheets, and incompatible conductor exponents on the same vertical sheet were re-derived. The one-component graph alternative was checked separately. A literal quad typo in the source is a typesetting correction only}{\ref{night:P4.55}}

In the remaining degree-two horizontal sector after P4.55,
suppose that the generic transverse curve along the other
image component $B$ has at least three normalization branches.
The unique image-graph cycle cannot be concentrated at one
point with three local critical-curve branches, namely the
smooth branch $A$ and two local branches of $B$.
\end{proposition}
\begin{proof}
\noindent\textbf{Step 1.} Write these two local branches as $B_+,B_-$ and the
point as $p$. Let $W_1,\ldots,W_c$ be the normalized
conductor components mapping onto $B$, with degrees
$d_i>0$ over $B^\nu$. Put
\[
                         r_B=\sum_i d_i\ge3.
\]
This is the number of generic normalization branches
of the full transverse plane curve along $B$.
The normalization conductor has components
$H,R,W_1,\ldots,W_c$ and a tree graph. Its total
point-identification contribution is $c+1$.
The horizontal section $H$ maps with degree two to
$A$, while the rational vertical curve $R$ maps
with degree one. Distinct points of $H$ are not
identified in the normalization $N$. Distinct
points of $R$ are not identified either: a
contracted vertical tree effecting such an
identification would create a cycle in the ruling
fibre. All conductor branches below are counted
on their normalization.

\noindent\textbf{Step 2.} We first record the exact spectral bound used in
the proof. Let $E=N_{A/Y}^*$, so $\deg E=2$ by
adjunction and $c_1(Y)\cdot A=0$. On $\mathbb P(E)$
use the quotient convention, tautological class
$\xi$, and fibre class $F$. Thus
\[
          \xi^2=2,\qquad \xi F=1,\qquad F^2=0,
          \qquad K_{\mathbb P(E)}=-2\xi.
\]
The actual normal derivative maps of P4.55 identify
$H$ with the normalization of the degree-two
spectral component $\Gamma_H$, and identify the
third sheet with a section $\Gamma_R\cong A$.
Remove the common zero divisor of the cubic tensor,
of degree $b\ge0$. For an integer $t$ write
\[
       \Gamma_R=\xi-tF,\qquad
       \Gamma_H=2\xi+(t-b)F.
\]
If $\delta_H$ is the normalization colength of
$\Gamma_H$ and $n=\Gamma_H\cdot\Gamma_R\ge0$,
adjunction and the intersection product give
\[
 \begin{split}
       g(H)+\delta_H&=p_a(\Gamma_H)=1+t-b,\\
       n&=4-t-b,\\
       g(H)+\delta_H+n&=5-2b.
 \end{split}                                                   \tag{R20.1}
\]
These are the actual residual components of the
cubic, with their common zero divisor retained.
Since $g(H)\ge1$, they imply $b\le2$.

\par\medskip\noindent\textbf{The exact graph when $c\ge2$.}
Let $r_A$ be the number of normalized conductor
points over the branch $A$ at $p$. It is two if
$H\to A$ ramifies there, and three otherwise.
Let $r_+,r_-$ count the points over $B_+,B_-$.
Every $W_i$ is surjective to $B^\nu$, so
$r_+,r_-\ge c$. Let $n_p$ be the number of
normalization points of $S$ over $p$.
Their nonempty supports in $\{A,+,-\}$ pairwise
intersect: distinct local hypersurface components
intersect in a curve.

\noindent\textbf{Step 3.} A singleton support $\{A\}$ forces every point to
carry $A$. Then $n_p\le r_A$ and the graph
contribution is at least $r_++r_-\ge2c>c+1$.
A singleton $\{+\}$ forces every point to carry
$+$ and gives contribution at least
$r_A+r_-\ge c+2$; the same applies to $\{-\}$.
Thus every support has at least two elements, and
\[
 \text{contribution at }p
       \ge\frac{r_A+r_++r_-}{2}\ge c+1.
\]
The whole global budget is $c+1$. Equality is
forced throughout:
\[
                 r_A=2,\qquad r_+=r_-=c,\qquad n_p=c+1.
                                                               \tag{R20.2}
\]
Every point over $p$ has exactly two normalized
conductor branches, with different image labels,
and all graph identifications occur at $p$.
The horizontal cover ramifies over $p$.
Each $W_i$ has precisely one point over each of
$B_+,B_-$ and has valence two; $H$ and $R$ have
valence one. The connected tree is therefore
the path
\[
                    H-W_1-W_2-\cdots-W_c-R.             \tag{R20.3}
\]
Choose the label $+$ at the $H$ endpoint. The next
point on $W_1$ has label $-$. Every internal point
has support $\{+,-\}$, so it joins that point to
the $+$ point of $W_2$. Induction gives the $+$
point on the left of every $W_i$ and its $-$
point on the right. The $R$ endpoint carries $-$.

\noindent\textbf{Step 4.} The endpoint hypersurface factors, denoted
$f_H,f_R$, thus have supports $\{A,+\}$ and
$\{A,-\}$. The $c-1$ internal factors $f_j$
have supports $\{+,-\}$. These are all the local
irreducible hypersurface factors.

\par\medskip\noindent\textbf{The actual length-two ribbon.}
The set-theoretic intersection of $f_H$ and $f_R$
is $A$ alone. Along generic $A$, the first factor
supplies two ordinary-triple branches and the
second supplies the third, so their intersection
multiplicity is two. Their one-dimensional
complete-intersection ring is Cohen--Macaulay,
without a zero-dimensional embedded component.
A coordinate $z$ on the smooth $A$ is a
nonzerodivisor. The ring is finite free of rank
two over $\mathbb C\{z\}$, and its quotient by
$z$ has length two.

\noindent\textbf{Step 5.} If both factors had no linear term, their
restrictions modulo $z$ would belong to
$(x,y)^2$. The Artinian quotient would contain
the independent classes $1,x,y$, contrary to
length two. At least one factor is smooth.
It cannot be $f_H$: a smooth surface is its own
normalization and supplies only one generic
preimage of $A$, whereas this factor supplies
the ramified double cover. Thus $f_R$ is smooth.
Take the actual ambient coordinate $u=f_R$.
On $u=0$, the divisor of $f_H$ is exactly twice
the smooth $A$, so in coordinates
$A=(u=v=0)$ it has restriction $v^2$ times a
unit. Weierstrass preparation in $v$, followed
by depression and using $f_H\in(u,v)^2$, gives
the exact expression
\[
                    f_H=\text{unit}\,(v^2+u^2g(u,z)).     \tag{R20.4}
\]
All changes preserve $A$. The full cubic normal
tensor is therefore, up to a nonvanishing scalar,
\[
        \left(\prod_j f_j(0,0,z)\right)
                     u\bigl(v^2+u^2g(0,z)\bigr).          \tag{R20.5}
\]
The coefficient of $uv^2$ in the last factor is
one, so the common zero order is exactly the
sum of the orders of $f_j|_A$.

\par\medskip\noindent\textbf{The pole on the smooth rational branch.}
On the smooth branch $f_R=0$, its conductor is
the restriction of the product of the other
factors. The factor $f_H$ contributes $2A$.
Each $f_j$ cuts out only $B_-$ there. This
normalization branch is already smooth, and
the pole support on its minimal resolution is
SNC. Hence $A$ and its $B_-$ curve are smooth
and transverse. Choose the latter's equation
$z=0$ on this surface, with $z$ still a
coordinate on $A$, and extend it to the ambient
threefold. Then
\[
             f_j|_{f_R=0}=\text{unit}\,z^{e_j},
                  \qquad e_j\ge1.
\]
Their sum $e=\sum e_j$ is the generic conductor
exponent of the global component $W_c$ on
this branch. By (R20.5) it is also the local
common cubic-zero order. The branch on
$f_R=0$ is smooth; the plane-curve conductor
formula therefore gives
\[
          e=\sum_{\text{other branches}}I
                         \ge r_B-1\ge2.
\]
Combining $b\ge e$ with (R20.1) forces
\[
     b=e=2,\quad r_B=3,\quad
     g(H)=1,\quad \delta_H=n=0.                         \tag{R20.6}
\]
Thus $\Gamma_H$ is smooth and disjoint from
$\Gamma_R$, and there is no further common
zero of the cubic.

\par\medskip\noindent\textbf{The exact normalization at the ramified endpoint.}
The cover $H\to A$ ramifies at $p$ by (R20.2).
Equation (R20.4) presents its spectral cover as
$w^2+g(0,z)=0$. After translating its double
root, smoothness from (R20.6) gives
$g_z(0,0)\ne0$. Hence
\[
 \begin{split}
        N_H&=\{w^2+g(u,z)=0\},\\
        (u,w,z)&\longmapsto(u,v=uw,z)
 \end{split}                                                   \tag{R20.7}
\]
is the actual normalization near that point.
It is smooth, finite, and birational onto the
branch (R20.4); off $u=0$ its inverse is
$w=v/u$. Coordinates $(u,w)$ solve
$z=\psi(u,w)$, with $z|_H$ of order two in
$w$. The minimal normalization resolution
has no exceptional curve over this smooth
point, and $H$ is $u=0$.

\noindent\textbf{Step 6.} The other conductor branch $W_1$ meets $H$
transversely by SNC. Thus $u$ is a coordinate
on $W_1$ and has order one on its image
$B_+$. By (R20.2), $W_1$ has only one point
over the normalized $B_+$ point, so its
global degree is $d_1=1$. The same smooth
SNC argument on $f_R=0$ gives $d_c=1$.
Since $\sum d_i=r_B=3$, the case $c=2$
would have sum two, while $c>3$ is impossible.
Therefore
\[
                 c=3,\qquad d_1=d_2=d_3=1.              \tag{R20.8}
\]
There are exactly two internal factors, and
$e_1=e_2=1$.

\par\medskip\noindent\textbf{The conductor mismatch.}
Each $f_j$ intersects $f_H$ only along $B_+$.
Its pullback to the actual smooth normalization
(R20.7) consequently has divisor supported
only on $W_1$. Its restriction to $H$ vanishes
to order two: the order on $A$ is one and
$H\to A$ has ramification index two.
Transversality of $H$ and $W_1$ makes the
divisor exactly $2W_1$. Thus the generic
intersection multiplicity of $f_H$ with
each internal factor along $B_+$ is two.

\noindent\textbf{Step 7.} The path (R20.3) identifies the same global
component $W_3$ on the smooth $R$ branch
over $B_-$ and on an internal factor over
$B_+$. At $B_-$ its conductor exponent is
$e=2$. At $B_+$ it has intersection
multiplicity two with $f_H$ and positive
intersection multiplicity with the other
internal factor. Its conductor exponent
is at least three. These cannot be two
values of the coefficient on one global
irreducible conductor component.

\noindent\textbf{Step 8.} The full marking is also explicit. Over
$B_+$ the factors $(f_H,f_1,f_2)$ carry
$(W_1,W_2,W_3)$; over $B_-$ the factors
$(f_1,f_2,f_R)$ carry the same ordered
global components. If $p_+,p_-\ge1$ are
the intersections of the internal factors,
the conductor vectors are
\[
         (4,2+p_+,2+p_+),\qquad
         (1+p_-,1+p_-,2).
\]
There is no missing intrinsic branch-delta
term: the exponent two of the smooth $R$
branch forces its two intersections to
equal one, which makes both other generic
$B$ branches smooth as well. The third
entries already give the contradiction.

\par\medskip\noindent\textbf{The case $c=1$.}
Now the conductor graph has $H,R,W$ and
budget two. The points of $W$ over $B_+$
and $B_-$ cannot be identified in $N$,
since that creates a loop in the conductor
tree. Points carrying the two labels are
distinct; their supports intersect, so
both carry $A$. Each contributes at least
one, exhausting the graph budget. Each
has exactly one $A$ and one $B$ branch,
and no other identification occurs.

\noindent\textbf{Step 9.} If $H$ ramifies over $p$, these two
$A$ branches are $H$ and $R$. There are
only the factors with supports
$\{A,+\}$ and $\{A,-\}$. The length-two
ribbon argument makes $f_R$ smooth.
No other factor contains $B_-$, and this
smooth branch has no singularity there.
Thus $B_-$ cannot be a critical curve.

\noindent\textbf{Step 10.} If $H$ is unramified, the three normalized
$A$ preimages are distinct. Two occur
at the $B$ attachments and the third
is an $A$-only normalization point.
There are three factors, each supplying
one generic $D_4$ branch. Every pair
meets only along $A$, with multiplicity
one. Its complete-intersection curve
is the reduced smooth $A$: it is pure,
generically reduced, and has no other
component. Modulo $z$, length one forces
independent normal linear terms.
All three factors are smooth near $p$
and their only pairwise intersection is
$A$. Neither $B_+$ nor $B_-$ can be a
critical branch. This final contradiction
completes the proof.
\end{proof}

The exact spectral intersections, length-two
ribbon basis, normalization equation, scalar
cubic coefficient, endpoint-degree list and
marked conductor vectors are executed in
\path{code/part4/degree_two_triple_point.py}.
Its \path{.json} and \path{.out} certificates
report PASS.

